% Generated from Bliss Pages Final Edition.idml and 89 website articles; do not edit by hand.
\BookColophon{书　名\BookTab{}幸　福\BookLineBreak{}著　者\BookTab{}李葆春\BookLineBreak{}排　版\BookTab{}李葆春\BookLineBreak{}装帧设计\BookTab{}王　灏\BookLineBreak{}出　版\BookTab{}Duetone Corp.\BookLineBreak{}胶版印刷\BookTab{}Marquis Book Printing Inc.\BookLineBreak{}数字印刷\BookTab{}Blurb Inc.\BookLineBreak{}版　次\BookTab{}2019 年 4 月多伦多第一版\BookLineBreak{}\BookTab{}\BookTab{}2019 年 4 月多伦多第一次印刷\BookLineBreak{}规　格\BookTab{}Tradebook 6"×9" (152 mm×229 mm)\BookLineBreak{}标题字体\BookTab{}冬青黑体 | San Francisco Pro Display  \BookLineBreak{}正文字体\BookTab{}方正新书宋体 | Minion Pro  \BookLineBreak{}字　数\BookTab{}120 千字\BookLineBreak{}印　数\BookTab{}1 – 600 册\BookLineBreak{}国际书号\BookTab{}ISBN 978-1-9990777-0-9\BookLineBreak{}Copyright © 2019 by Baochun Li\BookLineBreak{}Printed in Canada and the United States of America\BookLineBreak{}版权所有，侵权必究}
\BookTitlePage{幸福}{李葆春}{幸\\福}{李\\葆\\春}
\BookForeword{序}
\BookParagraph{让我从哪儿说起呢？}
\BookParagraph{大概是二十年前吧，那时候我和李葆春都还在UIUC上学，有一回 UC Berkeley 的 Lotfi Zadeh 教授来讲学。如果你对学术圈不熟悉的话，这个Zadeh教授可不是一般人，他当年几乎是一个人独创了“模糊数学”这个领域，后来的什么“模糊集合”、“模糊算\BookLineBreak{}法”、“模糊逻辑”、“模糊语言”……模糊这模糊那的全是他一个人弄出来的。}
\BookParagraph{二十年后，我成了一个在大公司混饭吃的碌碌无为的前科学工作者，而李葆春却像Zadeh老师一样，成了名校教授、著名学者，IEEE Fellow。}
\BookParagraph{不过这不是我想说的。}
\BookParagraph{我想说的是，李葆春老师这些年其实写过的字远远不止这本书里的十几万字，但是李老师的那些文字，大部分看到这本书的朋友们都不会有机会看到 —— 就算看到了估计也看不明白 —— 那些论文、著作都涉及专业学术领域，就像Zadeh有关模糊集合的那篇论文，引用数高达九万多，但是我保证你没看过。}
\BookParagraph{你在这本书里看到的聪明的、有趣的、深邃的李葆春，还只是他的一小部分。不过这可能也是所有有趣的人的一个共同特征\BookLineBreak{}吧 —— 对自己的专业、工作以外的世界兴致勃勃，充满了令人惊喜的洞察和令人动容的热爱，而且，必须幽默。}
\BookParagraph{但其实这还不是我最想说的。}
\BookParagraph{那天Zadeh老师的讲座都讲了什么算法什么学术发现我一概不记得了——这可能也是我后来在学术界一事无成的原因吧。但是他在讲座最后讲过的一个段子我到现在还记得。}
\BookParagraph{他说我在学术界混了大半辈子，最后送给同学们一个人生忠告吧：}
\BookParagraph{Friends come and go, but enemies accumulate.}
\BookParagraph{我操没想到老头猛地来了这么劲爆的一句毒鸡汤，当时阶梯教室里很多快睡着的人一下子都醒过来了。}
\BookParagraph{回想了一下，这句话还真对，我从小到大 —— 从小学中学到大学到出国到工作到回国，前前后后有过很多朋友，但是随着身边城市的变化以及人生阶段的改变，一个个都从我的生命中离开了，就像模糊逻辑老头说的：}
\BookParagraph{朋友来来去去。}
\BookParagraph{想了半天，好像只想出极有限的几个反例 —— 李葆春是我第一个想到的。}
\BookParagraph{从燕南园到北大附小到美国小镇尚佩恩到芝加哥到后来的纽约、上海，好像李葆春从来没有离开过。}
\BookParagraph{一九九六年夏天，我和李葆春都在芝加哥实习，平时在芝加哥住，星期五下班回学校，星期天傍晚再回芝加哥。}
\BookParagraph{我们上学的大学城距芝加哥开车将近三小时，我们经常结伴一起开车。算起来每次三小时的路程，整整一个夏天我们应该说了很多话——用句套话说是：谈人生谈理想谈爱情。}
\BookParagraph{就在刚才，我试图从无数个从小到大有李葆春出镜的回忆场景中找到一个开始的地方的时候，不知道为什么就想到了九六年夏天我们坐在来去芝加哥路上的这个画面。}
\BookParagraph{可能是这本书闹的吧，有一种熟悉的、夏天的味道。}
\BookParagraph{印象里那是一个黄昏的画面，汽车穿过美国中西部一望无际的田野，路边是大片大片的玉米地，我们打开车窗，风呼呼地吹进来。}
\BookParagraph{我那辆别克旧车车顶内饰的布开胶松弛了，平时就那么耷拉着，但有风的时候会往下鼓起来，像是一张帆。}
\BookParagraph{像是我们要去远航。}
\BookDate{褚明宇}
\BookDate{2019 年 4 月 8 日}
\BookTOC
\BookPreface{前　言}
\BookParagraph{外公七十多年前写过一本书叫《龙虫并雕斋琐语》，收录了他当年在报纸上发表的一些小品文。外公从《姓名》开始，聊到《请客》，到《开会》，到《跳舞》，到《战时的书》，那些时而针砭时弊、时而风趣幽默、时而又引经据典的文字，有一种穿越时空的力量，这么多年以后，竟然还能令人沉浸于其中，击节称叹。这种只有文字才有的独到魅力，一直吸引着我，让我对它情有独钟。就像外公那篇《写文章》里说的：「看到兴高采烈、手舞足蹈的当儿，虽未至于五体投地，却已一心欲仙。」}
\BookParagraph{这本书收录了我十几年来写的几乎所有文字。当年开始写的初衷不外乎是依靠文字记下来一些感想和生活，留给多年以后的自己和家人，也可以赠送给朋友看。这些文章都不长，一方面是因为写作的时间不多，另一方面我偏爱超短篇幅的文字，因为这样的文字有点儿像浓缩咖啡，清香而醇正，有一种意犹未尽的感觉。}
\BookParagraph{两年前因为参加了褚明宇在分答开的向上移动社区，开始回答一些社区里的问题。这些文字既然写出来了，我就同时发在微博里。后来问题逐渐多了起来，我写得也更经常了，有时一个月可以写好几篇。既然是回答问题，不太可能压缩得过于精炼，所以文章写得也逐渐长了起来。不过一方面因为长微博有字数限制，另一方面因为写得越长就越难写得好，每篇文章还是不会超过一两千字。题材和问题有关，五花八门，什么都聊两句。不过似乎还是和情感有关的文字多一些，怪不得有的读者管我叫「情感博主」。好在情感是永恒的话题，也许多年后依然值得一看。}
\BookParagraph{既然已经费尽心力写出来了，我一直很想把所有的文字汇集成一本书。书可以印刷出来，泛黄的纸张让人成就感油然而生。既可以赠送给好朋友，算是一份特别的礼物；又可以留给家人将来永久收藏。几个月前女儿用 Blurb Bookwright 排版印刷了一本画册，图文并茂，令人羡慕。理论上来说我也可以步其后尘，把我自己这本书也印出来。可是由于我不会排版，这个似乎很浩大的工程就一直搁置在一边。}
\BookParagraph{恰巧这几天因为一个偶然的机会，我发现了一个排版的小窍门，可以直接把我平时存在 Day One 里的文字导入。找到这个窍门就像玩游戏找到了通关的秘诀，一发而不可收。我利用去香港开会的几天零星时间，从挑选中英文字体开始，对文字做了全面的修订，完成了这本书的编辑和排版。}
\BookParagraph{外公的《龙虫并雕斋琐语》虽然好看，但是有一个缺点——写作的年代和现在离得到底过于久远了。他如果「穿越」到今天，将这些文字发在微博里，不加个「V」什么的，说不定门可罗雀，读者还不如当年发表在报纸专栏上多。俗话说：「妻子是人家的好，文章是自己的好。」我的文字水准虽然无法与外公当年相提并论，但是我一直希望能融入一些当今时尚的元素，一点儿「modern twist」。如果几十年后翻出来再看，敝帚自珍，触景生情，还能有一点儿「遥想公瑾当年」，一点儿「nostalgic」，那就再好不过了。}
\BookDate{2019 年 2 月 27 日，多伦多}
\BookPart{幸福}{幸\\福}
\BookArticle{房　子}{房子}
\BookParagraph{又有朋友迁入了又大又新的房子，参观的人络绎不绝。不禁想起美国的劳动人民经常爱犯的错误，管买房子说成「买个家」，就好像中国人说买个花园别墅一样。其实，房子就是房子，是众多的投资中的一种。把房子说成家或别墅，不外乎是为提高价码而进行的感情上的诱导。投资都是有风险的，所以与投资有关的决定，感情应该中立。}
\BookParagraph{但是话又说回来，凡人都想住个又大又漂亮的房子，我更不能免俗。从没有房子变成有房子，从小房子喜迁大房子，从地上的房子换成船上的房子。有了房子，就要买上档次的家具；有了车库，怎么也得再配上一大一小的汽车。有了花园，要铺上绿得可爱的草坪；有了邻居，要比拼五颜六色的彩灯。房子成为上好的聚会时的谈资，收入的晴雨表，成功的象征，孩子的未来，和自我价值的实现。当孩子们欢呼雀跃上上下下追打嬉戏时，因为有了房子，我们确实仿佛找到了那梦寐以求的「家」。}
\BookParagraph{房子象青春的岁月，只能前进，不能后退。住八个人一屋的宿舍的时候，点着蜡烛，各自诉说着心中的女孩。和女朋友校园里散步的时候，盼着遥远的将来那能落脚的小屋，和自己亲手缝制的窗饰。在两室一厅的公寓里十几个帅哥美女喝酒狂欢的日子里，房子成了可有可无的虚设。房子象绊脚石，石头越重，那火红的日子就越遥不可及，追也追不上。总想有那么一天，能走出房子，去体会那转瞬即逝的幸福。}
\BookDate{2005 年 3 月，多伦多}
\BookArticle{幸　福}{幸福}
\BookParagraph{一九九一年清华的暑期小学期，摆脱不了影响心情的干燥和炎热。哥儿几个商量，好不容易这么闲的大好时机应该充分利用起来。也不知是谁忽发奇想，要组织我们零三班自己的舞会。可惜清华男女生比例之离谱远近闻名，我们班三十四个同学，就算六个女生全到了，舞会也没法儿开。一哥们儿出了一招：直接邀请全计算机系所有三十二个女生参加我们班的舞会。}
\BookParagraph{计算机系的所有女生都住在七食堂旁边的六号楼，平常未经允许不准自由出入。我带着请柬，第一次走进六号楼，一个房间一个房间地直接表达来意。每次说完离开，就能听到一片大笑的声音。来到四五一房间零五班女生宿舍，敲门进去，发现靠窗上铺的女生，穿着一件乳白色休闲的短袖衫，阳光映在她的脸上。在我介绍来意的时候，她回报以微笑，让人感到温暖。}
\BookParagraph{去年在拉斯维加斯的四季酒店庆祝我们结婚二十周年纪念日那天，我想起了发请柬这事儿。心情还是有点儿紧张，到底绝大多数的女生我都不认识。可是，在那个没有手机没有微信、每月还在发粮票的时代，这大概是唯一的办法。我们组织的舞会出乎意料地成功。来了很多系里其他班的女生，在音乐的渲染下，场面轰动，高潮迭起，气氛热烈。她也来了，披肩的头发，穿着一身白色的短袖衫。远远看到，朴素自然，落落大方。}
\BookParagraph{暑期小学期结束后，接着是在京郊防化兵工程学院的三周军训。在灯光篮球场的一次舞会，再次见到她，清丽的脸似乎比原来晒黑了一点儿，头发为了军训打成两个小辫儿披在肩上。两三次舞会以后，我们也算认识了。军训结束的那一天返校，我一个人到七食堂门口等她。等了好久也不知她来不来的时候，看到她从六号楼出来，穿着一条白色的连衣裙，印象里太阳正好快要落山了，显得特别轻盈透亮。站在食堂门口，我问她有没有时间一起在校园里走走。她明显犹豫了一下，不过还是答应了。}
\BookParagraph{走到清华学堂边上，几个男生在大草坪上弹着他们的吉他，享受着放假前这个轻松的晚上。她忽然对我说：「我有男朋友了。」}
\BookParagraph{回到六号楼的时候，我问她有没有通信地址，她爽快地告诉了我。暑假我住在北大燕南园，夏天的晚上知了叫得睡不着。我给她写信用的是一种只有清华有卖的天蓝色横线四百字稿纸，每张都特别的薄。}
\BookParagraph{开学了。晚上背着书包上自习的时候，我一个教室一个教室地站在后窗找她坐的位置，找到了就过去和她打招呼。有时候上完自习，推着我的二八永久自行车，在回宿舍的路上走得慢一点儿，聊一会儿天。}
\BookParagraph{九月快过完了的一天，她从校园里彻底消失了。该上的大课，平时上晚自习的一教到五教，还有清华学堂和图书馆全都找遍了，没有见到她的踪影。一年后她告诉我，男朋友从上海来北京找她玩儿，她课不上了陪他，景山，天安门，圆明园，到处游览北京。等到她重新开始回来上课以后，我就再也没有见过她。}
\BookParagraph{我二十岁生日那天请客，在大学生之家请哥儿几个大吃大喝。晚上十点多钟，她下自习路过时隔着栏杆看到我们，明显停了几秒钟。大学生之家的大喇叭里，正放着黑豹乐队的《无地自容》：人潮人海中又看到你 / 一样迷人一样美丽 / 慢慢地放松慢慢地抛弃 / 同样仍是并不在意 / 你不必过分多说你自己清楚 / 你我到底想要做些什么 / 不必在乎许多更不必难过 / 终究有一天你会离开我。}
\BookParagraph{新年快要到了的一个晚上，九号楼十点四十五准时熄灯。大家忙着点蜡烛，映着摇曳的烛光，我用纯蓝墨水的钢笔给她写了一张卡片：「这几个月跟你在一起的时候特别高兴，也许今后不可能再有机会了。」}
\BookParagraph{新年的钟声快要敲响的时候，她突然出现在九号楼我宿舍房间的门口。什么都没说，只是送给我一个她自己亲手做的小钱包和一张小卡片。钱包是打开金属暗扣里面好几层的那种，用一种质地厚实而光滑的纸做的，纸上是一幅油画，油画里有一个小男孩和一个小女孩，小男孩手里拿着一朵玫瑰花。打开小卡片，上面写着：我想会有的。 那一年，她十八岁。 打开卡片的那一瞬间，我一个人坐在宿舍上铺的床上，感到了一种时间凝固了一般的幸福。}
\BookDate{2011 年 7 月，2017 年 2 月，2017 年 7 月，多伦多}
\BookArticle{生　命}{生命}
\BookParagraph{曹建农老师在Facebook上首发的第一篇文章，是在他自己几天前发的一张图片的评论里，而图片显示的是一本名为《The Big Questions: A Short Introduction to Philosophy》的书的扉页。图片有十八个「Like」，而其评论只有五个。主要的原因大概是，如果用 Facebook 的手机客户端来浏览 newsfeed，图片自然非常抢眼，但几乎没有人会在几天以后去仔细翻看一张旧照片新添的评论。}
\BookParagraph{可惜的是，一段简短文字的力量和震撼，远远超过了Facebook里俯首皆是的照片。曹建农老师在他的短文中提到了几天前在香港理工大学因为工作压力而跳楼自杀的一位高级管理人员。他说：}
\BookParagraph{———}
\BookParagraph{校长哽咽着用慢于平时一倍的语速悼念我们离世的同事，我们站立默哀。已经过去一个星期了，心里仍然不时总在想，不知道他当时是以怎样的心情从自己的办公室走出来，落两层楼，徘徊于露台，然后跃落而去；再在这之前，他又是以怎样的心情写下一封以致全体同事开头的遗书。}
\BookParagraph{一位优秀善良的资深同事就这样走了。如果在他深度困惑的时候，有人可以觉察，倾听分担一下；更如果他自己可以觉察，向别人表达述说一下，这一切或许都可以不发生。如果连失去生命都不可惧，那么还有什么不可以言说。可是，这就是我们面对的现实，就是那么困难去相信能有人可以倾听、理解，那么全然依靠我们自己的觉悟。}
\BookParagraph{这也是我为什么会再翻出这本书来读。所有的哲学问题其实都是最平凡而基本的问题，都是从古希腊开始被问了两千多年却还没有答案的问题，但是我们还是必须不断再问。从什么是自我，什么是生活的意义，再问到什么是勇气、正义、英雄，什么是好的生活、好的教育。我们能觉悟吗？我们对悲剧事件的反问应该得出积极的结果，使我们更加努力地去生活，更加勇敢。}
\BookParagraph{今天在回香港的飞机上，葆春拿出新买的书在读，对我说，还有一本，你拿去读。是我看到很久却一直没买的龙应台的一本书，题为《目送》。读到里面的一篇「跌倒」，这样写道：「谁教过我们，在跌倒时，怎样的勇敢才真正有用，怎样的智慧才能度过？跌倒，怎样可以变成远行的力量？失败，为什么往往是人生的修行？何以跌倒过的人，更深刻，更真诚？」有人教过我们吗？我们教过别人吗？}
\BookParagraph{———}
\BookParagraph{这段文字是曹老师在大连返回香港的航班上手书在便笺纸上的。因为我们同行，我看着他当时不假思索，一挥而就，一气呵成，然而却写得逻辑清晰，非常精彩，读来酣畅淋漓。我忽然想起前天与曹建农老师、贾晓华老师一起在大连滨海西路的木栈道上边走边聊的情景。记得当时曹老师提到一个段子：「中国人上学时忙着各种考试，大学毕了业忙着结婚，结婚后忙着生孩子带孩子，孩子大了忙着催孩子结婚，退了休忙着给孩子看孩子。」 虽是笑谈，颇有冯唐那篇《活着活着就老了》之风范。而贾晓华老师则忆起当年熟读的小说《钢铁是怎样炼成的》里的主人公保尔 · 柯察金的一句名言：\BookLineBreak{}「人最宝贵的就是生命，生命对于每个人来说只有一次。人的一生应该这样度过：回首往事，他不会因为虚度年华而悔恨，也不会\BookLineBreak{}因为碌碌无为而羞愧；临终之际，他能够说：“我的整个生命和\BookLineBreak{}全部精力，都献给了世界上最壮丽的事业 —— 为解放全人类而斗争。”」}
\BookParagraph{也许是出于一种好奇，我读了一段曹老师提到的那本名为《大问题》的书。这本书是一部西方哲学简史，深入浅出，厚积薄发，发人深省。从苏格拉底到柏拉图，从亚里士多德到康德，全书阐述的是几个基本而又艰难的哲学问题，其中最重要的一个就是「生命的意义」。}
\BookParagraph{在我所受的教育里，哲学一直处于一种可有可无的地位，甚至有被妖魔化的趋势。「哲学」这两个字，似乎永远与「马克思主义哲学」联系在一起，永远与政治课如出一辙，枯燥无味。而看了一段这本英文书，却有豁然开朗的顿悟感，因为其实哲学是一种开宗明义的思辨，它与现实生活水乳交融，无处不在。哲学是对极为现实的各种问题的看法 （ideas），这些看法每个人都有，只不过有的浅薄、有的深刻而已。就像书的引言里所说：}
\BookParagraph{「The difference between philosophy and the popular alternatives is ultimately one of quality — the quality of ideas, the thoroughness of understanding. Because we all live by our ideas anyway, the choice becomes not whether to do philosophy or not do philosophy, but whether to accept a cheap and unchallenging substitute or to try the real thing. The aim of this book is to give you an introduction — to the real thing.」}
\BookParagraph{也许，对于那位留下遗书跃身而下的同仁来说，他需要的是可以推心置腹的好朋友，而非哲学。也许，身处娱乐八卦豪宅名车满天飞的新时代，我们需要的是出人头地飞黄腾达，而非哲学。也许，在高考即将取消英语、而网络词汇大行其道之际，我们需要的是基本的读写能力（literacy），而非哲学。但是，一个普通的周五晚上，我却莫名其妙地较起真来，任性而天真地憧憬着一种假象。在我虚无缥缈的想象中，我们不仅爱看《小时代》，更爱看《大问题》；不仅爱看高圆圆黄海波的《咱们结婚吧》，更爱看Michael Sandel 教授在哈佛大学苏格拉底式的演讲 《Justice: What’s the Right Thing To Do?》。因为后者如同千年古刹里的得道高僧，如同深红的原装可口可乐，如同崔健的《南泥湾》，是「The Real Thing」。}
\BookDate{2013 年 12 月，香港}
\BookArticle{The Cranberries}{The Cranberries}
\BookParagraph{The Cranberries 的歌里，我最喜欢的还是那首「Dreams」。每次听到这首歌，都会有一种特别浪漫、特别开心的感觉，大概是因为每次都能想起二十年前的一部浪漫喜剧「You’ve Got Mail」。电影里当这首歌响起的时候，正好是 Meg Ryan 演的 Kathleen 自言自语的旁白：}
\BookParagraph{「What will NY152 say today, I wonder. I turn on my computer. I wait impatiently as it connects. I go online, and my breath catches in my chest until I hear three little words:}
\BookParagraph{You’ve got mail.}
\BookParagraph{I hear nothing. Not even a sound on the streets of New York, just the beating of my own heart.}
\BookParagraph{I have mail. From you.」}
\BookDate{2018 年 1 月，多伦多}
\BookArticle{空中小姐}{空中小姐}
\BookParagraph{王朔的小说里，我最喜欢他的第一部作品《空中小姐》。很多人说这部爱情小说笔法略显生涩稚嫩，我则偏偏喜欢文字里那种强烈的时代感，情景交融，简洁明快，丝毫不拖泥带水。}
\BookParagraph{——}
\BookParagraph{「你知道我现在最大的愿望是什么？」}
\BookParagraph{「什么？」}
\BookParagraph{「临死前，最后一眼看到的是你。」}
\BookParagraph{「小傻瓜，那时我早老了，老得不成样子。那时，也许你想看的是孩子。」}
\BookParagraph{「不会的不会的。」}
\BookParagraph{——}
\BookParagraph{这样连中学生都能写得出的伏笔，这么多年以后再看，和后来作品中那种北京侃爷的玩世不恭相比，竟然更加令人感动。}
\BookParagraph{朔爷三十多年前动笔写《空中小姐》的时候，正好二十六岁。和小说里写的一样，那时他刚从海军复员，还不认识舞蹈学院的沈旭佳，更没见过北京电影学院的徐静蕾或者四川妹子王子文。就像航班在跑道上缓缓滑行准备起飞，就像刚刚改革开放八十年代初期的中国，一切的变数才刚刚开场。}
\BookParagraph{那时候我特爱看《当代》和《收获》，看余华和王安忆，看得如火如荼。但是我却没有看过《当代》这篇只有三万字的短文，题目叫做《空中小姐》：}
\BookParagraph{「我没再见过王眉，也没得到过她的音讯。有一年，我在北京火车站看见一个女孩背影很像她，我没追上去看，因为她决不可能出现在北京站。即使是休假、公出，民航也给她们飞机乘的。还有一次，我坐缓缓出站的火车和一列天津方向开来的火车相错而过时，有个从车窗往外看的女孩和我对视了半天，直到递次而过的车窗远去。我真的以为那是王眉了，但由于如上的原因，我最终认定是自己看错人了。」}
\BookDate{2018 年 5 月，多伦多}
\BookArticle{快乐至上主义}{快乐至上主义}
\BookParagraph{朋友间请客吃饭出去玩怎么做才不会尴尬？}
\BookParagraph{各位老师好。我是一个女生，毕业一年多，吃喝玩乐特别不上道。不管是女性朋友或是男性朋友一起出去，给人的感觉就是这人小气，爱纠结。唉，好吧！自己确实挺心疼钱的。在农村长大，对城市生活和城市朋友间见面的消费不太清楚，对这些吃法玩法也不感兴趣，怎么做才能不花很多钱又能让大家高兴的相聚呢？（问题来自「向上移动」）}
\BookParagraph{说实话，我虽然不是农村长大，但是和你性格很相似：小气、爱纠结，挺心疼钱的。这大概是因为我的性格从小就比较谨慎小心，可能是小时候父母言传身教的结果吧。}
\BookParagraph{如果允许我做一个自以为是的比较，我觉得巴菲特虽然非常成功，性格其实跟我也差不多。他一生积累了八百亿美元的财富，但除了每个孩子给了十亿美元之外，大多数捐给了盖茨基金会和他自己的基金会。他至今还住在五十年代花了三万美元买的一栋普通的房子里。别人问他为什么不搬到好一点儿的房子里住，他回答：「I’m happy there. I’d move if I thought I’d be happier someplace else.」}
\BookParagraph{这句话的要义在于，有时候说不定少花点儿钱，反而更加快乐。}
\BookParagraph{英文有一个词儿叫「mindfulness」。这个词有点儿佛教的味道，大意是无论是工作还是玩儿，尽量「往心里去」。有的时候跟朋友一起吃饭，会觉得钱花了不少但是却很无聊，一点儿都不开心快乐。说不定少花点儿钱，把一起吃饭改成一起骑车或者跑步，反而能更加享受其过程，更加「往心里去」。}
\BookParagraph{所以说啊，又小气又心疼钱不一定是一件坏事儿。逛逛书店再来顿三鲜饺子，开心的体验说不定远远超过去米其林三星的粤菜馆吃饭。无论是提议「断舍离」还是倡导「极简主义」，骨子里都有种「少花钱少买东西反而更快乐」的意味。}
\BookParagraph{当然了，什么事儿都不能走极端是吧。有的时候斥巨资购买设计出色、质量上乘的物品，去环境优雅、海景无敌的馆子吃饭，或者入住服务周到、体贴入微的酒店，绝对是一种享受。}
\BookParagraph{所以我一直相信，没什么意义的消费尽量避免，但是该花的钱一定要花。一方面「斤斤计较」，另一方面「一掷千金」，两者绝不矛盾。而衡量花钱还是不花钱的「性价比」，看的是它是否能带来快乐。虽然有的朋友可能不一定理解，但我想时间长了，总能找到可以一起分享快乐的好朋友。}
\BookParagraph{这就是我的「快乐至上主义」。}
\BookParagraph{祝大家圣诞新年快乐。}
\BookDate{2017 年 12 月，多伦多}
\BookArticle{决　斗}{决斗}
\BookParagraph{李老师，您好。相比您和郭芳老师，微博里有些人的言论总是喜欢反复奚落一些说过蠢话或者犯过错误的人。我是一个戾气很重的人，想问您的问题是：如何避免残忍和不成熟的行为？（问题来自微博私信）}
\BookParagraph{你听说过十九世纪初法国有一个叫伽罗华（Galois）的数学天才吗？}
\BookParagraph{他十几岁时仅用了几年时间就提出了群和域的理论，用这套理论彻底解决了是否可以用根式求解高次方程的充分必要条件这个困扰了数学界三百多年的问题，奠定了抽象代数的根基，并普遍应用于现代编码理论中。}
\BookParagraph{有一次因为一个偶然的机会，他遇到了一个向他倾诉衷肠的舞女。正是因为这个舞女，他接受了和一个枪法很准的对手要和他决斗的邀约。}
\BookParagraph{伽罗华在决斗前一天晚上彻夜未眠。他连夜写了一封长长的信，试图把自己所有未曾发表的数学发现悉数写到信里。天快亮了，他有点儿着急，写道：「不行，我没有时间了，没有时间了。」}
\BookParagraph{他明明知道自己赢不了这场决斗，还是要去。}
\BookParagraph{伽罗华终年仅仅二十岁，却是人类历史上最伟大的数学家之一。现代数学几乎半壁江山，都是基于他首创的抽象代数理论。}
\BookParagraph{相比而言，在社交网络上用一个匿名的账号和别人唇枪舌剑，就跟在公共汽车上几个陌生人为了一点儿小事儿就吵得不可开交一样，既没有什么必要，又没有什么意义。逞一时口舌之快，并不能带来什么「快意恩仇」的感觉，反倒会令看客对你刮目相看。而知书达理，温文尔雅的人，照样可以外柔内刚，笑傲江湖，指点江山。}
\BookParagraph{我有一个缺点：胆儿小。我从小到大最怕各种动物，无论是趾高气扬的公鸡，还是街头流浪的小猫小狗，都能让我绕道而行。因为胆儿小，与其反唇相讥，我认为不如让步隐忍，全身而退。}
\BookParagraph{金庸笔下姑苏慕容家「以其人之道，还治其人之身」的武功虽非绝世，至少会让对手想起「己所不欲，勿施于人」，自然而然地忌惮三分。}
\BookParagraph{您要是确实血气方刚，咱要玩儿就向伽罗华或者普希金学习，玩儿点儿大的。自古英雄出少年，明知山有虎，偏向虎山行。}
\BookParagraph{不过「决斗」之前，最好先提出一套代数理论，或者写几部传世的诗作。这样维基百科也能榜上有名，不枉此生。}
\BookDate{2018 年 6 月，西安}
\BookArticle{Say No}{Say No}
\BookParagraph{李老师，您好！关注您的微博有一段时间了，常看您回答一些年轻人的问题，今天我也有个小问题想请教您。}
\BookParagraph{我即将与父亲同事的女儿去上海玩几天，还有另一个她自己的朋友同行，与我更是初次相识。客观地说我们的关系属于几年未见的、非常普通的朋友，在旅行中做个伴。}
\BookParagraph{由于上海某乐园正值旺季，票价比平时贵，算上一些附加消费大约会达到700元一天。我旅行是用的自己的存款，考虑到研究生开学还会花一大笔钱，我不太愿意花这个钱，所以与另两女生商量前一天晚仍住在那个乐园附近，但去的那天分头行动，她们俩去乐园玩，我回城内看艺术展，晚上再汇合。但她们不放心我一个人，就退掉了房间说不去了，但又说非常遗憾。我感觉自己很扫她们的兴，因为她们确实很想去，所以目前我花这笔钱跟她们一起去乐园仿佛成了维持她们旅行心情的唯一方式。}
\BookParagraph{我感到不解的问题是，我应该如何处理自己的需求和别人（情感）需求的冲突？}
\BookParagraph{1. 在我的价值观里，她们愿意去就去，我不太愿意去就不去，分头行动再汇合即可，大家都开心，不需要为彼此牺牲什么。但我没考虑到她们希望三个人总是在一起，甚至愿意因此不去乐园了。如此一来，她们将去与不去放在了利益的对立面，让我不知道该如何处理类似的事。}
\BookParagraph{2. 父母是普通工薪阶层，700一天肯定觉得贵的，但他们听说此事后，表示愿意他们出钱让我去。我感觉此举意在让我合群和学会迁就别人。多出700我自己也是可以负担得起的，只是基于对存款的规划，不愿意而已。}
\BookParagraph{总结为一句话就是，我大概太贪心，既不太愿意打乱自己的规划，又不想伤害别人的感受。我想这种矛盾或许来自我自己强烈的自我意识，它告诉我应该遵从自己的内心；但是周围环境（例如父母在此事中的举动）又时刻提醒我应该考虑别人的感受，甚至牺牲自己去顺从别人，以维持一种「快乐的大氛围」，所以我感到很纠结。}
\BookParagraph{如果是事实需求的话，我相信总能找到折衷的办法。但当我面对的是情感需求的时候，我就不知道怎么办了。有时候我觉得如果随意洒脱一点谁牺牲迁就一下都不是什么大事，但是有时候这种事情发生得多了，就容易改变一个人的风格。虽然此事只是小事，但我感觉它反映出了我在人生某种大原则上的摇摆不定。希望听听您对此类事的看法，谢谢您！（问题来自微博私信）}
\BookParagraph{十几年前我还算是年轻教授那会儿，每年有一次机会可以和系主任聊一会儿天。系主任「面授机宜」，算是我当时的「mentor」。}
\BookParagraph{有一次他看到我在十几个会议里担任了各种这组委会那委员会的成员，不但没有因为「没有功劳至少还有苦劳」而表扬我，还竟然建议我少接受几个这方面的邀请。}
\BookParagraph{「You should learn to say “no” to people.」他说。}
\BookParagraph{这是当年系主任给我的所有建议中，我印象最深刻的一句。}
\BookParagraph{因为「say no」是一种「桃花潭水深千尺」的哲学，一种「牧童遥指杏花村」的境界。我即使在多年以后的今天，依然觉得这种境界犹如「南阳诸葛庐，西蜀子云亭」般遥不可及。}
\BookParagraph{「Say no」之所以遥不可及，是因为我们从小到大耳濡目染的所有待人接物的道理，什么「在家靠父母，出门靠朋友」啊，什么\BookLineBreak{}「为朋友两肋插刀」啊，什么「毫不利己，专门利人」啊，什么\BookLineBreak{}「个人利益要服从集体利益」啊，什么「人生的意义在于奉献，而不是索取」啊，都与「say no」的哲学南辕北辙。}
\BookParagraph{我之所以一直在用英文的「 say no 」来代替中文的「说不」，是因为前者我觉得很自然，而后者却因为这些从小到大的道理，觉得很唐突，很没有礼貌。}
\BookParagraph{比如我平常跟小女儿建议她早点儿弹会儿琴什么的，她每次只需要一句「thanks but no thanks」，就一定可以把我凌厉的武功招数自然而然地化解于无形。我不但一点儿都没觉得生气，还觉得有点儿俏皮可爱。}
\BookParagraph{而换了我呢，每一次没有狠下心来「say no」，到最后都差点儿沦落到「鞠躬尽瘁，死而后已」的境地。因为在我的「公理系统」里，一旦答应了别人，我绝对不会中途变卦，而且一定要把答应别人的工作做得尽善尽美。}
\BookParagraph{一旦答应了别人，不仅仅可能像你的故事里一样多了一笔不菲的花费，而且还花了自己的时间。几百块钱省下来可以用在很多其他地方，但至少将来还可以再挣；时间则是一种更加无情的资源投入 —— 一旦过去就再也回不来了。}
\BookParagraph{这就是「say no」哲学的最大好处：最初别答应别人，一切就都不会发生。}
\BookParagraph{所以，理直气壮地、不计后果地、自私自利地、不顾集体利益地、「做一天和尚撞一天钟」地说「no」吧！洒脱地牺牲一下别人，迁就一下自己，这就是你「踏破铁鞋无觅处，得来全不费功夫」的人生哲学。}
\BookParagraph{但是我不得不再次强调，这种境界实在太难以企及了。}
\BookParagraph{比如今天有一个陌生人通过微博问我问题，我本来完全可以对她说「no」，可还是一念之差，想着「说不定将来可以成为一个朋友呢」，花了一个多小时时间，给她写了如上的文字。}
\BookDate{2018 年 7 月，多伦多}
\BookArticle{命　运}{命运}
\BookParagraph{李老师您好，我最近两年人生走得特别不顺利，这两天又在人生选择上焦虑得不行。您能帮我指点一下吗？先感谢您！}
\BookParagraph{我本科毕业于上海一所二本院校（我是山西省一本线），大四考研考的是北京师范大学，初试过了，面试没过，难受了两个多月。因为自己心疼父母，不想让父母有太大压力，我就边工作边二战，还是考的北师大。但结果未遂人意，我被调剂到了陕西省委党校。开学几天了，我并不是特别开心，因为感觉自己不是特别适合党校的氛围，还特别想念上海。所以我又想再考了，这次想选个大城市比较好的学校，压力不是特别大的。跟中部城市比，还是想念上海。}
\BookParagraph{您觉得我这样想奇怪吗？身边的学长学姐说时间耽误不起，可以把时间好好利用起来考博。我有点犹豫，不知道自己是不是过于执念。麻烦您帮我指点下吧，谢谢您！（问题来自微博私信）}
\BookSubtitle{一}
\BookParagraph{前两天跟小女儿闲聊，不知道怎么就聊到了一个假想的场景。我问她：在荒无人烟的野外，一只灰熊发现了你，向你虎视眈眈地慢慢走来，你怎么办？}
\BookParagraph{这运气确实没法儿再糟糕了。}
\BookParagraph{但要按我从小到大受到的教育，正确的心态应该是去「挑战命运的安排」，向张海迪同志学习，学习她在残酷的命运挑战面前，没有沮丧，没有沉沦，以顽强的毅力与命运抗争，去做一个社会主义新人。}
\BookParagraph{毛主席说的好，人定胜天嘛。}
\BookParagraph{至少可以跟电影《阿甘正传》里Jenny的名言说的：Run, Forrest, run!}
\BookParagraph{女儿的回答却完全在我的意料之外：}
\BookParagraph{「I will accept my fate and just die, right there.」}
\BookParagraph{「那你为什么不跑啊？」}
\BookParagraph{「No.  I will just die.」}
\BookParagraph{这让我想起了电影《Legends of the Fall》的结局：主角都在年轻时去世了，只有 Tristan 活到了老年。他遇到了一只灰熊，拔出匕首和熊搏斗，画外是 One Stab 的声音：「It was a good death.」}
\BookSubtitle{二}
\BookParagraph{按常理出牌的话，我想一般人遇到这种「命运」，都会跑的。无论如何，「跑」是以顽强的毅力去跟命运抗争，是不屈服于命运的安排。不跑必死无疑，跑还有一线希望，为什么不跑？}
\BookParagraph{所以我对女儿的回答感到非常震惊。}
\BookParagraph{「I will accept my fate.」}
\BookParagraph{这是什么人生哲学？从来没听说过，跟我受到的教育格格不入，然而又令人耳目一新。}
\BookParagraph{在你的故事里，当命运像贝多芬的交响乐一样来敲门的时候，你收到的那张纸上写着：陕西省委党校。}
\BookParagraph{为了回答你这个问题，我还专门在网上找到了这个学校的地址，发现是在陕西省西安市市区。}
\BookParagraph{确实，跟上海相比，我也更喜欢上海。}
\BookParagraph{但是这可不是什么新疆维吾尔自治区吐鲁番市，或者西藏自治区林芝市，这座城市总人口一千二百万，机场可以直飞法国巴黎，有喜来登，有香格里拉，丽思卡尔顿就要开业。}
\BookParagraph{如果考虑买房的性能价格比，西安甚至可以说还超过了上海。}
\BookParagraph{If I were you, I will accept my fate.}
\BookParagraph{And just die right there.}
\BookSubtitle{三}
\BookParagraph{但是，如果不跑必死无疑，跑还有一线希望，为什么不跑？}
\BookParagraph{好吧，那我们就来仔细分析一下因为「跑」而成功的概率。}
\BookParagraph{咱们假设你复习准备了一年，成功考上了上海一所普通高校。研究生学成毕业后，工资是八千一个月。}
\BookParagraph{而你如果继续在咱省委党校学习，因为不需要准备考研，可以提早一年毕业，工资低一点儿，是七千元一个月。}
\BookParagraph{你复习准备一年的因为工资损失而带来的沉没成本，是八万四千元。你需要在上海工作七年，才能追上西安工作的收入，刚刚持平。}
\BookParagraph{但如果西安房价每平米一万、上海每平米五万呢？那按照咱们假想的收入水平。在西安每年的收入可以买八平米住房，而在上海只能买两平米。也就是说，要是想和西安的住房水平持平，上海学校毕业需要找一个一年税后四十万，税前六十万的工作。}
\BookParagraph{当然了，大城市生活质量不能简单地用买房面积来衡量。}
\BookParagraph{那考研还指不定能不能考上呢。}
\BookParagraph{阿斯顿马丁确实好，我也和大家一样喜欢啊。可惜喜欢归喜欢，上班儿还得继续开我那辆开了十七年的本田。}
\BookSubtitle{四}
\BookParagraph{据说，野外碰到灰熊千万别跑。}
\BookParagraph{跑了它一定会像追猎物一样追你，而灰熊的短距离冲刺速度远远超过奥运会百米冠军。}
\BookParagraph{而以静制动，以不变应万变，倒也不一定就会束手就擒，必死无疑。}
\BookParagraph{四面楚歌之时，背水一战，也许还有一线「置之死地而后生」的希望。}
\BookDate{2018 年 9 月，多伦多}
\BookArticle{二十而折腾}{二十而折腾}
\BookParagraph{你好，李老师，我是一名大三的医学生，最近一段时间很迷茫，不想学习，但又十分愧疚，心里的小杂事太多。比如感情，以及减肥什么的，还有说不清的，但要是让我坐在那儿学习我真的学不下去，很焦虑。对我来说，考研是一个非常重要的事情，我爸妈说我要是能考上一个比现在好的医学院就可以了，可是我还是想拼拼去考我自己理想的医学院，但又十分担心自己的实力不够，在科研这方面我更加不如别人（心里自己想的能考上理想医学院的人），不知道跟谁求助，觉得自己心里有个潜在的心结不能打开，迷茫的很，以前自己喜欢的东西都不能坚持做下去，感觉我真的是很糟糕的一个人。希望李老师能点拨一二，谢谢李老师。\BookLineBreak{}（问题来自微博私信）}
\BookSubtitle{一}
\BookParagraph{你看过几年前一部叫《重返二十岁》的电影吗？}
\BookParagraph{哎呀呀，这里请允许我跑一下题。一想起这部电影，我就想到了杨子姗。从《致我们终将逝去的青春》开始，我就是杨子姗的铁杆儿粉丝——她实在是太漂亮了。尤其是她在《重返二十岁》里的一颦一笑，让我深刻地明白了为什么古人会写出「北方有佳人，绝世而独立，一顾倾人城，再顾倾人国」的词句来。}
\BookParagraph{但是我更有共鸣的是，电影里的老太太忽然回到二十岁的那一瞬间，杨子姗表现出来的那种中了彩票般狂喜的心情。}
\BookParagraph{哈哈，因为我也特想重返二十岁啊。}
\BookParagraph{而且如果不是知道这不可能实现，也许很多人都特想重返二十岁。}
\BookParagraph{如果我没有猜错的话，你今年二十一吧？}
\BookSubtitle{二}
\BookParagraph{我们常听到的那句名言「三十而立，四十而不惑，五十而知天命，六十而耳顺」， 可是你发现没有，没人记得哪位圣贤说过二十「而」干什么。}
\BookParagraph{那怎么没人说二十而谈恋爱，二十而运动减肥，或者二十而考研啊？}
\BookParagraph{原因很简单：二十多岁你有的是时间折腾，爱干什么干什么。}
\BookParagraph{哲学家常问，生命的意义在于什么？}
\BookParagraph{生命的意义就在于「折腾」啊。}
\BookParagraph{「早晨八九点钟的太阳」——当然要可着劲儿玩儿可着劲儿折腾了。}
\BookParagraph{想吃就吃，想玩儿就玩，想恋爱就多找几个试试，想学个新乐器，就立马学起来。}
\BookParagraph{那父母老是说要学习怎么办？}
\BookParagraph{没事儿，一天的时间有的是，学习和折腾不矛盾。就好像你准备考研期间，不是照样每天都在刷淘宝玩微信？}
\BookParagraph{不但不矛盾，学习之余你不玩儿不折腾，多半儿学习的效率也高不到哪儿去。所谓「事半功倍」，说的就是多玩儿多折腾，学习的时候反而效果更好的道理。}
\BookParagraph{或者这么说吧，玩儿得开心也是一种能力，越能玩儿能折腾的人，学习就越轻松，将来多半儿也「混」得越好。}
\BookSubtitle{三}
\BookParagraph{我二十多岁的时候有一次坐长途飞机，邻座正好坐了一个和我年龄相仿的女生。}
\BookParagraph{她长得什么样儿、我们在飞机上都聊过什么，我都完全没有印象了。}
\BookParagraph{只记得十几个小时的航班。我们一刻不停地整整聊了一路。}
\BookParagraph{就像你的问题一样，每一句话永远是逗号，话永远说不完，筵席永远不会散场。}
\BookParagraph{我不知道她叫什么名字，下飞机以后，我也再也没有见过她。}
\BookParagraph{也许那十几个小时，就是我全部的青春。}
\BookParagraph{就像我喜欢看《致我们终将逝去的青春》那部小说一样，我也很羡慕你。}
\BookParagraph{因为无论你考研是否顺利，减肥能否成功，自己喜欢的东西能不能坚持做下去，你此时此刻都真的很幸福。}
\BookDate{2018 年 10 月，多伦多}
\BookArticle{信　教}{信教}
\BookParagraph{李老师，您好！最近有一个问题纠结了很久，想听听您的看法。}
\BookParagraph{最近通过朋友接触到了基督教，去过几次教会的活动。首先，教会里面人之间的那种关系是我很少见过的。他们防备心很少，很容易就分享自己的私事。不过另外一位朋友说本质原因是因为他们没有利益之争，而不是因为他们是基督徒。第二，牧师是我见过的唯一一个没有私欲的人。他是一个美国传教士，我很敬佩他，也很欣赏他做事的方法。第三，我无法认同圣经里面的很多东西。比如上帝拿男人的一根肋骨造了女人我就很不服气，又比如爱神应该超过爱自己的家人。牧师认为上帝和圣经的一切都是对的，人应该服从，我就觉得这样的话那神和一个古代的君王有什么区别呢？我为什么不能质疑他？第四，我不相信神的存在。教会里面有的人提到基督教和圣经的历史之长，并以此来向我证明神的存在，这完全不能说服我。牧师说好好研究圣经和在教会服侍，就会受感动，会信。但是我觉得我根本不信，对圣经也不认同。在他们那里我像一个怪人。}
\BookParagraph{有朋友说，没有思辨能力的人才容易被洗脑，去信教，但是我热爱的褚明宇老师就是个基督徒。我真的不知道自己应该怎么做——是不再去接触这件事，还是捏着鼻子读圣经期待会有不一样的想法产生？非常想知道您对基督教的看法。期待您的回复，谢谢！（问题来自微博私信）}
\BookSubtitle{一}
\BookParagraph{最近流行讨论《英雄联盟》和iG战队，你平时一定认识几个那种特别喜欢打游戏的男生吧？}
\BookParagraph{小说《Ready Player One》里说过，一个极其简单的游戏——比如Atari 2600里的经典——之所以引人入胜，是因为它简单到了极其粗暴的地步（brutally simple）。你的手指在机械地对机器开战，你的目标就是活着，活得越长越好。}
\BookParagraph{当你在全神贯注地打游戏的时候，你可以彻底忘了时间，忘了失恋，忘了孤单，忘了这个阴雨绵绵、令人沮丧的真实世界。可以沉迷于游戏当中，在游戏里欢呼雀跃，流连忘返。}
\BookParagraph{你绝对不会去质疑：游戏里的人和故事是真实的吗？}
\BookParagraph{因为游戏最吸引人的地方，正是可以投入到它的怀抱里，可以去享受那种只有全身心投入才有的温情。就好像一个无所不知、无所不能的男朋友，阳光、睿智、温暖、勇敢，让一个女生彻底放下所有的猜疑和不安，私定终身，一生一世地依赖于他。}
\BookParagraph{直到所有的装备弹药消耗一空，所有的「命」战死疆场——直到屏幕上出现「GAME OVER」。}
\BookParagraph{信一个宗教，就像打一个游戏一样。}
\BookParagraph{只有一个显著的区别：信仰宗教永远没有「GAME OVER」。}
\BookSubtitle{二}
\BookParagraph{你去质疑一种信仰，就好像去质疑那个无所不知、无所不能、阳光睿智的男朋友一样。可以当然是可以，只是大煞风景，完全没有了「沉迷于其中」的快乐，没有了「托付终身」的温暖——多半很快你就不爱他了。}
\BookParagraph{不过这不赖你，你把科学和宗教混为一谈了。}
\BookParagraph{科学的魅力，在于它有一套完美无缺的逻辑体系，可以在不断地被人质疑中发展。而科学研究的快乐所在，正是在不断质疑前人的结果的过程中，得到突破性的进展。}
\BookParagraph{十九世纪初研究欧氏几何的故事，正好可以完美的诠释这种质疑。数学家对欧几里得的几何原本中的第五条公设（postulate）早就有疑问，觉得它不该是公设，而应该是可以证明的定理。}
\BookParagraph{著名数学家高斯就曾希望，能不依靠第五公设就推出三角形内角和一百八十度。他发现，如果假设三角形内角和大于一百八十度，逻辑上可以推出矛盾。但是如果假设小于一百八十度，却无论如何也推不出矛盾来。}
\BookParagraph{无独有偶，匈牙利数学家 Bolyai 推了一辈子也没推出矛盾，建议他儿子：「Leave the science of parallels alone.」结果儿子不听，继续推导，终于推出了一整套自成体系而又和欧氏几何完全不同的理论，和俄罗斯数学家 Lobachevski 几乎同时发现非欧几何。}
\BookParagraph{更有甚者，另一位数学家黎曼（Riemann）发现，如果不假设直线的长度是无限的，连三角形内角和大于一百八十度都无法推出矛盾，终于创立了黎曼几何。}
\BookParagraph{你瞧，研究数学和信仰宗教或者打游戏竟然有异曲同工之妙：都有一种能让人沉迷于其中、竭尽一生之心力的魅力。唯一的区别是，研究数学的快乐来自质疑，而信仰宗教的快乐则来自笃信。}
\BookParagraph{这么说吧，你就把基督教当成欧氏几何里的一条公设好了。}
\BookParagraph{那到底应不应该信仰基督教呢？}
\BookParagraph{其实答案跟应不应该打游戏一样，非常简单：跟你的男朋友\BookLineBreak{}（或女朋友）保持一致就行了。}
\BookParagraph{当然如果你正好单身，跟你热爱的褚老师保持一致，也不是一件坏事儿。}
\BookSubtitle{三}
\BookParagraph{我有一回去一个叫「Canada’s Wonderland」的游乐场坐了个叫「Behemoth」的过山车。上了车系上安全带，车子缓缓地升得越来越高，夕阳越来越刺眼。}
\BookParagraph{科学地说，当过山车升到70米时，将会以75度的角度下坠，3.9 秒内达到每小时 125 公里。}
\BookParagraph{我当时说的是：}
\BookParagraph{Oh my god.}
\BookDate{2018 年 11 月，多伦多}
\BookArticle{优　秀}{优秀}
\BookParagraph{李老师，人应该如何正确看待比自己优秀的人呢？（问题来自微博私信）}
\BookSubtitle{一}
\BookParagraph{咱们来假设一个场景：公司年会上有一个抽奖活动，几百个人可以有一个中奖的机会，如果中了奖，可以得到一部你一直梦寐以求的苹果（或者华为或者锤子或者小米）手机。}
\BookParagraph{你的一位同事抽中了，非常高兴。你会是什么心情？}
\BookParagraph{我想你大概跟我差不多。如果不认识这位同事，大概没觉得什么，甚至会为他高兴。如果恰好和这位同事比较熟，大概会有一点点羡慕，但是一小会儿以后也就忘记了。}
\BookParagraph{一个人比自己「优秀」，就好像她抽奖时运气比自己好一点儿。也许会有点儿羡慕，但一小会儿之后最好就忘了。}
\BookSubtitle{二}
\BookParagraph{其实我觉得，中奖的同事我会心生羡慕，但是「优秀」则连那一点儿羡慕都没太大必要。}
\BookParagraph{主要是因为「优秀」这个词儿很难精确定义。}
\BookParagraph{离我住的地方不远，有一家有五十年历史的自行车店，叫「Mariposa Bicycles」，专门手工制造铝合金自行车。店主七八十岁了，造了一辈子自行车，还把自己的手艺传给了儿子。虽然自行车早就可以大规模生产，材料也早就演变成碳纤维或钛合金，但是只要还有顾客愿意买这个老店的自行车，他们父子俩就可以凭手艺继续把店开下去。}
\BookParagraph{如果跟一个天天「空中飞人」全世界出差开会的企业高管比，我甚至觉得这父子俩更加「优秀」。几十年后，高管的「PPT」早就烟消云散了的时候，父子俩手工制造的自行车还可以骑着上路。}
\BookParagraph{我一直有点儿羡慕像这父子俩这样的人。羡慕他们可以一辈子就专心做造自行车这么一件事，把这件事做得尽善尽美，而且还有买主认可他们的手艺。}
\BookParagraph{我家附近还有一家汽车修理店也是如此。日出而作，日落而息，在同一地点开了三十多年了。我开的一辆十七年前买的车，又锈又旧，但每次坏了就放心地送到这家店去修理，「新三年，旧三年，缝缝补补又三年」，竟然还能天天开着上班。做为一个连换个灯泡都要上网研究半天型号的我，对这家汽修店的老板也一样羡慕不已。}
\BookParagraph{所以说啊，「优秀」这个词儿真的很难定义。有人大概觉得\BookLineBreak{}「富二代」啊明星啊网红啊大公司老板啊天天吃香的喝辣的更加\BookLineBreak{}「优秀」；而我呢，相比「治大国」，我更加偏爱「烹小鲜」；相比「待到理想化宏图，咱重摆美酒再相会」，我则更向往「小小竹排江中游，巍巍青山两岸走」。}
\BookParagraph{所以啊，在你因为别人更「优秀」而「羡慕嫉妒恨」的时候，他说不定也在暗地里羡慕你呢。}
\BookSubtitle{三}
\BookParagraph{我之所以不怎么羡慕富二代、网红和老板们的「优秀」，是因为这些跟在抽奖活动里中奖差不多，需要有很好的运气。}
\BookParagraph{我更偏爱那种不怎么靠运气，而只需要有一点儿像样的手艺和特长，每天认真而专业地做好自己的事情的「优秀」。}
\BookParagraph{外科医生每天做一台手术，钢琴老师每天教几个孩子，专车司机每天跑几趟机场，饭店厨师每天做几道拿手菜——我觉得都特别优秀。}
\BookParagraph{想起周围有这么多比自己优秀的人，我感到日子过得很踏实，人气十足，心情愉快。}
\BookDate{2019 年 1 月，多伦多}
\BookArticle{Institutionalized}{Institutionalized}
\BookParagraph{亲爱的李教授，最近刚和男友分手。他遇见我时，似乎以为我在思想、为人等方方面面都和他很匹配，给我套上了一个理想化的框。后来渐渐发现我没他想的那么理想，甚至远远低于他要求自己的标准，没有他定义的那么「优秀」。两人在一起太多不开心，最终他选择了和我分开。我们俩似乎有些价值观上的差异，他野心更足，而我只想在小范围内对于我爱的人和爱我的人造成一定影响。他说我还没搞清楚自己要什么，而他很清楚。我觉得遇见他其实也是一种幸运，虽然相处真的很累，但我确实看到了自己的矛盾和不足。}
\BookParagraph{这位前男友是一位高学历、非常focused、有崇高追求的人。但我总觉得他在象牙塔里太久，似乎对社会上其他的和他不太相似的人没有那么多包容和耐心。可能因为我们真的不是一路人吧。}
\BookParagraph{李老师，您是在人生中的什么时刻知道自己的「天命」，清晰地意识到自己擅长并热爱做的事情呢？您在选择伴侣时又是怎么知道对方是正确的人呢？（问题来自微博私信）}
\BookParagraph{你这一句「亲爱的李教授」，我觉得挺温暖的，感觉自己颇像一位慈祥的老爷爷（快了快了）。好吧，就冲这句话，即使是绞尽脑汁，我也要把你这篇命题作文写出来。}
\BookSubtitle{一}
\BookParagraph{你看过电影《肖申克的救赎》（The Shawshank Redemption）吧？}
\BookParagraph{如果没看过，我建议你立刻找来看一看；如果看过了，再看一遍也是个不错的主意。}
\BookParagraph{IMDb上，依靠广大人民群众投票，这部电影的排名二十年以来一直稳居世界第一。看看排名世界第二的是著名电影《教父》，跟它同年上映并获得六项奥斯卡奖的《阿甘正传》只排到第十二，你就可以想象这部电影有多牛。}
\BookParagraph{但是这部电影当年在美国上映以后票房惨淡，差点儿连成本都没回来。而《阿甘正传》则卖了三亿美元票房，大红大紫。}
\BookParagraph{也难怪，这部电影除了最开始几分钟，没有一位女性角色。一个讲监狱的普通故事，片长甚至还跟《阿甘正传》完全一样——两部电影都是两小时二十二分钟。}
\BookParagraph{主角 Andy 本来当年导演想让 Tom Hanks 演的，结果连Tom Hanks都不太看好这部电影的剧本，决定去演《阿甘正传》了。}
\BookParagraph{就是这么一个普通得不能再普通的「鸡汤励志」故事，没有女主角，没有枪战追车的大场面，竟然凌驾于《教父》之上，排行世界第一。}
\BookSubtitle{二}
\BookParagraph{而这部电影里最感人的一段情节，我一直认为不是什么关于\BookLineBreak{}「希望」的鸡汤，而是 Brooks 假释出狱的那一段儿。}
\BookParagraph{Brooks 在监狱里关了五十年，出狱后没多久就自杀了。}
\BookParagraph{Red 听到这个消息，说：「The man’s been in here fifty years, Heywood. Fifty years! This is all he knows. In here, he’s an important man, an educated man. Outside he’s nothing!」}
\BookParagraph{「I am telling you these walls are funny. First you hate them, then you get used to them. Enough time passes, you get so you depend on them.」}
\BookParagraph{「That’s institutionalized.」}
\BookParagraph{我就不替你翻译成中文了，主要是因为我不会翻——尤其是最后那个画龙点睛的词儿：institutionalized。}
\BookParagraph{你的那位前男友，我一点儿也不喜欢。不要说他因为你不够\BookLineBreak{}「优秀」提出了分手，就算是他不提，你差不离儿了也可以主动提出来。}
\BookParagraph{因为他已经被他周围的环境「institutionalized」。}
\BookSubtitle{三}
\BookParagraph{一个人其实挺容易就会被「institutionalized」，不需要五十年的时间。}
\BookParagraph{这是因为人不但适应环境的能力极强，而且一旦适应了以后，突然让他跳出这个「圈子」的思路，非常困难。}
\BookParagraph{比如一位同事本来发言时不会打官腔，当了领导，讲话时的语气就和别的领导一模一样，反而不会正常说话了；一位教授本来还能写写论文，拿到了大项目，就只会到处做做报告，反而不会做实验写论文了；或者一位风险投资圈子里的哥们儿，你跟他聊天千万别聊什么赵丽颖啊张雨绮啊的八卦，人家压根儿没听说过。}
\BookParagraph{没错儿，我们的领导、教授，还有著名风险投资人都是高学历出身——不是清华北大至少也是个什么九八五——而且都有正确的人生观世界观，崇高的追求，远大的理想。他们逢人便问『你最近在「做」什么（研究或项目）？』}
\BookParagraph{但是我就不能什么都不「做」？什么都不想「要」？}
\BookParagraph{我就不能随波逐流，做一天和尚撞一天钟？就不能时不时关注一下赵丽颖的电视剧，或者发几条微博朋友圈？}
\BookParagraph{他们认为这难以理解，认为这是人生观世界观价值观的大问题。}
\BookParagraph{这就是「institutionalized」。}
\BookParagraph{回到你的问题上来，我从来都没搞清楚自己擅长或热爱什么事儿。不过这也没什么关系，即使是牛顿，不是也觉得自己只是个海滩上捡贝壳儿的孩子；晚年把科学「做」得差不多了，不是也全心全意地去研究神学了嘛。}
\BookSubtitle{四}
\BookParagraph{那如何去选择「正确」的人生伴侣呢？}
\BookParagraph{我记得二十多年前在北京办理结婚登记手续时，婚前体检的一部分，是看一部二十分钟的教育片。}
\BookParagraph{教育片的内容我基本全都忘了，只有一小段记得异常清晰。这段教育广大未婚男生，结婚以后，在女生生病不舒服时，要悉心照料，给她递一杯热水。}
\BookParagraph{乔布斯当年试图说服百事可乐的 John Sculley 出任苹果 CEO 时说过一句名言：「你是想卖一辈子糖水，还是想跟我改变世界？」}
\BookParagraph{你要找的男友，首先应该是一个善良而重感情的好人，别动不动就要去改变世界。}
\BookParagraph{当你感冒发烧时，他无论在世界上任何地方，都会奇迹般地飞回来出现在你的眼前，给你倒好一杯热水，放在你的床前。}
\BookParagraph{人和人之间的感情，永远比改变世界重要得多。}
\BookDate{2019 年 1 月，多伦多}
\BookArticle{优　点}{优点}
\BookParagraph{李老师，您好。我是从褚老师那认识的您。看了您不少微博，听了您不少分答，特别喜欢、敬佩您。您的善良、真诚以及没有架子，让我特感动。想听您夸夸自己。（问题来自分答「在行一\BookLineBreak{}点」）}
\BookParagraph{我其实还真没啥优点，不过倒是有一特点。}
\BookParagraph{我的特点就是运气一直很好。}
\BookParagraph{就比如说高考吧——我的物理和生物一直是弱项，而语文是高考强项。结果我那年的全国统考试题，正好语文题极难，而物理非常简单，结果轻松过关，在人生第一大关口上展示了极其强悍的好运气。}
\BookParagraph{上了清华一年过去了，忽然注意到一个同年级的女生绝对是小美女啊，怎么以前没发现？后来才知道她刚刚换了一副新眼镜——以前那副太难看了。我学习不如小美女，体育也不太灵，除了学校作文比赛得过名次，还真没什么突出的特长。小美女不但有人追，还有男朋友，结果后来竟然被弱不经风其貌不扬的我追到了，在人生终身大事上一锤定音，我再次展示了运气好的重要性。}
\BookParagraph{后来出国了，运气更是越来越好。比如说褚老师吧，竟然在校园旁边的一个小馆子撞运气偶然碰上，后来一起吃香喝辣，无酒不成席，成了很好的朋友。我一直觉得，一个人的性格特色是很容易受他周围的朋友影响的。而我正好在上学的时候遇到了褚老师，吃喝玩乐之间，少花了多少2198啊——IPO时买了后来的蓝筹股，这不是运气好是什么？}
\BookParagraph{虽然说运气不可能一直好下去，人生如赌局，有时也会有运气很糟糕的时候。可是金榜题名、洞房花烛、他乡故知几场大局都赢了，绝对算是难能可贵了吧？中学班主任大陈儿教导我们，「世上没有单纯的运气，机遇只偏爱有准备的头脑」，也许是我从小都在「时刻准备着」呢？}
\BookParagraph{虽然没「创业」发过财，没走红当过领导，没戴过劳力士的手表，没玩儿过徕卡的相机，可从小第一批就戴上了红领巾，上中学第一批就玩儿上了苹果原装电脑，眼前有可爱的家人孩子，远方有难得一见的朋友，此般运数，夫复何求？}
\BookParagraph{这不，今天一位从未见过的陌生人，到好久都没玩了的「在行一点」上——这名字也太拗口了，还是叫分答吧——把我小小地表扬了一下。这样的「机遇」，像是走在胡同里忽然飘来悠扬的琴声，令人心情大好，怎么「有准备」的头脑都准备不来——我们还是把它算在「运气好」的头上吧。}
\BookParagraph{谢谢你。}
\BookDate{2019 年 2 月，多伦多}
\BookArticle{完　美}{完美}
\BookParagraph{李老师，很抱歉打扰您，想到您可能会偶尔看到这封私信，觉得无比荣幸。跟很多网友一样，我是通过褚老师认识您的，就这么关注了这好些年，包括加入「向上移动」社区。虽然我们并不认识，但对我来说，能看到您和褚老师的文字和分享，是相当有启迪意义的。通过您俩，我更深入和全面地去了解自己的生活和内心。}
\BookParagraph{我今年26岁，去年看到了个别亲人的骤然去世。想到再也看不到的人，无法为对方做的事，想到将来会让深爱的家人和朋友伤心和牵挂，心理和身体受到了相当大的打击，我意识到健康是如此的可贵。我生性内向，但我一直努力着去更主动地用含蓄的方式向身边的人表达爱。花了一年多的时间我才缓和了一点儿，从那以后我常常思考生死的问题。有时候能释然，相信只要自己过好每一天，不管遇到任何事都能坦然和乐观面对；但更多的时候，感受到的是恐慌、巨大的失望、伤心和无助。生而为人，到底怎么样才能坚强和快乐呢？}
\BookParagraph{如果您不介意的话，我想听听您关于生老病死的看法，虽然这个话题很沉重、很空泛。谢谢老师。（问题来自微博私信）}
\BookParagraph{哈哈，我这个「情感博主」看来要拓展业务范围了——绝对是业务能力一直很强的体现啊。}
\BookSubtitle{一}
\BookParagraph{要聊「生死离别」的话题，得从有关生死的「鸡汤」开始聊起。因为只有想明白了自己的生死，才能继续去琢磨有关亲朋好友生死的话题。}
\BookParagraph{这样的鸡汤简直俯拾皆是啊——比如最近我刚看到一篇，中心思想是「别把好东西留给特别的日子，你活着的每一天都是特别的日子。」}
\BookParagraph{这样浅显的道理，理应每个人都明白，不是说了跟没说一样嘛。}
\BookParagraph{你还别说，大多数人还真的没怎么想过这些问题。前一阵子和一个年轻朋友聊聊起加拿大的退休金减税制度，他说：「退休不是还早着呢？」}
\BookParagraph{是啊，很多事儿看来遥遥无期，什么时候发现竟然近在眼前了，自然有点儿让人措不及防。}
\BookParagraph{还是让我们来看看鸡汤界的一哥乔布斯怎么说的吧：「Remembering that I’ll be dead soon is the most important tool I’ve ever encountered to help me make the big choices in life.」}
\BookParagraph{是啊，其实平时想想如果今天是生命最后一天该干什么这个问题，不但不沉重，有时候反而还挺好玩儿的。}
\BookParagraph{比如如果你的答案是立马到车行提一辆特斯拉或者保时捷爽一把，哈哈那是因为你钱还不够多。有统计数据表明，真正的富人最常开的车是丰田。什么时候你只开丰田了，你就基本不在乎车这个玩意儿了。别玩了，抓紧时间努力挣钱吧。}
\BookParagraph{又比如如果你的答案是要趁最后一天还能说话向心仪已久的暗恋男生或者女生表白，那你可千万别当真啊，假设仅仅是假设，不是真的。你要是一冲动不顾我「千万不要表白」的一贯立场真的表白了，到时候惨遭人家断然拒绝可别怪乔帮主的鸡汤。}
\BookParagraph{著名二十世纪数学家 Paul Erdos 也喜欢玩儿这个，曾经畅想什么才是「最完美的死法儿」。他最为钟爱的死法儿是在一堂数学课上正在讲一个证明，学生举手提问：「What about the general case?」他回答：「I will leave it to the next generation.」然后去世。}
\BookParagraph{我怎么觉得他这死法也太高大上了吧——要是我的话，估计是这么一个场景：跟往常一样，字斟句酌写好一条微博，里面并没有提到那些「啊朋友再见」式的煽情文字，发出去一会儿看到还颇有十几个赞。这时正好Instagram上收到了一条负面评论，跟小女儿学习，登录小号和他针锋相对展开论战：「Yo! Yo! Chill!」最后再加上一句「You’ll have to pry it from my cold, dead fingers!」然后欣慰地微笑去世。}
\BookParagraph{哈哈，作为爱玩微博的小市民，我也就这么点儿追求啦。}
\BookSubtitle{二}
\BookParagraph{完美的死法虽然并没有那么恐怖，但也不能天天想啊。如何\BookLineBreak{}「完美」地活着，倒是值得多琢磨琢磨。}
\BookParagraph{我认为其实目标可以很简单明确：活得越长越好。}
\BookParagraph{我越琢磨，越觉得「活得长」是个很「给力」的概念。你想啊，要是有个你恨之入骨的老板或者情敌，又拿他没有办法，你只需要「君子报仇，十年不晚」，比他活得长不就得了。英文里还专门为「比他活得长」创了一词儿，叫「outlive him」——「If there’s nothing you can do to him, just outlive him.」哈哈，解气！}
\BookParagraph{不过要真想活得长，还得是一个小抠门儿才行。在加拿大，政府或者单位会发一种叫「defined benefit」的退休金。什么叫「defined benefit」呢？就是说你每月领取的退休工资是固定的，一直发到去世为止。这种退休金的大意是，把活得短的人那里结余下来的钱，发给活得长的。这么一来，如果你是个小抠门儿——英文叫「cheapskate」—— 那你一定会为了多领退休金也要活得长点儿。你越抠门儿，活得长就越有动力。}
\BookParagraph{哈哈，别懒着了，赶紧每天早上出门跑步做操，「为祖国健康工作五十年」吧。}
\BookSubtitle{三}
\BookParagraph{系主任最近发了一封邮件，一位早已从系里退休了的教授最近去世了，终年八十多岁。他的亲人为他在加拿大最大的报纸「The Globe and Mail」里发了讣告。}
\BookParagraph{我有功夫的时候很喜欢看报纸上的这些讣告，因为它们大多行文优美，发人深省。我跟大女儿说，等将来我去世了，一定要把一篇行文优美的讣告发表在「The Globe and Mail」上，有点儿坐飞机坐了一等舱的感觉。}
\BookParagraph{那位教授的讣告里说，他一生和他妻子共同度过，没有子女。他妻子和他都酷爱天鹅，经常长途跋涉去天鹅出没的野外去看天鹅。}
\BookParagraph{八年前，他的妻子去世了。他一个人在养老中心住了八年，中心的服务员跟他混熟了，被他风趣幽默的语言所折服，平时就亲切地称呼他为「The Professor」。}
\BookParagraph{「The Professor」—— 完美。}
\BookDate{2019 年 2 月，多伦多}
\BookArticle{Get Busy Living}{Get Busy Living}
\BookParagraph{李老师你好。我们如何能从平淡的生活中获得快乐？虽然觉得自己还有很多风景没看很多东西没有接触，但有的时候会觉得人生不过如此，生活每天都是平凡乏味，就是感受不到快乐！抱歉都是负能量，希望得到您的建议，谢谢！（问题来自微博私信）}
\BookSubtitle{一}
\BookParagraph{不知道你发现没有，做一件自己以前从来没有做过、或者做过但是从来没有成功过的事儿，可以给人一种非常快乐、非常开心、现在讲话叫「幸福感爆棚」的感觉。}
\BookParagraph{比如说，我上小学时上过一堂现在还记忆犹新的「自然常识」课。之所以一直记得，是因为老师要求我们把圆周率背到小数点后五十位，作为当晚的作业。起因是当时正好有一则新闻，说科学家茅以升偶遇一个孩子，一老一小聊了一会儿竟成了忘年之交，原因是两人都发现对方可以把圆周率背到小数点后一百位。}
\BookParagraph{结果第二天上课时老师检查作业，班里除了我以外，很少有人能把五十位完整背下来。}
\BookParagraph{现在想想，这个老师也是够难为人的 —— 这五十位的圆周率，我背了一个晚上。可是当我发现竟然真的可以一字不错完整背熟时，成就感油然而生，觉得所有的痛苦都是值得的。}
\BookParagraph{这五十位圆周率，我现在还会背。甚至女儿五六岁的时候，还在我耳濡目染之下，也把五十位圆周率背得滚瓜烂熟。而且后来我和女儿都发现，其实并没有白费功夫：五十位里的任意一小段儿，都可以作为一个完美的账号密码。}
\BookParagraph{这样的例子还有很多。比如很多人热爱的极限运动马拉松，其实跑起来是一件非常痛苦的事儿。那为什么还要动不动就去跑呢？我想，和喜爱蹦极、喜爱跳伞、喜爱登山等玩法差不多，很重要的一部分原因是想体验完成了一件自己从来没做过的事儿时的那种欣喜若狂的感觉。}
\BookParagraph{所以啊，想从平淡的日子里获得快乐太容易了。只要做成一件你一直没有机会做的事儿 —— 比如跑一个五千米、把《红楼梦》看完、完成一个在线课程、临摹一份字帖、或者写一本书出版，你会发现，一旦完成，快乐就会像老虎机中大奖时的钢蹦儿，源源不断地哗哗落下，数都数不清。}
\BookParagraph{我一直不会自己用针线钉扣子。如果将来有朝一日我真的学会了，说不定一高兴网购一台Singer的缝纫机开始裁剪布料也未可知呢。}
\BookSubtitle{二}
\BookParagraph{我还有一个建议：一定要「忙」起来，最好忙得不可开交，别给自己太多闲功夫，去想有关「人生」的大问题。}
\BookParagraph{我喜爱的电影《肖申克的救赎》里有一句名言：「 Get busy living, or get busy dying.」}
\BookParagraph{你还别说，这句话我想了一会儿，还真的很难「信达雅」地翻译成中文。}
\BookParagraph{所以我特别羡慕翻译家的水平 —— 上中学时我看过一本小说叫《春月》，原著是英文写的。但是在吴世良这样的著名翻译家笔下，文字于柔和中流露出一抹凄美，竟然比英文原著更好看。}
\BookParagraph{你要是难以体会这句话的含义，建议你继续努力提高英文水平，别跟有的不懂英文的「作家」一般见识。}
\BookParagraph{「忙」很重要，但是忙着「living」还是忙着「dying」，比仅仅是「忙」还重要得多。「Living」不仅仅是柴米油盐，还有对不可预知未来的憧憬，还有平常日子里满心的希望。}
\BookParagraph{比如你如果恰好在酒吧偶遇一位高挑漂亮的姑娘（或者彬彬有礼的帅哥），厚着脸皮上去跟人家攀谈几句，即使是铩羽而归，也总算是「get busy living」，比偷拍几张照片再发到朋友圈儿——或者压根没去过酒吧——强了不少。}
\BookParagraph{鸡汤里人家是怎么说的来着——「You only live once.」}
\BookSubtitle{三}
\BookParagraph{话又说回来，在平常的日子里体会到快乐也是一种能力。不同的人对快乐的「敏感度」不一样，就好像同一幅名画，有的人感动得掉眼泪，有的人无动于衷一样。}
\BookParagraph{不过我们都特别幸福地正好生活在一个高科技的时代，随时随地可以感受到快乐。}
\BookParagraph{比如最新一期的《Wired》杂志里讲到一个故事。说一个计算机高手特别爱吃拉面，去日本玩儿的时候，他编了一个程序把沿途所有的拉面馆按照好评度计划了一条最优路线，然后和女朋友一起，按照程序输出的最优结果，一家一家吃了一圈儿。}
\BookParagraph{我坚信，他编完这个程序时，一定有一种快乐得要大喊「I’m the king of the world」的冲动。}
\BookParagraph{不过后来他女朋友实在受不了他这个「优化」狂人，还是跟他分手了，哈哈。}
\BookParagraph{我自己呢，远远没有他这么热爱计算机——高科技时代能让我写的文字多几个人看，多得几个「赞」，就满心欢喜、心满意足了。}
\BookDate{2019 年 3 月，多伦多}
\BookArticle{Chill}{Chill}
\BookParagraph{李老师，祝您青年节愉快。我也来私信求解惑了。是一个关于「原则」的问题。比如，最近和朋友玩游戏喝酒，喝一杯我就是喝一杯，倒满。朋友们喝着喝着就会倒少点儿，最后几乎成半杯了。褚老师引用过一句 GQ 的话，曾对我有所帮助，大意是「在一个俱乐部，两类人比较讨厌，一是没规矩的，二是太坚持规矩的」。所以想问下李老师怎么看。一，我是否应该坚持「满杯」。二，我是否应该坚持自己且不用在意别人的「鸡贼」做法。求指教。（问题来自微博私信）}
\BookSubtitle{一}
\BookParagraph{你知道褚老师说的这个「GQ」杂志的名字是哪两个英文单词的缩写吗？}
\BookParagraph{没错儿，是「Gentlemen’s Quarterly」——「男士季刊」。}
\BookParagraph{可是「GQ」杂志其实不是季刊，也不是双月刊，而是一份月刊杂志。}
\BookParagraph{这不是丢人吗？作为一份美国最有名的男性时尚杂志，连名字都在误导读者，还怎么取信于人？}
\BookParagraph{你可能会说，这不是钻牛角尖儿吗？也许这个杂志创刊时是季刊，后来变成了月刊却没有改名，这似乎也没什么关系吧？}
\BookParagraph{去研究 GQ 杂志的名字是否误导确实有点儿钻牛角尖儿。这就好像你第一次去和一个不认识的女生到一家正式的西餐厅约会，发现她竟然先拿那一把靠近盘子的叉子吃开胃沙拉，于是认真地给她介绍了一下西餐刀叉的摆放顺序一样。}
\BookParagraph{大煞风景之余，还完全没有必要 —— 用哪一把叉子吃沙拉，和 GQ 到底是不是季刊一样，是无足轻重的一个细节。}
\BookParagraph{而你所谓的「原则」性问题，竟然是喝酒到底是半杯半杯还是满杯满杯喝的问题。}
\BookParagraph{你所信任的 GQ 杂志还说过，力量训练之后不要喝酒，有氧训练之前也不要喝酒，前一天晚上喝酒还会影响后一天早上运动的成效。总的来说，喝酒无益于健康。}
\BookParagraph{那你是否应该和朋友建议，咱俩要不然别喝酒了，改喝可乐？}
\BookParagraph{等等，可乐也同样含有大量糖分和咖啡因，应该改成喝冰水。如果你朋友的健康不算原则性的问题，那什么是原则性的问题？}
\BookParagraph{你的朋友这时一定会想：这人整个儿傻到家了啊，下次不能再和他喝酒了。}
\BookParagraph{饭局不欢而散，这就是你在「原则」问题上钻牛角尖儿所希望得到的结果吗？}
\BookSubtitle{二}
\BookParagraph{前一阵子我在法国开会，碰到了几位学术界的同行，聊到了我的新书《幸福》。我跟他们「宣传」这本书用到的一句名言是：如果搞科研是让你精神紧张的「毒药」，那我的这本书就是让你心情放松的「解药」。}
\BookParagraph{而对于你这样完美主义到了钻牛角尖儿地步的人，我也有一份「解药」—— 美国青少年爱用的一个单词，叫「chill」。}
\BookParagraph{「Chill」这个词儿还真不太容易翻译成中文。从发音来看，让我想起了小时候最爱的「清凉油」；从词义来看呢，又让我想起了一句北京话叫「哪儿凉快哪儿歇着去」。}
\BookParagraph{与其试图翻译成中文，咱要不然还是继续假想那个跟一个未曾谋面的女生第一次约会的故事吧。}
\BookParagraph{第一次约会去一家正式的西餐厅，这个选择本身就明摆着是「chill」的反义词儿 —— 超级「un-chill」。}
\BookParagraph{你想啊，吃完了这么贵的饭你如果说「我来我来」抢着买单，给她一种「欠了你」的心理暗示，她对你的印象能和「chill」沾上边儿嘛。}
\BookParagraph{正确的约会地点是一家路边儿的小饭馆儿，一边儿喝着大扎的啤酒吃着新鲜出炉的烤串一边聊天，既可以半个钟头买单走人谁也不欠谁的，又可以「他乡遇故知」，海聊到馆子打烊。聊着聊着如果正好聊到了她特想去的一家网红甜品店，那不正好是下次一起「chill」的好去处嘛。}
\BookParagraph{吃饭时既不能一句话不说各自玩儿手机，又不能吹上一晚上自己从前光荣的奋斗故事，因为这两者都一点儿也不「chill」。约会结束以后即使再喜欢她也要「chill」，别动不动就给她发微信，让人家觉得你平时无所事事，没有朋友，也无法自己独处。}
\BookParagraph{没错儿，「chill」代表的是一种人生哲学。你越计较的事，越喜欢的人，就越有可能说出钻牛角尖然而却于事无补的昏话，做出南辕北辙的傻事儿来。「船到桥头自然直」，你越不在意，大家的心情就越愉快。}
\BookSubtitle{三}
\BookParagraph{当然，如果不是第一次约会，聊天聊到了真正原则性的问题，你一定要就事论事，旗帜鲜明地表达自己的观点。}
\BookParagraph{不过，在清晰地表达你的观点的同时，你还要说明「为什么」：你得出这个结论依据的是什么可靠的论据。}
\BookParagraph{这就像中学时语文课写作文一样：议论文需要明确的论点和论据，逻辑清晰，「晓之以理」；而散文则「动之以情」，形散而神不散。}
\BookParagraph{如果这位和你约会的女生还是固执己见，无法认同你的观点，那说不定你们将来也不太可能互为对方的「解药」，一起过日子也不一定能「幸福」。}
\BookParagraph{最终分手的时候，你如果实在想要完美主义，还是可以找个机会善意地提醒她：西餐里的沙拉，应该是用左手最外面的那把叉子来吃。}
\BookDate{2019 年 5 月，多伦多}
\BookArticle{泯然众人}{泯然众人}
\BookQuote{李老师好！我本来是一个超级活泼开朗、无忧无虑的人。但是在经历了很多个突然长大的瞬间之后，我感觉现在的自己丧丧的，有种底色日渐悲凉的感觉，似乎快乐很难且很短暂。对于这种悲观，您有什么建议吗？（问题来自微博私信）}
\BookSubtitle{一}
\BookParagraph{日子过得不开心，自然有很多可能的原因。甚至有的时候，亲人或者朋友无心说出的一句话，听了都会让人觉得生活索然无味，世界黯然失色。}
\BookParagraph{但是既然你的问题这么笼统，有一种可能的原因，我打算单拎出来掰碎一点儿聊聊。}
\BookParagraph{我觉得身边的很多人 —— 也包括我自己 —— 从小到大受到的教育，都有点儿过于「正能量」。比方说，我们要有崇高的理想，远大的抱负，从师要「青出于蓝而胜于蓝」，无论将来做什么工作，都要「干一行爱一行」，「三百六十行，行行出状元」。我们生怕自己变成了「伤仲永」中的那个「泯然众人矣」的天才少年，生怕「做一天和尚撞一天钟」。}
\BookParagraph{等到我们为人父母的那一天，带孩子则更是如此：巴不得把自己未竟的理想和抱负，将来让孩子来替我们实现。我们「可怜天下父母心」，为孩子找最好的学校、上最好的特长班，同时永远在羡慕别人家的孩子：怎么教育的，这也太优秀了吧。}
\BookParagraph{但是很多时候我们发现，从上学到上班，从谈个恋爱到结个婚，真实的世界其实极其残酷：工作不顺心，领导不待见，甚至谈个恋爱都谈不成。「待到理想化宏图」似乎永远没戏，只有时间，在一年一年没有悬念地慢慢流逝。}
\BookParagraph{偏偏这时候，「头条」里充斥着各种「大牛」的成功经历，朋友圈里则永远是「哪壶不开提哪壶」—— 我这儿还没男朋友呢，闺蜜都结婚了；我这儿论文不断被拒，人家的论文永远榜上有名，搞不好还能得奖。就连机场的书店里，随处可见的都是马云、冯唐、李开复这样的「人生赢家」，令人相形见绌。}
\BookParagraph{这大概就是你所说的「突然长大的瞬间」吧。}
\BookSubtitle{二}
\BookParagraph{二十多年前我在清华上学的时候，同宿舍有一个室友，而且还恰巧睡在我的下铺，可以算是「睡在我下铺的兄弟」。}
\BookParagraph{我发现我们班里的「神人」都有一个离奇的姓名 —— 我下铺的这位兄弟姓「银」，我们都管他叫「老银」。}
\BookParagraph{老银虽然睡在我的下铺，但跟我完全不是一路人。我属于比较「追求上进」的同学 —— 每天晚上一定要去自习教室上自习，晚上十点准时下自习的那种。老银则是我的「镜面反射」：他从来不上自习，甚至不经常参加打牌或者踢球等一系列人民群众喜闻乐见的课外活动。}
\BookParagraph{他的爱好时不时会变，但却永远别具一格。比如大二有一阵子，他喜欢修电视机。一位同学捡回来了一台没画儿了的黑白电视，他靠着一把老旧的电烙铁，愣是把这台电视给修好了。害得我十几年前买过的一台背投影的彩电后来颜色坏了，我第一反应竟然是自己打开后盖去修。}
\BookParagraph{临近大学毕业时，班里的同学全体下馆子聚餐。像电影《致我们终将逝去的青春》或者《芳华》里的散伙饭场面一样，同学们频频举杯，依依惜别，男同学大多老泪纵横。据我观察，老银像没事儿人似的，冷眼旁观，不屑一顾。我的毕业纪念册上，「睡在我下铺的兄弟」没有留言。}
\BookParagraph{我和老银明显不是「一条线上的蚂蚱」，毕业以后也就自然而然地失去了联系。每次老同学聚会时打探老银的消息，消息灵通人士寥寥。二十多年来，只是听说老银没怎么工作过，曾经迷过几年炒股，没有女朋友，后来辗转移民了加拿大。在加拿大也没怎么工作，据说这几年靠救济生活，白天喜欢在咖啡店闲着看看人来人往的风景，晚上则住在政府给无家可归者提供的住所（homeless shelter）。}
\BookParagraph{据说有同学曾经问过老银，问他为什么不找个工作。他的回答是：为什么要工作？}
\BookParagraph{一个清华的毕业生，在遥远的加拿大名副其实地「流落街头」，竟然过得还不错。有的人倡导「断舍离」，倡导「极简主义」，老银则把「极简主义」活到了极致，连个自己的家都给「断舍离」了。}
\BookParagraph{心情「悲凉」的你一定会开心地发现，即使你将来「混」得再差，也比老银混得好些。}
\BookParagraph{不过我一直觉得，老银像一位少林寺里扫地的和尚，武功卓绝，日子过得很牛逼。}
\BookSubtitle{三}
\BookParagraph{前几天去香港，跟一位香港理工大学的教授朋友吃饭。酒过三巡聊起香港时局，朋友说这几天经常会觉得一个人就像大海中的一滴水，像沙漠里的一捧沙子，风暴来临之际，只有随波逐流的份儿。}
\BookParagraph{其实不需要挂什么「十号风球」，即使是风平浪静，我们绝大多数人也不可能出人头地，在历史的长河中，充其量也只不过是一朵浪花罢了。}
\BookParagraph{每天买个菜，做个饭，朋友圈里点几个赞，在咖啡店喝一杯咖啡，看着窗外行色匆匆的人们，「比上不足、比下有余」，不是也挺让人满足的么？}
\BookParagraph{那年我竟然把家里那台背投影的彩电修好了，后来又看了好多年。}
\BookDate{2019 年 11 月，多伦多}
\BookArticle{追星}{追星}
\BookSubtitle{一}
\BookParagraph{周末参加了的一个学术会议的晚宴。晚宴上群星荟萃，「大佬」云集，舞台上美丽的主持人还在致谢，舞台下觥筹交错，敬酒环节已经提前悄然开场。}
\BookParagraph{「李老师您好。我先干为敬，您随意！您用微信吗？」}
\BookParagraph{「瞧您这话问的 —— 这跟问 “你会说中文吗？” 差不离儿了吧？」}
\BookParagraph{「那我能加一下您的微信吗？」}
\BookParagraph{「没问题没问题，不过我是真的没有朋友圈儿，不是把您屏蔽了啊。」}
\BookParagraph{正当我习惯性地拿出手机打开微信二维码的时候，舞台上忽然传来了亲切熟悉的歌声。我转过头去，看到舞台上一位女歌手，身材高挑，披肩长发，手持麦克风，正在开始唱黑豹的《无地自容》。}
\BookSubtitle{二}
\BookParagraph{这种感觉很神奇：在一场规模盛大的晚宴上，在一片嘈杂的敬酒祝辞中，竟然听到了一首二十多年前清华的校园广播中常常播放的摇滚歌曲。歌声高亢浑厚，动人心弦。}
\BookParagraph{它让我想起了我的《幸福》这本书里的同名压轴作品「幸福」中的一小段儿：}
\BookQuote{我二十岁生日那天请客，在大学生之家请哥儿几个大吃大喝。晚上十点多钟，她下自习路过时隔着栏杆看到我们，明显停了几秒钟。大学生之家的大喇叭里，正放着黑豹乐队的《无地自容》：人潮人海中又看到你 / 一样迷人一样美丽 / 慢慢地放松慢慢地抛弃 / 同样仍是并不在意 / 你不必过分多说你自己清楚 / 你我到底想要做些什么 / 不必在乎许多更不必难过 / 终究有一天你会离开我。}
\BookParagraph{节目结束后我问主持人：「能帮我介绍一下儿这位唱歌的老师吗？」}
\BookParagraph{见到这位从未谋面的女歌手的一瞬间，我想起电影「A Star Is Born」里 Lady Gaga 刚刚唱完《Always Remember Us This Way》以后的一段儿台词：「That was unbelievable. That was unbelievable what you did out there.」}
\BookSubtitle{三}
\BookParagraph{我们后来简短地聊了一会儿。从当年的窦唯和张楚，聊到 Lady Gaga《I’ll Never Love Again》最后忽然切换到 Bradley Cooper 在钢琴边上唱的尾声：「Don’t want to give my heart away to another stranger...」}
\BookParagraph{她说：「我无数次感叹，如果没有音乐，我的快乐将会少一大半儿。」}
\BookParagraph{聊到《无地自容》和当年的校园广播，她说：「每个人的记忆中都有这样一首歌，这也正是音乐的魅力。」}
\BookParagraph{聊到酣畅之时，我想起了《Always Remember Us This Way》里的歌词：「When the sun goes down, and the band won’t play, I’ll always remember us this way.」}
\BookParagraph{「刚才刚好录了一段您唱这首歌儿时的视频 —— 我能加一下您的微信吗？」我拿出手机说。}
\BookDate{2019 年 5 月，成都}
\BookArticle{五星红旗}{五星红旗}
\BookQuote{李老师您好，我想了很久还是想要给您发这样一个私信。我是您在香港理工大学当访问学者时教过的 e-commerce 的研究生。读书的时候深深折服于您的学术成就以及任何时刻都能够化繁为简的教学技巧及非常具有感染力的个人魅力，微博有时看您给其他人答疑解惑也觉得十分受用。}
\BookQuote{众所周知，香港最近的种种令人感到痛心，特别是昨天今天理大发生的事件更是让人觉得匪夷所思，我觉得已经发展到了失控的地步。看到曾经生活学习过的香港变成如今的样子，感到十分难过。同时也不解，为什么香港会变成这样？香港还能好吗？十分想听听您对香港最近发生种种的看法，如果您能回复我会十分开心！但是如果有不便，您不能或不想发表意见，我也十分理解。无论如何，还是非常感谢您看完我写的这些，一直觉得人生路上能遇到过您这样的老师，是我的幸运，谢谢！（问题来自微博私信）}
\BookSubtitle{一}
\BookParagraph{收到你的来信觉得特别亲切。六年前在香港理工大学访问期间，我教的这门 e-commerce 学位的研究生课是我第一次在香港教书，所有教学内容都需要从头精心准备，可以算是别具一格的一次体验。课虽然是英文教学，我发现大多数的学生都是内地生，课上课下互动很多，教起来非常开心。整个学期认识了好几个经常来问问题的学生，有的渐渐成了朋友，甚至现在还在保持联系。}
\BookParagraph{不仅仅是学生，访问期间我也结识了不少同系的教授，很快就和他们成了聊得来的朋友。有内地背景的，有香港本地的，也有加拿大人。大家平常一起在学校的餐厅海聊，参加系里组织的活动，去羽毛球馆打球，去星巴克咖啡店喝咖啡，去铜锣湾吃上海菜，去离岛区「行山」（hiking）。和内地背景的朋友说普通话，和香港本地的教授侃英文，大家都对我非常友好。听不懂广东话，我居然逐渐喜欢上了这个完全陌生的城市。}
\BookParagraph{其实不仅仅是我自己，那位加拿大来的教授在香港理工大学工作了二十多年，既不会说普通话也不会说广东话，也一样很喜欢这个城市。他告诉我，他最喜欢香港的是这里中西方水乳交融的文化。就拿吃来说吧，香港无论西餐还是粤菜都相当地道，甚至有巴西烤肉和日本料理拼在一起吃的日料小饭馆。}
\BookParagraph{我闲下来的时候，有时会在理工大学的校园里漫步。我很喜欢这里的红砖建筑，有点儿像南加州大学，建筑风格各异，但是颇有异曲同工之妙。很多楼都相连在一起，下雨天去什么地方多半儿可以室内到达。每栋楼本来各有自己的名字，但是大概是为了方便标注，一概以二十六个字母命名。我最喜欢的地方还是图书馆，一楼有大量可以随手翻阅的图书，内容涵盖天南地北。图书馆里还有一家咖啡馆，我经常坐下来边喝咖啡边干活儿。}
\BookParagraph{如果说理工大学的图书馆一直是我的「三味书屋」的话，那我住的地方 —— 九龙塘的又一城就是我「百草园」。在这里可以尽情地「拔草」—— 从奢侈品店、名牌时装店到苹果店，一应俱全。没事儿的时候我喜欢在又一城电影院看电影，在滑冰场看别人滑冰，或者在一家叫 Log On 的商店徜徉 —— 那里各种小玩意儿一应俱全，设计精奇，特别考验自律能力。圣诞期间，苹果店旁边的圣诞树大得出奇，上面挂满了各种饰物，玲琅满目，节日气氛浓郁。}
\BookParagraph{前几天，听说香港理工大学的图书馆被水淹了，星巴克早就给砸了，红砖建筑起火了，一所大学的校园成了攻防的战场。}
\BookParagraph{又一城也未能幸免。据说里面的圣诞树给烧了，只得停业。}
\BookSubtitle{二}
\BookParagraph{就好像欧氏几何所有的定理都是基于五大公理推出的一样，做人也必须有一些最基本的底线和「公理」。所谓「公理」，就是无论每个人的立场和观点如何，全社会都一致同意它们是正确的。}
\BookParagraph{而小学生守则里「遵守国法校纪，爱护公共财物」这条，则是做人的「公理」之一。}
\BookParagraph{违背了做人最基本的底线的人，无论其初衷和目的，都不值得同情。}
\BookParagraph{如果你说为了达到自己的目的，可以不顾做人最基本的底线，不用遵纪守法，那你和杀人犯强奸犯又有什么区别？}
\BookParagraph{就如同在电影院大声打电话会影响其他人的道理一样，世界上从来没有过绝对的「自由」—— 在你强调自己打砸抢烧的「自由」的时候，也必定会剥夺其他人的自由。}
\BookParagraph{当一座城市里说普通话都可能危及人身安全的时候，当一所大学里所有的内地学生需要集体紧急撤离的时候，这座城市和这所大学，不是变得更「自由」了，而是失去了自由。}
\BookParagraph{当一家咖啡店、一辆小巴、一家银行、一个过街天桥、一个地铁站、甚至一所大学的校园都可以随时被当做「死物」砸烂烧光的时候，这个社会不可能安居乐业，只能人人自危。}
\BookParagraph{当一个社会里有「黑衣人」可以藐视法律、可以嚣张到无故当街打人却无人敢管的时候，这个社会不是在进步，而是在迅速倒退。}
\BookParagraph{我最喜欢的电影《肖申克的救赎》里，Red 谈起 The Sisters 时说过一句名言：「You have to be human first. They don’t qualify.」}
\BookSubtitle{三}
\BookParagraph{我经常试图理解「抗议者」的逻辑，但却无论如何也想不明白。}
\BookParagraph{比如说，假设「抗议者」大获全胜，香港警察全线落败，再也出不来了。抗议者自己若是遇到人身威胁，打电话报警却没有任何响应，下一步再去找谁？}
\BookParagraph{再比如说，假设「三罢」极其成功，几百万市民被迫全都上不了班上不了学。无人给超市送货，无人维护无线电话和网络，「抗议者」们如何生活、又如何通讯联络？}
\BookParagraph{归根结底，大概「抗议者」们追求的理想社会，不外乎是他们自己可以不用上班、不用上学，但其他人则需要给他们保证供水供电、保证无线通信、保证衣食无忧。而警察则需要各司其职，给他们保障社会治安良好，但又不能管他们自己的违法行为。}
\BookParagraph{我相信，如此「在电影院大声打电话」的「抗议者」们，自己一定觉得比内地生高人一等，一定觉得说香港广东话比说普通话「高级」。}
\BookParagraph{也许他们应该去找一本美国六十年代民权运动历史的书 —— 对，就是他们热爱的那个美国 —— 看看什么才是真正的自由，什么才是真正的平等，而他们玩儿的，又是历史上什么人曾经玩儿过的。}
\BookSubtitle{四}
\BookParagraph{前几天去香港访问，正好路过香港理工大学的红磡湾校园。红磡湾校园离主校园有一小段距离，似乎没有受到破坏。那天天气特别好，蓝蓝的天空让人心情愉快。车等待红绿灯准备左拐时，隔着路边的树林，我远远地看到校园楼顶上有一面五星红旗，在迎着风飘扬。}
\BookParagraph{没错儿，无论「抗议者」们的目的是什么，无论香港的局势将来有什么变数，香港永远是中国的一个特别行政区，飘扬的国旗也永远是五星红旗。}
\BookParagraph{那首五十年代曾经脍炙人口的经典《歌唱祖国》所唱的，几十年过去了，还是这么令人心动：}
\BookParagraph{\hspace{1em}\hspace{1em}\textbf{宽广美丽的土地}\BookLineBreak{}\hspace{1em}\hspace{1em}\textbf{是我们亲爱的家乡}\BookLineBreak{}\hspace{1em}\hspace{1em}\textbf{英雄的人民站起来了}\BookLineBreak{}\hspace{1em}\hspace{1em}\textbf{我们团结友爱坚强如钢}\BookLineBreak{}}
\BookDate{2019 年 11 月，多伦多}
\BookArticle{运数}{运数}
\BookParagraph{今天周日，照例喝完咖啡，来到办公室完成一些需要整块时间才能完成的工作。工作间隙在网上无意中看了一篇批评 Apple Music 的文章，提到它对古典音乐的支持很差，应该好好向柏林爱乐乐团一个叫 Digital Concert Hall 的 app 学习。}
\BookParagraph{看到 App Store 上好评如潮，有点儿好奇，就注册了个帐号试试。你还别说，无论界面风格还是用户体验的设计都无懈可击，刚用了一小会儿就相见恨晚，喜欢得不得了。正打算交费，网站检测到我的 IP 地址是从多伦多大学来的，提示我说多伦多大学的师生可以免费使用。}
\BookParagraph{我琢磨着，今儿个运气还不错啊。发现今天其实不是一个普通的周日，而是 02/02/2020，一个八位数的「Palindrome」—— 正反来读完全一样。其实八位数的 Palindrome 并不少，比如 11/02/2011 也是。但它不是特别完美，因为如果按照英国人的日期写法，应该把月份写在后面，写成「02/11/2011」，那这一天就不是 Palindrome 了。}
\BookParagraph{如果这么完美主义地严格要求，那上一个八位数的 Palindrome 要回朔 909 年，到 11/11/1111；而下一个八位数的完美 Palindrome 不用这么长时间，不过也得等到 101 年以后，到 12/12/2121。}
\BookParagraph{我打开一首普罗科菲耶夫的钢琴协奏曲 —— 是王羽佳和柏林爱乐乐团演奏的 —— 开始浮想联翩：一百零一年以后，如果我有个外孙女，她估计会像包柏漪笔下的春月一样，九十岁高龄，五世同堂。那时候她的子女们，估计已经长住在火星上了吧。}
\BookDate{2020 年 2 月，多伦多}
\BookArticle{模拟游戏}{模拟游戏}
\BookQuote{亲爱的李老师，关注您很久了，我很喜欢看您写的文字：言之有物，温和、有力，闪着带有理性的感性，让人看了很安心。}
\BookQuote{有个问题我挺好奇的，就是每次看您写的内容的时候，我总感觉仿佛您的生活里充满的都是积极的、快乐的、温馨的内容，好像那些负面的社会新闻从未吸引您的注意力，选择地隔离它们，也从不会因它们产生任何情绪的波动和思绪的停留。好像您的眼睛只看得到光明，看不到阴暗。比如最近几度上热搜的杭州杀妻案、各种少女失踪情杀案件、蒙冤入狱的人生等。好像那些骇人听闻的苦难、令人揪心的他人经历跟你的生活井水不犯河水，不会产生任何的联系。我有些困惑，是这样的吗？如果是，想知道您的观点。如果不是，想知道您的处理方式。}
\BookQuote{过去有段时间我想清楚一个道理，每个人的生活都有各自的忙碌，大部分的新闻热搜其实都跟我个人毫无关系。我可以毫不关心。这种不置可否的「自私」确实节省了很多精力。但是人是社会的动物，一个人的经历和见闻，不可能与社会割裂开，而且长此以往，会不会让人丧失共情力，变成一个冷漠的人呢？可是看这些新闻，会不会产生不必要的戾气？}
\BookQuote{身为二十来岁的年轻人，在这个信息爆炸的时代，我们该如何筛选过滤信息源，该如何管理自己的注意力和精力，平衡社会与我们个人的关系？如何处理那些负面新闻给我们带来的心理冲击（或者压根儿不关心？）（问题来自微博私信）}
\BookSubtitle{一}
\BookParagraph{假如你想买辆车，无所谓「颜值」，无所谓品牌，也不差钱，唯一的希望就是这辆车质量靠得住，可以开很长时间不用怎么修，应该如何选择？}
\BookParagraph{无论如何，我建议你千万别去什么「汽车之家」之类的论坛。}
\BookParagraph{因为如果你到这样的论坛上去调研一款车的质量，你会发现那里充斥着各种「吐槽」：比如刚开了两个月发动机开着开着就开始冒烟啦、变速箱动不动就「顿挫」啦等等，不一而足。你一定会得出一个结论：无论是丰田马自达还是保时捷法拉利，质量都相当不靠谱，非常容易坏。}
\BookParagraph{甚至看得多了你还会感慨：十几年前买的车结实得很，没想到这年头汽车行业的质量越来越不靠谱了。}
\BookParagraph{但是实际上恰好相反，汽车行业总体来说质量比十几年前提高了不少。现在无论你买什么车，即使是直接挑一辆好看的买下，质量也差不到哪儿去。}
\BookParagraph{这是典型的样本选择偏差：绝大多数人平常开车开得好好的，完全没有必要到论坛上去郑重宣布「我刚买了辆丰田，开了三年没有半点儿毛病。」就好像到了医院你只能看到病人一样，到了汽车论坛上你也只能看到负面的「吐槽」。不能依靠研究汽车论坛来评估汽车品牌的质量，就好像不能因为流感有可能造成极小比例的病人死亡，就认为流感是一种极其可怕的传染病一样。}
\BookParagraph{看社交媒体上的新闻，自然也是同样的道理。坏消息在新闻中延绵不断滚动播出，不等于世界末日即将来临。就像绝大多数新车不用担心质量问题一样，绝大多数人过的日子都早出晚归，买房买车，上学上班，几十年如一日，平淡无奇。其实即使是轰动一时的新闻，也一样会像火车车窗外的小站站台，昙花一现，一晃而过。试想，现在还有谁会再去关心马航 370 或者马蓉翟欣欣的下落呢？}
\BookParagraph{几年前的一年年末，我去加勒比海坐了七天的邮轮，没有新闻，没有上网。}
\BookParagraph{回来以后发现，没人找过我，也没有错过任何值得追溯的新闻。}
\BookSubtitle{二}
\BookParagraph{劲爆的新闻就像火锅 —— 越吃就越爱吃，爱吃的人多了，火锅店自然会雨后春笋般地冒出来。同样，越是惊世骇俗的新闻看的人越多；而为了看的人多，作者就要添油加醋推波助澜写成「爆款」，尽量让你也看到。再加上用手机看新闻和以前的报纸杂志不同，只要标题引人注目，「重磅」、「无敌」、「震惊」，再加上若干惊叹号，即使我自己也无法抵御它的诱惑，想点进去看个究竟。}
\BookParagraph{本来这样的「标题党」「爆款文」并没有什么问题：人家也没有强迫我去看，我不看不就行了？哈哈，还真的由不得你不看：「爆款文」被亲朋好友转发「刷屏」，绝大多数人为了跟上这个时代的步伐，总还是会看一些的。看完了未必不会感慨万千或者义愤填膺，继续转发到自己的「票圈儿」。}
\BookParagraph{久而久之，由于经典的「劣币驱逐良币」效应，写高质量文章的人只会越来越少。高质量的文章需要基于事实，心平气和，有理有据，逻辑清晰，几经修改而成稿，本来就很难写。好不容易写出来了，即使冠名为「深度好文」也一样不可避免其门可罗雀的命运。天下虽然也不是没有免费的午餐，但俯瞰外滩全景的午餐，是无论如何不会免费的。这样恶性循环下去，即使是朔爷复出，金庸转世，也无力回天。}
\BookParagraph{那如何在「辣子」里找到那些宝贵的「鸡丁」呢？}
\BookParagraph{很遗憾，由于大多数标题已经无法代表文章的质量，我还真的没有什么迅速有效的好办法。唯一一个我能想到的途径，是去关注极少数我觉得靠谱的偶像或朋友的微博。比方说如果褚老师公开发表了一篇文章，即使是广告，我都会先睹为快。}
\BookParagraph{可惜芸芸众生之间，无论是偶像还是朋友，这样的人越来越少了。将来你要是偶然发现了好看的文字，发现了当年的辛夷坞、慕容雪村和醉钢琴，一定要推荐给我。}
\BookSubtitle{三}
\BookParagraph{你前几天有没有听说这么一条普通的新闻：微软发布了一款全新的游戏，叫 Microsoft Flight Simulator 2020。}
\BookParagraph{这款飞行模拟游戏逼真到了什么程度呢？它真实地再现了全世界两百万个城镇，十五亿座建筑，还有三万多个机场。所有的天气、地貌和建筑都是根据真实情况实时渲染的：无论是晴间多云还是暴雨如注，在飞机座舱内的景色都和真实场景一模一样。高速公路上跑着模拟的汽车，雨后还会惊喜地看到彩虹。}
\BookParagraph{这个游戏哪里还是游戏啊，简直就是一个全新打造完全真实的虚拟世界。}
\BookParagraph{它让我想到了斯皮尔伯格的电影「Ready Player One」：在三维的虚拟世界里，我们一起慢慢长大，成为同生共死的亲密爱人和战友。}
\BookParagraph{也许有一天，会有那么一个模拟游戏，可以让我全然忘记所有那些骇人听闻的苦难，那些令人揪心的场景，全心全意地进入一个幸福的虚拟世界。在那里，就像王朔在《一半是火焰，一半是海水》的结尾里说的，一路乘船、火车回家。在那里，我穿过广袤的国土，看到连绵起伏的著名山脉，婉蜒数千公里的壮丽大川，看到成千上万、随处可遇的开朗的女孩子。}
\BookDate{2020 年 8 月，香港}
\BookArticle{答案}{答案}
\BookParagraph{李老师，越来越多的人开始提到或者接受「世界不是非黑即白，看问题不能二元对立」的看法，我也很赞同这个看法，但一直找不到合适的语言描述世界不是非黑即白，那到底是什么样的呢？我们又该在这样的世界里以什么样的勇气和信念生活下去呢？ 希望听听李老师的看法。（问题二月二十七日来自微博私信）}
\BookParagraph{李老师我又来问问题啦！想请问李老师，在和父母长辈交流有不同观点时，有比较可行又不伤害对方的方法吗？我知道很难改变一个人的对某件事情的看法。如果是其他人，我根本就不在乎，反而还是一个筛选出那些不同想法的人的机会。但和亲密的人长期有思想上的隔阂又让我很难过。（问题九月十三日来自微博私信）}
\BookSubtitle{一}
\BookParagraph{看问题时是该非黑即白，坚定立场，还是该「萝卜白菜各有所爱」，那得看这提的是个什么样儿的问题。}
\BookParagraph{比如说，如果你的问题是「地球是不是圆的」或者「吸烟是否对健康有害」，那我当然应该斩钉截铁地给出一个正确的回答，因为可以用科学的方法明确得出正确的结论。}
\BookParagraph{如果对方辩友反驳说，我认识个老爷子抽了一辈子烟，还不是活了八十多岁？我自然可以告诉他，不能用具体的个例来推翻统计上的规律 —— 保险公司在卖人寿保险时抽烟的人保费更高，正是直观地表现了这种统计规律。}
\BookParagraph{但如果你的问题是「刘亦菲是不是比 Angelababy 更漂亮」，那当然就仁者见仁、智者见智，没有正确错误之分，没什么可争的。不但没什么可争的，这样的话题还是闲聊中万一断了话题可以用到的谈资：「哎你最近看《摩天大楼》和《花木兰》了吗？你觉得刘亦菲更好看还是 baby 更漂亮？」如果大家正好都是刘亦菲或者 Angelababy 的粉丝，话题跟香槟酒似的正好续上，火热继续。}
\BookParagraph{这两种问题性质完全不同，有点儿像我小时候看过的报纸。那时报纸不多，经常能看到的有两种：《参考消息》和《北京晚报》。}
\BookParagraph{《参考消息》虽然经常看到，但是我不怎么爱看。主要是因为那里面每一篇文章都很简短，不外乎就是些国际国内的大事儿，没有任何带一点儿感情色彩的评论。这就有点儿像第一种问题，黑的就是黑的，白的就是白的，泾渭分明。}
\BookParagraph{而《北京晚报》则有点儿像现在的今日头条，社会万象，娱乐花边新闻，可以看得很过瘾。夏天搬个小马扎儿拿上一份儿《北京晚报》到阳台，凉风习习，越看越爱看，惬意的很。这就有点儿像第二种问题，「站」不「站」肖战无关紧要，只要知道《陈情令》中有个叫魏无羡的人，咱就大有可以聊的话题。}
\BookParagraph{这第二种问题，「友谊第一，比赛第二」，无所谓谁赢谁输。就像照着菜谱炒菜，「放盐、糖、味精各少许，翻炒至八成熟」—— 不用这么精确，模棱两可，开心就好。而第一种问题则不然，关乎自然科学或者社会科学中的素养，不懂当然没有问题，也可以拿来争论不休，但是无论如何，客观上答案只有正确的和错误的两种，没有中间路线。把它「佛系」地当成第二种问题来含糊其辞，号称「看问题不能二元对立」，是完全错误的。}
\BookSubtitle{二}
\BookParagraph{那如果是第一种问题，正确答案从何而来呢？}
\BookParagraph{Netflix 上最近上映了一部纪录片叫《The Social Dilemma》，如果你没看过，值得花时间看一看。很多人看了觉得触目惊心，可我看了没觉得很意外啊，说的不就是今日头条嘛 —— 你爱看什么，社交媒体平台就有更多的类似内容会推送给你，你就会越来越相信你对这个问题的见解是正确的。}
\BookParagraph{比如「地球是不是圆的」这样的问题，如果你在社交平台上看到论证「地球是平的」的文章，会发现里面摆事实、讲道理，证据确凿，言之有物，看多了一定会越来越相信地球其实是平的。你还别说，相信地球是平的的人还挺多 —— 美国每年都专门召开个会议，就叫「平地球国际会议」(Flat Earth International Conference)，有点儿粉丝后援团聚在一起开展活动的感觉。}
\BookParagraph{我认为，与其到网上随手「刷」一些支持自己立场和观点的推送文章，不如去「特别关注」一些高水平的作者写的文字。高水平的作者虽然平时写的文章不一定很多，但是一旦下笔，如果回答的是第一种问题，其文字中的观点一定逐字逐句地仔细推敲琢磨过，力求客观、符合历史规律。}
\BookParagraph{高水平的作者不经常在他们不熟悉的领域发表见解。他们会抱着学习的态度，不会不懂装懂，对每一种观点都充分怀疑，去伪存真。著名哲学家罗素 —— 就是《围城》里褚慎明亲切地称为 Bertie 的那个罗素 —— 就曾经说过：「这个世界上的一大问题是，傻子和狂热的人永远认为自己正确，而聪明人则满满地全是怀疑。」(The whole problem with the world is that fools and fanatics are always so certain of themselves, and wiser people so full of doubts.)}
\BookParagraph{所以说，即使是高水平的作者写出来的文字，我们也不能照单全收，而是应该自己静下心来仔细想想他说的是否真的有道理，逻辑上有没有问题，甚至用科学的方法从侧面调查研究一下支持其观点的论据是否可靠。}
\BookParagraph{我所敬仰的经济学家 Thomas Sowell 曾经说过：「意识到自己的无知，需要不少知识。」(It takes considerable knowledge just to realize the extent of your own ignorance.)}
\BookParagraph{遗憾的是，在这个「人工智能」的大时代，往往越是罗素或者 Thomas Sowell 这样相当靠谱的作者，写出来的文字看过的人就越少。}
\BookSubtitle{三}
\BookParagraph{那如果和父母或亲友对一个问题的看法完全不同，该怎么办？}
\BookParagraph{最近金斯伯格 (Ruth Bader Ginsburg) 去世了。这位金斯伯格大法官是位传奇般的人物，一生非同凡响。两年前有一部叫《On the Basis of Sex》的电影非常好看，刻画过她的生平。我还发现老太太非常热爱运动健身，写作风格则既浅显又简练，这一点比 Thomas Sowell 强不少：Sowell 的著作篇幅往往都很长，最著名的《Basic Economics》这本书看到后来，都有点儿车轱辘话来回说的感觉了，浅显易懂没有问题，但完全谈不上简练。}
\BookParagraph{金斯伯格婚礼那天，她的婆婆给了她一条忠告：「In every good marriage, it helps sometimes to be a little deaf.」（一桩好的婚姻中，稍微有点儿耳背是好事儿。）}
\BookParagraph{金斯伯格说，她的婚姻五十六年，一直都是这样做的。即使是在最高法院的工作中，她也将这句话奉为信条：就算别人说的不对，也最好「这耳朵进，那耳朵出」。}
\BookParagraph{我想，老太太可能是在说，每个人都在走他喜欢的那条路，而即使是至亲的人，你也很难改变他自己选的那条。}
\BookParagraph{这让我想起了我喜爱的电影《花木兰》里的一句台词：}
\BookParagraph{Take your place, Mulan.}
\BookDate{2020 年 9 月，香港}
\BookArticle{理想国}{理想国}
\BookParagraph{葆春老师好。经常看您的微博，您的回复真诚实在，给即将毕业步入社会的我很多启发。我原本打算留学美国，托福 GRE 也已经考完，现在虽然拜登赢面大，但是现在考虑到特朗普期间急转直下的中美关系，美国大幅取消中国留学生签证，心里还是打鼓。想请教下您怎么看待 2021 年留学美国的形势呢，是不是现在对于留学生而言，放弃美国考虑其他国家才是明智之举。非常感谢您！（问题来自微博私信）}
\BookSubtitle{一}
\BookParagraph{中美关系我虽然不懂，但是我以前提起过我的舅舅王缉思，他研究中美关系数十载，是国内当之无愧的中美问题专家。所以但凡跟中美关系有关的问题，我都去读他的文章，希望从中找到答案。}
\BookParagraph{他十年前曾经写过一篇文章，题目是《如何看待美国地位的相对下降》，后来也经常重申其中的观点：和其鼎盛时期相比，美国虽然有衰落的迹象，但是「瘦死的骆驼比马大」，对于我们这些普通人，离真正需要担心是否要去美国留学这样的问题的时候还很遥远。我们还是来摘抄一段王缉思教授的原文吧：}
\BookParagraph{「我们的问题在于下结论太草率，往往凭一时一事。如九一一事件发生，有人认为美国从此就不行了。现在利比亚的危机出现，美国束手无策，于是乎美国不行了。对于金融危机的判断也是这样。关于美国兴衰的研究，我想有许多问题可以讨论，但我这里想说的是，国内的研究者也好，我们的整个舆论也好，几十年来低估美国、唱衰美国的时间更多一些，而有很多次已经被事实证明是错误的。有的时候就是一种情绪被掺杂在研究或者判断里面：我不喜欢美国，就不喜欢说美国强大。」}
\BookParagraph{同样的道理，现在虽然财政危机、政治两极化等问题还在更为集中地不断出现，但是如果我们把眼光放长远一点儿，很难就此认定美国已经从此衰落、一蹶不振了。既然如此，如果我们假设美国的学术水平无论人文还是工程学科还是世界一流的，那我们就没有任何理由不去这样全世界最好的一流大学留学。这些大学有最好的师资力量和科研实力，留学又是「过了这村儿就没这店」的事儿，为什么要舍本逐末呢？}
\BookParagraph{我听说，很多国内的学生都考虑放弃去美国，而改为去其他国家和地区的大学留学。如果这是因为客观的因素比如疫情或者签证问题，当然无可厚非；如果是因为「急转直下的中美关系」而自己改了主意，那我就不明白了：中美关系跟我们普通人去美国留学有什么本质或必然的联系？不仅如此，我这个人还天生有一种「逆水行舟」（contrarian）式的叛逆心态：如果大家都因为微信公众号看多了而不去美国留学，不正好是我去美国最好的机会么？}
\BookParagraph{我盲目地相信，那些号称不去美国的同学，如果收到一份哈佛法学院的录取通知书，一样会兴奋地跳起来。他们不去美国，多半儿不过是嫌自己申请到的学校不够知名而已。}
\BookSubtitle{二}
\BookParagraph{不过，不知道你发现没有，很多人但凡选择了在一个国家或地区留学，毕业后会自然而然地定居在那儿。}
\BookParagraph{这就意味着，如果你是一个普通人，去一个地方留学很有可能也就意味着将会在那里长期工作定居、抚养下一代。}
\BookParagraph{如果真的有这样的「惯性」存在，那去什么地方留学这个问题就变得复杂得多了。我们不仅需要关心未来几年，更要关心未来几十年这个地方将会如何发展，是否适合在这里长期居住、儿孙满堂。}
\BookParagraph{把这个问题换一种问法就是，有没有这样一些「理想」的国家或地区，可以对于一个普通家庭来说尽量保证未来几十年的幸福？}
\BookParagraph{王缉思教授在其近著《世界政治的终极目标》中，恰好也讨论了「理想国的标准」这个问题 —— 看来要想把中美关系研究得深入，涉猎真的得称得上是「博大精深」才行，涵盖政治、经济，甚至哲学。我还是来摘抄一段他的原文吧：}
\BookParagraph{「一个成功的、令人向往的国家应当是：第一，没有严重的外部和内部的安全威胁，国内政治稳定，暴力犯罪率低；第二，国家和民众都比较富裕，经济稳步增长；第三，国家有相对统一的信仰体系、道德准则和主流价值观，同时包容一部分公民所奉行的其他信仰，公民对国家认同度高。第四，公民之间贫富差距较小，公民平等在教育和社会保障体系中得到较好体现，社会不公能够通过法律和政策调整得到矫正，抑止官员腐败；第五，公民的自由权利得到充分保障，个人自由同民族、国家的自由相一致。」}
\BookParagraph{如果以这样的高标准严要求来衡量，北欧国家 —— 比如丹麦 —— 一定可以金榜题名。我的偶像刘瑜老师为美国斯坦福的政治学家福山的著作《政治秩序与政治腐败》写的导读「如何到达丹麦？」，讲的就是如何到达丹麦这种国家如童话般的境界 ——「有法治、又民主，政府还高效而廉洁」。}
\BookParagraph{丹麦或其他北欧国家虽然不错，可惜我有点儿怕去一个不说英语的陌生国家，也懒得花时间再去学一门外语 —— 真的，我觉得那些可以在非英语国家长期定居的中国人都非常令人敬佩。}
\BookParagraph{这么说来，如果出国去一个英语国家留学继而长期居住，选项就所剩无几了：英国、美国、加拿大、澳大利亚、新西兰。}
\BookParagraph{这些选项中，很多人最终还是去了美国。}
\BookSubtitle{三}
\BookParagraph{二十八岁那年，我在美国留学，周围的人基本毕了业都留在了美国。我则反其道而行之，决定去加拿大，一个我从来没去过的地方。}
\BookParagraph{为什么要舍近求远呢？}
\BookParagraph{我曾仔细研究过美国和加拿大的区别：两国的移民政策、医疗系统、还有退休制度，都完全不同。最后得出了一个结论：虽然两个国家都远远不能算是「理想国」，但是从我这样一个普通人的视角来看，加拿大比美国更适合长期定居。}
\BookParagraph{我们拿退休保障制度当一个简单的例子来聊聊为什么吧。}
\BookParagraph{美国联邦政府的退休金制度叫 social security。如果是公司的雇员，每年需要扣除工资的 6.2\%，直到 137700 美元的收入上限，公司还要同时再支付同样一笔钱。也就是说，如果年薪超过了 137700 美元，每年要从自己应得的收入中拿出一万七千美元交给政府。六十六岁退休年龄一到，如果工作了三十五年，可以最多每月领到三千美元左右的退休金。}
\BookParagraph{如果我们简单一算，会发现这个退休金制度不是很合算：三十五年交了六十万美元，退休每年只能拿回来三万六千美元，需要至少活到八十三岁才能把本金收回。可是如果自己拿六十万美元去投资，总是会有投资分红收入的，怎么感觉给了政府就没有利息了呢？}
\BookParagraph{其实还真是这么回事儿。年轻人每年交给 social security 的钱，美国政府会先给退休后的老年人发退休金，如果有剩余则去买国债。也就是说，收入减去支出余下的钱，并没有去投资，而是借给政府花掉了。}
\BookParagraph{随着平均寿命越来越长，人口逐渐老化，十年前支出就已经超过了收入，预计再过十几年，就算可以把以前借给政府的钱全拿回来，也都会尽数花完。那时候这个退休金制度是否会发生变数，就不好说了。}
\BookParagraph{反观加拿大的退休计划 Canada Pension Plan，把雇员和公司交的钱加起来，每年最多只需要给政府交五千八百加元。而六十五岁退休年龄一到，最多每年可以拿回一万四千加元，活到八十一岁就可以把交过的本金收回。}
\BookParagraph{两者相比，我更愿意给政府少交点儿，到退休后也少拿点儿。}
\BookParagraph{更为重要的是，加拿大把收到的退休金交给一个叫 Canada Pension Plan Investment Board 的公司去投资，在全世界买了不少现金流还不错的产业，比如机场、输电线和高速公路。这样基本可以保证无论人口如何老化，任何时候都不会没钱发退休金。}
\BookParagraph{您可能会想，这一个月才一千多块钱，不够花啊？}
\BookParagraph{其实政府的退休金，平时虽然只不过是「锦上添花」，可一旦到了山穷水尽身无分文的时候，还是可以保证最基本的生活的。而恰恰是因为有了这样的基本保障，一辈子就不需要特别节省或者投资，可以有多少钱就花多少钱，尽情当一个「月光族」，不怕「outlive my money」。}
\BookParagraph{有多少钱全花光 —— 这是多幸福的一件事儿啊。}
\BookSubtitle{四}
\BookParagraph{头一次从美国长途开车去加拿大的那天晚上，一路上基本漆黑一团，没有路灯，也没有导航和手机。直到远处渐渐看到了万家灯火，我才知道多伦多已经快到了 —— 这里没有亲人，也没有朋友。}
\BookParagraph{其实，无论是去哪里留学，还是到哪里定居，说不定答案都非常简单。}
\BookParagraph{哪里有你喜欢的馆子，有你多年的友情，有你爱着或者想追求的人，就去哪里。}
\BookParagraph{如果那儿正好酷似童话般的丹麦，安全、富裕，可以安安稳稳地做个「月光族」，那当然就更好啦。}
\BookDate{2020 年 12 月，香港}
\BookArticle{银婚}{银婚}
\BookParagraph{今天是毛主席诞辰，也是我和郭芳的结婚纪念日。二十五年前在北京结的婚，今年可以算是「银婚」了。}
\BookParagraph{记得当时结婚前几天在中关村的一个商场花了两百多块钱买了一对普通戒指，后来多年来也没买过什么金银首饰，应该算是「裸婚」吧。姥姥当时批评我结婚不能如此简单，戒指一定要纯金的，还当场给了我一小块黄金作为礼物。}
\BookParagraph{燕南园六十号姥姥的卧室，从小到大看惯了外公亲笔写的一首送给姥姥的诗，是他们结婚四十五年时写的：「甜甜苦苦两人尝，四十五年情意长。七省奔波逃猃狁{\BlissSymbols ①}，一灯如豆伴凄凉。红羊{\BlissSymbols ②}溅汝鲛绡{\BlissSymbols ③}泪，白药医吾铁杖伤。今日桑榆晚景好，共祈百岁老鸳鸯。」诗从三十年代的抗战写到六十年代的文革，绝对可以算是「甜甜苦苦两人尝，四十五年情意长」了。外公和姥姥五十年金婚纪念那年，照了一张黑白合影。后来这首诗和这张合影，就永远嵌入了他们位于万安公墓的墓碑之中。}
\BookParagraph{希望我们将来也可以有机会纪念金婚，最好还能青出于蓝而胜于蓝，再纪念一下钻石婚 ——「今日桑榆晚景好，共祈百岁老鸳鸯」。}
\BookDate{2021 年 12 月，多伦多}
\BookFootnote{{\BlissSymbols ①} 猃狁 (xi{\BlissCJKFallback ǎ}n y{\BlissCJKFallback ǔ}n)：古代族名，是匈奴族的旧称。这里指日本侵略者。}
\BookFootnote{{\BlissSymbols ②} 红羊：从古代纪年方式来看，“ 红羊 ” 指 1967 年，是文化大革命期间。}
\BookFootnote{{\BlissSymbols ③} 鲛绡：薄纱。}
\BookPart{将爱情进行到底}{将\\爱\\情\\进\\行\\到\\底}
\BookArticle{结　婚}{结婚}
\BookParagraph{\textbf{Trade-off }n.}
\BookParagraph{1. a balancing of factors all of which are not attainable at the same time.}
\BookParagraph{2. a giving up of one thing in return for another.}
\BookParagraph{— Merriam-Webster Dictionary}
\BookParagraph{事业生活双丰收的时候，大凡也就该结婚了。婚结了自然是好事儿，远的三大件嫁妆不说，近的也有减税绿卡种种甜头。结婚是意义深远的里程碑，或一锤定音，或一辈子就那么两三回，终于可以「雄姿英发，羽扇纶巾」，双双出没于县城集市之间，鸡尾酒会之上。为了结婚，咱即使再拿不出款子，也要八抬大轿，张灯结彩，大红喜烛，含糊不得。老外更不能免俗，又得求婚，又要买钻石戒指，庆典之上还要好多帅哥美女盛装陪衬，争先发言，专业摄影师全程追踪，如同一场自编自导自演的故事片，非同小可。}
\BookParagraph{结婚其实是一种投资，日子久了，才能看出水平的高下。遗憾的是，它多半儿是一辈子风险和注金最高的一次赌注，赌的是一生的幸福。当然了，不喜欢可以换别人，但是谁能知道换了是否就更理想呢？离婚再结是有代价的，人往高处走，凡聪明人不到万不得已不费这劲儿。结婚又如同大国间的外交，牵一线动全局，无论政治经济，如履薄冰，务必就共同关心的问题广泛地交换意见，达成共识，发联合公报，才或可长治久安。兵临城下之时，或三十六计，或围魏救赵，或远交近攻，八仙过海，力挽狂澜。一个小名儿，一个电话，肩膀上轻轻地一拂，真情流露，远胜那情人节时的新款手机。在这罗马史诗般的乐章之后，或可达到那世外桃源、男耕女织、悠然见南山的境界。要是一辈子也理会不得这柴米油盐中波澜壮阔的言外之伏笔，恐怕他也当不了咱外交部门的领导。套用洋文字典里的单词儿，这结了婚的方程没有最优解，有的只不过是所谓的「trade-off」。}
\BookParagraph{朋友的朋友结婚了。新娘子预先参加了糕点制作培训班，亲手制作了三层有九十九朵玫瑰的结婚蛋糕，人见人爱，精彩绝伦。又有朋友青梅竹马，结婚七年，却孜然一身流连于美国美丽的农村，爱人奔波于中美之间，数月未曾鸿雁传书，不知其所。姑妄揣测之，朋友的朋友如中朝人民伟大的友谊，唇亡齿寒，舍身取义，却不免危机四伏。朋友则如中苏忘年之交，精神支柱消匿于无形，如何才能力挽乾坤？拙文要是让朋友的朋友读到，一个电话，我就又得回家好好研习马克思主义哲学唯物辩证法了。}
\BookDate{2011 年 8 月，多伦多}
\BookArticle{永恒的爱情}{永恒的爱情}
\BookParagraph{最近把周末所有的闲暇时间，全都用来编程。}
\BookParagraph{发现编程有点儿中国古代封建社会「自给自足的自然经济」的感觉，守着自己的一亩三分地，辛勤耕作，先苦后甜。不需要关心达沃斯的时局、比特币的兴衰、或者高晓松的八卦，「采菊东篱下，悠然见南山」，颇为自得其乐。}
\BookParagraph{昨天在设计实现一个全新的聊天留言板。用户界面里每一条留言的背景颜色，尝试了不同程度的浅灰、浅灰里暗含淡紫、橘红、青绿、湖蓝等颜色。长方形背景边框的圆角，还有字体的大小和粗细，也尝试了好几种不同的方案。既别太圆润，也别太棱角分明，才能恰到好处。}
\BookParagraph{设计界面里好几个细节的时候，都和郭芳讨论了一下，就像二十多年前在清华学堂的自习教室里，和她在讨论一道线性代数的习题。}
\BookParagraph{也许就像我几年前写的一条微博里所说，永恒的爱情，其实就是随时随地可以聊天儿。而关注的话题不是达沃斯，不是比特币，也不是高晓松，而是一个聊天界面设计的背景色。}
\BookDate{2018 年 2 月，多伦多}
\BookArticle{异国恋}{异国恋}
\BookParagraph{怎么维持异国恋，大家给我出出招吧。（问题来自「向上移动」）}
\BookParagraph{去年端午节的时候，忽发奇想，一个人去看了《北京遇上西雅图之不二情书》。汤唯一如既往的漂亮，不过最喜欢的台词儿，还是那句「那时没有一秒钟就可以到达的电邮，等一封信，漫长得如同一生，但是慢一点儿又有什么不可以呢。慢一点儿，才能写出优雅浪漫的话语；慢一点儿，才能仔细寻觅盼望的爱情。」}
\BookParagraph{那年离开北京去美国读书时，行李箱打上了五彩尼龙绷带，她专程来机场送行。一切都是那么笃定，就好像只是一场短暂的别离，很快就要重逢。}
\BookParagraph{我们每三天给对方写一封信，每一封信都大致能在七天左右到达。每一封信里提到的，都是上上封来信里的事儿。但久而久之习惯了以后，也就觉得非常自然，理应如此。我们靠来往书信约定下次打电话的时间，她会专程骑车去北大燕南园姥姥家，一边和姥姥闲聊，一边等我的电话。那时的国际长途电话，一美元一分钟，因为舍不得钱，其实说不了几句话，就是想听听对方的声音。}
\BookParagraph{那会儿清华大学刚刚连上国际互联网，她的教研组只有一台主机可以收发国际电子邮件。那台主机是实验室所有人共享的，所以她就只有晚上十点钟以后才能收发邮件。负责机房的师兄很照顾她，经常会多等一会儿再锁门。为了多一点儿私密性，我专门编了一个加密和解密的程序，她收到的加密邮件如果没有解密，别人无法随意阅读。}
\BookParagraph{我所在的小镇离芝加哥机场大约三小时车程。为了能顺利地到机场接她，我很早就开始学开车。她到达美国那天的前一个晚上，我到超市买饮料，为了她到来做准备。我平常买橙汁为了省钱，买的都是「from concentrate」。那天，我毅然决然地买了「not from concentrate」，仿佛这是全世界最奢侈的品牌。}
\BookParagraph{从那以后，我们搬过家，但是再也没有分开过。就像童话故事里说的：「Lived happily ever after.」}
\BookParagraph{两个女儿长大了，我们给她们每人一个iPhone，然后建了一个群，叫「FamiLEE」，中文大概可以叫「老李他们家」。通信太方便了，就开始怀念当年写信的日子。大段大段手写的文字，特别真实，特别简单，特别自然。就像《不二情书》那部电影说的：「信是千里之外的眼睛，它能注视人心。」}
\BookParagraph{花一点儿时间写信吧。贴上邮票寄出，将来再用橡皮筋绑好所有的信，权当青春的纪念。}
\BookDate{2017 年 6 月，西安}
\BookArticle{心　态}{心态}
\BookParagraph{李老师好，我很喜欢您的文字，也喜欢您的生活方式，每次看您的微博都觉得您的生活就是我一直追求的生活。但是我目前还在迷茫的阶段，在复旦读大二，对于自己的未来，对于以后的生活伴侣都有诸多焦虑，害怕自己做出错误的选择。请问老师年轻的时候遇到过这样的焦虑吗？该如何调整心态呢？（问题来自微博私信）}
\BookParagraph{无论什么时候、什么事儿，我都经常会焦虑啊。比如刚刚收到你这个问题那会儿，我就想，完了，命题作文来了，写不出来怎么办？你的问题就差「字数不少于八百字，文体不限，诗歌除外」了。}
\BookParagraph{一个人的焦虑，是因为未来的阴错阳差，完全无法预料。我们都希望一个好的结局，越是重要的事儿，越害怕事与愿违。在焦虑这件事儿上，我对你最多只是「五十步笑百步」而已。}
\BookParagraph{王朔是北京一比较贫嘴的作家，多年前写过一部小说，名字叫《顽主》。里面给「五十步笑百步」造句：「新娘上轿，前五十步笑，百步以后哭。」}
\BookParagraph{你瞧，同样六个字，换个方式来看，结果完全不同。}
\BookParagraph{有些事儿吧，我们还真得想个办法，换个角度来看。}
\BookParagraph{比如寻找「生活伴侣」这事儿吧，什么才叫「错误」的选择？为什么一个「错误」的选择就满盘皆输？}
\BookParagraph{一个滴酒不沾的数学家，天天不是在餐馆打工，就是在公园演算数学证明；另一个酷爱喝酒、撸串，天天不求上进，小富即安，哪个选择是「错误」的？}
\BookParagraph{如果运气不好正好碰到了一个「错误」的选择，为什么不能立马离开这个选择，一切重新开始？}
\BookParagraph{我们系里有一个葡萄牙裔的秘书，经历过四任丈夫，一共抚养了四个儿子。她自豪地将他们的照片摆在办公室的书架上，我看着这些照片，个个光彩照人。有一瞬间，我觉得他们的妈妈才是真正的「人生赢家」。}
\BookParagraph{如果一个选择虽然非常重要，但是没有对错之分，甚至选与不选都无所谓对错，是不是心情就一下子轻松好多？如果新娘上了轿，前五十步笑，百步以后可以换一个轿子坐，是不是上哪个轿子都差不多？}
\BookParagraph{刚才给你写这篇作文的时候，我一旦明白了写得怎么样其实无关紧要，一下子就豁然开朗，如释重负。于是草草成文，一气呵成。}
\BookParagraph{写完一看，好像也没差到哪儿去啊。}
\BookDate{2018 年 2 月，多伦多}
\BookArticle{克　夫}{克夫}
\BookParagraph{李老师，我是在《向上移动》认识您的，包括微博每次看到您写的内容都有如沐春风的感觉，是不同褚老师的另一种广博和包容。}
\BookParagraph{最近有件事让我翻来覆去睡不着，请您空闲时指导指导。按说我是不信「命」一说的。可是开始谈恋爱后告诉家人，家人了解到女孩是「铁扫帚命」后对我很是担心。他们的依据是我姑姑也是这种命，结果克夫改嫁，自己不能寿终正寝；一个表嫂同样出生于不吉利的月份，表哥十年前三十岁出头就死于白血病。}
\BookParagraph{开始我觉得无所谓，但我爸已病故，我妈特别害怕再失去我，现在我也开始犹豫。让我犹豫的是姑娘对我很好，我不忍心仅仅因为可能会被克死或终生不幸而说分手，但被身边这些事例搞得不能坚定自己。我妈也陷入了是信命还是不能这么辜负女孩的两难处境。作为一个科学工作者，李老师一定是不会迷信「命」这一说的，能不能说说您的看法？（问题来自微博私信）}
\BookParagraph{看完你的故事，我真为那个一意孤行要嫁给你的女孩担心。如果我有机会见到她，听她讲了一个一模一样的故事，我一定会坚决地跟她说，千万别嫁给这个男生。}
\BookParagraph{假设这个故事将来是这么发展的：你们俩冲破重重阻挠，历尽千辛万苦，不顾你母亲的坚决反对，终于走入了婚姻的殿堂。结婚之后，但凡出现任何波折——孩子虽然聪明但是长得普通、工作虽然稳定但是升迁无望、买房虽然坚决但是只跌不涨、父母虽然行动如常但是小病缠身——你第一时间会怎么想？}
\BookParagraph{没错儿，无论真正的原因是什么，你百分之百会认为，这是女孩「克夫」的结果。}
\BookParagraph{关键问题在于，出现让你产生这种想法的波折的概率，也是百分之百。因为这样的波折，是生活的一部分。}
\BookParagraph{而你的想法其实也不能全怪你。你在婚前婚后收到了大量的心理暗示，致使「克夫」这个词儿和各种生活中不尽人意的故事一直在紧密关联着，在你的心里根深蒂固。大多数人都有这个倾向：一个错误的观点重复了很多遍，就不由自主地变成了事实。}
\BookParagraph{那一个「科学工作者」是不是会迷信「命」这个东西？绝对不可能，因为这是一个逻辑错误。所有和「命」有关的逻辑命题，都认为如果事件A和事件B有关联（correlation），就可以推出它们之间有因果关系（causation）——事件A导致了事件B的发生。}
\BookParagraph{比如你观察到看电视看得多的孩子比较熊，于是得出了看电视容易让孩子更熊的结论。但是这个结论是不对的，正确的事实是熊孩子比正常孩子更爱看电视。}
\BookParagraph{比如你发现风车转得越快，风就会越大，于是得出因为风车旋转而形成大风的结论。而正确的事实自然是因为有了风，风车才会旋转。}
\BookParagraph{在中世纪的欧洲，人们发现虱子一旦离开人体，人就会很快生病，所以一度认为虱子是对健康有益的。但其实真正的原因是虱子怕热，所以人发烧的时候不会长虱子。}
\BookParagraph{这些例子里的逻辑错误还算有情可原，因为至少依据的事实是基于统计数据的关联性。而说张三家的大门「路冲」所以儿子考不上大学，说李四家的寡妇是天蝎座的所以天蝎座的女人「克夫」，王五家全家都是属猪的所以没法儿不胖等等，则连基于统计数据的关联性都不存在，纯属令人啼笑皆非的无稽之谈。}
\BookParagraph{所以说，如果那个一意孤行的女孩听完了我的建议，不顾我的坚决反对，还是一心一意地要和你结婚的话，用一句不科学的话来形容，你的「命」简直好得令人惊叹。将来你们婚后如果孩子聪明健康、工作轻松稳定、房子早早买入、父母行动如常，你一定不要忘了感谢她。}
\BookParagraph{这叫什么来着？好像叫做「旺夫」的命，是吧。}
\BookDate{2018 年 4 月，多伦多}
\BookArticle{天下大乱}{天下大乱}
\BookParagraph{老师，我和前女友认识四年了。我们在一起三年，因为读研异地了一年。在这一年中因为一点点小摩擦，在她身边她觉得有个人更懂她，而我说话可能更加现实直白，一个星期前我们分手了。我觉得好痛苦，为什么四载不如几月珍，为什么直接宣判我死刑，而选择和别人在一起？}
\BookParagraph{真的特别痛苦 —— 如果我做错了，可以告诉我啊，给我一次机会去改。但是她一直冷战逼着我说分手，上周我跑了一千多公里偷偷去找她，结果等来的却是她和那个男生一直在微信上聊，并且两个人互相都有了好感。}
\BookParagraph{我觉得有时候如果她能想想我们三年的感情，何必会成为现在这样子？（问题来自微博私信）}
\BookParagraph{如果你去知乎搜索「失恋」这个关键字，第一个跳出来的问题是「怎样快速走出失恋？」这里面有四千六百三十四个回答等着你去读，最高赞的回答有六千七百一十七个赞同。}
\BookParagraph{可是我还是应该负责任地告诉你两个事实。}
\BookParagraph{一、千万不要抱有幻想，无论你怎么做，一切已经无法挽回；}
\BookParagraph{二、「快速」地走出失恋是不可能的。}
\BookParagraph{毛主席一九六六年给江青写过一封信，信里说「天下大乱，达到天下大治。」失恋给人的感觉，就是一种「天下大乱」，一触即发，声泪俱下。}
\BookParagraph{有人建议去旅行、去运动、去读书、去恋爱，但是你会发现，「天下大乱」的时候，其实这些都不管用。即使你真的试图去再爱上别人，每一次微信提示音响起，每一次手机在身边震动，你都会问自己「是不是她？」}
\BookParagraph{我高中有一段时间，曾经一度发疯似的开始给她写信。信的内容我不记得了，但是不外乎和你现在一样，在问她「为什么」这样的问题。说是写给她看的，但是其实只能算为了自己而写。因为我无法知道信寄出以后，她到底会不会打开信封，把信看完。}
\BookParagraph{有一次妈妈发现了我在写信，说你还想不想好好准备高考提高学习成绩了？我非常不争气地哭了。}
\BookParagraph{我自始至终没有收到过她的回信。也无从知道所有那些悬而未决的问题的答案。}
\BookParagraph{三十年后，我没有加她的微信。因为我不能肯定，如果我发出请求成为她「新的朋友」，这个请求会不会过期。}
\BookParagraph{不过你放心，时间是一付万灵药。时间会在你痛苦时、反思时、声讨时、回忆时，一年年不经意地过去。直到有一天，你会发现自己的记忆越来越模糊，如果不去翻看以前的照片，你会连她美丽的模样都记不清了。}
\BookParagraph{在这段痛苦的时间里，不知道哪里来的力量，你也许会通读萨特、叔本华、尼采的著述，迷上西方哲学，成为一名中青年学者，著作等身；你也许会沉迷于各路论坛知乎微博，横刀立马，舌战群儒，成为一名「十万加」「大V」作家；你也许会云游四方，看破荣辱成败，大彻大悟，成为一代高僧。}
\BookParagraph{「天下大乱，达到天下大治。」}
\BookParagraph{毛主席说得对。}
\BookDate{2018 年 6 月，西安}
\BookArticle{两　难}{两难}
\BookParagraph{尊敬的李老师：您好。我想冒昧咨询您一个感情问题。我是一名大三学生，和男友交往了大概八个月。现在到了升学时期，我的成绩比较好，可以保外校；男友的差一些，可以保本校，考研考上我所中意的学校很难。但是在这几个月的相处中，我发现我们非常合适，他很聪明，为人很好，对我也很好，然而对于升学，他看重专业想留在本校，我则更想去更好的学校，我们总是因为这个闹得不愉快。我的问题是：我应该为了男友留在本校，还是去更好的学校发展？另外，您在我这个时期遇到过类似的情况吗？那时您是怎么处理的呢？（问题来自微博私信）}
\BookParagraph{你出的这个作业，我想了将近一个月才动笔。}
\BookParagraph{大概是因为这有点儿像一个经典的两难问题（dilemma）。无论哪一个选择都有它的坏处，英文叫做「damned if you do, damned if you don’t」。}
\BookParagraph{这就好像在问：女朋友生日那天晚上老板要求加班，我到底是去还是不去？}
\BookParagraph{不过其实你的问题和我举的这个例子又不太一样。我的例子里，「两难」表现在要不然得罪女朋友要不然得罪老板；而你的问题里，选择去更好的学校读研，「得罪」了男友，却无法保证对将来的事业发展更有利，这只是一种期望而已。而不去更好的学校读研，不会「得罪」任何人。}
\BookParagraph{而我想了半天，给你的建议恰恰是尽量争取去更好的学校读研。}
\BookParagraph{这是因为我觉得两个人在读书时分开住一段时间，虽然有点儿辛苦，却是一件非常浪漫的事儿。}
\BookParagraph{我说的「辛苦」，不是大晚上该学习的时候，你们俩微信来微信去、语音视频聊个没完的那种辛苦。}
\BookParagraph{而是每逢节假日，他提上一个双肩包，坐上全中国最慢的绿皮火车，去看望十几个小时甚至几十个小时以外的你的那种辛苦。}
\BookParagraph{而我说的「浪漫」呢，当火车徐徐进站，他在车窗眺望而你在站台上等着他，你们终于互相发现了对方，不等下车，笑着遥相挥手示意的那一刻，也许你会感觉到。}
\BookParagraph{短暂的分别，本来就是一件非常浪漫的事儿。正好又能就读一所更好的学校，岂不是锦上添花，何乐而不为？}
\BookParagraph{你也许会问，读研要好几年时间啊，这可不是短暂的分别。这么好的男友，将来因为这个跟我分手了怎么办？}
\BookParagraph{这个问题我是这么理解的：如果计算「分手」的条件概率，你们一直在一起说不定比分开几年还要高一些。}
\BookParagraph{只是因为在车站每一次拉着旅行箱为他送别时的那份不舍，还有拎着一塑料袋零食终于等到火车进站时的那份开心，很多人能清晰地记一辈子。}
\BookDate{2018 年 6 月，西安}
\BookParagraph{（后记：去年我乘坐绿色的特快T109次火车从北京站出发历经十五个小时到达上海站的时候，发现似乎站台上已经没有人等着接人了，大概站台票早已停售。）}
\BookArticle{表　白}{表白}
\BookParagraph{李老师，您好！关注您的微博已经有些日子了，每次读您的文字正如大家评论的那样，散发着一种理性睿智的温暖感。}
\BookParagraph{想要请教您的是一个偏个人情感方面的问题。在异性交往这方面，我发现我有一个让其他人无法理解的观念：从心底里我会认为身边女生只会有女朋友、普通女生两类，当然同学、实验室伙伴那种特殊关系的除外。}
\BookParagraph{这个问题一直困扰着我。昨天刚和心仪的女生大大方方地当面表白了，但因为女生即将面临今后一年甚至几年的出国生活而被拒绝。回到寝室我半夜醒来，把所有和她相关的信息和联系方式都删了，想着今后应该也不会再联系，完全没有一点那种「买卖不成仁义在」的情谊。}
\BookParagraph{对此，我自己分析的原因有以下两点：首先我认为可能是我自己太功利了，导致和女生分手、表白失败，从此两人就天各一方了，无需再有任何联系；其次，我确实在男女交往这块儿缺少必要的认识，甚至是男女间普通的朋友关系的认识上也有问题。}
\BookParagraph{如果能有幸看到李老师的回复，不胜感激！谢谢您！（问题来自微博私信）}
\BookParagraph{咱们来换位思考一下：如果你是这个女生，你应该如何选择？}
\BookParagraph{一个男生冷不丁跟你说，他对你早已心仪已久，并且觉得你们两个人再合适不过了，愿意跟你发展更进一步的亲密关系。「合同」已经草拟好了，你是签还是不签？}
\BookParagraph{这有点儿像一桌酒席：男生的这杯酒已经高高地举起来了，还振振有词地说：来，给个面子吧，干了？}
\BookParagraph{你肯定会想，这男生以为他自己是谁啊？}
\BookParagraph{无论这男生以前是不是个好朋友，你的最优策略一定是：果断地、坚决地、不留任何余地地拒绝他。但最好能找一个相对客观的理由，让他觉得是因为不可控的客观因素。}
\BookParagraph{你瞧，你心仪的女生完全是按照这个策略来处理这件事儿的。即使换成我是这个女生，也一样会这么做。}
\BookParagraph{如果你的「表白」再高调一点儿，搞个什么公开场合九十九朵玫瑰啊拉个横幅啊什么的，相当于八抬大轿敲锣打鼓威逼人家签这份「合同」，那被拒绝的概率就更加接近百分之百了。}
\BookParagraph{马克思主义哲学里说：「事物的发展是螺旋式上升的过程」，「从量变到质变」。其实恋爱也是一样。}
\BookParagraph{「表白」有点像高手过招里亮出自己的底牌（play the trump card），风险极大，又完全没有任何好处。「质变」之前「表白」，注定会失败；「质变」之后「表白」，又显得画蛇添足，多此一举。}
\BookParagraph{这时你一定会问：我就是想确定一下我们正式男女朋友的关系啊，为什么不行？而且我们俩以前是很熟的朋友加同学啊，又不是萍水相逢。}
\BookParagraph{关键的问题是，为什么一定要强迫一个女生答应跟你确立正式男女朋友关系？确立与否，照样可以一起吃饭、喝酒、看电影、唱歌啊。}
\BookParagraph{与其去「表白」，不如拉上她一起到酒吧看世界杯，一起到上海法租界吃零食，一起看「创造101」，或者一起到黄埔江边练习人像摄影。}
\BookParagraph{当然如果你们能一起学习，一起准备英语考试，互相抽背单词，一起讨论一道悬而未决的疑难习题，那绝对是天下最浪漫最铭心刻骨的爱情。}
\BookParagraph{与其大半夜的删这个删那个，给她打个电话聊聊天吧，说不定还有希望。}
\BookParagraph{就是千万别再「表白」了。}
\BookDate{2018 年 7 月，西安}
\BookArticle{心有灵犀}{心有灵犀}
\BookParagraph{李老师，看了您关于党校的问答，真的忍不住想跟您说一下我的情况。我可能就是那种就地躺倒的人，也这么做了，每天乖乖巧巧按部就班，现在最大的困扰是找对象的问题。我的生活习惯和爱好都比较独，接触不到新朋友，介绍的又很难聊得来。想找志趣相投的，可我真的很像日本小说中独居的人，不喜欢吃不喜欢逛不喜欢凑热闹玩（可能也有经济原因）。如果您有兴趣的话，也开导开导我吧。（问题来自微博私信）}
\BookParagraph{找终身伴侣这件事儿，我比较相信高科技。}
\BookSubtitle{一}
\BookParagraph{斯皮尔伯格最近拍的一部叫《Ready Player One》的电影你看过吗？}
\BookParagraph{电影里的男主角 Wade 和女主角 Samantha 就是在一个叫做「OASIS」的虚拟现实游戏里相识的。游戏里 Wade 的好朋友对他的网恋不以为然，警告他：「You’ve got to be more careful about who you meet out on the OASIS.」他斩钉截铁地反驳：「She just gets me. We even finish each other’s sentences.」}
\BookParagraph{Finish each other’s sentences —— 两个人聊天儿不但不尴尬不紧张，还越聊越 high，完全停不下来，只嫌时间过得太快。聊到酣畅淋漓之处，如同一个人左右互搏一样。这已经不仅仅是「心有灵犀」了，简直就是爱情的最高境界啊。}
\BookParagraph{而如此境界，仅仅靠亲朋好友介绍相亲，估计多半儿体验不到。我一直觉得靠相亲来寻找终身伴侣效率相对比较低：本来天各一方的两个陌生人，突然坐在一张桌子上喝咖啡，这得多能说才能聊得起来啊。}
\BookParagraph{但是如果依靠高科技，在虚拟现实的游戏里，为了一个共同目标的队友，还真有可能萌生出一种「两情相悦、心意相通」的浪漫情怀。你想啊，你和一个从未谋面的女生边语音聊天边一起打「吃鸡」，跳伞和射击时听到的是她略带娇滴滴的笑声和欢呼声，会不会跟斯皮尔伯格的电影里一样，整个人彻底地爱上了她？}
\BookParagraph{没错儿，你其实爱上的不是一个人，而仅仅是一种声音。但是这种声音像Scarlett Johansson在电影《Her》里的声音一样，让人深陷其中，难以自拔。}
\BookSubtitle{二}
\BookParagraph{不仅仅是虚拟现实游戏，其实靠谱一点儿的婚恋网站和交友app，依赖的也是高科技的算法。}
\BookParagraph{普通的婚恋网站需要你去填写自己的兴趣爱好，然后用推荐算法来帮助你找和你最相似的陌生人。信息填得越详细，推荐给你的人就越准确越靠谱。}
\BookParagraph{比如说如果你喜欢黄章晋或者高晓松，算法就可能推荐那些和你一样热爱黄章晋或者高晓松的粉丝。}
\BookParagraph{如果你将来有办法把自己在知乎、微博或者朋友圈的所有内容和浏览记录都分享给这样的「人工智能」算法，我敢肯定，算法对你的了解会远远超过你的家人朋友——因为算法不会忘记任何历史信息，它甚至可能超过你自己。}
\BookParagraph{你想象一下，这样的算法推荐给你的人，成功率是否比亲朋好友随便介绍的相亲对象高很多？}
\BookParagraph{不喜欢下馆子不喜欢逛商店不喜欢凑热闹都没关系，但凡你有喜怒哀乐七情六欲，「人工智能算法」都能把你映射为一个高维空间中的向量，用一个向量之间的余弦相似度，就能在这个世界的某个角落里找到和你「心意相通」的那个人。}
\BookSubtitle{三}
\BookParagraph{可是在这个柴米油盐的世界上，真的存在这样一个人吗？}
\BookParagraph{即使真的存在，她真的还在等着你的出现吗？}
\BookParagraph{也许你踏破铁鞋最终找到的，是一个和你志趣大相径庭，天天有事儿没事儿心情不好就跟你吵上一架，磕磕绊绊地也和你过了大半辈子的人。}
\BookParagraph{不过这也没有关系，很多年以后孩子长大了，几个哥们儿叙完旧，英雄不减当年，再度开玩儿「绝地求生」的时候，你如果还能怔怔地记起当年一起打游戏时她的声音，就心满意足了。}
\BookParagraph{就像斯皮尔伯格那部电影里的台词说的：「As terrifying and painful as reality can be, it’s also the only place that you can get a decent meal. Because, reality is real.」}
\BookDate{2018 年 9 月，多伦多}
\BookArticle{一瞬间的爱情}{一瞬间的爱情}
\BookParagraph{李老师您好，您是我的偶像。我现在有个很好的女朋友，我很爱她，想和她走过一生。想听您聊聊爱情这事。（问题来自微博私信）}
\BookParagraph{爱是一瞬间的事儿。}
\BookParagraph{像《罗马假日》——在黑白的罗马自然而然地发生了一段朴素至极的爱情：「Rome, by all means, Rome. I will cherish my visit here, in memory, as long as I live. 」也像《致我们终将逝去的青春》——当青春飞扬逐渐挥霍一空的时候，余下的只有刻骨铭心的爱情：「我不哭，阮阮，我愿赌服输。」最像的还是——你听说过吗——张爱玲二十三岁那年写过一篇文字，极短、极简，题目就叫《爱》：}
\BookParagraph{——}
\BookParagraph{这是真的。}
\BookParagraph{有个村庄的小康之家的女孩子，生得美，有许多人来做媒，但都没有说成。那年她不过十五六岁罢，是春天的晚上，她立在后门口，手扶着桃树。她记得她穿的是一件月白的衫子。对门住的年轻人同她见过面，可是从来没有打过招呼的。他走了过来，离得不远，站定了，轻轻的说了一声：「噢，你也在这里吗？」她没有说什么，他也没有再说什么，站了一会，各自走开了。}
\BookParagraph{就这样就完了。}
\BookParagraph{后来这女人被亲眷拐了，卖到他乡外县去作妾，又几次三番地被转卖，经过无数的惊险的风波，老了的时候她还记得从前那一回事，常常说起，在那春天的晚上，在后门口的桃树下，那年轻人。}
\BookParagraph{于千万人之中遇见你所要遇见的人，于千万年之中，时间的无涯的荒野里，没有早一步，也没有晚一步，刚巧赶上了，那也没有别的话可说，惟有轻轻的问一声：「噢，你也在这里吗？」}
\BookParagraph{——}
\BookParagraph{愿你们将来的生活永远朴素、简单、圆满。}
\BookParagraph{中秋快乐。}
\BookDate{2018 年 9 月，多伦多}
\BookArticle{暴　爱}{暴爱}
\BookParagraph{你好李老师。}
\BookParagraph{我遇到了一个人。他身上有我喜欢的全部品质。而在此之前，我不知道那是我在伴侣身上需要的全部品质。}
\BookParagraph{我们之间的交集只有短短五天，这之中的每一天，和他的每一个互动，都让我快乐。而在此之前，我不知道我可以这样快乐。}
\BookParagraph{可是在我们第一次接吻第二天，他告诉我他不能和我在一起。他觉得 doomed。}
\BookParagraph{我当着他的面，念了我这五天的心路历程，念了我这五天写下的诗，哭了两声，自己下车走了。}
\BookParagraph{我不知道我接下来该怎么办。我指的不是要不要删掉他的联络方式，也不是要不要再争取一下，我觉得都没必要了。只是我的喜欢突如其来，又被迫胎死腹中，我被自己的感情吓到了，也被它被迅速扼杀的事实吓到了。}
\BookParagraph{我记得《欲望都市》里有一集，提出分开的作家忽然捧着花回来，和女主角一夜春宵后，留下纸条说自己做不到就走了。女主角愤怒地砸了桌上的花瓶。可我不愤怒，我现在无比孤独。我不知道原来，在我对一个人感觉这么好的时候，他却觉得没有结果。我想抱着我的妈妈哭一哭。}
\BookParagraph{想听您说点什么。}
\BookParagraph{（故事来自微博私信）}
\BookSubtitle{一}
\BookParagraph{一九一二年四月十日，泰坦尼克号首次从英国南安普敦启航，前往美国纽约。}
\BookParagraph{一九一二年四月十五日凌晨，在北大西洋与冰山相撞，一千五百多人遇难。}
\BookParagraph{仅仅五天时间，竟然演绎了一场头等舱的阔小姐和四等舱的穷画家之间惊心动魄的爱情故事。}
\BookParagraph{这五天时间里，和 Jack 的每一次互动，都让 Rose 觉得快乐。}
\BookParagraph{而在此之前，Rose 从来不知道她可以这样快乐。}
\BookParagraph{没错儿，你的故事让我迅速想到了二十年前那部家喻户晓的电影——《泰坦尼克号》。}
\BookSubtitle{二}
\BookParagraph{同样是二十年前，这部电影刚开始公映的时候，褚明宇曾经大笔一挥，写了一篇檄文发表在 MIT BBS 上，题目就叫《大骂Titanic》。这篇文章虽然短小，但极其精悍，全文毫不拖沓，行云流水，锋芒毕露。我先来摘抄一段：}
\BookParagraph{——}
\BookParagraph{《Titanic》给我的感觉只有一个字：俗。其实俗也不一定就是坏事儿，毛主席他老人家早就说过我们的文艺是要为工农兵服务的。就象当年流行过的《妈妈再爱我一次》还有《喜福会》之类的破烂儿，骗完了家庭妇女们的眼泪就可以老老实实的收摊儿走人了。不过俗得无与伦比的东西要是被捧得太高吹得太悬主题歌咿咿呀呀悲悲切切满大街的放，就显得有些讨厌，让人忍不住想站起来骂两句了。}
\BookParagraph{——}
\BookParagraph{你瞧，从旁观者的角度来看，任何闪电式的爱情，都有「这就俗了」的嫌疑。《Titanic》还算好的，至少是两情相悦，直到灾难来临，还立志携手共度难关。你的故事里明显是你很爱他，但是他却不怎么爱你；你觉得可以以身相许，他却不想私定终身。如果说《Titanic》俗，那从一开始就注定要短命的爱情，就更加不值一提了。}
\BookParagraph{但是我很羡慕你。}
\BookParagraph{我们几乎每个人都曾经「暴怒」过——就是气得七窍生烟、呼吸都费劲儿的那种感觉。}
\BookParagraph{但是我们不是每个人都「暴爱」过——就是爱一个人爱到昏天黑地，像一个刚打好了氢气的气球，越飞越高，飞得快的令人动容，忽然一声巨响，一切戛然而止。}
\BookParagraph{这种欲哭无泪、心如刀割的心情，不是每个人一辈子都能体验一次的。}
\BookParagraph{你的运气简直好得让人难以置信，好就好在他五天就把你甩了。这过山车一般的几天时间太短，就像热水里刚出来去冰水里淋浴，从「I’m the king of the world」堕入万丈深渊，把你所有的痛苦体验加倍放大，很少有人能亲身经历。}
\BookParagraph{而闪电式的爱情，多半儿早晚是要以分道扬镳收场的。即使美丽得像张雨绮一样倾国倾城，几任男友，闪电恋爱闪电结婚，双胞胎都生了，不是一样要离婚？所以他提出「咱们没有将来」越早越好，「Jack」沉入大西洋以后，把你的微信删了，你早晚可以跟我们的阔小姐一样——「my heart will go on」。}
\BookParagraph{如此完美的「暴爱」，能不让人心驰神往吗？}
\BookSubtitle{三}
\BookParagraph{二十年前在美国的一个周末，我去看了一场叫《泰坦尼克号》的特价电影。可能是想再看一遍，刚好褚明宇也去了。}
\BookParagraph{电影的结尾有一个细节。老太太Rose在船上安详地睡着了，房间里摆放着她年轻时在沙滩上骑马的照片，正是 Jack 说要带她去玩儿的。画面转到一个梦幻般的场面，破旧的沉船慢慢地变成了当年的泰坦尼克，Rose一袭白衣，美丽得令人窒息，从豪华的旋转楼梯拾级而上，Jack西装笔挺，微笑地在上面等着她。两人终于接吻时，泰坦尼克上的所有人全体起立，热烈鼓掌。}
\BookParagraph{褚明宇后来说：这个电影真的挺好看的。}
\BookDate{2018 年 9 月，多伦多}
\BookArticle{三枚金针}{三枚金针}
\BookParagraph{李老师，您好，不知道您能不能看见我的私信，想请教您一个问题，就是一个平凡的人会遇见的情情爱爱的问题。}
\BookParagraph{喜欢了一个很特别的男孩子，留着飘逸的长发，一双眼睛很忧郁，像个流浪诗人，身上也有一种说不出来的气质，非常特立独行的一个人，也是一个很自我的人。分手以后一直走不出来，一直沉溺在这种喜欢里，因为觉得在大街上再也找不到第二个这样的人了。现在的他，在我眼里实在是特别得无可取代了。可能因为年轻，喜欢这样的「怪人」，喜欢他的不合群，甚至想保护他的自我和他尖锐的性格。我不知道怎么办，真的很难走出这种喜欢。不是失恋的痛苦，真的是发自内心的太欣赏这个人了，很害怕自己再也遇不见这样的人。}
\BookParagraph{我想知道您怎么看待年轻女孩欣赏怪人的美这件事。那种欣赏常常掺杂着模糊不清的崇拜、向往甚至是保护欲。会在与之交往的过程中，自觉扮演起相对适应规则的一方，好去衬托他们对于社会关系的迟钝和漠然。也会因为怪人们每每释放出的「目中无人」，而高估他们的才华。在被他们抛出的谜面彻底击溃后，还要反思是否自己的理解能力过于低下。您觉得这种崇拜是一时的，还是会随着成长改变呢？（问题来自微博私信）}
\BookSubtitle{一}
\BookParagraph{哈哈哈，你大概忘了，我的专业是理工科，不是恋爱心理学。}
\BookParagraph{不过即使是我这个外行，都可以斩钉截铁地告诉你，确实如你所说，以后你在「大街上」也好，小胡同里也罢，再也不会遇到他这样的人了。}
\BookParagraph{这就好像我多年前第一次看《围城》，自从那个大雨的下午方先生从唐小姐家里离开之后，一直热烈地期盼着唐晓芙再次出现在小说中。}
\BookParagraph{这让我看《围城》后一半儿时，有一种屏住呼吸，一直期待着小学生作文中的首尾呼应、交响乐里的华彩的感觉。}
\BookParagraph{可惜直到那只祖传的老钟从容自在地打起来，唐小姐再也没有出现。}
\BookParagraph{这种莫名的失落感，也许正是《围城》这本小说的精华所在，正是钱老先生想讲给我们听的人生道理。}
\BookSubtitle{二}
\BookParagraph{不过我看完你的问题，心里倒是隐隐有一种穿越到了金庸小说里的感觉。而此时此刻正跟我聊天的，是《神雕侠侣》里我最喜欢的女孩儿——郭襄。}
\BookParagraph{郭襄在黄河北岸风陵渡口初遇神雕侠，自此无论是三枚金针，还是生日的三件大礼，直至绝情谷纵身一跳，十六岁的郭襄对她的「大哥哥」生死相依，情深至此，读来荡气回肠。}
\BookParagraph{「在绝情谷底，郭襄从怀中取出最后一枚金针，说道：“大哥哥，当日你给了我三枚金针，曾说过凭着每一枚金针，我可相求一事，你无不允。今日我来求恳：不论杨大嫂是否能和你相会，你千万不可自寻短见。”」}
\BookParagraph{我如果也有金针，一定会跟郭襄说：神雕侠确实独一无二，世上只有一个。}
\BookParagraph{不过人生苦短，甜点有很多种，各有各的甜法儿，为什么一定要苦苦寻觅以后不会再出现的那一种呢？}
\BookParagraph{不论以后是否还会再遇到这样的人，不论唐小姐是否还会再出现，千万要多尝试几种不同的甜点，才能不枉此生。}
\BookSubtitle{三}
\BookParagraph{我最爱的英文电影之一 —— 其实是挺小众的一电影 —— 叫「Eyes Wide Shut」。}
\BookParagraph{虽然这部电影没什么人看过，它的导演却是大大有名：他的名字叫 Stanley Kubrick。}
\BookParagraph{Kubrick 一辈子导演过好多部超级牛的电影，每一部都足以让一位导演名垂青史：「The Shining」、「A Clockwork Orange」，还有最著名的「2001: A Space Odyssey」。}
\BookParagraph{「Eyes Wide Shut」是他的最后一部作品。他对拍摄的细节极其较真儿，一部电影竟然拍了将近两年时间。电影的主题你说不定感兴趣：它是一个关于婚姻和出轨的故事。}
\BookParagraph{这部电影后期剪辑刚刚结束没几天，Kubrick 心脏病突发去世。}
\BookParagraph{Tom Cruise 和 Nicole Kidman 是两位我喜爱的明星，联袂主演了这部电影以后，很快就离婚了。}
\BookParagraph{电影里我最喜爱的一句台词是这么说的：}
\BookParagraph{Life goes on. It always does, until it doesn’t. But you know that, don’t you?}
\BookDate{2018 年 10 月，多伦多}
\BookArticle{处女情结}{处女情结}
\BookParagraph{李老师您好，我想请教一下处女情结的事。}
\BookParagraph{大学曾经谈过一个女朋友，差不多交往了一年的时候，我知道她不是处女了（我和她至今没有发生关系）。我的天塌下来了，我开始对她冷淡，发脾气，我开始心里恨她，恨上帝，心里想的是将来要把那个男的阉割掉。我的心理产生了严重的矛盾，深爱她又接受不了她，整日处在抑郁当中，天天失眠被吓醒，天天在纠结到底爱不爱她。甚至一度有了轻生的念头，一下爆瘦20斤，真的走到了人生的尽头。}
\BookParagraph{终于她也看出我的状态。为了不伤害她，我告诉她我抑郁了，她每天给我开导，逗我开心说会走出来的。我也看了很多这方面的书籍，我决定坚持我的爱，我要把精力放在学习上，我要考研，每天早出晚归。渐渐的我看到了希望，我突然不再想这件事，我好像更爱她了，每天学习到十一点也要给她打电话聊天，不时会看看她的照片。}
\BookParagraph{可是考研结果失败了，工作也没有着落，整个人垮了。所有的事接踵而来，朋友也伤害了我，这次彻底完蛋了，我又恢复到先前，不是冷淡，就是发脾气。最后她提出分手了，我也没有过多解释，我们就这样分开了。几次疼痛难忍，伤心绝望，我都没有联系她，我不确定我能有勇气坚持下去。}
\BookParagraph{故事还没有完。现在我研究生也毕业了，工作还不错，终于几年之后我又碰到了一个我喜欢的人。我追了她很久，就在快要确定关系的时候，她告诉我她谈了一个三年多的男朋友（之前她说没有谈过），我的天又一次塌了下来。开始追她之前我就想过她如果不是处女我是可以接受的，可是到了现在又活不成了。}
\BookParagraph{我真的很讨厌这个情结，我也不知道为啥我这么强烈。两件事让我彻底对爱情绝望了，我想了很多很多，是不是我的占有欲，控制欲，还是我变态？如果我不爱她们，为什么会如此绝望？如果爱她们，为什么还接受不了？李老师，希望您能聊聊这个事，我不怕观点有多犀利，我想勇敢面对，谢谢。（问题来自微博私信）}
\BookSubtitle{一}
\BookParagraph{好好认清形势吧，同学——你现在只有两个选择。}
\BookParagraph{要不然你单身一辈子。这就意味着将来你老了，很可能没有人陪你聊天，没有人陪你吃饭，没有人为你掉眼泪。你没有子女，没有亲人，独守空房，房子里只有冷冰冰的家具和电器。你玩儿着手机，潜意识里试图在一个叫 Soul 的 app 里找陌生人说话。}
\BookParagraph{要不然你全心全意地爱她，将来说不定可以一起过日子，一起抚养你们的孩子。你老了退休了赋闲在家，可以和外孙子外孙女闲聊，万圣节可以带着孩子们出门要糖吃，和她们一起大喊「Trick or \BookLineBreak{}treat!」，听她们弹琴、唱歌，给她们讲你当年热爱的金庸小说，还有你当年如何追到她们的姥姥的故事。}
\BookParagraph{现在就决定吧，我这里在线等着你。}
\BookSubtitle{二}
\BookParagraph{你的女朋友以前有过男朋友，不但不是坏事儿，还是一件值得庆贺的好事儿。不仅如此，她以前的男朋友越多，你就越值得开心。}
\BookParagraph{你想啊，如果她以前交往过不少男生，什么人都见识过了，结果现在竟然会爱上了你，这是不是至少说明你长得还不算难看，见识还不算浅薄，能力还不算平庸，品格还不算低俗，至少比起她以前的男朋友，还有可圈可点之处？}
\BookParagraph{人往高处走，水往低处流。恭喜你在竞争中脱颖而出，立于不败之地。}
\BookParagraph{那如果反过来，她以前从来没有过男朋友，你是她的初恋又如何呢？}
\BookParagraph{如果你们是青梅竹马、两小无猜时就相爱到现在，那可喜可贺，咱们可以拿这个素材写一篇感人的短篇小说。}
\BookParagraph{如果不是，那隐含的风险就有点儿像买一支 IPO 的股票：你买的是谷歌的话自然大赚，但买的如果是 Snapchat，很可能血本无归。}
\BookSubtitle{三}
\BookParagraph{Kubrick 的电影「Eyes Wide Shut」开场的圣诞酒会里，一个风度翩翩的陌生男士邀请 Alice 跳舞，边跳边聊：}
\BookParagraph{「You know why women used to get married, don’t you?」}
\BookParagraph{「Why don’t you tell me?」}
\BookParagraph{「It was the only way they could lose their virginity and be free to do what they wanted with other men. The ones they really wanted.」}
\BookParagraph{「Fascinating.」}
\BookSubtitle{四}
\BookParagraph{我姥爷和我姥姥刚结婚的时候，每天工作时都要抽烟。姥姥也不跟他理论，只是每买一包烟，就另把一份烟钱单独存起来，很快就积攒了很大一笔钱。}
\BookParagraph{姥爷看到了这笔钱，什么也没说，很快就戒了烟。按咱现在讲话，叫「秒懂」。}
\BookParagraph{姥爷和姥姥后来培育了五个子女，晚年儿孙满堂，其乐融融。}
\BookParagraph{作为一个男生，想要表达对你女朋友如火如荼的爱，最好的办法不是送给她九十九朵玫瑰表白，不是陪她到香港「血拼」，也不是带她去见你的爸爸妈妈，准备丰厚的彩礼。}
\BookParagraph{而是慢慢地彻底地改变自己。}
\BookDate{2018 年 11 月，多伦多}
\BookArticle{爱与不爱}{爱与不爱}
\BookParagraph{李老师您好，我想请教一下关于我喜欢他但他不喜欢我的事。}
\BookParagraph{平时和他聊天，是我主动开始的多。可能是我不够有自知之明，没有察觉到他有不想聊下去的意思，感觉良好。他跟我讲过他高中时他父母帮他申请晚上不去上晚自习，为了他能玩好休息好；大学时吃两份饭，担心别人很奇怪地看他，对面也放一个餐盘，假装对面还有一个人；讲他研究生在美国去很多地方玩的趣事；还有他的婚姻观，培养孩子的观念，等等。有一次还叫我去他家坐了一会儿。}
\BookParagraph{但是，在叫他出来玩了几次以后，跟他说了我爸爸去年生病去世的事情，那以后他就再也不愿意出来和我吃饭玩了，但还是会跟我在微信上聊很多。}
\BookParagraph{直到最近，我问他是不是讨厌我，为什么不愿意和我一起玩了，他说为了我好，要保持距离，他替我保持距离。我问他原因，他说对我聊天的绝大部分内容没有兴趣，不想浪费时间；说要是和他聊篮球聊游戏，他就很感兴趣（他知道我不玩游戏）。我就说，那你可以带我玩，或者我看也行，他就说要和我保持距离。}
\BookParagraph{平时聊天的感觉挺好，所以他说这些我挺惊讶的。我好不容易喜欢上一个人，对发生的这些事实在想不明白为什么，问他他也不是十分愿意说清楚。朋友说，其实就是他不喜欢我。我想可能也是这样，但我想他很明确地告诉我，他说话又总是很含糊很温柔。请问我应该怎么办？（问题来自微博私信）}
\BookSubtitle{一}
\BookParagraph{我真的希望能告诉你点儿振奋人心的好消息，比如你还有一线成功的希望，比如如果你真心对他好，认真学习篮球技术认真打游戏，他一定会义无反顾地爱上你。}
\BookParagraph{可惜我只能跟你说一个既成事实：你好不容易爱上的男生，他并不爱你，将来也不会爱上你的。}
\BookParagraph{注意，我说的是他不爱你，并没有说他对你反感。}
\BookParagraph{恰恰相反，多半儿他还挺喜欢你的某一方面的，比如喜欢你这个忠实的听众，喜欢跟你在微信上讲述他自己从前的故事。}
\BookParagraph{可喜欢你这一方面并不代表喜欢整个的你。}
\BookParagraph{你想啊，你对其他男生是不是也经常会喜欢他们身上某一两个细节——比如做什么事儿都靠得住，比如性格特别阳光——但是如果他们向你表白，你一定不会考虑？}
\BookParagraph{没错儿，喜欢一个人身上的一个细节，不等于你就要嫁给他。}
\BookParagraph{那你又要说了，他为什么不跟你明确说清楚啊？}
\BookParagraph{等等，我怎么觉得他跟你说得简直没法儿再明确再清楚了呢？}
\BookParagraph{他明明知道你不玩儿游戏，不懂篮球，还专门拿这两样儿举例；明明知道跟你出来玩儿才能进一步发展和你的感情，却坚决拒绝走出这一步；他都说了要「保持距离」了，你还要怎么更明确——「对不起，我不爱你，将来也永远不会爱你」？}
\BookParagraph{他在不伤感情的基础上对你已经仁至义尽了——这阵子估计心里正想：这姑娘怎么一意孤行、没完没了啊？}
\BookSubtitle{二}
\BookParagraph{要是说他「很含糊很温柔」，你大概没见过真正「很含糊很温柔」的男生。}
\BookParagraph{我在美国上学的时候有个小师妹，到美国的时候本科刚从国内顶尖名牌大学毕业，既冰雪聪明又如花似玉，按常理说，在美国应该有大把的男生追她才对。}
\BookParagraph{可惜的是，她刚到美国一下飞机，在机场就被一个其貌不扬又特别不爱说话的男生接走了，根本就没别人什么机会。}
\BookParagraph{这个男生是真的不爱说话。我们大家跟他聊天，他一般就是听着，最多也就回一两句，从来没见过他高谈阔论，讲过自己的什么故事。让人感觉有点儿像一个久经战阵的老兵，「欲说还休，却道天凉好个秋」，什么都经历过了，什么都不会再新奇。}
\BookParagraph{那时中国学生会每年活动的重中之重，就是去机场接新生。当时他在机场第一次见到小师妹时是怎样的喜出望外，怎样「蓦然回首，那人却在灯火阑珊处」，就不得而知了。大概跟《致青春》里接新生接到郑微的老张心情差不多吧。}
\BookParagraph{和郑微不尽相同的是，小师妹很快就和我们这位不爱说话的幸运儿恋爱了。几年以后他们在学校附近的一个风景秀丽的公园结婚，婚礼大宴宾客，还请了我去参加。}
\BookParagraph{两人机场初次相识之后的第十个年头，在纽约离婚了。据说是男生婚外爱上了别人。}
\BookParagraph{不过我们的小师妹也不是吃素的。离婚后立即开展一项庞大的工程，寻找下一位终生伴侣。连Facebook上的状态都写着：looking for a relationship.}
\BookParagraph{功夫不负有心人，一年后她就找到了如意郎君。}
\BookSubtitle{三}
\BookParagraph{「请问我该怎么办？」}
\BookParagraph{你应该向小师妹学习她百折不挠、全力以赴「looking for a relationship」的精神，迅速组织力量，「放长线，钓大鱼」，把找男朋友作为一个重大研发项目来贯彻实行，争取在最短的时间内找到你最爱的人。}
\BookParagraph{但是不要等到十年以后才明白一个道理：}
\BookParagraph{你爱的人也许并不爱你。}
\BookDate{2018 年 11 月，多伦多}
\BookArticle{多线程}{多线程}
\BookParagraph{李老师您好！我想跟你请教有关前男友的问题。}
\BookParagraph{我和他分手是高考完第二天，他对我冷暴力，逼我说出「分手」二字（我觉得男生很擅长这样）。分手后一周后他有了新女朋友，那个女孩我也认识，我很伤心。大学我和前男友在一所学校。他和那个女孩在一起的时候私底下还和我保持联系，原因是因为不能彻底忘掉，还是会想起我。我当时还很爱他，他偶尔会来找我我从不拒绝。我们也一致对外保密保持着联系，但一直是他找我，因为他有对象所以我没主动找过他。}
\BookParagraph{后来他分手了，分手后喝醉酒了来找我倾诉，说没珍惜我等等。就在我还以为我们会复合的时候，他在大学交了个新的女朋友。事情发生的很突然，后来我就劝自己放弃。他也消失很久了没有联系过我。我和他就这样拉拉扯扯到大三。}
\BookParagraph{他最近又找我了，说梦到我了。他来找我我还是没拒绝。我清楚他有女朋友，从他说的话里感觉他对自己的现任女友不是很走心。他还叫我宝贝。原本我以为自己已经彻底放下了，但现在又开始在意他了，期待他的早安和晚安，也还是和以前一样，他不找我我绝对不找他。我该怎么办？我感觉他心里有我，这三年我一直没开始新的恋情，也是因为对他抱有一丝希望吧。他一找我，我就开始幻想各种关于他的未来。真的不知道该怎么办了。朋友劝我放手，也都不耐烦了。我表面上表现出对他不在意，但我心里明白自己的想法。}
\BookParagraph{老师我就想问问你，我有必要再为他在心里留位置吗？谢谢老师，麻烦您了！（问题来自微博私信）}
\BookSubtitle{一}
\BookParagraph{你在写这个问题的时候，心里一定门儿清。}
\BookParagraph{你清楚地认识到他和你的距离越来越远，认识到你们成功的机会越来越渺茫，认识到你早就应该放弃他。}
\BookParagraph{但是微信里收到那句「早安」的瞬间，你还是身不由己，想着「反正他不找我，我绝对不会去主动找他的」，然后秒回。}
\BookParagraph{这就好像一个痴迷炒股的股民，眼睁睁地看着手里的股票一年年地大跌，知道应该早点儿把股票清仓止损，可是却总抱有一线希望，想着「说不定已经跌到底了，很快就会大涨」，一直不卖，直至血本无归。}
\BookParagraph{所以说如果真要短线炒股，不能手工去操作交易，而是应该设计一个算法——也就是一个计算机程序——由算法去交易，而需要手工设计调整的，仅仅是算法的参数而已。}
\BookParagraph{人都有感情，有七情六欲，容易情不自禁、身不由己，所以炒股这种事儿，就应该采用没有感情的算法去做。将来自动驾驶算法成熟了，开车也可以尽量让算法来开，不至于把车开到桥下去。}
\BookParagraph{可惜的是，人工智能算法还不成熟，至少现在还没法儿在微信替你和男友聊天。男友的「早安」来了，还是忍不住想回，怎么办？}
\BookSubtitle{二}
\BookParagraph{想回就回呗。}
\BookParagraph{不过如果你学过一点计算机科学，你会发现里面经常会提到一个「多线程」的概念。}
\BookParagraph{有了「多线程」，一个计算机程序可以同时运行好几个算法，每个算法都以为自己拥有整个CPU。通俗一点儿说，就是一心二用甚至一心多用，好几件事儿同时进行。这就好像《射雕英雄传》里周伯通的双手互搏之术，心无杂念，分心二用，左手画方，右手画圆一样。}
\BookParagraph{你想回他的微信，觉得希望虽然渺茫但终究还有一线希望，行，你随时可以和他开聊。不过回的同时你还得再多开几个「线程」，同时跟其他几个男生聊天，每一个都聊得火热，让他们以为是唯一在和你聊天的人。}
\BookParagraph{在计算机科学里，这叫做「虚拟机」，你虽然只有一个「物理机」，但是可以分身出多个「虚拟机」，同时和好几个男生「谈恋爱」。}
\BookParagraph{可这不是骗人吗？}
\BookParagraph{准确地说，不是和好几个男生「谈恋爱」，而是同时「date」好几个男生。「date」这个英文词儿的好处是没有排他性，你可以中午 date 这个男生，晚上去 date 另一个男生，大大提高「谈恋爱」的效率。}
\BookParagraph{而提高计算机程序执行的效率，正是当年发明「多线程」的初衷。}
\BookSubtitle{三}
\BookParagraph{英文里分得很透彻：「boyfriend」你可以「date」很多，但是「fiancée」你只能有一个。}
\BookParagraph{而什么时候一个男生就从「boyfriend」升级到「fiancée」了呢？}
\BookParagraph{在他实在受不了你的「多线程」你的「虚拟机」，于是在一个灯光昏暗的西餐厅，酒过三巡，忽然对你单膝跪下，拿着三个月工资买来的钻戒，对你说「Will you marry me?」的那一刻。}
\BookParagraph{「梦到你了」、「早安」、「晚安」？}
\BookParagraph{这些统统不算。}
\BookDate{2018  年  11  月，多伦多}
\BookArticle{男生的品格}{男生的品格}
\BookParagraph{李老师，您好，我想请教您一个问题，关于男生应该具备怎样的品格。}
\BookParagraph{这段时间一直在和对象闹矛盾，经历了吵架到后面吵都不想吵了，直接冷战的局面。后面双方都爆发了一次，我竟然发现我有点儿不认识她了，有点儿陌生。她告诉我，有时候觉得我是一个自负、并且气量小的人。我想了想，她说的没错，我对于某些事会认为很简单，不值得一提，那些事我没有去做过，就直接认为很简单。我活在自己营造的主观世界里，气量小也让我很困惑，有时候我和她沟通时常会被她呛到，气得说不出话。我只是想让她知道一下我那个时候的心情，我想我那样做肯定是错误的吧，但我不知道别的情侣是怎么去处理的。}
\BookParagraph{前几天我们分手了，因为处理不好两人沟通的方式。我们在一起也不算短，算算有两年啦。可是那种堆积起来的无力感，真的没有办法呀，明明双方都是爱彼此的。也希望能通过李老师感知一下其他情侣吧。想问一下李老师，男孩子该具备哪些品格呢？（问题来自微博私信）}
\BookSubtitle{一}
\BookParagraph{你听说过博弈论（game theory）里有一个经典问题叫「囚徒困境」（Prisoner’s Dilemma）吗？}
\BookParagraph{警察逮捕了两名嫌疑犯，但没有足够证据指控两人有罪。于是警察给两个人提供了相同的选择：认罪或者保持沉默。如果两人里只有一人认罪，认罪的人无罪释放，另一人获刑十年。如果两人都认罪，两人各获刑五年。如果两人都保持沉默，则两人各获刑一年。}
\BookParagraph{哈哈，估计五十年代把这个故事写到论文里的那哥们儿警匪片儿看多了——里面动不动就有逮捕嫌犯的场景：「You have the right to remain silent. Anything you say can and will be used against you in a court of law.」}
\BookParagraph{在囚徒困境的故事里，如果两名嫌犯都自私地为自己争取最大利益（刑期尽量短），他们都会认罪。即使两人都明明知道如果两人合作双双保持沉默的话，刑期只有一年——两人还是都会认罪。}
\BookParagraph{永远为自己争取最大利益得到的均衡点，叫纳什均衡（Nash Equilibrium）。这个博弈论的故事告诉我们，从全局来看，纳什均衡往往并不一定是最好的结果。}
\BookParagraph{如果你和你女朋友出现矛盾，吵架升级到了冷战的局面，你是认错道歉呢，还是保持沉默僵持不下？}
\BookParagraph{如果两人都认错，结果是和好如初；两人都保持沉默呢，结果是分手。但是如果一个人认错，另外一个人保持沉默呢，结果是认错的这方丢尽了脸，还不如分手的好。}
\BookParagraph{按照博弈论的理论，如果你自私地为自己争取最大利益，你应该保持沉默，让冷战继续僵持下去。}
\BookParagraph{可惜的是，你们俩都太聪明了，明明知道合作就可以和好如初，最后还是要分手。}
\BookParagraph{而证明任何博弈游戏里都存在纳什均衡的 John Nash，却依靠他的天才证明得了诺贝尔经济学奖。}
\BookSubtitle{二}
\BookParagraph{你一定要问，那也不能一直忍让啊，时间长了，女朋友岂不是更要蹬鼻子上脸，欺人太甚了？}
\BookParagraph{好吧，现在我们又要用到一个理论。不过这回咱不玩儿博弈游戏，而是去解一个优化问题。}
\BookParagraph{计算一个优化问题的最优解，乍一看似乎很简单——只需要每一步都去挑选这一步里所有选择最好的一个，慢慢地不就逼近全局最优了吗？}
\BookParagraph{其实不然：这样的算法叫做贪心（greedy）算法，大多数情况下，贪心算法会让自己陷入「局部最优解」中。就好像你去爬山，如果只上山而不下山，就有可能陷在一座小山的山顶上，雾气缭绕之中，看不见远处还有更高的山峰在等着你。}
\BookParagraph{那怎么办？五十年代初有一个叫 Nicholas Metropolis 的物理学家写了篇论文提出了一个算法：你还是像以前那样爬山，每一步如果是上山就迈出这一步，如果是下山呢，那你根据一个概率随机抽一个奖，如果抽中就迈出这一步。这样我们利用一个随机算法，就可以离开小山头，向大山头进军了。}
\BookParagraph{这个算法后来以他命名，属于一种蒙特卡洛算法（Monte Carlo method）。蒙特卡洛算法的精髓就在于它的随机性，后来广泛应用于很多科学研究领域。比如谷歌鼎鼎大名的AlphaGo，所有人都知道是基于机器学习，但其实它的算法的核心是蒙特卡洛树搜索（Monte Carlo Tree Search），同样是利用算法的随机性去试探搜索树中的一些未知节点。}
\BookParagraph{回到你的问题上——你的最优策略是，在明知吃亏时掷出一枚硬币，如果正面朝上，即使吃点亏也认了。}
\BookParagraph{你瞧，我完全不认识你，还在这里一个字一个字地回答你的问题，不正是「吃点亏也认了」嘛。}
\BookSubtitle{三}
\BookParagraph{那男生应该具备什么样的品格呢？}
\BookParagraph{哈哈，光《小学生守则》就有十条，可以列举的优秀品格太多了。}
\BookParagraph{不过我印象最深的还是最后一条：「诚实勇敢，不说谎话，有错就改。」}
\BookParagraph{勇敢地跟你的女朋友认个错吧，说不定还能和好如初。}
\BookParagraph{当然如果能再懂一点儿计算机算法，就更完美啦。}
\BookDate{2018 年 11 月，多伦多}
\BookArticle{最优停止}{最优停止}
\BookParagraph{上上周六，我给喜欢的师兄发微信告白了，但师兄最近有女朋友了，拒绝了我。我后半夜发的消息，他早上回复的我，并且说要和我聊聊。但我当时被拒绝感觉很丢脸，和他见面也很尴尬，就没有去。我当时想着就是告白失败就放弃了。}
\BookParagraph{上周还好，这周不知道怎么回事，每当下午四点多以后，看他就离开办公室，我猜是和女友约会去了，我就会难过。真的，很难过，从我表白之后，我们再也没有说过话。}
\BookParagraph{我俩一个办公室，座位挨着，中间一块挡板，其实我难过的也有一点是觉得他根本一点都不关心我的感受。告白那天我还看他给另外一个小师妹的朋友圈评论了，最近有时候他在办公室也会很开心的和别人聊天什么的。我因为觉得尴尬，现在在办公室他参与的聊天我都不会参加了，他是不是一点都不关心我的感受呢？}
\BookParagraph{我其实一直想要表现得落落大方的样子，和他聊一次。我也不想在毕业前和他一直不再说话，毕竟是我喜欢的男生，我希望我们还能像从前那样聊天什么的。可每次一想约他聊天，就想人家会不会要去约会，没时间呢？那多尴尬，我就说不出口了。我也讨厌自己为什么别人不喜欢我我还总挂念着他，甚至会想万一他们将来分手了……李老师，我脑子挺乱的，该怎么办？（问题来自微博私信）}
\BookSubtitle{一}
\BookParagraph{你有没有想过一个问题：如果你三十岁前命中注定会遇到三个男生，希望能和其中你最爱的最优秀的那个结婚，那你在遇到第一个男生时应该怎么做呢？}
\BookParagraph{是直接表白，私定终身呢——还是守株待兔，静等生命中你下一个爱上的男生出现，看看是不是比第一个更加出色？}
\BookParagraph{正确的最优方案是：不要跟遇到的第一个男生表白，而是静等第二个男生出现。第二个男生如果比第一个男生优秀，就直接表白和他结婚；如果不如第一个，就等到第三个男生出现跟他表白结婚。}
\BookParagraph{在全部六种的出现顺序里，用这个方法，可以有三种情况——50\%的概率——能和最优秀的那个男生结婚。}
\BookParagraph{如果有四个男生，最优方案还是略过第一个男生，从第二个男生开始，一旦发现比第一个更好，直接表白。这时你有46\%的概率找到那个最优秀的男生。}
\BookParagraph{如果有五个男生，最优解是略过前两个男生，从第三个开始，一旦发现比前两个更好，直接表白。这时你有43\%的概率找到那个最好的男生。}
\BookParagraph{你发现没有，这个算法按顺序分成两个步骤。第一步只去发现和收集有关男生的「数据」，而不去表白；第二步和第一步发现的最优秀的那个男生比较，一旦发现一个更好的，马上表白。}
\BookParagraph{即使你遇到的男生很多，一样有最优的算法：第一步这个收集「数据」的阶段应该是所有男生总数的百分之三十七。}
\BookParagraph{这么说吧，如果你从二十岁开始你的「数据收集」阶段，三十岁之前必须确定婚姻人选，那你在二十三岁（零八个月）之前，应该对男生持「收集数据」的观望态度；但一过二十三岁，就随时准备表白结婚。用这个算法，你找到最佳选择的概率，大约也是37\%。}
\BookParagraph{这就是著名的「最优停止」（optimal stopping）理论。}
\BookParagraph{你瞧，向你一辈子喜欢上的第一个男生表白，很有可能是一个错误的决策。要是以后遇到更好的呢？}
\BookParagraph{所以说啊，他把你婉言拒绝了，避免了你因为一个错误的选择而和最优的结果失之交臂，你应该感谢他才对啊。}
\BookSubtitle{二}
\BookParagraph{你一定要说，可是表白不等于就一定能成功啊！}
\BookParagraph{你说得没错儿，咱们前面的理论做了一个不太现实的假设，表白的确不能保证结婚。说到底，不是人人都是刘雯刘亦菲。}
\BookParagraph{那如果表白的成功率只有百分之五十怎么办呢？}
\BookParagraph{答案很简单：提前结束对男生的数据收集工作，进入算法的\BookLineBreak{}「表白」阶段——从37\%提前到25\%。用这个最优算法，你和最优秀的男生结婚的概率，也是25\%。}
\BookParagraph{那如果以前收集过「数据」的男生，你发现比后来的男生都优秀，希望吃「回头草」去跟他表白是否可以呢？}
\BookParagraph{咱们假设你美得倾国倾城，他恰巧还没找到女朋友而且接受了你的表白的概率是50\%，那么你的最优策略是推迟进入「表白」阶段：从37\%推迟到61\%。这时你找到最佳选择的概率恰好还是61\%。}
\BookParagraph{哈哈，下回如果父母再催你找男朋友，你就说现在正处于「最优停止理论」里的「观望」阶段。}
\BookSubtitle{三}
\BookParagraph{好吧，咱们不聊算法了——你太喜欢这个男生了，希望不惜一切代价追到他，「optimal stopping」be damned.}
\BookParagraph{你虽然平时经常和他聊天，但是表白却以失败告终，主要是方法不得当：你给他发微信，还留给他大量思考如何回复的时间，这大大降低了成功的概率。}
\BookParagraph{表白最好的方式，是语言上含糊，行动上明确。}
\BookParagraph{比如聊得正high的时候，讲一个事先准备好的故事，引入入胜，风情万种。你们一起开怀大笑，然后你挽住他，慢慢而随意地靠在他的肩上。}
\BookParagraph{也许——仅仅是也许啊——他说有女朋友只是一个拒绝你的借口。也许你如果有足够的勇气，还有希望。}
\BookParagraph{人生如戏。就当是练兵场上，为命中注定的那下一个男生而\BookLineBreak{}「实战演习」吧。}
\BookDate{2018 年 11 月，多伦多}
\BookArticle{一厢情愿}{一厢情愿}
\BookParagraph{李老师，不知道我的私信您看不看得到。女生暗恋该表白吗？我暗恋上了一个不常见面的同事（一个单位，不一个部门，单位太大，不常见面）。注意到他之后，就开始觉得自己好卑微，好自卑，觉得自己哪儿哪儿都不好。每天都期待在单位食堂看见他，所以每天都把自己打扮的美美的，但是看见他的时候又怕他看见我，赶紧低下头或者把头扭过去……我都开始讨厌现在的自己了。已经半年了，偶尔跟他微信聊一下，因为采访过他一次。我办公室的同事想撮合我俩的时候，我却极力掩饰自己对他的好感，不让同事看出来，但过后又后悔，我是不是有病……这半年我上百度上知乎找女生暗恋男生的经验，想尽一切办法想知道自己下一步该咋办……未果。我该怎么办？（问题来自微博私信）}
\BookSubtitle{一}
\BookParagraph{不知道你发现了没有，英文里竟然没有跟「表白」对应得贴切的词儿。}
\BookParagraph{「Propose」明显不对：表白是想说「我喜欢你」，不是想说\BookLineBreak{}「嫁给我成不成？」}
\BookParagraph{「Confess」有点儿「承认错误」的意思——但爱一个人怎么能说是错误呢？}
\BookParagraph{这么普通这么简单的词儿竟然没有对应的英文，不符合英文的惯例啊。一般来说，同一种意思英文的说法比中文只多不少。比如「快乐」，英文里就有「joyful, cheerful, merry, jovial, jolly, gleeful, delighted, grinning, radiant, blithe, joyous, exuberant, elated, ecstatic, blissful」等数十种说法。中文估计把成语全算上也没有这么多。}
\BookParagraph{这只能说明一个问题：西方文化里不怎么「表白」—— 因为\BookLineBreak{}「表白」被拒绝的概率实在太高，风险实在太大，所以完全没有必要。}
\BookParagraph{注意，「表白」这个词儿虽然没有，「暗恋」还是有的。普通的说法是「I had a crush on him」，正式一点儿的说法叫「infatuation」。}
\BookParagraph{可越是暗恋，越不能表白——因为暗恋时表白就像买彩票一样，不是不可能中奖，只是中奖极其罕见。}
\BookParagraph{你换位思考一下：假如一个跟你只有几面之缘、微信里聊过几句的同事，突然对你说：「你知道吗？我一直喜欢你。」你是多半儿答应他呢，还是多半儿被他吓了一大跳，下意识地拒绝他？（当然如果是易烊千玺另当别论。）}
\BookParagraph{不过你表白不成功说不定也不是一件坏事儿。你可以从暗恋迅速而直接地过渡到失恋，为将来找下一个男朋友铺路搭桥。从这个角度来看，与其花半年时间去暗恋，还真不如刚开始就直接「表白」，让自己彻底失望，心灰意冷。}
\BookParagraph{爱情其实就像一个电子游戏：游戏里你有好几条命，每玩儿死一次，只不过是少了一条命而已，分分钟又满血复活，东山再起。}
\BookParagraph{而你的游戏，现在才刚刚开始，后面还有好多好多条命在等着你呢。}
\BookSubtitle{二}
\BookParagraph{不过在还没有一败涂地之时，我们都曾一厢情愿地幻想过爱情。}
\BookParagraph{如果你特别希望一直暗恋的男生也喜欢自己，应该怎么办？}
\BookParagraph{答案很简单：与其去表白，不如去培养自己和他的共同语言。}
\BookParagraph{比如如果他是一个足球迷，那你就应该花一些时间，仔细研究一下英甲、德甲、意甲和欧洲杯近十年来著名的球星和比赛，研究一下足球比赛中的基本战术和打法，不能浅尝辄止，而是深入到随时可以出口成章的地步。}
\BookParagraph{你暗恋着的男生时间有限，真正感兴趣又花了大量时间精心研究过的东西多半儿并不多，你如果真能静下心来速成，聊天时不经意间提起，和他的共同语言可以迅速爆表。}
\BookParagraph{很多人都憧憬能有那种心意相通、心心相印的爱情，他说不定也不例外。}
\BookParagraph{你所有的努力，就像足球比赛里在对方半场里各种密集的传球配合，希望能找到一线可乘之机，撕开对方防线，杀入对方禁区。}
\BookSubtitle{三}
\BookParagraph{那如果你和你喜欢的男生共同语言早已「爆表」，他却似乎一直没有明白你的心意，怎么办——是不是最终还是要表白？}
\BookParagraph{如果他近在眼前，你可以握住他的手，但什么也不要说。}
\BookParagraph{如果他远在天边，你可以手写一封长长的信。信里絮絮叨叨地聊着很多家长里短的闲事，像写给奶奶的家信，而一点儿也不像情书。}
\BookParagraph{唯一和家信不同的是，信封背后的封口上有一张贴纸。贴纸上印着「I like you」。}
\BookParagraph{而在「I」和「like」之间，发信人用红色的水笔添上了\BookLineBreak{}「really, really」两个词儿。}
\BookDate{2018 年 11 月，多伦多}
\BookArticle{生成对抗}{生成对抗}
\BookParagraph{李老师你好，看了您解答的很多问题，被您的智慧和温情深深折服。我也想请教李老师个感情问题。也算不上感情问题吧，但最近为这事儿倒是挺烦心的。最近在图书馆学习，对面坐了个很漂亮的女生，由于都是在考研，位置也相对固定，每天她几乎都坐在我对面。所以每次我都会不由自主的看她，但是我经常发现她也在看我（哈哈，自夸一句，我颜值还不算差），每天都会有那么几次目光撞在一起，然后迅速移开，当做什么都没发生。我真的特别特别想认识她，只是怎样才能和一个陌生人认识呢？怎样迈出第一步呢？我该做什么说什么呢？像我这样特别想认识别人但又仿佛有人际交往障碍的人该怎么办呢？（问题来自微博私信）}
\BookSubtitle{一}
\BookParagraph{我有一个坏消息和一个好消息要告诉你。}
\BookParagraph{先说坏消息吧——你能成功地和她认识直至成为朋友的概率很低，所以你要做好彻底失败的心理准备。}
\BookParagraph{注意，我说的是认识她直至成为朋友，不是仅仅能和她说上几句话，或加上她的微信。}
\BookParagraph{我甚至反对动不动就来一句「咱们加个微信吧」这种做法。且不说人家不一定愿意加你，就算加了你，几天后被她删了的可能性也不小——因为你一定会不等下次见面就主动用微信联系她，一言不合的风险很大。不加微信，反而显得你超然大度，与众不同。}
\BookParagraph{你一定会问：等等，为什么会失败——不是有专门讲搭讪心理学的书吗？}
\BookParagraph{书确实是有，但是你觉得看弹钢琴的书能学会钢琴演奏，看炒股的书能避免炒成股东吗？理论永远是理论，咱改革开放的总设计师当年说过，「黑猫白猫抓到耗子的才是好猫」——实践永远是检验真理的唯一标准。}
\BookParagraph{那如果心动女生总是会时不时地看你，是不是希望大大增加呢？}
\BookParagraph{实际情况是，她看你不一定是因为你长得高富帅特有吸引力，而很有可能是因为其他原因——比如发现你总是看她觉得很烦人。}
\BookParagraph{即使她确实是因为你高富帅而注意你，除了这一个事实之外，她对你来说完全是个陌生人。你不了解她就无法投其所好，以她的兴趣爱好展开话题，很难成功地认识她并和她成为朋友。}
\BookParagraph{但是你又特别特别想认识她，怎么办？}
\BookSubtitle{二}
\BookParagraph{别忘了，我这儿还有一个好消息要告诉你。}
\BookParagraph{虽然你成功的概率很低，但是确实有一个行之有效的方法可以提高这个概率。至于你是否愿意采纳我的建议，要看你对这件事儿有多重视。}
\BookParagraph{比如你毕业了要考研或者找工作，这些都是非常重要的事儿，你一定会花大量时间和精力准备。}
\BookParagraph{没错儿，只有当搭讪一位让你心动的女生这件事儿在你心目中也像考研找工作一样重要时，你才有可能抓住这个百年不遇的历史机遇，一战成名，成为一个幸运儿。}
\BookParagraph{那怎么准备呢？当然不是去看书，而是在不断的实战演习中提高你的水平。}
\BookSubtitle{三}
\BookParagraph{你听说过机器学习里有一个叫做「生成对抗网络」的东西吗？}
\BookParagraph{听着好像挺玄乎，但其实原理很简单。有两个神经元网络，一个叫「生成器」，一个叫「鉴别器」。生成器尽量以假乱真，生成一些假的样本；而鉴别器则尽量明辨真假，把假的样本以较高的概率分辨出来。}
\BookParagraph{我们让生成器和鉴别器来玩儿一个游戏：用鉴别器的输出去训练生成器，让它生成的样本更接近真实；再用生成器的输出去训练鉴别器，让它长一双火眼金睛，去伪存真的能力越来越强。如此反复成百上千次，直到最终生成器生成的样本，鉴别器已经无法分辨真假。}
\BookParagraph{没错儿，你就是这个「生成器」。而你需要找的是一个熟识的女生去当那个「鉴别器」——她需要从你的一言一行中看出，到底你是个可以托付终身的「蓝筹股」呢，还是一个「满纸荒唐言」不学无术三观不正的侃爷？}
\BookParagraph{比如你如果上来就问「咱们加个微信？」，她说「回头再说吧」，你如何应对？}
\BookParagraph{你跟她说「不好意思同学，我找不到手机了，你能给我手机打个电话吗？」她要是说「你找别人吧」，这句话你怎么接？}
\BookParagraph{如此反复和你的「鉴别器」激烈交锋，直到你的语音、语调、口才、幽默，一切都不卑不亢、恰到好处、炉火纯青、像个外交部发言人为止。}
\BookSubtitle{四}
\BookParagraph{当然了，搭讪这件事儿头一两句话至关重要，既不能落了俗套，又必须立竿见影，也是你「生成对抗」训练的重要议题之一。}
\BookParagraph{比如你在图书馆坐在心动女生对面，看到我的这篇文字，情不自禁地笑出了声儿。}
\BookParagraph{你抬起头，正好和女生目光交接，她好像在问：「你笑什么呢？」}
\BookParagraph{你自然而然地跟她解释：「微博上有个叫李葆春的，哈哈哈哈，实在是太能吹了，笑死我了。」}
\BookParagraph{她问：「李葆春是谁？」}
\BookParagraph{·  ·  ·  ·  ·  ·}
\BookParagraph{希望有朝一日能够在这条微博下面看到一条评论：「我就是故事里的那个心动女生。」}
\BookDate{2018 年 11 月，多伦多}
\BookArticle{相　亲}{相亲}
\BookParagraph{李老师，我碰不到喜欢的人了，老是回忆以前，去接触新的圈子新的朋友都觉得累，不是真的自己。是不是女生到了年纪都要草草相亲结婚？现在这种时代怎么没有所谓的爱情呢？（问题来自微博私信）}
\BookSubtitle{一}
\BookParagraph{Al Gore —— 就是克林顿时代美国的副总统后来跟小布什竞选没选上的那个Al Gore —— 十二年前曾经拍了一部纪录片，是关于全球气候变暖的，片名叫「An Inconvenient Truth」。}
\BookParagraph{这部纪录片基本就是 Al Gore 自己做的一个演讲，似乎很枯燥无味，但 Rotten Tomatoes 打了93\%的高分，值得一看。光是Gore的演讲技巧就可圈可点，比近几年国内企业家那种乔布斯式的演讲风格高明多了。}
\BookParagraph{比如有一张幻灯片讲到二氧化碳浓度最近二十年直线上升，Gore自然而然地站上了一台升降机，一按按钮，慢慢随着升降机腾空而起。}
\BookParagraph{Al Gore 的这部电影当年就得了两项奥斯卡奖，他本人第二年还得了诺贝尔和平奖。}
\BookParagraph{哈哈，聊得一高兴差点儿跑题了。}
\BookParagraph{是不是女生到了年纪都要草草相亲结婚？}
\BookParagraph{虽然「草草」两字略显悲观，到了什么「年纪」也不是很明确，但是你说的基本没错儿，对大多数人来说还真是这么回事儿。}
\BookParagraph{借用Al Gore的话来说，it is 「an inconvenient truth.」}
\BookSubtitle{二}
\BookParagraph{不过我一直盲目地相信，世界上总是存在这么一个人，是你一旦认识就能爱上的人。}
\BookParagraph{我喜爱的电影「When Harry Met Sally」里就有这么一句经典的台词儿：}
\BookParagraph{「Somewhere out there is the man you are supposed to marry. And if you don’t get him first, somebody else will, and you’ll have to spend the rest of your life knowing that somebody else is married to your husband.」}
\BookParagraph{看了这段儿有没有一点儿「大干快上、只争朝夕」的感觉？那个你还不认识但早就属于你的人，你要不赶紧第一时间拿下，他就是别人的了。}
\BookParagraph{所以说啊，这个时代还是有「所谓」的爱情的，可惜这爱情只偏爱那些早已准备好、随时可以冲出起跑线去抢的人。}
\BookParagraph{抢到爱情的算你的，抢不到的也别怨天尤人。}
\BookParagraph{「各就各位。预备。」}
\BookParagraph{赶紧起床去相亲吧——发令枪就要响了。}
\BookSubtitle{三}
\BookParagraph{不过，和「草草相亲结婚」相比，你真的更向往更喜欢那种两情相悦、心心相印的爱情吗？}
\BookParagraph{Harry在电影「When Harry Met Sally」里说：}
\BookParagraph{「Right now everything is great, everyone is happy, everyone is in love and that is wonderful.  But you gotta know that sooner or later you’re gonna be screaming at each other about who’s gonna get this dish. This eight dollar dish will cost you a thousand dollars in phone calls to the legal firm of That’s Mine, This Is Yours.」}
\BookParagraph{你瞧，爱情小说的结尾，永远是「live happily ever after.」就像一场平衡木体操表演永远不让你看到下法一样，完美的爱情故事里也不会让你看到为了一件鸡毛蒜皮的小事儿大打出手闹离婚的结局。}
\BookParagraph{爱情就像不带氧气就去爬八千米的雪山，如果能成功登顶自然很 high 很有成就感，但要是雪崩了呢？}
\BookParagraph{所以说啊，草草相亲结婚说不定也没什么坏处。李开复说「买车是一生最坏的投资」——其实为了找到爱情而「投资」，说不定比买车还要糟糕。}
\BookParagraph{因为钱花了还能再挣，爱情一旦变成渺无音信甚至反目成仇，就很难再回来了。}
\BookDate{2018 年 11 月，多伦多}
\BookArticle{怀　疑}{怀疑}
\BookParagraph{李老师，很抱歉再次打扰到您，我和她碰上了一个问题，这个问题让我们都不知道该怎么办了。}
\BookParagraph{女友认为我出轨了。我和她异地恋，我那天去看她和她待在一起的时候，一位我曾经的女性朋友发消息给我，内容是上次聊天聊到的键盘的事情。这位朋友换了一个蓝牙键盘觉得很可爱，发消息和我说了一下。当时我的女友就在我边上，她看到了有其他女生发消息给我。我因为怕她误会，删除了记录，那天我女友看见就直接崩溃了，导致我们的关系变成了现在这样。}
\BookParagraph{我和那位女性朋友是和女友在一起前认识的。有时候会联系，平常聊天我都是和她聊一些我和我女友的事情。因为我自己朋友很少，很多时候和女友闹矛盾我都不知道怎么办，又想解决问题就会去问问她，以她的的视角看看有没有办法。对于那些我做的事情我都是担心我女友会多想才删除的，没想到会变成现在这个样子。我可以拿我这辈子任何重要的事情发誓，我要是有什么出轨的想法，我都不得好死。我不知道怎么去处理这些问题，我明明很爱我的女友，为了见她，我各种省钱，为了回她消息，我会时不时看手机，都已经出现了幻听。现在我和我女友的关系变成这样，我真的不知道怎么去做。（问题来自微博私信）}
\BookParagraph{哈哈，谈恋爱嘛，总是要吵架的——我都可以想象你们之间的对话：}
\BookParagraph{「我可以发誓，绝对从来都没有过出轨的想法。」}
\BookParagraph{「那你没事儿干嘛要删微信聊天记录啊！」}
\BookSubtitle{一}
\BookParagraph{大概是七八年前的圣诞节前几周吧，我运气不太好，收到了法院的一封信：「Jury duty call」。}
\BookParagraph{什么是「Jury duty call」？哈哈你放心，不是传票，而是要求我去法院履行每一个加拿大公民都必须履行的义务：参加一个「Jury panel」。}
\BookParagraph{「Jury panel」又是什么呢？就是一两百人在一个大房间里待命，时刻准备着，哪一个案子开庭了需要选择陪审团成员（juror）时，就进去几十个人备选。}
\BookParagraph{我那天的运气真的是超级糟糕。一个强奸案件开庭审理，随机抽签竟然把我抽中了，双方律师发现我的职业是教授，欣然同意，不由分说，我立即成了十二个陪审团成员之一。}
\BookParagraph{案子审理了五天。审理过程中，不允许带手机，不允许上网，必须专心认真地听讲，每个人只给了一支笔和几张纸用来做笔记。}
\BookParagraph{但后来我发现，这五天时间是我印象里过得最充实的五天。}
\BookParagraph{因为这段时间，法官给我们陪审团的每一个成员，详尽而又清晰地讲解了加拿大法律的要义。给人的感觉，就像张无忌在武当山顶跟张三丰学太极拳，新鲜热辣，现买现卖，学了马上就可以拿来用。}
\BookParagraph{每一个陪审团成员，刚开始可以完全是法盲——需要的只是你公平、能独立思考、听得懂英文。}
\BookParagraph{而双方律师总结发言结束之后，对案件的判决就由陪审团单独决定，法官无权过问。}
\BookSubtitle{二}
\BookParagraph{在法官给我们讲过的法律要义中，我记得最清楚的两条，是「presumption of innocence」和「beyond a reasonable doubt」。}
\BookParagraph{「Presumption of innocence」好懂：在控方（Crown prosecutor）证明有罪（guilty）之前，必须假设辩方（defendant）是无罪的。}
\BookParagraph{而「beyond a reasonable doubt」法官则花了很长时间给我们解释。控方必须提供足够证据（burden of proof），证明辩方是「guilty beyond a reasonable doubt」。也就是说，不能是「可能有罪」（probably guilty），但也没必要完全不可能无罪（proof to an absolute certainty），而是这两点之间的一个状态。}
\BookParagraph{而这里的怀疑（doubt）不能基于同情或偏见，而是要基于逻辑和常理（reason and common sense）。}
\BookParagraph{正确理解「reasonable doubt」，可以援引著名的「Blackstone’s formulation」：宁可放走十个坏人，也不冤枉一个好人。}
\BookParagraph{我记得我当时想的是，这都是谁发明出来的——发明这玩意儿的人简直是个天才。}
\BookSubtitle{三}
\BookParagraph{我想说的是，男女朋友或夫妻之间若是有所猜疑，我同样提倡「beyond a reasonable doubt」。}
\BookParagraph{而你的这个案例，你的女友作为控方律师，完全无法提供足够强有力的证据，可以「beyond a reasonable doubt」。}
\BookParagraph{和别的女生聊天不等于在和她谈恋爱，而删除微信聊天记录，则很可能是因为你害怕她会多想，会误解。}
\BookParagraph{如果我把你这个案子交给一个十二人陪审团，他们一定会一致认定，没有足够证据证明你出轨。}
\BookParagraph{「Beyond a reasonable doubt」这道理听着很高大上，但其实很朴素：咱不能捕风捉影，而是要有真凭实据。}
\BookParagraph{其实我很能理解你女友作为控方的心情：万一辩方出轨了呢，心里过不去这个坎儿。}
\BookParagraph{但是换个角度想，如果因为捕风捉影而和一个有志青年—— 现在流行叫高富帅——失之交臂，多不合算啊。所谓「过了这村儿就没这店儿」了。}
\BookSubtitle{四}
\BookParagraph{那次审理的强奸案件，双方证人证词陈述花了五天时间。陪审团却只用了一小时就达成了共识：罪名不成立（not guilty）。}
\BookParagraph{谁也不知道强奸案是否发生过，只是把「beyond a reasonable doubt」的原则用上了而已。}
\BookParagraph{希望你的女朋友看了我的这篇文字，可以回心转意，和你和好如初。}
\BookParagraph{就是别再删微信聊天记录啦。}
\BookDate{2018 年 12 月，多伦多}
\BookArticle{前女友}{前女友}
\BookParagraph{李老师，我们是在他有女朋友的时候认识的，那时候我们仅仅相处几天就迅速喜欢上了对方。我知道他有女朋友就一直保持着距离，直至后来发生了一些事我拉黑了他。上个月他跟女朋友分手后开始找回我，我们最近聊天以及相处模式都跟情侣没什么两样。可是我一直疑心很重，我不知道他是不是已经忘了他的前女友，他跟我在一起的时候会不会在想她。我不知道该不该等，等他完全忘记她，我只希望别人心里是完全只有我的时候我才会交付真心。我也给自己很大压力，怕他会拿和我的相处模式跟她前女友做对比，毕竟他们刚分手不久。李老师我该怎么办？（问题来自名为「意大利」的微博私信）}
\BookParagraph{意大利同学：}
\BookParagraph{谢谢你的来信。因为最近这一段时间比较忙，拖了将近一个月才有机会给你回信，实在不好意思。不过你就当是自己用信纸手写的信，不远万里从中国寄到加拿大，到现在才收到对方的回信，大概也可以算是一份惊喜吧。}
\BookParagraph{你问我你该怎么办，可其实你心里多半儿早已经有主意了：你希望和新任男友这份感情能够简单、纯粹，没有任何杂质；你希望两个人的心里都完全只有对方。为此你甚至愿意再等一段时间。}
\BookParagraph{我衷心为你理想中的爱情点赞，但是恰恰相反，我认为任何一个有情有义重感情的人，都不应该忘记他的前女友（或者前男友）。}
\BookParagraph{其实不用说是曾经发自内心地深爱过的人，就算是你的老同学、老领导、老战友，难道不应该多年一直铭记在心，逢年过节互致问候，每隔几年拎上一瓶儿茅台，找个临街的馆子小聚一下吗？}
\BookParagraph{你可以想象一下，一个可以把前女友忘得一干二净的人，将来万一和你分手，是否也有可能把你也忘得一干二净呢？}
\BookParagraph{我的立场是，他越是忘不掉以前和前女友的感情经历，越是不愿意告诉你他内心深处那些刻骨铭心的场景和故事，他这个人就越靠谱，越值得你去真心地爱。}
\BookParagraph{相反，他越是满口向你保证前女友早就忘了，从来都不曾想念她，那这个人就越像一个感情上的骗子，就越需要多加小心。}
\BookParagraph{当然，如果你不会动不动就盘问他和前女友的事儿，自然也就不会谈论这个话题，不会让他处于两难的境地。}
\BookParagraph{我甚至认为，任何一段感情中，都应该对自己在对方心目中的地位有一种绝对的自信。比如他前女友的生日聚会或者婚礼，大可陪他积极参加，或者建议他自己去参加。他前女友生活中遇到难处，也大可替他买一张火车票，鼓励他去她的城市看望她。}
\BookParagraph{那这么放任他和前女友单独接触，会不会旧情复燃？}
\BookParagraph{当然有可能。可这旧情复燃的概率，和他在同学聚会上爱上一个女同学、战友聚会爱上一位老战友的概率差不多。你不能因为这些小概率事件，就不让他去参加老同学和老战友聚会，是吧？}
\BookParagraph{我要是他，你这样一位超级自信、通情达理、重感情的女朋友，打着灯笼都难找，前世积了德才修得来，怎么会傻到丢了西瓜去捡芝麻呢？}
\BookParagraph{当然，如果真碰到这样的人，让他早点儿用实际行动向你表明心迹，也不是一坏事儿。放心，今后还有的是机会。}
\BookParagraph{希望你将来和男友的感情发展顺利圆满，生活简单幸福。}
\BookParagraph{新的一年就快到了，提前祝你新年快乐。}
\BookDate{葆春}
\BookDate{2018 年 12 月 24 日，多伦多}
\BookArticle{罗马假日}{罗马假日}
\BookParagraph{我遇到了我喜欢的人，在喜欢他的过程中，我确实变成了更好地自己。因为之前他没有明确拒绝我，让我有了一丝希望。但是现在他明确拒绝我了，同时也因为是某方面的合作伙伴，我们达成了做好朋友的关系。我很舍不得他这个好朋友，但是我不知道该如何与自己喜欢的人做好朋友？（问题来自微博私信）}
\BookSubtitle{一}
\BookParagraph{你一定看过六十五年前拍的那部老电影《罗马假日》吧？}
\BookParagraph{也许是奥黛丽 · 赫本和格里高利 · 派克在电影里合作得太出色了，当时的报纸上有传言，赫本和派克在拍电影的那一段时间曾经相恋过。}
\BookParagraph{对于这一传言，赫本是这么说的：}
\BookParagraph{「Actually, you have to be a little bit in love with your leading man and vice versa. If you’re going to portray love, you have to feel it. You can’t do it any other way. But you don’t carry it beyond the set.」}
\BookParagraph{中文大意是：「如果你想把爱情演好，你就得有一点儿爱的感觉。但是电影拍完了，你也没必要继续爱下去。」}
\BookParagraph{派克则是这么说的：「Everyone in the set was in love with her.」（剧组里的每一个人都爱她。）}
\BookParagraph{《罗马假日》是赫本拍的第一部好莱坞电影，第二年她就靠这部电影里的出色表现得了奥斯卡最佳女主角奖。派克后来也非常喜欢罗马，每次带孩子去罗马都要去电影里著名的「Mouth of Truth」去玩儿，重温他当年和赫本稍带临场发挥的那段儿表演。}
\BookParagraph{赫本是在派克为了庆祝电影公映举行的一场鸡尾酒会上认识了她的第一任丈夫 Mel Ferrer 的。后来赫本和 Ferrer 一起生了一个孩子，婚姻持续了十四年。}
\BookParagraph{无独有偶，在这个电影拍摄过程中，派克也认识了一个去采访他的年轻女记者。电影上映两年之后，派克和他原来的妻子离婚，和这个女记者结婚了。}
\BookParagraph{赫本和派克因为这部电影，成了一生的好朋友。}
\BookSubtitle{二}
\BookParagraph{几年前网上有一个流传得很广的故事。故事里说拍电影时派克其实是很喜欢赫本的，赫本结婚时，派克还送给她一枚蝴蝶胸针。这枚胸针赫本非常喜欢，一直珍藏着。}
\BookParagraph{到赫本去世后，派克亲自参加了她的葬礼，在葬礼上老泪纵横，对着赫本的棺材说：「你才是我今生唯一的真爱。」十年后的一场拍卖会上，年迈的派克亲自前去买回了那枚胸针，去世时还把胸针放在身边。}
\BookParagraph{这个故事听起来挺感人，我记得五六年前微博上还有不少人转发过。}
\BookParagraph{但这段故事其实极有可能是后人杜撰的。谷歌搜索「Hepburn Peck Butterfly Brooch」，没有什么值得信赖的消息来源。}
\BookParagraph{唯一值得信赖的一篇是《人物》杂志在赫本去世那年发表的文章。文章里没有提到胸针，仅仅提到了赫本葬礼中用到的鲜花是派克和其它朋友们一起送的。}
\BookParagraph{更何况故事实在是有悖常理，赫本的葬礼上她前夫 Ferrer 等人都在场，派克不太可能突然跟大家宣布，其实赫本才是他一生的真爱。}
\BookSubtitle{三}
\BookParagraph{无论这段感人的故事是否真实，《罗马假日》公映后赫本一直住在欧洲，派克则一直住在美国。}
\BookParagraph{他们各自都很忙，很少联系。就像我们大家也很少跟远方的朋友联系一样。}
\BookParagraph{我觉得，你越是喜欢这个男生，越想和他做好朋友，你和他就越做不成好朋友。}
\BookParagraph{我建议你远走高飞，别动不动就和他微信联系，该忙什么就忙什么，联系越少越好。}
\BookParagraph{只有这样，你们才有可能真正成为一生的好朋友。}
\BookParagraph{就像赫本和派克那样的好朋友 —— 一个人去世了，葬礼上还有另一个人送的鲜花。}
\BookDate{2018 年 12 月，多伦多}
\BookArticle{分　手}{分手}
\BookParagraph{15点17分:}
\BookParagraph{李老师，您好呀！首先提前祝您新春快乐！之前一直看您在微博上给大家解答各种问题，每次都能有很多感悟。最近我也遇上了一个问题，如果能有幸获得您的解答我一定开心的原地转圈圈。}
\BookParagraph{我从去年开始交往了一个男朋友，但因为各自的学业，我和他之间存在七个小时的时差。我们平时也会用打字聊天，视频通话等方式保持联系，所以我本来觉得问题不是很大。但最近男朋友家里发生了一些事情，让他很难过，精神压力也很大。他因为不希望我为他担心，总是说的有所保留。而让我感到无力的是即使我知道了事情的全部，在这个时候也没有办法立马出现在他身边陪他度过难关，语言上的安慰又很单薄。想请问李老师对异地恋甚至异国恋的这种状态怎么看？谢谢您！}
\BookParagraph{18点33分:}
\BookParagraph{李老师，就在两个小时前，我的男朋友向我提出了分手，一切猝不及防。我在一个多小时里表达了我所有的想法，他仍然坚持原来的立场，所以最后仍然是结束了。如果您看到了这里，我真诚地感谢您愿意花时间听我的这些碎碎念，如果您仍决定回答我上面的那个问题，我也真诚地希望它能对每一对正在进行异地恋的情侣有所帮助，希望大家都能收获幸福。再一次，新春快乐。}
\BookParagraph{（问题来自微博私信）}
\BookParagraph{芝子同学：}
\BookParagraph{恰巧在你给我写第一封信之后不久，你男友就提出了分手，这实在是有点儿巧合，令人惋惜不已。}
\BookParagraph{我一直觉得恋爱的最高境界是：你们两人并不需要对方，但却都想要对方。}
\BookParagraph{「需要」和「想要」，一字之差，天壤之别。}
\BookParagraph{你的来信里，我注意到了一个细节。}
\BookParagraph{你提到了和男朋友两地的时差，是七个小时。}
\BookParagraph{我似乎都能感觉得到，那无数次加上七个小时来计算对方时间的夜晚，还有微信里无数条来来往往的消息。}
\BookParagraph{但这么多的消息里，有多少仅仅是「需要」（need），有多少是「想要」（want）？}
\BookParagraph{「需要」是仅仅想有个人柴米油盐地陪着你，「想要」则是一种无端的热爱，一种掐不灭的激情。}
\BookParagraph{就好像没有人「需要」一部莱卡的微单，但是很多人都「想要」。}
\BookParagraph{与其一条一条地发微信，表达「想要」最好用让人看得见摸得着的东西——写一封「碎碎念」的长信，寄一大摞刚刚冲印好的彩色照片，每一张的背后都有一行亲笔写的脚注。甚至寄一部最便宜的 GoPro，里面录满了你一天一天的快乐时光。}
\BookParagraph{其实，心里觉得特别难受无助的时候，第一个想要聊天的人，才是你内心深处最「想要」的——你平常说不定没怎么想起这个人，从来没有「需要」过他。}
\BookParagraph{我想起了 Ernest Cline 写的科幻言情小说《Ready Player One》，就是后来被斯皮尔伯格拍成了电影的那部小说。}
\BookParagraph{小说里的男女主人公 Parzival 和 Art3mis 在一个叫OASIS游戏里不打不相识，是针锋相对的竞争对手。但是时间长了，Parzival发现，Art3mis和自己共同语言竟然多到了「we finish each other’s sentences」的程度。两人天天从早到晚在游戏里泡在一起，在游戏里相爱相惜。}
\BookParagraph{他们的「网恋」正发展到如火如荼的节骨眼上，Art3mis突然消失了。}
\BookParagraph{Parzival不惜一切代价，到处找她，可是无论如何也无法联系到她。}
\BookParagraph{后来的情节峰回路转，精彩纷呈，而小说和电影完全不同的是，直到大结局时，男女主人公才第一次在真实世界中见到对方。}
\BookParagraph{「We sat there awhile, holding hands, reveling in the strange new sensation of actually touching one another.」}
\BookParagraph{你看，千辛万苦才得来的，不是令人一见倾心的声音，不是动人心弦的文字，更不是无穷多条微信，而是「touching one another」的真实。}
\BookParagraph{我希望你喜欢的人即使消失了也终将归来，希望你能最终找到「想要」你的、愿意在最难过时和你长聊的、你可以「touching one another」的人。}
\BookParagraph{新春快乐。}
\BookDate{2019 年 2 月，多伦多}
\BookArticle{Bittersweet}{Bittersweet}
\BookParagraph{李老师你好。去年考研后期一个朋友把你的微博推荐给我，那时候学习累了就看看你的微博，很幸运能遇到你。}
\BookParagraph{说一下我最近的心事吧。考研差不多一年的时间，在自习室每天都能遇到的女孩子，考完试意外地刷到了她的微博，可开心了。加了她的微信，每天都能聊很久。她很可爱，和她聊天也很开心。前几天出了考研成绩，我考的还不错，最近在学校准备复试。她说她明天来学校，也来自习室看书。明天就能见到她了，真开心。她现在没有男朋友，我就很纠结要不要追她。很想能和她在一起，但是考虑到毕业后我要去另一个城市去上学，距离她要一千多公里，又想干脆和她做个普通朋友吧。}
\BookParagraph{李老师，我该怎么办呢？（问题来自微博私信）}
\BookParagraph{哈哈，「情感博主」强势回归啦。}
\BookSubtitle{一}
\BookParagraph{你这绝对是「身在福中不知福」啊，不是已经成功地追到她了吗？}
\BookParagraph{最完美的恋爱，都是两个人自然而然开始的。比如六七十年代那会儿，自行车是一特「长脸」的大件儿。如果你骑着一辆永久牌二八自行车去人家家里接一位女生，她两三步助跑跃上你的自行车后座，轻轻地搂着你腰间的那一刹那，你们的恋爱就已经悄然开始了。}
\BookParagraph{从当年的自行车快进四十年到现在的微信，每一场愉快的聊天，都是恋爱悄悄开始的契机。虽然说微信可以同时和好几个人聊天，但是一心很难二用，这好几个人中，多半儿只能有一个人能聊得开心。这个人正好就是你。}
\BookParagraph{一起去学校自习教室学习，更是一种非常浪漫的恋爱方式。你想啊，如果一起去吃饭，很难有机会非常深入地讨论一件事儿；一起去看电影呢，只有看完了才能三言两语交流对电影的看法。吃饭看电影更像是试探性的约会，是否能成还真的不好说。}
\BookParagraph{只有一起去自习，才能你给我讲一道题，我帮你抽几个单词，把心思用到极致。当你和她讨论了半天终于把一个问题弄明白了的那一瞬间，心意相通，相视而笑，颇有「琴箫合奏，笑傲江湖」之感，竟是比她跳上你的自行车更上一层楼了。}
\BookParagraph{她要是看到了你的这个问题，估计一定会感慨，这还叫「普通朋友」，那到底什么才是真正的爱情？}
\BookSubtitle{二}
\BookParagraph{不好意思啊，我这篇命题作文写着写着又有点儿跑题了。}
\BookParagraph{你真正关心的问题，是明知道几个月以后会天各一方，相距甚远，是否还应该继续发展这场恋爱？}
\BookParagraph{那咱们就来做一个最坏的假设：你们一旦异地，很快就要分手。}
\BookParagraph{可即使如此，不是还有眼前这几个月时间吗？}
\BookParagraph{能天天在一起开心地度过几个月时间，没有老板天天催，没有「deadline」之前必须做完的事儿，是一件非常幸福的事儿。文青范儿可以到云南小憩身心或者到西藏体验高原反应，馋嘴猫可以把全城一家家网红馆子吃个遍，至不济也可以宅在家里一起打游戏看电影吧。几个月的闲暇，对于恋爱中人来说，简直就是人间天堂啊。}
\BookParagraph{那几个月之后的分手呢？}
\BookParagraph{英文里有一个词儿叫「bittersweet」，字面上看好像是「先苦后甜」，有点儿「把最甜的葡萄留到最后再吃」的意味。但其实不然，这个词儿既不是要教育劳动人民「苦尽甘来」，也不是要先吃最甜的葡萄，而是在形容一种混在微微的甜味儿中的苦——比如那种加了橄榄有一点点儿苦的martini。这种甜比纯粹的甜口感更佳，实属甜中的上品。它最牛的一个地方在于，不是所有人都喜欢。}
\BookParagraph{明知有可能要分手的爱情，是「bittersweet」的。像是Canadian Whisky和gin混合的鸡尾酒一样，越是悲剧式的结局，越令人陶醉于其中，难以忘怀。}
\BookParagraph{这杯酒已经调好了，你愿意尝一尝吗？}
\BookSubtitle{三}
\BookParagraph{未来永远是无法预测的。}
\BookParagraph{比如我现在正在飞往香港的航班上，到达之后的一两天中，会和哪几位朋友一起吃饭聊天，心情是愉快还是难受，我现在都不知道。}
\BookParagraph{更不要说几个月或者几年以后你们是否会分手这种事儿了。}
\BookParagraph{可是我就喜欢这种随意和不确定，这种「serendipity」——要是一切都可以预期，那日子过得还有什么意思？}
\BookParagraph{正是因为这个原因，我其实一点儿也不羡慕什么公司高管啊电影明星啊那种有钱人——他们的日程早就被助理排得满满的，想闲下来去麦当劳吃一份薯条可乐巨无霸都没戏。}
\BookParagraph{可是有一件事儿却可以非常靠谱地预测——你一辈子再次遇到一个可以一起开心地聊天一起自习而又年龄相仿的女生，非常不容易。}
\BookParagraph{就像那句名言：「结婚太重要了，一辈子就那么两三回。」}
\BookParagraph{也许有一天，你们真的会分手。但是就像迪斯尼电影「The Nutcracker and the Four Realms」里说的：「When you miss someone, you remember them. And one day that’ll make you smile.」}
\BookParagraph{你说你该怎么办呢？}
\BookDate{2019 年 2 月，多伦多}
\BookArticle{王牌}{王牌}
\BookQuote{一月二十八日：李老师新年好！有个情感问题求助您。我和我的好搭档发生了地下情，他有女朋友我有男朋友。在最近一次热烈的出差后，在这些相处的过程中，我日渐感觉我想和男朋友分手而想和他正式在一起。但他也许不能和女友分手。我也不知道我的想法和感受是否正确，因为半年后我们就会结束搭档关系（我们说好地下情关系到搭档结束为止），我该怎么做？喜欢您的文字.....}
\BookQuote{四月一日：李老师您好，很抱歉再次打扰。之前向你请教的问题，其实已经算是解决了，当然也继而出现了新问题。原先的问题因为到了需要他担负责任的时刻，我也极不理性地告知了他的女朋友，所以他「逃跑」了，是真实的那种逃避。我心理上的痛苦是我现在的新问题。感谢李老师做了我的树洞。（问题来自微博私信）}
\BookSubtitle{一}
\BookParagraph{咱今儿个不聊爱情，咱们来学点儿英文。}
\BookParagraph{两年前有一部电影叫做《Miss Sloane》不知道你看过没有，我看了不下五遍。主要是电影里面的英文对话密集得令人完全无法放松下来，像是上了一堂生动的英语听力课，看完了甚至觉得英语听力小有长进，还想再看几遍，「温故而知新」。}
\BookParagraph{影片中的主角，Elizabeth Sloane，是一位「lobbyist」，一个专门为了阻止或通过一份法案去游说国会议员的工作。她在影片开场就说了一句名言，后来成为点睛之笔：}
\BookQuote{Lobbying is about foresight, about anticipating your opponent's moves, and devising counter measures. The winner plots one step ahead of the opposition and plays her trump card just after they play theirs. It's about making sure you surprise them, and they don't surprise you.}
\BookParagraph{而这段儿需要细细品味的话，最重要的是那句「plays her trump card just after they play theirs.」}
\BookParagraph{什么是「trump card」呢？哈哈和特朗普没什么关系，英文里是「王牌」的意思，好像扑克牌中的「大王」，可以傲视群雄的那张。}
\BookSubtitle{二}
\BookParagraph{而你当初如果真切地希望那位男生为了你可以离开他的女朋友，也绝不能一厢情愿地空想，而是像「lobbyist」这个职业一样，把改变这位男生的想法作为一个重大专项来抓，做出一套切实有效的行动计划。}
\BookParagraph{不过恐怕你一直心里很清楚，这个「逆水行舟」的计划即使可以做出来，可以按部就班地完美实施，这位男生离开他的女朋友的概率也不太高。}
\BookParagraph{主要的一个问题是，他一直把和你的这段感情当成「地下情」来看，甚至还和你约定了到期的时间。你希望把「地下」情转到「地上」，就需要扭转他对你和你的感情的观点，螳臂当车，很不容易。}
\BookParagraph{这时候如果你还希望能达到自己的目的，手中就必须要有一张「王牌」。不仅仅如此，你还要尽可能地预见到对方的「王牌」是哪一张。}
\BookParagraph{「Lobbying is about foresight, about anticipating your opponent's moves, and devising counter measures.」你需要构想对方在你们俩这盘棋中所有可能的下法，然后一一构想应对之策。他无论如何布阵，都在你的意料之中。}
\BookParagraph{而你的「王牌」，一定要在对方亮出自己的王牌之后才亮出来。你这一招制敌的路数，就好像金庸小说里说的：「独孤九剑，有进无退。招招都是进攻，攻敌之不得不守，自己当然不用守了。」}
\BookSubtitle{三}
\BookParagraph{可惜我的这篇文字还没来得及写呢，你已经把你的牌打出去了。而你出的不但不是「王牌」，没有「晓之以理，动之以情」，甚至竟是最糟糕的那张牌，仿佛你希望把你爱的男生推得越远越好。}
\BookParagraph{不过，即使你很早就收到了我的建议，真的可以把这位男生当做一个假想的对手（opponent），把这个目标作为一个重大专项来制定一套详尽周密的计划吗？}
\BookParagraph{不仅仅是你，我猜即使换了任何一个人，都很难把自己的情人当作一个头号对手，即使这样对实现自己的目标有帮助。}
\BookParagraph{所以说，你和你心仪的男生再怎么「相见恨晚」，从一开始就无缘地久天长。}
\BookParagraph{至少，在这以后很多年，每次睡着之前无端地想起这个人的时候，你可以对自己说：我爱过他，却从来没有像对待敌人一样，对他祭出那张早已准备好了的「王牌」。}
\BookDate{2019 年 4 月，多伦多}
\BookArticle{爱情麻辣烫}{爱情麻辣烫}
\BookQuote{李老师您好，好久没见您在微博上分享问答了。最近也有一个困惑想听听您的看法：两个人在一起到什么程度就可以结婚了？什么原因可以促使一个人下定结婚的决心？（问题来自微博私信）}
\BookSubtitle{一}
\BookParagraph{你的问题让我想起了二十年前的一部电影，叫《爱情麻辣烫》。}
\BookParagraph{按理说，二十年前的电影绝对可以说是「老」电影了。可是在我的心目中，这部电影一点儿也不老，反倒是颇为入时，可以说是国产电影中的一部经典，和三十年前的《霸王别姬》有一拼。}
\BookParagraph{该片是张扬的处女作，可我觉得和《洗澡》并列，可以说是他最好的作品 —— 就好像慕容雪村的第一部小说也是他最好的一部小说。}
\BookParagraph{电影一开头是王学兵陪着未婚妻和未来的岳父母去吃麻辣烫：「今儿的火锅太辣了。」}
\BookParagraph{然后镜头一转，开始了五段互相之间似乎没什么关系的故事。「声音」里的男生爱上了同班一个身穿白裙子的女生；「麻将」里三个老头儿一起和老太太相亲；「玩具」里刚结婚的一对儿不打不成交；「十三香」里的中年夫妻在孩子的泪水中离婚；「照片」里两个人终于决定在一起时，竟然一分开就再也找不到回家的路了。}
\BookParagraph{我觉得张扬的厉害之处在于五个故事排列的顺序。少年时的相思开场，紧接着的是老年时的相亲，新婚时的相对无言，中年时的分离，最后是青年时的相爱。}
\BookParagraph{而五个故事的弓弦，是由王学兵演的这对准备结婚的情侣串起来，也是以他们结婚的婚礼收尾。}
\BookParagraph{张扬好像是在告诉你，按照他安排的这个顺序想明白了这五个故事，你就可以结婚了。}
\BookSubtitle{二}
\BookParagraph{你的问题有两种可能的解释：一种是打定主意要结婚，但是不知道什么时候结婚为好；另一种是还不太确定是否应该和一个男生结婚，不确定和他结婚以后会不会幸福。}
\BookParagraph{第一个问题我的答案是：结婚越早越好。}
\BookParagraph{你一定会问：那暂时还没钱买车买房办喜酒怎么结婚？}
\BookParagraph{简单啊。没有钱买房可以先租房住着，没有钱买车可以买辆二手车，没有钱办婚礼可以从简。}
\BookParagraph{总之两个人都没钱也完全没关系，照样可以早早结婚。甚至越穷的时候结婚，将来的婚姻就越可能源远绵长，即使真的分开了，也会友谊长存。}
\BookParagraph{因为没钱的时候结的婚，经常能体现两个人的一种特别牢靠的默契 —— 而婚后的天长地久，往往依靠的正是这种「怎么又想到一块儿去了」的默契。}
\BookParagraph{而第二个问题 —— 在说不上知根知底时如何下定决心和一个男生结婚 —— 就不是一两千字能聊的完的了。你看连张扬都说了：「相爱是容易的，相处是困难的。」而想「长期共存」，你得找一个优点多于缺点、比较靠谱的好人。}
\BookParagraph{那什么样的男生比较靠谱呢？我有一个自作主张不一定对的判据供你参考：平常愿意经常自己做饭的男生 —— 无论做得好吃不好吃。}
\BookSubtitle{三}
\BookParagraph{买三角钢琴的时候我听说，上好的钢琴用到的杉木多是阿尔卑斯山高海拔地区背阴处采伐来的，那里终年阳光不多，树木长得极其缓慢，每英寸需要长二三十年；但是又极其均匀，年轮之间的距离几乎一致。}
\BookParagraph{婚姻其实也和钢琴有点儿像：上好的婚姻稳扎稳打，极其缓慢而又极其均匀地成长。}
\BookParagraph{就好像西方婚礼的誓言里说的：「for better and for worse, for richer and for poorer, in sickness and in health, to love and to cherish, till death do us part.」}
\BookParagraph{电影《爱情麻辣烫》「声音」那一段里，有一段高中男生王艾和高中女生高圆圆的经典对话录音剪辑：}
\BookParagraph{「你喜欢和我在一起吗？」}
\BookParagraph{「喜欢呀。」}
\BookParagraph{「真的喜欢吗？」}
\BookParagraph{「真的喜欢呀。」}
\BookDate{2019 年 9 月，多伦多}
\BookArticle{离婚}{离婚}
\BookQuote{李教授你好，看了很多你的微博文章，觉得你都能从平凡生活中体会到人间喜乐。最近婚姻生活遭遇严重打击，夫妻缘分耗尽在最后一刻，签了分居协议，关于小孩的抚养以及为数不多的财产的处置。在我看来即使是分手也希望好聚好散，不知道如何去劝慰孩子爸爸。孩子父亲是统计学 PhD，性格上比较固执也不善交谈。无法接受他的条件，就是不想分手但是也不想让你好过，无法心平气和。如何去面对这样一个尴尬的过程，希望你能给一些建议。（问题来自微博私信）}
\BookSubtitle{一}
\BookParagraph{「好聚好散」看起来挺简单，但实际上是很难做成的一件事儿。}
\BookParagraph{难就难在你觉得这么「散」是「好散」，他却觉得那么「散」是「好散」。你们俩如果分别到我这个「街道居委会」来哭诉，我一定会觉得你说得确实很在理儿，他也确实有他的难处。}
\BookParagraph{比如说吧，你多半儿会拿「孩子的抚养」说事儿，抱怨他再怎么着总得多想想孩子吧？而他想着的则可能是你们俩多年的感情，说结婚这么多年了，总不能无情无义到了直接扫地出门的地步吧？}
\BookParagraph{而「街道居委会」作为观众，「是非曲直未有所定」，立场也很难说向着谁。就好像三五好友一起看多年前的那部《中国式离婚》，「吃瓜群众」有的同情陈道明演的宋建平，有的却反过来去支持蒋雯丽演的林小枫：「谁叫宋建平没事儿就跟娟子眉来眼去的？」}
\BookParagraph{俗话说「当局者迷，旁观者清」，现在连旁观者都有的「站」「李总」，有的「撑」「俞总」，长篇大论连篇累牍地先掐起来了，当局者还怎么能「好散」呢？}
\BookParagraph{所以依我看啊，你们夫妻俩结婚时间越长曾经的感情越深，「好聚好散」就越没有希望。}
\BookParagraph{能「散」就已经不错了。}
\BookSubtitle{二}
\BookParagraph{我外公在六十多年前写的《龙虫并雕斋琐语》中，曾写过一篇短文，叫《夫妇之间》，正好就提及了夫妻打架 —— 现在时髦的语言叫「开撕」—— 的问题。他在文中写道：「夫妇反目，也是难免的事情。但是，老爷撅嘴三秒钟，太太揉一会儿眼睛，实在值不得记入起居注。甚至老爷把太太打得遍体鳞伤，太太把老爷拧得周身青紫，有时候却是增进感情的要素，而劝解的人未必不是傻瓜。」}
\BookParagraph{你瞧，其实夫妻打架也不完全是一件不幸的事。即使不能增进感情，至少也给自己积攒了一点儿阅历，给日子增添了一点儿颜色。将来老了动不了的时候，大概和丈夫当年的相见恨晚含情脉脉早已忘了一干二净，却竟然还记得「缘分耗尽」的那些最后一刻，然后不忘了再为「男的都不是好东西」这样的永恒真理叹一口气。}
\BookParagraph{外公的文字笔锋一转，还写了另外一种离婚：「倒反是有些相敬如宾的摩登夫妇，度了蜜月不久，突然设宴话别，揽臂去找律师，登离婚广告，同时还相约常常通信，永不相忘。」乍一看来，这种摩登的离婚方式堪称完美；但是再一琢磨，你大概会觉得这种「好聚好散」的离婚既不打架又不「开撕」，如同白开水煮挂面一般，缺少了「中国式离婚」的精彩回合，还略微有那么点儿遗憾不是？}
\BookParagraph{所以说啊，能好好聊聊好好劝当然很好，否则大大的打上一架，歇斯底里，快意恩仇，哭得昏天黑地，岂不也是一番不同的人生体验？}
\BookParagraph{虽然将来不太会经常通信，但一定会永不相忘。}
\BookSubtitle{三}
\BookParagraph{什么事儿都有好坏两面 —— 离婚的好处，是从唯一选择回复到无穷多的可能性，就像手机恢复了出厂设置，新生活从此开始。}
\BookParagraph{再次体验一把到麦当劳跟陌生网友相亲的过程，今天跟这个见见，明天再跟那个聊聊，喜欢就穷追猛打，不喜欢就微信拉黑，「滚滚长江东逝水，浪花淘尽英雄」，多来劲儿啊。}
\BookParagraph{离婚有点儿像从现在的公司辞职，就算还没找好下一个工作，也不能算是「失业」，而只不过是暂时「in between jobs」而已。}
\BookParagraph{就是下次再结婚，别忘了签一份婚前协议。}
\BookDate{2019 年 10 月，多伦多}
\BookArticle{Mr. Right}{Mr. Right}
\BookQuote{李老师，我是一个 25 岁的女生，研究生毕业工作一年多，单身了四年左右。我发现，找对象太难了，喜欢上一个人并且让这个人喜欢上你也太难了。我的好闺蜜们都陆陆续续的找到了稳定的交往对象，并且向结婚迈进，像是约定好了一样。所以这件事好像成了我一个缺点，她们也经常催促。这期间也有男孩喜欢我或者表达好感，我也迫于压力逼迫自己去积极相处，但是结果也都不尽如人意。和不喜欢的人在一起，实在没法快乐。}
\BookQuote{最近我们部门的一个入职不久的新员工，也是小学弟，时不时找我聊天，大概持续了两三周。晚上下班以后我们微信上总能聊的很开心，并且小学弟本人也高高白白俨然就是我喜欢的类型，所以每天都很开心。}
\BookQuote{最近因为有喜欢的明星的电影上映，我就开玩笑发消息给小学弟说让去支持支持，他回复的也很令我开心，想我带他一起看。真的很开心，心里面甚至忍不住狂喜，觉得这次爱情可能真的要降临到我头上了。今天晚上我们一起出去看了电影，可是真的和真人去交流，却发现火花少的可怜，他似乎也并没有对我很上心，真的是在吃着东西开开心心看电影，和我的沟通交流也完全和之前不太一样。总之看完电影也就各回各家了，我非常非常难过失望。同时，他也没有在像之前那么热情在微信上和我聊天了。}
\BookQuote{我好困惑迷茫，现在的我要怎么心理建设，我还会去恋爱么，还有能力还有机会去恋爱么？我好担心我的全世界只剩我一个人孤孤单单。（问题来自微博私信）}
\BookSubtitle{一}
\BookParagraph{刚刚读到你提到自己二十五岁、「单身四年」的时候，我隔着手机的屏幕，都能体会到你那份有点儿焦虑不安、想「与时间赛跑」的心情。}
\BookParagraph{可你发现没有，找个合适的男生谈恋爱，跟到健身房运动减肥差不多：你再着急也没什么用。就如同减肥不是一天两天就可以见效一样，谈恋爱也不是一天两天就能谈成的事儿；就如同减肥会有反复一样，谈恋爱也会分手、会痛苦、会和好如初。}
\BookParagraph{不过你这种焦虑的心情我非常可以理解。估计你的心里总是惦记着，在多少岁之前一定要完成这件事儿，就像是一个「deadline」。时不我待，只争朝夕，离你心中的这个「deadline」越近，你就越觉得希望渺茫。}
\BookParagraph{但其实就像健身减肥是一辈子的事儿一样，谈恋爱也是一辈子的事儿，没有什么「deadline」。就像健身减肥可以随时开始一样，谈恋爱如果结局幸福圆满自然好，如果不尽如意，也一样可以随时开始，重新再来。}
\BookParagraph{很多人喜欢说「人生苦短」，我完全不同意：人生一点儿也不短 —— 人生很长很长，将来有的是时间，甚至长得有点儿让人不知道怎么打发。你想啊，假设你二十五岁开始工作，六十岁退休，能活到九十五岁，工作三十五年还算有事儿干，退休三十五年要打发时间，得玩腻多少游戏、看腻多少网红直播啊。}
\BookParagraph{想明白了「着急其实没有用」、「时间有的是完全没必要着急」这两条「定理」，咱们就在「战略」上成功地藐视了「敌人」。}
\BookParagraph{不过别忘了，要想好好地恋爱（或者实实在在地减肥），咱们还得在「战术」上重视「敌人」，把这个任务多快好省地完成。}
\BookSubtitle{二}
\BookParagraph{我不是很相信「一见钟情」。}
\BookParagraph{除非这「一见」是指一上飞机偶遇就开始聊天，越聊越投缘，一直聊到越洋航班落地，还觉得有点儿意犹未尽。}
\BookParagraph{我觉得谈恋爱要想谈到「合二为一」的境界，必须从了解对方开始。就连日常的小事儿 —— 譬如说买辆新车 —— 我们都倡导先做一些调查研究，更何况谈情说爱这样关乎幸福与否的大事儿。}
\BookParagraph{而了解一个男生越多，你就越可能因为他的一些小小的不足之处把他拒绝：「不行不行，这位同学有点儿口吃。」「妈宝男坚决不能要啊。」「凤凰男也不行。」「这人也太能说了吧，有点儿油嘴滑舌。」「这人也太木讷了吧，我说十句他只回一句。」}
\BookParagraph{我最近一直在准备买一辆新车，就差点儿犯了这个错误。刚开始买车的时候，我非常关注车是不是好看 —— 也就是所谓的「颜值」—— 但凡不好看的坚决不买。结果研究了一段时间我发现，好看的车都很贵，从买车到日常保养，贵到了让我觉得其实性价比很成问题。贵倒还在其次，很多豪华品牌其实可靠性不怎么地，要是开在路上出了问题，既不方便又不安全。而可靠性比较出众的车型，绝大多数都不好看。}
\BookParagraph{等到我穷尽所有品牌的车型时才发现，我想找的那辆既好看又可靠、既物美又价廉、既加速快又超级省油的汽车，其实根本就不存在。}
\BookParagraph{没有十全十美的车，也不会有十全十美的人。不擅长言辞的男生也许将来更适合相依为命，苦孩子出身的凤凰男也许前途不可限量。很可能啊，你好多年之后才发现，最初的那个不善言辞有点儿口吃的男生，其实才最适合你。}
\BookParagraph{所以啊，其实谈恋爱不可能永远一帆风顺。我们要列一份详尽的清单，把每一位男生（或者女生）的优势和劣势仔细拆解一番。综合评估分数最高者，我们要本着「敌疲我打、敌退我进、敌驻我扰」的精神，百折不挠，排除万难，去争取胜利。}
\BookParagraph{我终于选定了要买的新车 —— 是一辆超级难看的丰田。}
\BookSubtitle{三}
\BookParagraph{前一阵儿在飞机上偶然机会看了一个十几年前的老电影。这是一部浪漫喜剧，名字叫「27 Dresses」。}
\BookParagraph{故事情节其实很简单，不用剧透都能大致猜到。Jane 非常喜欢当朋友婚礼上的伴娘，曾经当过二十七次，但自己却从来无缘「Mr. Right」。}
\BookParagraph{当她第一次见到 Kevin 时，非常不喜欢他。而她一直暗恋着她公司的老板，却偶然结识了她的亲妹妹，两人打算火速结婚。}
\BookParagraph{其实浪漫喜剧是非常不容易拍好的一类电影。主要是因为不同人对「浪漫」的感觉完全不同，有的人觉得感人至深就会有人觉得俗不可耐 —— 按我女儿的话说叫「sappy」。众口难调，想把电影拍得雅俗共赏着实不容易。像斯皮尔伯格「The Terminal」这样的上乘之作，难得一见。}
\BookParagraph{我在飞机上完全没有想到，这样一部关于伴娘在婚礼上偶遇「男神」的传统爱情喜剧，我居然看得感动不已。}
\BookParagraph{就像电影里面一样，无论早晚，我们每一个人总是能过上属于自己的那最特别的一天。}
\BookParagraph{这一天的那个人也许和当初自己想象的完全不一样 —— 没有「颜值」、脾气倔强、不善辞令。}
\BookParagraph{但是这一天总是会到来。}
\BookParagraph{新年快乐。}
\BookDate{2019 年 12 月，多伦多}
\BookArticle{一半是海水}{一半是海水}
\BookParagraph{王思聪和孙一宁的故事，开始只不过就是一个有钱人爱上了一个美丽的女生追了半天最后满盘皆输的故事。故事虽然很平常，但还是可以复盘看看他到底输在了哪里。}
\BookParagraph{我觉得，王思聪没学过计算机科学，一定不懂什么是「面向对象」的程序设计。「面向对象」虽然有点儿复杂，但大意是说每一个「对象」都有公有和私有两种变量，公有变量外界可以随意访问，而私有变量则「仅自己可见」。}
\BookParagraph{比如说吧，如果你跟一个心仪的女生聊起密室逃脱或者肖战，多半儿她会很自然地跟你聊聊自己的看法；但如果你跟王思聪学习，去问人家「你现在在哪儿呢」「你约了谁啊」，那人家不但不一定有心情告诉你，反而会自然而然地产生反感的情绪，也许本来还有的一点儿好感也烟消云散了。}
\BookParagraph{所以老话里「谈对象」这个词儿挺中肯的：无论你的「对象」跟你熟到什么程度，人家都有自己的「私有变量」。}
\BookParagraph{当然了，王思聪可能确实有点儿陷入了单相思，于是做出了一个极其错误的决定 —— 临时去杭州找我们的女主。去就去了呗，可以说去办别的事儿顺便看看她，可以在杭州先逍遥几天，慢慢再约不迟。可我们的男主急得像热锅上的蚂蚁，一定要当晚跟我们的女主见上一面，这就不是恋爱，而是要挟了。}
\BookParagraph{人的本能都是要保护自己的，我们的女主自然是这方面的专家，决定立刻提出跟男主「做朋友」。所谓「做朋友」，就是把所有的公有变量全部转为私有，将来坚决不打算做朋友了。事已至此，便是华佗再世，也回天无力，后面死缠烂打的剧情，就毋需再看了。}
\BookParagraph{四十年前，在王朔的代表作《一半是火焰，一半是海水》里「海水」的那一半里，一个叫胡亦的女孩就跟书里的男主说过：}
\BookParagraph{「你这人怎么这样无礼 —— 我们不过是萍水相逢，一块玩了几天，我又没花过你一分钱，从始至终就是旅伴关系。别说没有什么，就是真有过什么，我想走你也管不着！难道你碰到对你热情一点的女孩子，就都以为她们一门心思要嫁你？」}
\BookParagraph{朔爷门儿清啊。}
\BookDate{2021 年 6 月，香港}
\BookArticle{生命的意义}{生命的意义}
\BookQuote{葆春老师，您好！我出身于偏僻的农村，是一个 “小镇做题家”，现在国内一所 985 大学的念研究生。一直以来，按部就班地学习，转眼就要 25 岁了。遗憾地是，我现在连个对象都没有。我也想谈恋爱，但时常感到自卑，不敢去追身边的女生。即使有女生找我聊天，我也是畏首畏尾，感觉自己不配。}
\BookQuote{一方面是，自己很普通，基本上没有什么优点。我没有侃侃而谈的口才，有些 “直男”。待人真诚，但没有很高的情商，说话有时会得罪人。虽然每周都坚持跑步，但体育不敢恭维。作为工科生，我更喜欢数学，动手能力却很差，甚至连游戏都不怎么会打，经常被同学笑话。我空暇时会看一些人文社科方面的书籍或者上网学一些自己感兴趣的课程，比如我经常看您的《幸福》，真的受益匪浅！哈哈，说来惭愧，我好像都没有什么正经的爱好。总之，我感觉自己是挺无趣的一个人。}
\BookQuote{另一方面是，农村物质条件匮乏，对自己的原生家庭感到自卑。我非常感激我的父母，他们在 “计划生育” 紧张的年代，“死里逃生” 给予我和姐姐宝贵的生命。姐姐凭借自己的努力，现在是一名老师。父母淳朴善良，但没什么文化，眼界相对狭隘。我担心将来的女友不会接受我这样的原生家庭。按照常理，婚姻似乎要门当户对。}
\BookQuote{读研后，圈子也小了很多，加之自卑，对未来的感情和生活充满了迷茫。老师，您有空的时候能给我提一点建议吗？谢谢！我提前恭祝您及家人虎年快乐！（问题来自微博私信）}
\BookSubtitle{一}
\BookParagraph{「What’s the meaning of life?」}
\BookParagraph{十年前，也就是 iPhone 4S 刚刚发布那年，全新的 Siri 简直可以说是人见人爱。不知道为什么，大家最爱问的问题就是这句「生命的意义是什么？」}
\BookParagraph{这个问题 Siri 当时最爱用的一个答案是：「All evidence to date suggests it’s chocolate.」}
\BookParagraph{这个回答自然是来自我喜爱的电影「Forrest Gump」里面的那句：「Life was like a box of chocolates. You never know what you’re gonna get.」}
\BookParagraph{说起巧克力，我小时候最爱吃的甜品当属酒心巧克力了 —— 这是小孩儿能喝点儿酒的唯一办法。有的巧克力里的酒多一些，沁人心脾；有些则几乎没有酒，正如世事无常，令人遗憾。}
\BookParagraph{但是现在想来，其实这个比喻还是不够「刺激」：生命的意义倒是更像坐过山车：爬得越高，俯冲得越狠；忽而春风得意、踌躇满志，忽而又五味杂陈、如入冰窟，这样的生命才有了意义。}
\BookParagraph{明明心里喜欢的女生不去追，主动来找你聊天的女生也不趁热打铁、火热开聊，生命还有什么意义？前怕狼后怕虎，因为自己想象的什么「门当户对」就不去恋爱，生命还有什么意义？因为自己无趣怕早晚有一天被女生拒绝就一潭死水不去开始，生命还有什么意义？}
\BookParagraph{失恋怕什么？咱谈恋爱就是奔着失恋的滋味去的！要是正好没失恋，就像过山车没玩儿上，还不够刺激呢！}
\BookParagraph{我中学同班同学中有个小美女，名字叫蓝宁。蓝宁不仅是绝对的美女，还博览群书，大家闺秀，书香门第。初二那年班里流行传看手抄本，蓝宁也要看。有男生说「这本书有点儿黄啊。」蓝宁斩钉截铁，寸步不让：「越黄越要看！」}
\BookParagraph{十年前的那个十月，如果我对 Siri 说「Siri，我爱你」，Siri 会用英文回答：「你还不了解我，葆春。」}
\BookSubtitle{二}
\BookParagraph{其实不仅女生，男生也应该多多结交。}
\BookParagraph{我一直认为，人一辈子认识不了几个有趣的朋友。大多数的人如过眼云烟，而只有极少数人，他们风趣幽默，个性鲜明，时间长了，像黑白照片里的一抹亮彩，成了生命的一部分。}
\BookParagraph{而不去多多结交男生，就很难认识这样有趣的人。}
\BookParagraph{比如聪哥，就是我多年的一位朋友。聪哥在香港当教授，是青年学者中绝对的翘楚，无论是质量还是数量，各种奖项多有斩获，学术水平几乎无人能敌。}
\BookParagraph{聪哥最爱问的问题是「什么时候约个饭？」然后到点儿就开着他的特斯拉来接我。我们的饭局天南地北，无所不聊，就是很少聊学术研究。我们大碗喝酒，再打上三两斤牛排，绝对属于「酒肉朋友」。}
\BookParagraph{有一次我旅居香港，独自居住时发高烧。想到有没有当地朋友的电话号码可以应急，我第一个想起的就是聪哥。}
\BookParagraph{发消息过去，聪哥秒回：他和他的特斯拉随叫随到。}
\BookParagraph{生命的意义是什么？}
\BookParagraph{生命的意义就在于，无论你的「圈子」是大还是小，有这么三两好友，在你危急的时候，赴汤蹈火，在所不辞。}
\BookSubtitle{三}
\BookParagraph{乔布斯最后一次去日本，是他去世前一年。}
\BookParagraph{那一年有一次他带家人去京都的一家日料店吃晚饭，懒得点菜，直接让老板推荐他们最好的寿司。吃完了一份意犹未尽，问老板能不能再来一份。}
\BookParagraph{离席时老板礼貌地欢迎他在不久的将来再来。乔布斯告诉老板，以后很可能没有机会了。老板问，那能不能给他女儿的照片签个名。乔布斯平时很少给人题词签名，这次也许是寿司过于惊艳绝伦，他欣然应允，一蹴而就：}
\BookParagraph{「All good things. 」}
\BookParagraph{是啊。所有美好的东西都像彗星划过夜空，像生日蛋糕上的蜡烛，像酒心巧克力，早晚有一天会消逝。}
\BookParagraph{All good things come to an end.}
\BookParagraph{在这之前 ——「Siri，我爱你。」}
\BookDate{2022 年 4 月，多伦多}
\BookArticle{相亲}{相亲}
\BookParagraph{今天跟朋友吃晚饭，聊到了相亲。}
\BookParagraph{我没相过亲，但是我一直觉得这是一种很有意思的体验。茫茫人海之中，三教九流各色人等，因为一个机缘巧合，两个本来完全陌生的人，就这么在咖啡厅见面认识了还聊了一会儿天。聊完又极有可能从此分道扬镳，再无往来。}
\BookParagraph{相亲的最大问题，我想大概是目的性太强。有了为了目的的审视，就少了一种自然而然的心动，就很难一见钟情。}
\BookParagraph{比如，我很难想象有目的的相亲中能聊什么。「请问你平时下班有什么兴趣爱好？」「我爱打游戏，你呢？」这样的聊法儿多土多无聊啊。}
\BookParagraph{不过，现在信息时代有了手机的一大好处是，天聊到聊不下去，不用再继续找新的话题，而是可以低头玩手机给别人发微信。}
\BookParagraph{不知道有没有那样一种两个人的相亲，没有那么多目的，就是多认识一个朋友，喝杯奶茶，聊几句天。聊完多半儿还是会沉没于通信录中，老死不相往来。}
\BookParagraph{可十几年后，回首往事，有没有可能还能忽然想起那天，一个清凉的午后的咖啡馆里，一个陌生人讲的故事中的一个小片段？}
\BookParagraph{如果微信发过去，对方竟然也能想得起来这段相亲的故事，那这种茫茫人海中毫无意义的相亲，也就有了特别的意义。}
\BookParagraph{所以，我觉得相亲还是挺好玩儿的。}
\BookDate{2022 年 9 月，多伦多}
\BookArticle{玩手机}{玩手机}
\BookParagraph{约会吃饭，女生开始玩儿手机怎么办？}
\BookParagraph{褚明宇的「霸道总裁」版本「不许玩儿手机了，听话」，理论上来说虽然可行，但是说话时的语气需要一点儿调侃，分寸很难掌握，我觉得不适合大多数男生的气场。}
\BookParagraph{如果约会吃饭时女生开始玩儿手机，有两种可能的情况。}
\BookParagraph{一种是对方确实有要紧的事务需要即时处理回复。这种情况，我认为最多也就是几分钟的事儿，正确的应对方式应该是耐心地等候一段时间。}
\BookParagraph{当然这段时间里，您不能自己也拿出手机开始玩儿 —— 对方可能是紧急事务，您就真的是在玩儿手机了。}
\BookParagraph{能不能继续吃饭呢？我觉得也不太合适，显得趁热把饭吃完比约会聊天更重要似的。}
\BookParagraph{我觉得比较少见但却可行的做法是：拿出一本英文版的《马斯克传》，从上次没看完的章节继续看下去。}
\BookParagraph{吃着吃着饭，拿出本书来看，虽然有点儿奇怪，但是并不令人反感，甚至还让对方边玩儿手机边产生一种天生的好奇 —— 看什么书呢？}
\BookParagraph{如果等候的时间超过了几分钟，对方还没有停下来或者解释一下的意思，那可能就真的是在玩儿手机了。}
\BookParagraph{这时候，我觉得应该自己反省一下，刚才到底是怎么把好好的天儿给聊到死胡同里去的。}
\BookParagraph{一般来说，这种情况说明对方不仅仅是对您不感兴趣，甚至有点儿反感 —— 基本上大势已去了。}
\BookParagraph{当然，如果您还不死心，想对女神做出最后的努力，可以一边看书一边自言自语、一语双关、进口攻退可守、听天由命地说：}
\BookParagraph{（注意，这里一个字都不能改 ——）}
\BookParagraph{哈哈哈，马斯克他妈真是太逗啦。}
\BookDate{2024 年 6 月，多伦多}
\BookArticle{生日快乐}{生日快乐}
\BookParagraph{时间和空间两个维度，偶尔想起来，都有点儿令人伤感。}
\BookParagraph{前一阵子忙了一两个月，终于在截止日期前最后一天，写完了一篇得意之作。现在闲下来，整理了一会儿近来发表过的文字。时间这东西过得真快，一年、两年的时间，如果不去看以前写的文字，都像是一晃而过，雁过无痕。}
\BookParagraph{前几天，和小女儿一起在机场候机时，她跟我说：你老了。}
\BookParagraph{没错儿，这就是时间让人伤感的地方。}
\BookParagraph{玩儿手机的时候，偶尔打开 Find My，看到此时此刻，全家都分别在世界上不同的位置 —— 和郭芳相距一万两千多公里，和小女儿相距两千多公里。}
\BookParagraph{我高一时去美国波士顿访问的时候，我的好朋友余晨同学给我寄的一封信结尾，写过一句至今难忘的话：「在地理上，我们比以往相距甚远；但是在心灵上，我们比以往任何时候都更加紧密。」}
\BookParagraph{拉回几十年时间的长卷，我清晰地看到的，还是那一年的初夏时节，那个穿着连衣裙，轻盈地跳上我的永久二八车，跟我一起从清华园到北大燕南园一游的女生。}
\BookParagraph{不能永葆青春，就让我们一起慢慢变老。}
\BookParagraph{郭芳同学，生日快乐。}
\BookDate{2025 年 8 月，香港}
\BookPart{芳华已逝 面目全非}{芳\\华\\已\\逝\\[18pt]面\\目\\全\\非}
\BookArticle{钢　笔}{钢笔}
\BookParagraph{幼时看外公伏案著书立说，从来只用毛笔。外公对毛笔择录甚严，小楷狼毫，产地出处，无一不亲自过问。只有在记日记的时候，外公才用钢笔，如毛笔般竖而握持，且只记流水帐，从不用之从事创作。我当时不懂为什么。上小学二年级的时候，我如愿以偿地得到了我的第一支在北大三角地商店六角钱买的钢笔。这支钢笔很短，但笔锋流畅，笔杆那端还有一个小熊猫。我小心翼翼的把钢笔别在铅笔盒里，每天随我排着路队放学回家。我用钢笔写周记，做应用题，是班上老师第一批批准用钢笔写作业的好学生。}
\BookParagraph{终于有那么一天，我心爱的熊猫钢笔从课桌滚落在地，再也不能流畅地写字。我偷偷地哭了，这才体会到一个人对一支笔会有那么依依不舍的感情，象旧日的故友。姥姥知道了，送给我一支英雄金笔，鼓励我学外公，写一手好字，做天下的文章。以后我长大了，姥姥却一直记着我喜欢钢笔，常将外公的金笔转赠于我。有几支我一来舍不得用，二则汗颜字写不成柳颜，文章又做得蹩脚，就一直留到今天。}
\BookParagraph{本以为信息时代，没有人在北美再用钢笔了。一年前参加与我做办公室邻居的一位在学术上德高望重著作等身的老教授的庆功会，会间基金会为表彰他的贡献，赠予一支钢笔。我们借此聊起钢笔，才得知老教授收藏钢笔，象多年的美酒，藏品弥足珍贵。老教授风趣地告诉我，其中有他几十年前的第一支钢笔，当年五元购得，后来却花了高价修复。当晚，我取出已故的姥姥当年相赠的派克钢笔，见笔如见人，只望有生之年，也能著书立说，做天下的文章。}
\BookParagraph{一个学生见到我的派克，喜出望外。原他也嗜钢笔，多年劝说别人改用钢笔。他兴趣盎然地娓娓道来，从 Waterman 到 Mont Blanc Jules Verne, 身价不菲，几百直至上千美金不等。我听得出了神，想起了我的第一支熊猫，想起外公的狼毫，想起姥姥赠的英雄，恍若隔世。}
\BookDate{2005 年 3 月，多伦多}
\BookArticle{钢　琴}{钢琴}
\BookParagraph{儿时没机会学钢琴，是因为家里没有钱买。其实，妈妈的乐器水准着实不低，是大学时文艺社团手风琴队的队长。妈妈自幼会弹钢琴，那是因为姥姥家多年以前有一架钢琴，陪伴妈妈长大。即使姥姥自己，当年年轻姑娘时在教会学校就读，钢琴自也是训练有素。文革开始的时候，妈妈积极要求入党。遗憾的是出身不好，于是亲率红卫兵上姥姥家抄家，所有值钱的首饰，字画，藏书，一并抄走，当然也少不了这架钢琴。据说，钢琴被红卫兵劈成了木头，当柴烧了。那以后，琴声就再也没有在姥姥家重现过。}
\BookParagraph{八十年代抢购风的时候，有一天，妈妈兴高采烈地告诉正上高中的我，她买了一架新钢琴。记得那时「星海」钢琴时兴，但妈妈图个省钱，买的是俄罗斯的琴，四千多块。那年头一月工资大多一百两百，对我来说这可是天文数字。妈妈平时最省吃俭用，但这一天她激动地告诉我，她买这琴是为了十几年后退休了自己弹。我从来没有见她这么高兴过。}
\BookParagraph{提货的那一天，妈妈特约了公司里几个精壮的小伙子，将崭新的钢琴直接运到了姥姥家的客厅。谁知姥姥知道了这事儿，怒目相向，严厉地批评了妈妈，说她铺张浪费。妈妈脸色铁青，难过了好多天。自此以后，妈妈经常轻抚着钢琴发呆，从没见她弹过一曲。翌年，妈妈确诊为癌症晚期，卧床不起。一次我陪着她看电视，林语堂的「京华烟云」主题歌响起的时候，妈妈的眼睛湿润了，紧紧握住了我的双手。我想，她是想起了她心爱的俄罗斯钢琴。妈妈的遗言里，没有提到她的钢琴。我没有哭，她一定是轻轻地把她自己的琴带走了，和她一起去了她生前最喜欢的地方。}
\BookParagraph{姥姥白发人送走了黑发人，就每年将收到的无数五颜六色的鲜花和贺年片，置于钢琴之上。琴多年没有人弹，却成了花的海洋。我出国了，姥姥每数月即修书一封，每每说起妈妈和她的钢琴，都笔锋生涩，话不成句。妈妈去世十二年后的一天，年迈的姥姥也走完了她最后的一程。那花丛中的钢琴则不知所终。身在异国的我，俗务缠身，也无缘再回姥姥家的客厅一趟，最后再在当年的琴上弹几个不成曲调的音符。}
\BookParagraph{女儿三岁了，我又想起了钢琴。倘徉于雅马哈钢琴行之中，从五十二寸的立式 U3 至六尺一的三角 C3，几千至几万美金不等。房子客厅边的一室特意空着，虚左以待，等着我梦寐以求的三角钢琴。爱人问，家里又没人会弹钢琴，买这么贵的琴干什么？我也犹豫着，翻开那满纸高低音谱号的谱子，想起妈妈的手和姥姥的信，或许今生与琴声是无缘再相见了。}
\BookDate{2005 年 6 月，多伦多}
\BookArticle{同　学}{同学}
\BookParagraph{刚到清华上学的时候，喜欢攀比学习。当然总知道「山外有山」的道理，也知道要经常比下有余而不比上不足，但年轻气盛，总要突破那高考文化强大的惯性，才能真正避实就虚。无意发现排学号的算法是同时「参考」高考分数和生源地区，来自北京的同学都在班上排到前列，我的更是排到第一，比北京第二的一个郑同学多了三分，不禁得意溢于言表，好像中了649。}
\BookParagraph{这一切在一两个学期之后迅速变成了美好的回忆，因为两个原因。第一是知道了很多牛人压根就没有参加高考，直接保送清华了，还随便爱上什么系就上什么系。第二知道了这些人实在不是吃素的，凡是跟数学有关的课程成绩都根本无法追随其后，就像玩赛车游戏时发现自己的车实在太烂，技术再好都没戏。赛车比赛变成了观赏性的比赛项目，看看这些牛人之间谁赢谁出局。当然这些牛人考试都接近满分，完全分不出高低来。}
\BookParagraph{这时江湖上开始传说一个一班的同学的轶事，他的名字叫廖成。传说中的廖成在众多的牛人中独领风骚，在众所周知的高难度课程《组合数学》中一学期从来不去上课，不记笔记，不上自习，就凭考前突击一下，立马满分。现在想想当然很可能他高中的时候把组合数学都学得差不多了，但是我想即使我也有这样的条件，也没有这样的魄力敢一学期不上课，尤其是在也没什么别的重要的事要干而且其他同学都不但去上课还提前占位子的情况下。}
\BookParagraph{真正认识廖成，还是在大四上学期准备 GRE 的时候。那会儿一门心思上海报贴到清华各个角落的新东方，觉得宿舍里天天打牌的人太多，就用打工的钱和廖成还有另外一个五班的同学一起到十四号宿舍楼租了个房间，正好我们三个都考大四下学期九四年四月的GRE。虽然租了房间，没怎么看到廖成好好学习来着，天天就我在那里笨鸟先飞。一来二去地跟廖成逐渐熟悉起来，晚上睡觉前有时也躺床上海阔天空地神侃一通，不再觉得他是遥不可及江湖传说中的人物。}
\BookParagraph{廖成长得挺帅，个头也挺高，但戴一副至少八百度的近视眼镜，好像没有女朋友。就像反着透过眼镜看人一样，廖成的脾气怎么也琢磨不透，有时候高兴了就侃到夜里两三点，但又有时不高兴的时候，一整天连一句招呼都不打。卧谈会的主要内容是他的高中生活，永远是各种高中时他在四川的吃喝玩乐，最神奇的是他的故事里从来不落俗套，没有提及一个女孩子，就好像他从来既不学习也不跟女孩套磁一样。和他自己讲的故事里无所顾忌地吃喝玩乐相反，他为人很绅士风度，每次我女朋友来找我，他总是彬彬有礼地接待她。他也相当风趣，GRE 单词里学到了一个 libido 的新词，立刻大喜过望，觉得用来当我的外号那是绝配，以致后来好长一段时间都非常顺口地叫我 bido，听着像是什么奢侈护肤品的品牌广告。}
\BookParagraph{一转眼九四年的春节到了，我抓紧时间利用春节多做几套真题，廖成却岿然不动，回家过春节去了。春节过完了也二月下旬了，廖成就好像刚刚回过味来，发现四月九号 GRE 就要开考。他拿了套新东方的真题一做，verbal 部分才四百多分。于是开始像换了一个人，成天从早到晚就看到他在房间里学习刚刚出版的俞敏洪的红宝书。那会儿的红宝书大概五六千个单词，他刚开始基本一个不认识。他一天学习十几个小时，进度是一天新背五百，复习一千五，方法是说英文答中文，说中文答英文，所有俞敏洪发明的助记法也一并背熟。比如你说「charisma」，他就立马反应「领袖魅力」和「China rises Mao」。就这样我眼看着他在两周之内verbal真题从四百多分揠苗助长到六百多分，真正见识了一回传说中大牛的「领袖魅力」。}
\BookParagraph{九五年八月我们俩同时出国了，我去了 UIUC，他去了普林斯顿。他是出国的同学里第一个有自己的个人主页的，小 icon 一个个搞得非常 fancy ，害得我们出了国的同学后来设计的个人主页都跟他那个长得差不多。看样子他在普林斯顿的实验室工作得很刻苦，我 finger 他的 Sun workstation，他永远不是 telnet 在上面挂着，就是在 console，idle time 不超过几分钟。后来我忙着赶论文，论文拒了又忙着修改；忙着结婚，结了婚又忙着过日子，联系越来越少了。只是还是经常到他的主页上看看是不是又多了一篇重量级会议的论文，看看他有没有「为清华争光」。}
\BookParagraph{零零年七月我刚到多伦多大学的时候，有一次晚上在办公室忙里偷闲给他发了一封惯例的 email，不外乎是「Haven’t heard from you for ages, how are you doing recently?」云云。他的email 回应快得出奇，上来就一句惊天动地的「heng! how dare u say this!  it’s u who mysteriously disappeared into the thin air!!!」 我跟他说我在多伦多大学开始工作了，他回答「so, Professor Li now, ehh? :-)))  sounds really nice.」我当场给他又回了封email，他在几分钟之内回答：「anyway, why u still stick to ur computer at this hr of night?  u’r not a grad student any more, go relax!  hv a life!!! :-)」}
\BookParagraph{零一年一月元旦刚过，北京第二的那个郑同学第一时间打电话到我办公室，告诉我廖成已经不幸去世的消息。廖成几个月前搬了一次家自己一个人单住，最近一段时间得了感冒，但可能没有及时让医生开西药治疗，等他的朋友发现找不到他去他家敲门时，他已经在家去世多时。廖成七三年的，去世的时候只有二十七岁。他去世的时候，博士论文还没有写完，据说是还差最后一章。}
\BookParagraph{廖成去世已经十年了，在办公室夜深人静的时候，我经常会莫名其妙地想起他来。想是否应该听他的话，go relax，have a life，再外加三个惊叹号和一个笑脸。时间长了，很多原本清晰的记忆逐渐地模糊淡忘了，也有很多年轻的希望逐渐变得渺茫。曲终人散的时候，特别想回到十四号宿舍楼的房间里，再听他讲一遍四川吃喝玩乐的故事，然后不忘了问他那里有没有能让人一见钟情的女孩子。}
\BookDate{2011 年 5 月，温哥华}
\BookArticle{语　言}{语言}
\BookParagraph{语言是会变的。像个小姑娘，越变越时髦，等到变得差点儿都不认识了的时候，看着她长大的人也许就老了。小时候去北京展览馆附近的莫斯科餐厅吃完饭到马路对面的出租汽车站乘出租车的时候，还未曾流行香港的「的士」一说。后来北京一块钱一位招手上车就近下车的面包车横行，全都改叫「面的」。坐出租车赶上香港的时髦变成了「打的」或「打车」，公共汽车也港味地演变成了\BookLineBreak{}「大巴」。本来还私下庆幸「没有」一词一时没来得及变成香港自创的「有」少两横的新字，近来似乎也无法幸免跳楼被演变成「木有」的命运。}
\BookParagraph{能否适应激流勇进的语言，变成了是否能够与时俱进的风向标。即便是像我这样对语言有一种与生俱来的兴趣的人，跟这个小姑娘几个月不见面，都难免刮目相看。看到新鲜的语言大珠小珠落玉盘，开始还只是被香港传染，继而受到拼音输入法的影响，再后来则因为手机短信输入的便捷，连标点都变成了空格。不仅仅是粉丝淡定神马童鞋泪奔坑爹等各路新词长江后浪推前浪，而且风水轮流转，老式的词儿像 Bob Dylan 的摇滚和王朔的小说，大浪淘沙，没人再敢跟饭馆把服务员招呼成小姐，问路的时候把大龄青年称呼成同志了。直到把所有情人节表达感情的语言全盘演进成 1314 的数字的时候，语言就像海风里的帆船，从模糊的记忆里呼啸而来，在夕阳中绝尘而去。}
\BookParagraph{中学的时候刚学会怎么写作文的时候，常常写点什么。影评书评也好，言情武侠也行，实在图个物尽其用之时，就写情书给那谁谁寄去。在那个护城河边儿上遛弯骑着自行车上学吃着西瓜高考继而打食堂的饺子的年代，从来没有想到会有这么一天，会「淡定」地把语言的强势丧失殆尽，会「out」成连打酱油都要百度一下的尴尬局面。愿赌服输之余，不知道等到女儿从 Webkinz 长大到 Justin Bieber 的时候，是否会有兴趣看看《收获》的小说，是否会回到物是人非的护城河，读读王安忆、陈村、余华和王蒙等人在没有网络的年代出版的文字，感悟让我怦然心动的语言的境界。}
\BookDate{2011 年 7 月，夏威夷}
\BookArticle{早　恋}{早恋}
\BookParagraph{高一下学期，我喜欢上了她。}
\BookParagraph{那年有个叫肖复兴的写了本小说，题目叫《早恋》。小说的情节早就不记得了，就记着当年忘了谁看了小说第一句话，说：「高二（5）班」？这也叫早恋？}
\BookParagraph{十年以后在UIUC褚明宇当主席的那届学生会组织看过一个电影，叫《爱情麻辣烫》。开篇的一段叫「声音」，讲的就是高中的早恋。穿着白色连衣裙的高圆圆在学校的操场转悠的那一刻，让我想起了她。}
\BookParagraph{她那会儿很喜欢我的文字。语文老师要求每周在自己的笔记本上随便找个主题写一篇文章，叫做「练笔」。写完了就在课间互相传看。高三时我写的练笔全让别的同学借走了，后来就再也没要回来，只记得写过最好的是一篇《最后的贵族》的影评。她也喜欢写练笔，字写得不如我圆润，但是绝对工整。我们似乎所有的共同语言都凝结在这些文字中，在她家午后的阳光下，感觉很暖和。}
\BookParagraph{她知道我喜欢她。但是她说，我们还是做朋友吧。}
\BookParagraph{高三毕业，她考上了北京一所财经学院。我后来去农学院找过其他同学，去轻工业学院找过友谊宿舍，却不知道为什么，没有再去找过她。}
\BookParagraph{十四年以后，我给一个很久没有联系过的信箱发了一封信，说我计划去香港开会。她竟然回信了。}
\BookParagraph{这绝对是一个意外，一个挺离奇的意外。感觉就好像二十多年前我坐飞机去北京，跟一个邻座的陌生女孩儿兴高采烈地聊了一路，十年后她突然给我来了一封电子邮件——那封邮件我没有回复。}
\BookParagraph{她在上海，但是她说正好届时要到公司在深圳的分部出差，我们可以在深圳见面好好聊一聊。我从香港坐摆渡抵达蛇口口岸，办好了一份深圳特区的落地签证后，终于打到一辆出租车，一个多小时后才到达罗湖她的工作地点。}
\BookParagraph{十四年后第一次见面，都说对方没怎么变。跟当年在她家一样那么聊得来，没觉得拘谨，一点儿也不生分。她笑起来还是那么美丽，跟我保存下来的当年唯一一张照片里一样。}
\BookParagraph{我上出租车离开时说，希望我们十四年之内还会再见。她说哪儿会再等那么长的时间。}
\BookParagraph{如果说这是个不成文的约定，那十四年之约快要到了。几年前她在微信上加我，说她还在上海。那以后我竟然再也没有碰到过去上海的机会。说不定哪一天，我要专程买好去上海的机票去找她聊天，然后问她哎你还记得高中时候有一电影叫《最后的贵族》吗？没错儿我那会儿还写过一篇练笔呢。}
\BookParagraph{但也说不定我到达浦东机场等出租车时给她发了微信，却再也等不来她的回复。说不定除了「哎你孩子钢琴考几级了」，我们已经无话可说。}
\BookDate{2017 年 2 月，多伦多}
\BookArticle{马拉松}{马拉松}
\BookParagraph{多伦多马拉松后半程相当冷清，和前半程欢呼雀跃的景象大相迳庭。跑到29公里左右时，路边看到一位中国人，孤伶伶一个人站着，在给每一位参赛选手加油。}
\BookParagraph{跑近了才发现，这个人怎么这么面熟啊，绝对在哪里见过。是不是我的研究生课上，总是坐在第一排课后还爱提问的那个男生？}
\BookParagraph{他看到我，明显非常高兴。在跑过去的那几秒钟里，我朝他微笑招手致意，甚至没来得及停下来聊上几句。}
\BookParagraph{平常的日子里，真的有那么些极少的瞬间，会有一种似曾相识的感觉。就像前一阵子去北郊的超市买菜，迎面碰上一个女生，觉得特别眼熟。停下来跟她聊了一小会儿才总算想起来，五年前常去的一家学校附近的日本寿司店，她在那里当服务员。熟悉了以后每次去吃饭，都有一种回到家里的感觉。最后她说快要辞职了，还特意互留了手机号码。}
\BookParagraph{再次见到她时，聊了几句她说，回头常联系啊，反正有手机号码。我努力地回忆了一下，已经想不起来把她的手机号码放在何处了。}
\BookParagraph{跑完马拉松全程，我在家睡着了，有点儿发烧。迷迷糊糊地，又想起路上碰到的这个人。茫茫人海中总是有人会就像这样和你擦肩而过，也许很快就想不起来他了，也许曾经见过几年，每次见到时聊过几句。时间一长，就把他们忘记了。直到重逢时感觉似曾相识的那一刻，发现对方还记得自己，竟然会有一点儿感动。}
\BookParagraph{睡醒了打开微信，看到一个群里的发言：「Saw Prof. Li in the waterfront marathon」，外加一个笑脸。点开他的头像，里面是一个跑步的人。我忽发奇想，破例点了「添加到通讯录」。}
\BookParagraph{微信随即提示：「由于对方的隐私设置，你无法通过群聊将其添加至通讯录。」}
\BookDate{2017 年 10 月，多伦多}
\BookArticle{同　事}{同事}
\BookParagraph{Bruce Francis 是前几年从系里退休的一位教授，今天去世，享年七十一岁。我和他认识十几年来，无论是学术研究，还是待人接物，他都是一个真正的学者（true academic）。}
\BookParagraph{我记得一次研究生硕士论文答辩会上，他问了半个小时的所有问题，都集中在论文第二页上的「Definition 1」。这仅仅是全篇的第一个定义，连模型描述都没有开始，更谈不上后来的定理证明。他问的所有问题，并不是要为难学生，而是确实想彻底明白这个定义的要义所在。}
\BookParagraph{还有一次我和他聊天，聊到在海外的中国人的姓名。他坚持认为，中国人无论走到哪里，都应该把姓放在名的前面，而不是颠倒过来。要让外国人去学习和适应中国人的文化习惯。}
\BookParagraph{正如系里一位和他合作多年的教授所说：}
\BookParagraph{「Bruce was a courageous intellectual, never taking for granted mainstream approaches and assumptions, always questioning their validity in search for the essence of each problem.」}
\BookParagraph{这个随波逐流的时代，也许我们正需要这种「打破砂锅问到底」和「挑战约定俗成」的精神。}
\BookParagraph{另一位同事对 Bruce Francis 则是这样评价的：}
\BookParagraph{「Bruce was a model colleague: a brilliant researcher, who wrote with clarity, taught with penetrating insight, and inspired everyone he met.  He faced the challenges of his illness with candor, courage, and grace. I will miss him.」}
\BookParagraph{短短一句「wrote with clarity, taught with penetrating insight」，写文章清晰易读，教书则直指人心，就像一把精良的瑞士军刀，给人一种豁然开朗的感觉。无论是作为一个学者还是作为一个教师，俱臻登峰造极的境界。}
\BookParagraph{这让我想起一句不知道谁说的名言：「教育不是往桶里灌水，而是将火把点燃。」}
\BookParagraph{不知道为什么，我竟然有些羡慕他。无论是正直、勇气还是优雅，都是难得的优良品质。而我自己，估计连几句纪念他的评价都写不了这么令人动容。}
\BookParagraph{What a life!}
\BookDate{2018 年 3 月，多伦多}
\BookArticle{成事在天}{成事在天}
\BookParagraph{李葆春老师，您好！关注您很久了，一直很喜欢您写的微博，以及回答微博时的文风。所以，今天犹豫了一下，斗胆跟您请教个问题，如有冒昧，敬请谅解。}
\BookParagraph{我是一个内心状态极其容易紧张、收缩和不舒展的这么一个人。融入新集体时总是遇到一些困难，很难自己独立地去解决，以一种大家都愉悦的方式来切入和收尾。每次都会把问题搞得很僵，乃至于接近于争吵。虽然最后目的是达到了，但心里难免龃龉。}
\BookParagraph{难道事情一定要搞成这个样子吗？人为了达到某种目的，例如职位的提升，亦或是收入的增加，一定要不惜一切手段吗？迫于「同侪压力」而不得不去做一些事情，值得吗？如果不做的话，那有一天万一自己变平凡了，被遗忘了，怎么办啊？内心好纠结啊，还望您能抽空为我答疑解惑，不吝赐教。不胜感激！（问题来自微博私信）}
\BookParagraph{你的问题让我想起了 Jennifer。}
\BookParagraph{我是二十年前在加州硅谷开会时认识的她。那是一个加州常见的那种暖洋洋的下午，Jennifer是一个台湾人，特别爱说台湾国语。一个北京人和一个台湾人第一次见面就像多年的朋友，聊得都快忘了时间。我当时还是一名博士生，她则刚当上教授没几年。}
\BookParagraph{Jennifer自视颇高，经常一针见血地批评我们组的一些同学的学术水平，认为他们的工作既没有深度，又谈不上创新。按现在时髦的语言，Jennifer 喜欢当面「怼」人，直来直去。}
\BookParagraph{谁知道她自己第二天在会议上做学术报告时全线告急，各种针锋相对的提问蜂拥而至。她在讲台上招架不住了，后来只有反复说最后一句话的份儿：「I am here to learn from you.」那激烈交锋的场面，完全不顾同行之谊，跟现在「锤黑」和「锤粉」在网上吵架也差不了多少。}
\BookParagraph{后来我毕业了，她也去了UIUC。我们联系得不多，只有开会时才能见到她。她犀利的风格毫无收敛的迹象，有一次和我打赌，说一位 UIUC 新来的教授肯定拿不到终身教职，因为她无法认同这位教授的学术水平。她事无巨细，事必躬亲，每次开会时住的酒店都要去Priceline上竞拍，似乎找便宜的酒店和写优秀的论文一样，一切都要完美无缺。}
\BookParagraph{一次我去UIUC开会，她还记得和我打赌赌输了的陈年旧账，愿赌服输，请我吃饭。席间她说起我导师的很多学生都去学术界当了教授，羡慕之情溢于言表。她说她自己的学生水平也不低，但是找学术界的工作很不顺利，她已经尽力而为在帮助他们去争取，还是差强人意。我安慰她说，谋事在人，成事在天，学生自己的事儿，就任由他们自己去折腾吧。}
\BookParagraph{和有的教授得绝症后到处演讲声名远扬不同，Jennifer知道自己的病情后，还在夜以继日地工作。她去世前一周还在认真地教一门研究生课程；前几天还在处理邮件，为学生安排将来的导师。Jennifer是单亲母亲，十年前她去世时孩子十几岁，据说学习成绩出色。二十年前的那一个下午之后，我再也没有机会和她畅聊过，但是我却一直把她当作我最好的老师。}
\BookParagraph{谋事在人，成事在天。无论工作中是否能心想事成，是否能完美无缺，尽量多和自己的家人和子女在一起吧。因为有了他们，你就不会「平凡」，也一定会有人一直记得你。}
\BookDate{2018 年 5 月，西安}
\BookArticle{乡　愁}{乡愁}
\BookParagraph{李老师，您好！每次从家回到工作的城市都要适应几天，止不住的伤感。请问李老师，您怎么理解「乡愁」？您是如何面对并排解这种愁苦的？谢谢。（问题来自微博私信）}
\BookParagraph{你的问题让我想起了北大燕南园。我的童年和少年大多数时间都是在北大燕南园的姥姥家度过的。}
\BookParagraph{燕南园里绿树成荫，尤其是夏天的午后，蝉鸣可以吵得睡不着午觉。一幢幢房子风格各异，很多房子轻巧地隐藏在树林里，不仔细找，不一定能注意得到。燕南园的地势比其他地方高出一点儿，所以有一个大下坡和一个小下坡出园。大下坡出来是校医院，小下坡出来是图书馆，过了图书馆就是未名湖了。还有一个小门，出来就是北大三角地的日用品商店，在那里我买了小学第一支熊猫钢笔。}
\BookParagraph{姥姥家在燕南园正中间。前院挺大的，夏天野草会肆无忌惮地生长，姥姥经常叫我到院子里拔草，却怎么拔也似乎拔不完。走进姥姥家，客厅挂着的是梁启超的集句联。上联是「人在画桥西，冷香飞上诗句」，下联是「酒醒明月下，梦魂欲渡苍茫」，笔墨平静端正，气势如虹。再向前走是书房和饭厅，一面墙的书架上，摆满了线装书。经过厨房，到达一个小巧玲珑的后院。我高三时就在对面的一间小房间里紧锣密鼓地复习准备高考。}
\BookParagraph{就像鲁迅在《从百草园到三味书屋》里写的那样，我想我们每一个人都有自己少年时代的「百草园」吧。你的百草园也许是一所机关的大院儿，一个农贸服装市场，一条北京的胡同，或者上海的弄堂。想回家，多半儿是想回到亲切的儿时记忆里；恋旧情结，也多半儿是在迷恋那时自己最爱的吃喝玩乐。无论是羊肉泡馍还是北京烤鸭，是南门儿外的长征食堂还是动物园的莫斯科餐厅，是小豆冰棍儿还是北冰洋汽水，总有那么几样儿东西，割舍不下。}
\BookParagraph{可是我不会经常想起燕南园。因为我发现，在一个陌生的城市住的时间长了，久而久之熟悉了起来，离开了一样会想念。「四海为家」这个词儿，大概形容的就是我这种到处漂来漂去，仍然不以为忤的「游牧民族」。「家」变成了一个随着时间慢慢改变的地方：少年时的家在北大，青年时在美国一个叫尚佩恩的小镇，而中年时则在多伦多这个北美二线城市。}
\BookParagraph{不过，我这样的「四海为家」一点儿也不稀奇。原因很简单：最亲近的人在哪里，哪里就是自己的家。说到底，北美的城市语言交流没有问题，很容易有家的感觉。我更佩服另一种人，竟然可以逆潮流而动，在不易旅行的年代，常年定居在举目无亲的地方。}
\BookParagraph{比如有一个名字叫「寒春」的美国人，早年是一个物理学家，曾参与过曼哈顿计划。她四九年和丈夫「阳早」在革命圣地延安结婚。新中国成立后的几十年长期定居在中国，养了一辈子奶牛，是这方面当之无愧的专家，设计建造了中国第一座机械化养牛场。至于她的英文名叫 Joan Hinton，和多伦多大学著名深度学习专家Geoffrey Hinton 教授是同一家族，则是后话了。}
\BookParagraph{我近年来很少再去北大燕南园。去年夏天去的时候，燕南园杂草丛生，少了蝉鸣，多了不少蚊虫叮咬。姥姥家已经改成了学校工程学院的办公室，隔着窗偷看客厅里的摆设，像是一个会议室。这时候「家」的概念，大概早已成了一种遥远的记忆。记忆永远是美好的，里面有熊猫钢笔，有上海手表，还有父亲留给我的永久二八自行车。我骑车带着她，从清华到北大燕南园，停在家门口。}
\BookParagraph{好好珍惜现在吧，还能回到少年时代自己记忆里的家。「酒醒明月下」之时，时过境迁，我小时候的家早就芳华已逝，面目全非。}
\BookDate{2018 年 7 月，西安}
\BookArticle{老毕}{老毕}
\BookParagraph{在这万众一心地关注每天又新增多少确诊案例的时局之中，不知道为什么，我竟然想到了老毕。}
\BookParagraph{也许是因为，再过几天，就是老毕去世一周年了。}
\BookParagraph{老毕是我在清华的大学同学。我们这些年的交情，既谈不上是「莫逆之交」，也不至于仅仅是「点头之交」，我想大概可以算是「合影之交」吧。我们的学术研究方向相仿，算是「小同行」，所以因为开会，每隔一两年总有机会见上一面。老毕有一个良好的习惯：每次开会见到我，都要跟我合影留念。最后一次见到老毕是他来多伦多参加一个学术会议，正好是我牵头组织的。我翻开当年的相册，轻而易举地找到了我们的合影，是与他同来的一个学生帮我们用手机拍的。照片中的老毕神采奕奕，笑容可掬。}
\BookParagraph{老毕大学时就有点儿胖。他在我的印象中永远显得和蔼可亲，大概是因为他爱笑的缘故 —— 微胖而又爱笑的人容易让人觉得亲切。老毕又是当年我们系的党支部书记，他慈祥的笑容中自然地多了一点儿领袖魅力。当年复习准备 GRE 时，俞敏洪在新东方给我们讲「charisma」是「中国升起了毛泽东」的那一刻，我正好就想起了老毕。}
\BookParagraph{老毕和我都可以算是土生土长的北京人，但是除了这一条，我哪一方面都没法儿和老毕相比。老毕是北京八一中学的高材生，品学兼优，高中时就入了党，并保送了清华大学计算机系。如果没有入选奥赛，保送清华当时非常难 —— 我所在的中学也号称算是全国重点中学，当年一个保送清华的名额都没有。大学毕业后，老毕留校读博士学位，仅用了四年就读完了，可以算是「书生意气，挥斥方遒」的样板。}
\BookParagraph{大学有一年，我在《马克思主义哲学》这门课中学得很不错，最后的学期成绩接近满分。这说明与其学工科，我其实早就应该去学文科 —— 大一的《中国革命史》我也相当擅长。可以看得出来，老毕很欣赏我，当年经常鼓励我写入党申请。有一次在三教的一个教室又聊起来这事儿，我很不好意思地跟老毕说，我觉得我自己还不够格，很多的东西还只知道一点儿皮毛，需要再进一步努力提高自己的认识水平。}
\BookParagraph{就像有的人性格外向有的人则内向一样，我觉得老毕一直是认真负责的「入世」性格，而我则说得好听点儿叫「出世」，说得不好听一点儿则是一个做一天和尚撞一天钟的「逃兵」。老毕博士毕业在美国进修了几年，继而顺理成章地在清华任了教职。自此，他在清华培养了一批又一批的学生，直到去世之前一两个月还在工作。按一篇纪念他的文章中的话说，「他是我目前见过对科研最拼命的人。」}
\BookParagraph{老毕去世之后，大学的同学们以他的家庭困难、孩子长期需要照顾为名义，发起了募捐。而纪念他的文章，讲述的故事则清一色的是他直到最后一刻，还坚守在教学和科研的岗位上，还在关心着论文是否能及时修改好，是否能投到一流的国际会议上发表。}
\BookParagraph{老毕和我是「合影之交」，但我们走过的路则大相径庭 —— 是一位优秀的党员和一个早早跑到国外的逃兵的区别。老毕去世转眼快一年了。时间飞快地流逝，时局则尽是变数，像万花筒一般，每一天都不同，令人目不暇接。一两年后，也许不会再有人想起老毕来，就好像那时也不会再有人想起现在眼前曾发生过的事儿一样。那时老毕大概早已成了百度百科的一个条目，而他的名字，渐渐地成为了历史百科全书中的一条脚注。很多年后他的子女和他曾影响过的学生依稀可以想起来的，不是他拼命写出来的科研论文，而多半儿会是他当年显得有点儿慈祥又有一点儿领袖魅力的笑容。}
\BookParagraph{老毕和我是「合影之交」。他是一位「够格」的党员，也是一个有益于人民的人，值得我向他学习。他让我想起了米特·罗姆尼前几天在弹劾案投票前演讲中所说过的一句话：「我会告诉我的子女和他们的子女们，我已经尽了我的最大努力，没有辜负我的祖国对我的期望。」}
\BookDate{2020 年 2 月，多伦多}
\BookArticle{时代}{时代}
\BookParagraph{最近我发现自己一没留神儿 —— 或者是没刷抖音 —— 又赶不上时代了。}
\BookParagraph{有朋友告诉我她有个朋友，不知道怎么，吃饭聊天儿开始用网络流行语言了。}
\BookParagraph{我有点儿好奇，网络流行语言是如何活学活用拿来吃饭的时候聊天儿的。是「今儿的火锅绝绝子了」吗？还是「太好吃了，拴 Q 栓 Q」或者「完了，芭比 Q 了」？}
\BookParagraph{她说都不是。人家是这么用的：}
\BookParagraph{「咱就是说能不能别放这么多油？」}
\BookParagraph{「这么辣，我真的会谢。」}
\BookParagraph{「我有个同事啥啥啥，我整个就是一个无语给到。」}
\BookParagraph{我一听，果然比我想象的更自然，融会贯通，炉火纯青。}
\BookParagraph{「好好学习，天天向上」之时，想起来十年多前有一次去夏威夷开会，好像是看到了「神马都是浮云」这个短语，开始浮想联翩，写了 \href{https://baochun.ca/post/芳华已逝面目全非/语言/}{如下的文字}，题目就叫「语言」。}
\BookDate{2022 年 7 月，多伦多}
\BookArticle{Brittle}{Brittle}
\BookParagraph{正在大唐不夜城游览的间隙，手机推送了一封系主任发给全系的邮件：系里和我共事过二十年的 Cristiana Amza 教授因病去世。}
\BookParagraph{Cristiana 比我晚来三年，年龄应该比我小一些。她是一位分布式系统领域的学者，十几年前她在一个国际会议拿到最佳论文奖，我们在 CN 塔顶层的旋转餐厅庆祝过。直到上学期，她还在我领导的一个委员会中共事过，线上开过几次会议。}
\BookParagraph{二十年来，我们聊得最多的话题是子女教育。Cristiana 是罗马尼亚人，她告诉我，罗马尼亚和中国一样，对子女教育极其重视。我的女儿和她的女儿是多伦多同一所女校毕业 —— 当年孩子还小的时候，我跟她「安利」说，据我的研究成果，这所学校是多伦多最好的女校。}
\BookParagraph{打开系主任的邮件时，正值大唐不夜城每天晚上九点的音乐喷泉开场。初秋微凉的夜色中，男女老少高举着手机，拍摄着这个号称亚洲最大的喷泉。我想，生命真的就像五彩的音乐喷泉一样，虽然五光十色，令人目不暇接，但是终究有曲终人散的一天。}
\BookParagraph{「Life is brittle.」我把邮件转发给组里的学生，加入了这句评语。说不定会有学生上过她的课，还有些依稀的记忆。}
\BookParagraph{邮件没有收到回复。}
\BookDate{2023 年 9 月，多伦多}
\BookArticle{家}{家}
\BookParagraph{今天凌晨一点一次线上会议开始前的闲聊中，系主任恰巧问我，「extended —— not immediate —— family」在哪里居住。}
\BookParagraph{那个点儿有点儿困，脑子转得有点儿慢。我反应了一秒钟，才明白她是在问我的父母。}
\BookParagraph{我回答她：我的父亲在北京。恰巧我已经订好了夕发朝至的动车卧铺，过几天从香港去北京看望他。}
\BookParagraph{好像也就是一瞬间的事儿 —— 父亲已经八十五岁了。}
\BookParagraph{而我的笔下，却极少写到父亲。}
\BookParagraph{我记得唯一写到父亲的一回，是大学三年级清华大学作文竞赛。那次作文竞赛是在主楼后厅举办的，几百人坐在一起，每个学生发一沓四百字的稿纸，限时三小时，爱写多少写多少。}
\BookParagraph{刚一开始，主持的老师回过头去，在黑板上写下了这次竞赛的题目：}
\BookParagraph{家。}
\BookParagraph{没错儿，就是这一个字。}
\BookParagraph{当时我写了什么，现在已经完全不记得了。我只记得梦幻般的故事中有一幅画面，其中有父亲的背影，还有他的永久自行车。我长大了去清华上学，曾经骑着这辆自行车，往返于宿舍和教室、北大和清华之间。而父亲的背影隐约而熟悉，有点儿像朱自清先生的散文。}
\BookParagraph{我还清晰地记得，那三小时里，我奋笔疾书，一气呵成，写了九页纸，将近三千六百字。题目就是《家》。}
\BookParagraph{我在那次作文竞赛中得了全校前三名，荣获香港校友奖学金。我想，大概阅卷的老师看到我和父亲的故事，眼眶湿润了吧。}
\BookParagraph{也许是天意吧：那一年，父亲只有五十三岁，和我现在同龄。}
\BookParagraph{我出国已经三十年了。每逢我生日，父亲都会来信，祝生日快乐。}
\BookParagraph{等将来那曲终人散的一天到来时，我会想念他的生日祝福，想念他的自行车，想念他的背影的。}
\BookParagraph{我可能想家了。}
\BookDate{2025 年 7 月，香港}
\BookArticle{酒公墓}{酒公墓}
\BookParagraph{一位朋友前些天提起，想筹备一个活动，聊聊「幸福」这个词儿，问我有没有什么好的建议。}
\BookParagraph{我第一时间想到的，竟然是这位朋友给我推荐的一本书：经济学人 2008 年左右出版的《讣告》（Book of Obituraries），书中收录了两百位各色人等的讣告。我想，幸福这两个字，可能只有在这本书中才能真实地感受到。至少，书中的人，知名也好不知名也罢，都早已仙逝，活着的人看看他们生前的故事，也会油然而生一种幸福的体验。}
\BookParagraph{朋友说：「那会不会有点儿丧？」}
\BookParagraph{我年轻时第一次读到一份感人至深的讣告，还是三十多年前看余秋雨所著的《文化苦旅》。这本书其他的文字都早已不记得了，只记得一篇《酒公墓》，寥寥数笔，记叙的是酒公张先生一生的故事。}
\BookParagraph{搜索了一下，发现百度百科就收录了全文。相隔三十年后，我又看了一遍这篇讣告。看到张先生一听到文中的「我」提及金岳霖先生的逻辑思想，马上就「抖抖索索地把我的手紧紧拉住，说自己将不久人世」，希望「我」能替他写一方小字碑文。}
\BookQuote{我问他小字碑文该如何写，他神情严肃地斟酌吟哦了一番，慢吞吞地口述起来：}
\BookQuote{“酒公张先生，不知籍贯，不知名号，亦不知其祖宗世谱，只知其身后无嗣，孑然一人。少习西学，长而废弃，颠沛流荡，投靠无门。一身弱骨，或踟蹰于文士雅集，或颤慑于强人恶手，或惊恐于新世问诘，或惶愧于幼者哄笑，栖栖遑遑，了无定夺。释儒道皆无深缘，真善美尽数失落，终以浊酒、败墨、残肢、墓碑、编织老境。一生无甚德守，亦无甚恶行，耄年回首，每叹枉掷如许粟麦菜蔬，徒费孜孜攻读、矻矻苦吟。呜呼！故国神州，莘莘学子，愿如此潦倒颓败者，唯张先生一人。”}
\BookQuote{述毕，老泪纵横。}
\BookParagraph{《文化苦旅》出版后已经三十多年，余秋雨先生都已八十高龄。但是重读他的文字，还是一如当年，鲜活而令人动容：}
\BookQuote{前年深秋，我回家乡游玩，被满山漂亮的书法惊呆。了解了张先生的身世后，我又一次上山在墓碑间徘徊。我想，这位半个多世纪前的逻辑救国论者，是用一种最潦倒、最别致的方式，让生命占据了一座小山。他平生未能用自己的学问征服过任何一个人，只能用一枝毛笔，在中国传之千年的毛笔，把离开这个世界的人慰抚一番。可怜被他慰抚的人，既不懂逻辑，也不懂书法，于是，连墓碑上的书法，也无限寂寞。}
\BookParagraph{「确实会有点儿。」我说。}
\BookDate{2026 年 2 月，多伦多}
\BookPart{好好学习 天天向上}{好\\好\\学\\习\\[18pt]天\\天\\向\\上}
\BookArticle{看　书}{看书}
\BookParagraph{褚明宇的「向上移动」里一人问了一个问题：大家平时还看纸质书吗？回答众说纷纭，大意是看书的好处是有仪式感和满足感。我倒是觉得，书最大的好处是可以借给别人。钱钟书在《围城》里说：「一借一还，一本书可以做两次接触的借口，而且不着痕\BookLineBreak{}迹。」就算是人家不还你，说不定多年后看见这本书，还能忽地想起你来。}
\BookParagraph{很少有其他东西有这个好处，可以随意借给别人，或者找别人借来看。一个著名的例子是十年前微软设计了一种叫「Zune」的MP3播放器，可以把正在听的歌儿用WiFi「借」给朋友听，但是最多听三遍。乔布斯看了演示说：「It takes forever. By the time you’ve gone through all that, the girl’s got up and left! You’re much better off to take one of your earbuds out and put it in her ear. Then you’re connected with about two feet of headphone cable.」}
\BookParagraph{钱钟书和乔布斯的借法儿不同，但是殊途同归。}
\BookDate{2017 年 6 月，西安}
\BookArticle{我的妹妹}{我的妹妹}
\BookParagraph{父亲还保留着我九岁时在北大附小三年级时的作文，题目叫\BookLineBreak{}《我的妹妹》。第一感觉是这个题目不好写，本来描写人物就很难把握写作时的分寸感，写亲人更甚。看了当年的文字，感觉还行。开头与故事呼应，有一场戏剧性的小高潮，结尾简洁不拖沓。妹妹现在在美国硅谷担任公司的高管，不但可以见到东西就买，而且也算成为了一个「有用的人」。}
\BookParagraph{————}
\BookParagraph{我的妹妹还不到四岁，胖乎乎的小脸上嵌着一双明亮的大眼睛。她活泼可爱，见到东西就新鲜，就要买。我非常喜欢她。}
\BookParagraph{一天早晨，妈妈带我到商店去买圆珠笔。妹妹也要去，我就同意了。妈妈带着我和妹妹来到商店，给我买了圆珠笔。正要走，忽然看见妹妹睁大眼睛，望着妈妈。妈妈看出了她的心思，就说：\BookLineBreak{}「以后再给你买，好不好？」 妹妹点了点头。}
\BookParagraph{一晃几个月过去了。妈妈早就把这件事给忘了。可妹妹还记得清清楚楚。一天，妹妹拉着妈妈来到商店，非要买圆珠笔。妈妈没有办法，只好依了妹妹。}
\BookParagraph{一个星期天的上午，我的表妹一家来看奶奶。带来了一个小钢琴。表妹还不到三岁。表妹刚刚弹了几下，我的妹妹就抢过来左弄右弄。因为她没见过。表妹很不高兴。等表妹走了以后，妈妈批评了她。妹妹哭了。她用手捂着脸，不时从手指缝里偷看，看大家是不是在注意自己。她看大家都不理她，有的在看报纸，有的在写作业，就一人拿着皮球出去玩了。}
\BookParagraph{我的妹妹爱学习，手也很巧。我真希望她长大能成为一个有用的人。}
\BookDate{2017 年 8 月，上海}
\BookArticle{大　V}{大V}
\BookParagraph{和十几位学校领导一起在校长办公室旁边的会议室开了一个正式的工作会议，讨论如何加强多伦多大学和东亚国家和地区的学术合作，尤其是如何提高在中国的品牌价值和影响力。}
\BookParagraph{聊到学校准备申请官方的微信公众号时，我想这个话题我强项啊，正跃跃欲试准备举手发言出谋划策，另外一位教授说：「微博也很重要。我在微博有五万粉丝，经常收到很多学生父母的来信，询问各种问题。」}
\BookParagraph{「大V」面前，从七年前开始认认真真写了好几百篇微博才几千粉丝的我，立马儿无话可说。粉丝的数目就像过日子，运气好的时候涨点儿，运气不好时再跌点儿，「事物的发展是螺旋式上升的过程」，「道路是曲折的，但前途是光明的。」按照现在的趋势，大概这辈子也涨不到五万粉丝。就好像中学里和同班的校花儿学霸一比，再如何发奋图强也无法与其相提并论。}
\BookParagraph{不过转念一想，比上不足，比下有余。郭芳的文字风格返朴归真，比我写得还强不少，还不是只有几百个读者？芸芸众生之中，创业也好，写微博也罢，真正成为「大V」的只有极少数人，和是否努力也没什么关系。}
\BookParagraph{但很多事儿努力一下儿说不定还真能做到。比如在安大略湖上开一条帆船；到德克萨斯养几匹骏马；孩子小时候给她们每天讲一个故事；把多年写成的文字用最好的纸张印一本书，女儿的婚礼上送给她。}
\BookParagraph{我终于举手发言：「I know a thing or two about Weibo as well — I’d happy to help. 」}
\BookDate{2018 年 2 月，多伦多}
\BookArticle{虚拟现实}{虚拟现实}
\BookParagraph{今天一天和几个学生一起批改期末考试的试卷。}
\BookParagraph{我们按惯例找了一间会议室，因为批改试卷是一件需要相当专注的事情。每一道题都需要仔细阅读学生的答案，同时迅速地衡量这个答案到底有多靠谱。}
\BookParagraph{可惜的是，改卷时的那种几乎有点儿虔诚的专注氛围，忽地烟消云散了。}
\BookParagraph{桌上放着一个学生的iPhone，震动了一下，过了一会儿又震动了一下，又过了一会儿再次震动了一下。}
\BookParagraph{也许是微信？要不然就是哪个新闻app的推送？我转念想着，然后不自觉地去看我自己的手机，虽然它并没有震动。}
\BookParagraph{曾几何时，我午夜梦回惊醒，第一反应是去摸我的手机；看今天的天气不是去打开窗帘，而是去找我的手机；走在路上在关注今天的股市，排队等咖啡则在刷新近的微博。手表可以震动，手机更可以震动，来电话可以用电脑接，签字可以用iPad。手机没电了，感到的不是悠然自得，而是深深的无助和无所适从。偶尔有时候漫无目的地抬起头看看周围，发现整个世界都已经沉迷在手机之中。}
\BookParagraph{然而扪心自问，上个月在手机里看过的「大V」里，最精彩的观点是什么？}
\BookParagraph{真的想不起来了。}
\BookParagraph{记忆犹新的那些故事，是电影，是小说，甚至是教材，没有一件是从手机里来的。「我的印象笔记」，却完全没有印象。}
\BookParagraph{手机是一个虚拟的现实世界，像电影《The Matrix》和《Inception》，一切都像极了真实的生活。直到有一天终于发现，它只不过是我们自己认为真实的梦境。呼啸而来，扬长而去，日记本空空如也，什么也没有留下来。}
\BookParagraph{没有手机，照样可以安然入睡；没有手机，照样可以知道时间；没有手机，照样可以看报纸和杂志；没有手机，照样可以学习，可以娱乐，可以开怀大笑，可以痛哭失声。我需要的不过是一块普通的石英手表，几本随身携带的纸质书，一付不需要充电的耳机，还有一辆引擎可以轰然发动的汽车。}
\BookParagraph{今天，我向手机宣战。}
\BookDate{2018 年 5 月，多伦多}
\BookArticle{演　讲}{演讲}
\BookParagraph{我一直觉得，无论是校庆还是毕业典礼，在众目睽睽之下做一个特别励志特别感人的演讲，是一件只有极少数人能做到的事儿。}
\BookParagraph{原因很简单：众口难调。阅历浅的人觉得励志的事情，阅历深的觉得小儿科；阅历深的人觉得感动的事儿，普通人又听不明白。语言过于书面过于正式，会觉得太空泛没有什么印象；过于随意过于市井，又会觉得太油滑。}
\BookParagraph{不过我认为，想做一个大多数人都觉得励志或者感人的演讲，至少要做到两件事儿。}
\BookParagraph{首先，一场激动人心感人至深的演讲，看上去轻描淡写，其实都是经过精心准备的。每一分抑扬顿挫，每一寸起承转合，都需要长时间的多次演练。即使是乔布斯这样久经沙场的演讲高手，在斯坦福大学毕业典礼上的著名演讲之前，还是紧张地准备了好几个月。直到最后他从家出发去斯坦福的路上，还在一遍一遍地练习。}
\BookParagraph{其次，演讲稿的用词越简单越好。从丘吉尔到马丁路德金，历史上最励志的演讲用到的语言，都是相对简单的。比如 J. K. Rowling零八年在哈佛大学毕业典礼上的演讲，就是一个简单到了极点的例子：「As is a tale, so is life: not how long it is, but how good it is, is what matters.」}
\BookParagraph{即使需要引经据典，也最好加以解释。比如，与其说「要立志，立鸿鹄之志」，不如说：「《史记》里有一句话，叫「燕雀安知鸿鹄之志」。这里的「鸿鹄」，是天鹅的意思。我们为自己的未来立志的时候，要向天鹅学习。」}
\BookParagraph{毛主席在著名的演讲《为人民服务》中是这么说的：「人总是要死的，但死的意义有所不同。中国古时候有个文学家叫做司马迁的说过：人固有一死，或重于泰山，或轻于鸿毛。为人民利益而死，就比泰山还重；替法西斯卖力，替剥削人民和压迫人民的人去死，就比鸿毛还轻。」}
\BookParagraph{「... but how good it is, is what matters.」}
\BookParagraph{这才叫气场，才叫水平。}
\BookDate{2018 年 5 月，西安}
\BookArticle{新概念}{新概念}
\BookParagraph{看到微信里有一个「每天20分钟120天新概念英语」学习班的广告：「仅需399元」，如果能每天「打卡」不缺勤就可以全额退款。}
\BookParagraph{这个英语速成班的理念是：全文背诵《新概念英语》一点儿用处也没有，需要用「拆解重构」的方法学习，学习一篇，相当于背诵一百篇。}
\BookParagraph{「不用背的新概念，听说读写全面提升。」}
\BookParagraph{你还别说，我小时候有一两个暑假，曾经认真自学过《新概念英语》。那是八十年代初，英语教材匮乏，我除了学过《新概念英语》，还学过陈琳英语和许国璋英语，仔细研究过十六种时态，从一般现在时到「过去将来完成进行时」（would/should have been doing）。}
\BookParagraph{我无法想象，三十年以后，还有人去学习《新概念英语》。}
\BookParagraph{无论是《新概念英语》还是陈琳许国璋英语，最关键的问题是学英语的理念是完全错误的。学习英语不能靠精读的办法，靠学习语法、强背单词、全文背诵、「拆解重构」的训练来学。而《新概念英语》每篇文章极短，文章之间没有语境的联系，第二册和第三册之间难度跨度过大，作为英语教材，早就该过时了。}
\BookParagraph{而是应该尽量脱离中文语境，靠大量快速地接触英文的办法来学。比如阅读，应该找大量报纸、杂志、小说，看完一遍坚决不看第二遍。单词不懂的可以根据上下文瞎猜，也可以查英英词典，阅读更简单的英文去理解更复杂的词义。}
\BookParagraph{那如果《华尔街日报》、《经济学人》什么的我都看不懂怎么办？}
\BookParagraph{学英语跟学数学学钢琴其实差不多，一口吃不成一个胖子，需要长期坚持，循序渐进。看不懂《经济学人》，可以看青少年小说（Teen fiction）；看不懂青少年小说，还可以看儿童读物。}
\BookParagraph{所以说，「120天每天20分钟」是学不会英语的。}
\BookParagraph{哈哈，写到这儿我都想开一个英语班儿了：「由学贯中西的北美知名大学教授领衔主演，让您在轻松愉快的氛围中享受英文的魅力。」}
\BookParagraph{这不，宣传文案都写得差不多啦。}
\BookDate{2018 年 7 月，多伦多}
\BookArticle{公众号}{公众号}
\BookParagraph{每天花大量的时间在阅读公众号上的文章，值不值？读多了觉得浪费时间，没读又怕错过了什么。（问题来自「向上移动」）}
\BookParagraph{英文有一个缩略语叫「TL; DR」，代表「too long; didn’t read」，中文大概可以译成「中心思想」或者「长话短说」。}
\BookParagraph{TL; DR: 别看微信公众号上的文章。}
\BookParagraph{微信公众号文章的最大优点是免费阅读。英文里常说「没有免费的午餐」，意思是说如果一个产品对用户免费，一定需要用其他方式来收费，对用户来说往往比收费的产品代价更高。}
\BookParagraph{微信公众号主要有两类：市场营销和文艺青年随笔。}
\BookParagraph{市场营销的主要目标是靠吸引人的文章、照片和视频来吸引读者的注意力和时间（attention span）。因为任何人的时间都是非常有限的，所以这一类公众号的目的就是抢占「时间」这种稀缺资源。因为写这类公众号的作者往往是专业从事这项工作的，他们可以穷尽一切可能的方式和一切当下最吸引人的话题去吸引你阅读转发，进而「引爆朋友圈」，抢占更多人的注意力和时间。}
\BookParagraph{文艺青年随笔的写作目标则不一样。它属于自娱自乐型，如果正好有爱好相仿的人来点个赞留个言，锦上添花，心情更加愉快。文艺青年随笔从来不对当前热门话题发表意见，阅读量一般很小，挺容易辨别的。可惜的是，文艺青年一般很难有大量的空闲时间从事长期的无偿创作，所以这类公众号就像当年的博客一样，容易开始却很难坚持，早晚要销声匿迹。}
\BookParagraph{不可否认，微信公众号里确实有少量质量很高的文章，一般都出自文艺青年的随笔。可惜的是，因为高质量的文章比例太低，需要「沙里淘金」的时间太多，从时间成本来看，往往得不偿失。}
\BookParagraph{那平时等地铁这种碎片时间里，不看公众号文章看什么呢？人的记忆往往偏爱长篇情节、故事有张力有跌宕起伏的文字。这么说吧，去年你在微信公众号看过的故事，现在还记得哪篇？跟朋友聊天时，与其靠「哎你看过咪蒙老师的那篇公众号文章吗」开场，不如用「莎士比亚老师写的 Macbeth 里讲了一个故事」、「乔治奥威尔老师在1984里说过」，或者至少「钱钟书先生在围城里有一句名言」来得引人入胜。}
\BookParagraph{当然了，如果你这篇微信公众号文章是别人在朋友圈里转发的，是「金」而不是「沙」的概率几乎为零。}
\BookParagraph{同样道理，我这篇回答也完全没必要看到结尾。因为不看没有错过任何值得记住的东西，而看这篇文字所花的时间，过去了就再也回不来了。}
\BookDate{2017 年 7 月，西安}
\BookArticle{学英语}{学英语}
\BookParagraph{大家怎样看待学习英语的事情？（问题来自「向上移动」）}
\BookParagraph{我在美国上学时有个好朋友。当年他有一个愿望：十年之内英语说得和美国人完全一样。其实说这句话时，他的英语已经非常出色了。出国前他为了准备GRE，当时没有新东方可上，于是自己一个一个词儿把《新英汉词典》背了下来。}
\BookParagraph{他为了练习英语口语，刚到美国就倾其所有买了一台崭新的彩色电视机，每天看电视里的广告。我开始没太明白，为什么看电视要专门去看广告。后来才知道，广告的好处在于它天天重复。他每天跟着广告练习，直到发音和语调都和广告一模一样。二十多年后，我还记得一则卖钻戒的广告词儿：「How could three months of salary last forever?」还有他现场认真跟读的神情。}
\BookParagraph{他用英文写作的时候，喜欢一边想一边嚼口香糖揪头发，常常陷入冥思苦想。他逐字地衡量着每一句话的句式和用法，一遍遍地修改同一个自然段，往往一晚上也写不了几段儿。他写得虽然慢，可一旦定稿就不再修改，所以跟我写得飞快但是经常大改的风格相比，颇有「龟兔赛跑」之妙。}
\BookParagraph{几年以后我们都当了教授，我去美国找他叙旧。走进他家的客厅，发现四壁皆空，只是正中间放着一个纸箱子，上面放了一台电视，正是他当年刚到美国时买的彩色电视机。纸箱子旁边放着一大摞录像带，是他一集一集亲自录下来的「Seinfeld」全集。他每天吃晚饭时，正好重温几集。每一个笑话，每一个包袱，都到了倒背如流的境界。}
\BookParagraph{学习英语没有捷径。无论听、说、读、写，你的付出远远超出常人，早晚可以「和美国人完全一样」。}
\BookDate{2017 年 7 月，西安}
\BookArticle{逆水行舟}{逆水行舟}
\BookParagraph{微信从来不发朋友圈的是什么样的一类人？或者说这种群体是什么心态？是不是跟不参加老同学聚会一样 —— 过得很好或很差？（问题来自「向上移动」）}
\BookParagraph{我就属于极少数从来没有开通过微信朋友圈的那一类人。主要原因是我不喜欢微信朋友圈这个产品的设计。这就好像一个人不喜欢武侠小说，所以从来没有看过一样。}
\BookParagraph{社交媒体平台的设计大体可以分为两类：阅后即焚式（ephe-meral）和永久储存式（permanent）。}
\BookParagraph{我偏爱风格纯粹、设计理念简单的产品。比如微博、知乎、YouTube和Twitter都属于永久储存式的产品，而Snapchat Stories和Instagram Stories则属于阅后即焚式的产品。}
\BookParagraph{而微信朋友圈这个产品，既可以阅后即焚（三天可见），半永久储存（半年可见），也可以永久储存。不仅如此，它还可以「分组可见」，可以「不看他的朋友圈」「不让他看我的朋友圈」，和当年QQ里的「对所有人隐身对他可见」和「对所有人可见对他隐身」有异曲同工之妙。再加上非好友评论不可见，到底我能看到他什么，他又能看到我什么，就显得扑朔迷离。}
\BookParagraph{所以才会有类似「他最近怎么不发朋友圈了，是不是对我屏蔽了」、「到底给不给同事和老板看朋友圈」、「哎这条朋友圈怎么你能看见我看不到啊」甚至「他怎么从来不发朋友圈，是不是混得不行」之类的问题。一个社交产品，如果连不去用它都有让人猜疑和误解的担忧，就有点儿危险了。}
\BookParagraph{所以我一贯的理论是，无论文字还是照片，要不然就完全公开，要不然就不发。而微信朋友圈支持这么多「是否可见」的开关，就是无法打开天窗说亮话，完全公开并永久储存所有文字和照片。}
\BookParagraph{我猜测很多人「看武侠小说」不是因为有多好看，而是因为别人都「看武侠小说」，所以他也要看。就好像别人都买房，所以他也要买一样。}
\BookParagraph{而我更钟爱的是一种逆水行舟式（contrarian）的风格：别人都买房我不一定买，别人都看武侠小说我不一定看，别人双十一都抢购我不一定抢，别人都玩儿朋友圈，我则自始至终「金盆洗手，淡出江湖」。}
\BookDate{2017 年 11 月，多伦多}
\BookArticle{写　作}{写作}
\BookParagraph{李葆春老师，特别喜欢您写的文章，也十分佩服。想请教您如何提升写作水平呢？这是您思维的具体化形式，通过思维锻炼出来的，还是依靠大量阅读，抑或是通过模仿和修改？（问题来自微博私信）}
\BookParagraph{其实很多有关写作的知识，我们小学时都学过，只是后来写文章时大概没有用上。我认为，写东西需要观点犀利鲜明，行文通透连贯，文风简单明了，修改大刀阔斧。}
\BookParagraph{小学时学的「中心思想」，其实说的就是要有鲜明的观点。你想啊，平时接触过这么多文字，真正印象深刻的，其实只有极少数「中心思想」清晰自然、入木三分的文字。有的人喜欢好几件事儿一起写，但互相之间又没有什么明确的逻辑关系，不如拆分成几篇文章来写。开门见山，首尾遥相呼应，虽然是小学生作文的路子，但正是「画龙点睛」式表达观点的好办法。}
\BookParagraph{小学时学的要有「段落大意」，则是行文通透连贯的要素。如果每一个自然段没有一个明确的大意，段落之间没有清晰的衔接关系，文字看起来就会令人百思不得其解。衔接关系最引人入胜的是「转折」，一停一顿之间，有一点儿小时候听故事时「那后来呢」的期待。最不该用的是「并列」，因为我们是在写文章，不是在做「PPT」。}
\BookParagraph{简单的文字，往往读起来反而显得雍容大度，底气十足。一段简单得不能再简单的标语，却最容易慢悠悠地爬到心里，再也不出来。无论是「团结紧张，严肃活泼」，还是「Think Different」，都占尽了「简单」的天时地利。简单的文风到了极致，平铺直叙之间，像是一道江浙菜，娓娓道来而又丝丝入扣，让人感到温暖。}
\BookParagraph{小时候学习《醉翁亭记》，据说欧阳修初稿开篇二十余字，但总觉得不满意，几经删改，最后收到五个字，就是著名的「环滁皆山也。」欧阳修起于呼之欲出、一气呵成，止于苦心孤诣、字斟句酌，文字既鲜活跳脱，又质朴无华，正是写一篇好文章的典范。}
\BookParagraph{一篇文章的初稿就像一幅画儿的铅笔底稿，需要我们离开安静的书桌，到嘈杂的星巴克点一杯咖啡，抽出一支普通的钢笔，坐下来一层一层地润色、渲染、修改。著名网站 Daring Fireball 的作者 John Gruber 曾经说过：}
\BookParagraph{I don’t use my favorite pen — which, of course, has black ink — but instead a pen with red ink. Editing is an angry, bloody act and therefore must be done in red.}
\BookParagraph{没错儿，修改就是一种「angry, bloody act」。对自己的文字刻薄一点儿，大刀阔斧，一针见血。让蘸满红色墨水的笔锋肆意地在纯蓝的字间跳跃，不破不立，像在健身房里挥汗如雨，畅快淋漓。}
\BookDate{2018 年 1 月，多伦多}
\BookArticle{莫名其妙}{莫名其妙}
\BookParagraph{李老师，冒昧打扰下。我最近想开个微信公众号，想尝试写类似「大象公会」这个号上形式的文章，和读者一起做认知上的升级，并借此了解和熟悉公众号管理一类的事。}
\BookParagraph{但是在公众号起名上我遇到了阻碍。我原本是想起「因果警察」，我想的是不判断是非，只追溯因果，这点跟我打算写的东西契合。但是由于该词汇过不了审核，我就开始困惑，一时不知道应该起个什么名字。所以我就很唐突地来问老师，您有没有什么起名字的建议？（问题来自微博私信）}
\BookParagraph{先给你泼一盆冷水啊。微信公众号和微博差不多，阅读量多半儿表现为 Zipf 分布。它是一种幂律分布，极少数公众号阅读量极高，而绝大多数没什么人读。读者寥寥，对写作热情没有什么激励，文字的质量就很难提高，也很难长期坚持下去。这些都和公众号的名字无关。}
\BookParagraph{不过新号开张儿，无论谁都想起一个好听点儿的名字。我认为，与其试图去体现公众号的主题，不如反其道而行之，起一个让潜在的读者莫名其妙的名字，去挑战他们的好奇心。}
\BookParagraph{就拿「大象公会」来说，我乍一看就完全不知道这是什么主题。而这种潜在的好奇反而可能让一位读者点开其中一篇文章，然后或开怀大笑，或感触良多，继而爱不释手。}
\BookParagraph{当然这种「莫名其妙」也还是要有个限度。要「莫名其妙」得似曾相识，而非「莫名其妙」得冷若冰霜。你可以先想想你的读者群大致年龄段，然后找一个这个年龄段的人十几二十岁时耳熟能详的名字。因为每一代人都有他们最钟爱的电影、电视剧、歌星和时尚，这些是藏在他们记忆深处挥之不去的。}
\BookParagraph{比如七十年代早期出生的读者，可以叫「铁臂阿童木」，叫\BookLineBreak{}「血疑」，叫「阳光灿烂的日子」，叫「花房姑娘」，叫「过把瘾就死」，叫「北京人在纽约」，叫「恋曲一九九零」，叫「新华文摘」，叫「芙蓉镇」，叫「排球女将」，甚至可以叫「小鹿纯子」。}
\BookParagraph{当然，如果你的读者群是四五十年代出生的，那你得考虑「红灯记」、「白毛女」、「喀秋莎」，或者「莫斯科郊外的晚上」。}
\BookParagraph{褚明宇曾经提到，他多年前在纽约曾经建议给一个乐队起名叫「新华书店」。他的套路就是试图去唤醒潜意识里对「新华书店」的一种热爱，因为那不仅代表着一家书店，更代表着一个年代。}
\BookParagraph{而管乐队叫书店的那种「驴唇不对马嘴」，正符合刘瑜笔下的\BookLineBreak{}「突兀之美」，有点儿像咬了一口的苹果商标，是起名字这门学问的最高境界。}
\BookDate{2018 年 3 月，多伦多}
\BookArticle{波特酒}{波特酒}
\BookParagraph{李老师您好，通过褚老师认识了李老师，经常默默地看您在分答、微博的回答，包括对于青年人问题的解答，特别喜欢您流畅、质朴的文风。所以我也想请教您一个问题，不知道能否有幸得到您的回答。}
\BookParagraph{我大学毕业于一所内地排名前十的985院校，翻译专业，毕业后在一家英语教育机构（不是新东方）任职，今年是第三年。其实当年毕业也考到了研究生，但是觉得国内研究生太水，没去读。国外文科类研究生学费略贵，虽然家境小康，但是父母赚钱不易，不想给父母增加负担。我是真的很享受当老师的过程，在工作中我发现我很适合做教研，也就是大量阅读原版读物，然后开发出适合不同年龄孩子阅读的材料。}
\BookParagraph{虽然工作上发展尚可，但是随着年龄的增长，愈发感受到办公室政治的残酷和累心，并不想参与其中。于是，萌生了去美国读教育类博士的想法。我知道一纸文凭对赚钱并没有什么用，只是想找到一个可以让自己沉浸于其中的事情去做，如果还能学有所成，就更是不负此生。李老师您自己是带博士生的，您觉得我这样的想法和背景 —— 我觉得自己研究背景还是浅了些 —— 适不适合直接申请读博士呢？申请博士在您看来有哪些需要注意的事项？（问题来自微博私信）}
\BookParagraph{你知道有一种叫「波特」（Port）的酒吗？}
\BookParagraph{波特酒是一种极甜的加强葡萄酒（fortified wine），适合饭后甜点时喝。它的主要产地是葡萄牙海边的一个叫 Porto 的城市，酒名也是因之而来。为了有足够的时间发酵，波特酒一般至少需要两年时间才能装瓶，装瓶后还需要再等至少十几年时间才能喝。它的酒精浓度比一般葡萄酒高，口感纯正、甜美，有一种「余音绕梁」的感觉。}
\BookParagraph{如果说上中学是开胃的鸡尾酒或者香槟酒，上大学本科是与主菜相配的红葡萄酒，那么读博士学位就是餐后的波特酒。}
\BookParagraph{读博士和波特酒一样，有其好处和坏处。如果你明白了这些好处和坏处，还愿意去读的话，我认为值得一试。如果能申请出国读自然最好，就好像要喝就应该喝葡萄牙出产的波特酒一样。}
\BookParagraph{按照西方人的偏爱，咱先从坏消息说起吧。和餐后的波特酒差不多，博士学位可有可无，读不读博士没有什么明显的区别。它无法帮你逃离办公室政治，无法帮你提高工资收入，也多半儿无法帮你结交「谈笑有鸿儒，往来无白丁」的朋友。}
\BookParagraph{更有甚者，读博士的过程其实相当痛苦。想写论文但是写不出来，想做实验但是结果不对，好不容易结果出来了投稿，又被审稿人拒绝，那是常有的事儿。}
\BookParagraph{读博士是一件只有「有志青年」才做的事儿。他们有点儿像数学家或者哲学家，不在乎收入，不在乎朋友，也不在乎「苦其心志」，就是喜欢研究一些普通人懒得去研究的问题。这大概是为什么「博士」英文里其实叫做「哲学博士」（Doctor of Philosophy）的原因。}
\BookParagraph{博士学位可有可无，又含辛茹苦，那它到底有什么好处呢？它最重要的好处，是它虽然苦，但是苦尽甘来时，有一种极为独特、浓郁的甜味儿。没喝过波特酒，你无法想象世上竟然还有这种和香槟或葡萄酒相比迥然不同的甜法儿，甜得发腻，甜得厚积薄发，甜得「一览众山小」，甜得「轻解罗裳，独上兰舟」。}
\BookParagraph{当然还有一个好处：一切尘埃落定之后，你会得到一份可以挂在墙上的证书，上面写着你的名字，还有「Doctor of Philosophy」。}
\BookDate{2018 年 4 月，多伦多}
\BookArticle{编程艺术}{编程艺术}
\BookParagraph{李老师您好！我是ECE开学大二的学生。昨天我们成绩登出来了，相比期中我感觉期末有了很大进步，但是我还是很不满意。我上大学前没有接触过编程，随便选了CE，我至今其实很满意，打算以后还是继续向编程方向发展。您有没有什么可以推荐的书或是一些怎样才能学好编程的建议可以和我分享呢？如果您能看到我的私信并回复的话，我在这里提前感谢您！另外祝您暑假愉快！（问题来自微博私信）}
\BookParagraph{我在清华上学那会儿，同学里曾经流行过一本书，叫《面向对象的程序设计》。那一年正好也是大学二年级，和你现在差不多大。}
\BookParagraph{这本书晦涩难懂到什么地步呢？}
\BookParagraph{我中学时代曾经多次参加过北京市和全国计算机竞赛，独立开发的教学辅助软件得过奖，家里甚至给我买了一台国产的苹果II兼容机。}
\BookParagraph{而看完这本书，我觉得我自己不适合搞计算机这行。}
\BookParagraph{这本书其实写得中规中矩，从头介绍了面向对象程序设计的各种基本概念。封装啊、继承啊、多态啊什么的，讲得头头是道。编程语言也非常前卫非常时髦，叫做 Eiffel。那个时代它和 Lisp 一样，大家都觉得是将来编程语言的大势所趋。}
\BookParagraph{可惜我无论如何认真学习，都搞不明白为什么要发明什么「多态」之类的概念。就连「面向对象」这个词儿，都让人不知所云。「找对象」这个词儿我明白什么意思；但是编个程序，为什么要面向什么「对象」？}
\BookParagraph{直到大学三年级时给一个公司打工完成了一个大型项目，我才找回了一点儿自信。}
\BookParagraph{学计算机编程有点儿像学弹钢琴：看书没什么用，需要真正动手去独立完成一些软件项目。刚开始可以从一些简单的项目做起，比如设计一款简单的游戏。就像钢琴曲目一样，任何一个项目都有更复杂的做法，需要循序渐进，一步一步由浅入深。}
\BookParagraph{比如可以给你的游戏设计一种人工智能算法，人可以和计算机玩儿。刚开始设计的算法可以用相对简单的 minimax 和 alpha-beta pruning，进而可以发展到蒙特卡洛树搜索（Monte Carlo Tree Search），甚至采用加强学习算法。}
\BookParagraph{无论是移动应用还是机器学习算法，一个项目要在 GitHub 这样的开源平台上开始。只有这样，任何一个小改动才有可能完全透明地公开展示在世界面前。}
\BookParagraph{编程是一种艺术。在 Github 上展示自己的编程作品，就像一个画家在画展里展示自己的画作。更为贴切的比喻是，编程就像在北京798艺术区街边儿的人像速写画家，一笔一画的素描，周围永远有路人和老外在围观。}
\BookParagraph{这才是真正面向「对象」的程序设计。}
\BookDate{2018 年 5 月，多伦多}
\BookArticle{逻辑思维}{逻辑思维}
\BookParagraph{李老师您好。我有一个问题想请教您：您平时有没有刻意训练您自己或者您学生的逻辑思维能力？如果有，能否分享一些行之有效的方法？期待您的答复！（问题来自微博私信）}
\BookParagraph{有个冷笑话是这么说的：博士生刚入学看不懂论文，抱怨「我怎么这么笨，连篇论文都看不懂？」等几年后再看不懂论文，开始抱怨「这个作者怎么这么笨，写篇论文都写不清楚？」这时他就可以毕业了。}
\BookParagraph{正好前几天有个学生来找我讨论她论文的草稿，我把第一页看了好几遍，还是没看懂。}
\BookParagraph{作为一个已经毕了业的博士，我说：「你的论文写得逻辑不清楚啊，我看了半天都没明白。」}
\BookParagraph{她反问：「我觉得逻辑很清楚啊，哪儿不清楚了？」}
\BookParagraph{这个问题很难，我一时竟然没有回答出来。}
\BookParagraph{一篇文章是否容易懂，直接反映了作者的逻辑思维和写作能力。不同作者的文字，只需要读一小会儿，高下立判。但是我一直悲观地认为，逻辑思维能力如果靠老师教，不是一时半会儿就能教得出来的。问题就在于我觉得难懂的文章，作者却并不这么觉得。而说清楚为什么难懂，比重写一遍还要艰苦卓绝。}
\BookParagraph{所以我认为，提高逻辑思维能力，基本上得依靠自学成才。}
\BookParagraph{那这个「自学成才」该怎么学啊？——如果不是靠看《罗辑思维》的话。}
\BookParagraph{首先，你工作里经常做的幻灯片，尽量不要用那种一行一个圆点儿的「bullets」。因为用「bullets」的潜台词是，你没有仔细想清楚它们之间的逻辑关系。无论是递进还是转折，全被「并列」这种最为枯燥单调的关系代替了。做幻灯片是为了讲一个引人入胜的\BookLineBreak{}「故事」，故事里有时「花开两朵，各表一枝」，有时「柳暗花明又一村」；起承转合之间，有时归纳，有时演绎，这样才能讲者眉飞色舞，听众茅塞顿开，意犹未尽。}
\BookParagraph{其次，你应该尽量多写原创的文章。写作可以强迫你整理自己的思路，是提高逻辑思维能力最有效的方式。小时候我们写议论文，老师告诉我们要写清楚论点、论据和论证方法。但其实我认为科学论文写作的三个要素在日常写作中也非常重要：「需要解决的问题是什么？」（What’s the problem?）「这个问题为什么重要？」（Why it’s important?）「解法有什么新意？」（What’s new?）}
\BookParagraph{最后，你需要培养独立思考的能力，而非人云亦云。这种批判性思维的特点，是不轻易认同别人的论点，而是要仔细研究他的论据和思路，同时想想相反的观点有什么道理。刘瑜写过一篇小说叫《那么爱呢》，里面有一个人特别爱抬杠，爱说「这话也不能这么说」。这么说话虽然是聊天的大忌，但精神可嘉，因为适度地批判，才能更好地继承和发扬光大。}
\BookParagraph{比如你如果快结婚了，在考虑什么时候贷款买房。无奈房价最近刚刚大涨，高不可攀，而且随时都有可能下跌。大家都说应该节衣缩食赶紧买，因为再等就再也买不起了。你这时不能盲目跟风儿，而是应该好好想想相反的结论是否逻辑上也说得通：说不定应该先租房子住，等到手头宽裕一点儿再买，因为理性的投资从来都是买跌不买涨。}
\BookParagraph{仔细衡量正反双方的立论，正是独立思考的体现。任何人都会「先入为主」地去寻找对自己有利的论据来论证自己偏爱的预设立场和观点，英文叫「confirmation bias」。独立思考的好处，就是靠经常去质疑似乎是正确的观点，来尽量纠正这个归纳推理中的系统误差（systematic error）。}
\BookParagraph{你看到这里，如果想说：「这个博主怎么这么笨，写篇微博都写不清楚？」}
\BookParagraph{那恭喜你，差不多可以毕业了。}
\BookDate{2018 年 6 月，西安}
\BookArticle{人来疯}{人来疯}
\BookParagraph{李老师好，我怀疑自己有社交恐惧。每次和别人面对面交流总会有脸麻的感觉，就像针扎一样。工作需要上台讲话，说话就结巴，恨不得低着头读完。老师我该如何练习，才能让自己生活更好一点？（问题来自微博私信）}
\BookParagraph{我小时候，有一拨儿孩子有一个共同的特点，家里来了客人就特别兴奋，又跳又闹的来劲儿。那会儿管这拨儿孩子叫「人来疯」。}
\BookParagraph{而我恰好和你一样，最怕的就是见到陌生人。我一见到陌生人就紧张得不得了，绝对属于「人来疯」的反面儿。就像你说的，现在好像管这叫做「社交恐惧」，英文叫「anti-social」。}
\BookParagraph{我想，如果用朋友圈儿里时不时冒出来的那种测试性格的小测验测一下，我的性格应该属于内向型（introvert）。我崇尚低调，从来不发朋友圈；不擅长参与集体活动；从来不踢足球，偏爱一个人长跑；与其去参加一个什么校友会活动，更喜欢在家看书写论文编程序。}
\BookParagraph{这样的性格，大概就是我害怕陌生人的原因。}
\BookParagraph{其实像你我这样内向兼「社交恐惧」的人，还有很多。有的调查发现，这样的人，至少有三分之一。}
\BookParagraph{然而现在的工作大多数都要和陌生人交流，开会要积极发言，随时要拿着「PPT」上台讲话，对外向性格的人极为有利。}
\BookParagraph{我刚开始也和你一样，而且越是重要场合的讲话，我就越紧张。不但上台紧张，上台之前也紧张，甚至前一天晚上准备的时候就开始紧张。我紧张得需要把要说的话全都死记硬背给背下来。}
\BookParagraph{日子就这么一天一天过去了，幸亏我只是上台说个话而不是表演钢琴独奏——那样我一定紧张得一个小节都弹不全——我发现其实也没有出过太多洋相。}
\BookParagraph{不知道你是否意识到了，其实这样的性格也没有什么不好。恰恰相反，内敛性格的人，更有可能利用自己一个人孤单的时间（solitude），深入地去钻研一个问题，精益求精，最终「无招而胜有招」，达到「独孤求败」的境界。}
\BookParagraph{苹果公司设计的第一台计算机，就是当年 Stephen Wozniak 足不出户，拿着烙铁和各式集成电路芯片在自己家里设计调试出来的。Wozniak这个人是典型的内向兼社交恐惧的性格，他虽然不擅长上台讲话，但是没有他四十年前设计的苹果计算机，就没有今天的苹果公司。}
\BookParagraph{在这个「直播」、「网红」、「大V」流行的时代，也许我们更需要的不是「人来疯」的那拨儿孩子，而是爱「宅在家里」的Wozniak。}
\BookDate{2018 年 6 月，西安}
\BookArticle{盲目自信}{盲目自信}
\BookParagraph{李老师，您好。请问如何驾驭或改变自己的思维习惯？最近在弗兰克 · 哈多克《意志的力量》中看到一句话，说一个人的脆弱或不幸，虽然可以追溯到出身和环境，但归根结底要看这个人的思维习惯。而我感觉自己很容易陷入没有意义且重复的想法或对未来的幻想，想听听您对这个问题的看法。谢谢。（问题来自微博私信）}
\BookParagraph{有一个英文短语叫「gut feeling」你听说过吗？}
\BookParagraph{世界杯淘汰赛最近比利时绝地反转赢了日本那场看了吧？伤停补时最后一分钟，从比利时守门员手抛球开始，一下儿就传到日本队禁区。第二次传球时，比利时前锋卢卡库跳起来假动作一漏，成功吸引了一位防守队员，球直接到了队友脚下，近距离射门绝杀。}
\BookParagraph{卢卡库跳起来漏球的假动作，当时是怎么思考的？}
\BookParagraph{他看不到自己后面的队友，没有时间充分地衡量这一漏到底风险有多大。}
\BookParagraph{他靠的就是「gut feeling」，是一种自信，一种直觉。直译为中文应该是「用肚子想问题」。}
\BookParagraph{充分相信你的「肚子」（trust your gut），其实可以解决很多问题。比如你作为公司领导面试了一位员工，有优点有缺点，到底给不给「offer」？比如你研究过程中有两个想法，时间不多了，到底继续去琢磨哪一个？比如你大学毕业出国，可以选择美国，也可以选择加拿大，你准备去哪里读书？}
\BookParagraph{这些问题都摆在岔路口上，决定了你到哪儿读书到哪儿工作，工作怎么做，跟什么人合作等等，按你的原话说，它们足以「驾驭和改变」你的生活。}
\BookParagraph{可是我们做这些决定之前，真的没有时间去列出具体的表格或推出具体的公式去优化最终的目标函数。我们能做到的，只有相信自己的「肚子」。}
\BookParagraph{我上大学时有一个室友，这位同学有两大业余爱好：申请入党和参加舞会。有一次开会时，他的发言一鸣惊人：班委是团委的左膀，团委是班委的右臂。他积极要求进步，头一次入党申请没有通过，一度伤心地哭了。去参加舞会时则像换了一个人，西服笔挺，红光满面，头发梳得锃亮。有一次从七食堂舞会回来，他非常纠结，因为他发现有两个女孩儿喜欢他，不知道该挑哪个了。}
\BookParagraph{我们都说：「这还不简单，谁跟你说话多你就挑谁呗。」}
\BookParagraph{他思考了一下：「她们都没说话啊，都在默默地注视着我。」}
\BookParagraph{你看，「trust your gut」其实也没那么容易，时机转瞬即逝，过了这村儿就没这店了。}
\BookParagraph{下回再遇到这种可以改变自己生命意义的事儿，别去问弗兰克 · 哈多克，而是盲目地相信自己。}
\BookDate{2018 年 7 月，西安}
\BookArticle{边　界}{边界}
\BookParagraph{李老师，您好。想请教您关于边界感的看法。最近发现在和朋友陈述自己的想法时，对方在我没有问她的意见时就已经对她的想法直言不讳。比如我说我想去打工旅游而不是现在就定下来，对方立马列举了各种原因告诉我打工旅游不好：耽误的时间成本难以弥补，打工旅游的工作都很简单而且工资也少，去外企是一个更好的选择。听到这些我很不舒服，因为我从始至终没有问过她的意见，只是想把我思考后的决定告诉她。我甚至会觉得对方这样是对我不尊重，虽然她是为了我好。想问问李老师，如何处理这种情况？谢谢！（问题来自微博私信）}
\BookParagraph{有一个GRE英文词汇叫「gerrymander」你听说过吗？}
\BookParagraph{中文大意是，把选举的选区人为划分得曲里拐弯儿的，为了自己党派选举时可以多得几张选票。}
\BookParagraph{英文作为一种交流的语言，绝对像Intel处理器一样是一个复杂指令集（CISC）：一个单词儿竟然能代表一大块儿复杂艰深的含义，而这个「gerrymander」就是一个很好的例子。而中文则是ARM处理器里的简化指令集（RISC），每一个词汇的含义简单有效，字字珠玑，组合起来又能形神兼备，行云流水。}
\BookParagraph{我们两个人之间的「边界感」，有点儿像「gerrymander」，似乎随便我怎么定义都可以。}
\BookParagraph{你作为反对党，如果说这样定义边界有问题，我一定会疑惑不解地、满怀幽恨地、恨铁不成钢地跟你说：「我这还不是为了你好？」}
\BookParagraph{为了你的前途，为了你的事业，为了你后一半儿人生的幸福，作为你的父母、你的好朋友、你的闺蜜，甚至你的战友、领导、老师、同学，我不告诉你残酷的事实真相，我不去关心你的个人问题，难道还会有别人像我这么好心吗？}
\BookParagraph{我吃过的盐比你吃过的饭还多，走过的桥比你走过的路都长，我参加红军长征时你还不会走路呢，你这真是「好心当作驴肝肺」啊。}
\BookParagraph{你现在不学习长大了哪儿来的饭吃？现在这也不行那也不行将来到哪找这么好的去？现在不结婚将来鸡飞蛋打了怎么办？越早生孩子孩子越聪明，这么简单的道理怎么就没明白？你看别人家都买房，你怎么还在犹豫？}
\BookParagraph{于是组织上要找你了解情况，居委会要找你促膝谈心，班主任辅导员要找你做思想工作，让你一定要记住：集体荣誉高于个人利益，不要自私自利，要有崇高的理想，在奋斗中实现自己的个人价值。}
\BookParagraph{对所有的这一切，你作为反对党，该怎么回应？}
\BookParagraph{你可以直接用上英文这个复杂指令集：「Why do you care? 」}
\BookParagraph{你可以温文尔雅有礼貌：「我已经没救了，天要下雨，娘要嫁人，您已经尽力而为了，将来不会后悔的。」}
\BookParagraph{你可以援引王朔《过把瘾就死》里的一句话：}
\BookParagraph{「我觉得你在思想上太关心我了！都快把我关心疯了！」}
\BookParagraph{你也可以用标准带北京口音的普通话反唇相讥：「有完没\BookLineBreak{}完？」}
\BookParagraph{最好再用阴暗的、叛逆的、让她觉得不打你一顿就不解恨的口气加上一句：「您管着吗？」}
\BookParagraph{祝你旅途愉快，一路顺风。}
\BookDate{2018 年 7 月，多伦多}
\BookArticle{英语口语}{英语口语}
\BookParagraph{李老师您好，看到您谈到学英语的问题，我作为一名英语专业的学生不免有点小激动，因为我就是那个还在抱着新概念四朗读的人，哈哈。}
\BookParagraph{谈到学英语，大多数中国人最津津乐道的，恐怕是如何能在短时间内掌握一口流利的英语口语。国内各种培训机构推崇的也是这种短时间、高强度（intensive）的学习方式，但学到的大多是一些类似于「how do you do?」之类的套话，长时间不练也就忘得差不多了。所以说白了还是得靠自学，国内学习口语缺乏语言环境，只能自己给自己创造环境，好在现在 podcast 之类的媒体软件很发达，想在国内练就一口伦敦腔也不是什么难事。但语言学习逃不掉面对面交流这个坎儿，自说自话时很顺，长时间不开口用英语和人交流等到再用时还是会语无伦次，尴尬百出。}
\BookParagraph{语言学上有预制语块（prefabricated chunks）的概念，大意是我们交流中所使用的词汇、句子大多是我们之前就储存过的，交流时就是自动提取、匹配这些语块。按照这个概念，我们可能需要大量积累各种场景里所需的表达方式，然后烂熟于心（proficiency），变成一种本能。}
\BookParagraph{这样我们能解决输入的问题，然而流利的表达就不止是语言了—— en、er、a等语气词就好像中文交流中有人特喜欢说然后，然后，然后……一样烦人。}
\BookParagraph{请问李老师，如何能做到交流中 clarity and grace 并重呢？谢谢！（问题来自微博私信）}
\BookParagraph{语言学（linguistics）我一点儿不懂，所以我今天要说的「理论」全凭朴素的直觉，在你这个英语专业毕业的专家面前，无异于班门弄斧。}
\BookParagraph{你提到的第一个问题是英文口语如何学习。我认为绝对应该从简单的对话开始。我自己的口语就是上小学时看电视跟着一个叫\BookLineBreak{}《跟我学》（Follow Me）的英国教学节目自学的。整个节目六十集，全部是情景对话，内容由浅入深，很多段落现在都能背出来。初中时又跟着一本叫《现代美国口语》的录音带学，还是情景对话，基本上奠定了我的英语口语基础。}
\BookParagraph{我觉得口语是否能学好，这些教学节目或教材的质量非常重要。很多口语教材里的英语不够地道，自然也就学不到标准的口语。比如两个人见面最简单的寒暄，常见以下的对话：}
\BookParagraph{Hi, Helen! How’s it going?}
\BookParagraph{Fine, thanks — and you?}
\BookParagraph{Just fine. Where are you off to?}
\BookParagraph{To the library. I’ve got a history exam next week and need to start studying. Ugh.}
\BookParagraph{Oh, no. Well, I’ll see you later then. Good luck!}
\BookParagraph{Thanks. See you later.}
\BookParagraph{这段对话里有几句，大多数中国学生说不出来。比如「Where are you off to?」，很多人会说成「Where are you going？」但是后者稍嫌正式，前者则更轻描淡写。又比如「Oh, no.」有一点儿「同情你」的含义，并不是要否定或拒绝什么事儿。这些都是表面很简单但实际上相当地道的英文。}
\BookParagraph{像这样的英文口语，难就难在用到的全都是极其简单的词，但是变化多端而飘忽不定，需要多听、多说，大量地重复练习，用起来才随心所欲。}
\BookParagraph{这里推荐美国国务院为外国人学英语而编纂的《Everyday Conversations: Learning American English》，上面这段对话就是摘自这本小册子。}
\BookParagraph{你提的第二个问题是如何在说话时尽量避免夹杂en、er等语气词。我先给你讲俩故事吧。}
\BookParagraph{看过最近 Google I/O 大会上 Google Duplex 的演示吗？计算机可以打电话给餐厅预约时间订座位，对话清晰自然，甚至可以说是通过了图灵测试。}
\BookParagraph{有趣的是，谷歌最近披露，早期的实验结果其实并不令人满意，预约电话经常失败，因为接听的人感觉像在和计算机说话。后来加入了en、er之类的语气词，让完美的句型里掺杂了一点儿停顿和犹豫，反而成功率大大提高，因为普通人都这么说话。}
\BookParagraph{我多年前在微软研究院访问时，有幸和清华当年最著名的少年天才、邓小平爷爷当面表扬「计算机要从娃娃抓起」的李劲合作共事过一段时间。他说英文最喜欢加「basically」，一段英文甚至要加进好几个。不过，从来没有任何人在意这一点，也没人抱怨听不懂他说的话。}
\BookParagraph{所以说，没必要过于追求完美。只要你能用英文清晰地表达自己的观点，多说几个en、er、basically，白玉微瑕，天才还是天才，没有人会在意的。}
\BookDate{2018 年 7 月，多伦多}
\BookArticle{看小说}{看小说}
\BookParagraph{李老师你好！看了你微博猜测向你咨询问题的人一定很多，我都怀疑你是否会看我的，好吧就算不看，我也得写啊问啊。}
\BookParagraph{李老师是个成功的人，而我却无所作为，每当想努力想学习，发现心怎么都静不下来。大道理小道理都懂，就是做不到耐着性子。一度怀疑自己是否记忆只有七秒，多次试验发现记忆不止七秒，自信重回，但是每次捧着书就是看不下去，想考什么证都有决心没耐性，甚为苦恼。}
\BookParagraph{以前也谈不上多么热爱学习，但不至于像现在这样成为一个不想看书、也看不进去书的人。拿着书很抗拒，不管是有兴趣还是交了钱要去考试的都不上心，可不学习不学点东西在这个城市真不知道能找到什么工作。已麻木的我以为别人的成功来得容易，真不知道别人是如何战胜困难，战胜自己的。我讨厌现在的自己。}
\BookParagraph{李老师啊，你有什么好的办法让我好好看书，让我下定决心付出行动吗？（问题来自微博私信）}
\BookSubtitle{一}
\BookParagraph{我一边做晚饭一边跟刚到家的郭芳说：「今儿个微博私信又收到了一个问题。」}
\BookParagraph{「什么问题啊？」}
\BookParagraph{「她问，大道理小道理都明白，就是不爱看书学习，怎么办？」}
\BookParagraph{「那她爱看小说吗？」}
\BookSubtitle{二}
\BookParagraph{你想没想过，不爱看书的一大原因，很可能是你手里的书写得一点儿也不好看。}
\BookParagraph{我们从小到大听过的都是似乎特别有道理的「励志」画外音：「读书破万卷，下笔如有神」啊，「人生能有几回搏」啊，「读万卷书，行万里路」啊，等等等等，不一而足。}
\BookParagraph{但是说真的，大多数的书，真的写得不怎么吸引人。}
\BookParagraph{比如说我家小女儿正在学的小学四年级语文课本，就有点儿令人捉摸不定。}
\BookParagraph{先看题目：《夜莺的歌声》、《小英雄雨来》、《一个中国孩子的呼声》、《和我们一样享受春天》、《触摸春天》、《永生的眼睛》、《生命、生命》、《花的勇气》……}
\BookParagraph{咱再随意打开一篇叫《触摸春天》的课文看看。}
\BookParagraph{看文章像看体操一样，要注意它的开场和下法。}
\BookParagraph{这篇文章开头写道：「邻居的小孩安静，是个盲童。」}
\BookParagraph{结尾写道：「我没有惊动安静。谁都有生活的权利，谁都可以创造一个属于自己的缤纷世界。在这个清香袅袅的早晨，安静告诉我这样的道理。」}
\BookParagraph{课后习题正好和这个结尾有关：「多美的课文……让我们结合课文和生活实际，说说对这句话的理解。」}
\BookParagraph{读了这样有点儿令人提不起兴致的文章，我要是小学生，说实话多半儿也会放下课本，去打《王者荣耀》。}
\BookParagraph{所以啊，不爱看书，很可能是书写得不好看；而你多半儿内心深处从小到大觉得不读书或者不学习是不对的，还要强迫自己看。就像背英语单词一样，越背越烦人，以前背的还都忘了个八九不离十。}
\BookParagraph{明知「不进则退」 还要逆水行舟，战斗还没打响，败局已定。}
\BookSubtitle{三}
\BookParagraph{那你爱看小说吗？}
\BookParagraph{比如钱钟书的《围城》、金庸的《天龙八部》、慕容雪村的\BookLineBreak{}《成都，今夜请将我遗忘》、艾米的《竹马青梅》、辛夷坞的《致我们终将逝去的青春》、王朔的《一半是火焰，一半是海水》，还有刘瑜的《那么，爱呢》？}
\BookParagraph{你也许会说，我说的看书学习不是看小说不务正业，是学点儿对将来事业发展有用的东西。}
\BookParagraph{你还别说，即使你将来的理想是当个数学家，你的当务之急还是放下微信，放下手机，放下所有的「大道理」和「小道理」，放下「学习」啊「成功」啊这些词儿，去亲身体验一把那种看一本书看得如痴如醉、如火如荼的感觉，一种唏嘘不已、物我两忘的境界。}
\BookParagraph{因为那时你会发现，在一本写得引人入胜的书里，竟然藏着一个和微信群朋友圈完全不一样的世界。}
\BookParagraph{我没有微信朋友圈，微信群里的红包总是抢不到，但却一直酷爱看小说，酷爱看一切令我陶醉于其中的书。所以说，我从来没有「成功」地「战胜」过自己。}
\BookParagraph{「谁都有生活的权利，谁都可以创造一个属于自己的缤纷世界。在这个清香袅袅的早晨，安静告诉我这样的道理。」}
\BookDate{2018 年 9 月，多伦多}
\BookArticle{捧　哏}{捧哏}
\BookParagraph{李老师你好，关注您微博也有一段时间了，最近有个问题想请教您，还望解答。}
\BookParagraph{怎样才能好好说话呢？我说话总是会无意地得罪对方，即使别人生气了还搞不清楚是怎么回事。怎样才能够提高那种语言上的分寸感？（问题来自微博私信）}
\BookParagraph{你这个问题问我，绝对问错人了。因为我说话和你一样，也经常说错，跟你相比最多只能算是「五十步笑百步」而已。唯一比你略强一点儿的是，我一般说错了话刚说完就马上能意识到，可惜永远无法收回了。}
\BookParagraph{因为我经常说错话，有时候也会有意识地想一个问题：怎么说才能少说错几句话？要是能把酒言欢，聊一晚上天儿，一句也不说错就好啦。}
\BookParagraph{后来我发现，其实做到这一点并不是特别费劲儿。}
\BookParagraph{你一定听过德云社的相声吧？聊天说话想少说错话，关键是别老去当那个逗哏的，说学逗唱样样精通；而是要多退居二线，去当那个捧哏的。}
\BookParagraph{捧哏的常说的话很简单：「嗯」「哟」「那是」「啊？」「别逗了」「等会儿等会儿」什么的，看过相声的都门儿清。}
\BookParagraph{你说话的任务是，听明白逗哏的那位讲的故事，然后想办法激发他继续讲下去的积极性。}
\BookParagraph{比如他故事讲到一半儿停下来，磕了一瓜子喝了一口啤酒，你这时正好可以接上话：「哟，那后来呢？」}
\BookParagraph{千万别跟刘瑜小说《那么，爱呢》里的「蜘蛛侠」似的，来一句「那也不能这么说」，呛了逗哏的一句，这相声也就说不下去了。真是活该小说里的「蜘蛛侠」找不到对象，谁会喜欢和这种人聊天吃饭啊。}
\BookParagraph{别较真儿，就能完美地演好一个捧哏的角色。}
\BookParagraph{其实捧哏的烘托场上气氛并不那么容易，人那也是门儿说话的艺术。我想起了高晓松几年前在优酷搞的那个叫《晓说》的节目。本来好好的看的人其实挺多的，就差家喻户晓了，结果高老师大概嫌自己一个北京人侃单口相声不过瘾，有一集找了一个优酷的美女主持人坐在对面儿沙发上，一起聊林徽因。}
\BookParagraph{结果这位美女有点儿辛苦地全程不停点着头，最多再加上几句「那然后呢」，捧哏捧得估计令高老师大失所望，过了几集就不再出现了。}
\BookParagraph{所以说啊，也不能净说「那后来呢」，还要时不时锦上添花，来一句「这话说得在理儿」、「您还挺明白」、「还真是这么回事儿」什么的，反正就是激励逗哏的继续变本加厉地猛说呗。}
\BookParagraph{那你一定会问了，逗哏的说的一点儿不逗怎么办？或者没人爱说话，都在玩儿手机？}
\BookParagraph{得，这时候就是您该「亮剑」出场的时候了。}
\BookParagraph{「我给大家讲一笑话啊，估计你们肯定没听过。」}
\BookParagraph{等你笑话讲完，一定会有人跳出来说「你这笑话一点儿都不逗，我都听过好几遍了。」然后继续口若悬河地开聊。}
\BookParagraph{这时你立马「敌进我退」，继续当个捧哏的就行了。}
\BookParagraph{捧哏到了最高境界，在座的各位竟然都有机会侃上一段儿，饭局在亲切友好的气氛中觥筹交错，酒过三巡，没人注意到你暗中\BookLineBreak{}「润物细无声」的起承转合。}
\BookParagraph{如此说话的功力，要是你将来有机会和高晓松「对谈」，估计他能从抗美援朝侃到《如懿传》，喜玛拉雅上再开一专栏儿，也未可知呢。}
\BookDate{2018 年 9 月，多伦多}
\BookArticle{昙花一现}{昙花一现}
\BookParagraph{老师，您好。我是一名直博生，一直很焦虑迷茫，压力大，感觉在崩溃的边缘。找不到能解决我困惑的人，想跟您聊聊，或许会打开我的心结。}
\BookParagraph{我的博士方向需要的知识，本科没有接触过。但导师总是说，这个还用学吗？那个还用学吗？直接用就会的呀！但从师哥师姐那里得知导师其实也不怎么会。就这样，导师让我做第一个工作，大概半年，仿真没有做出来，导师嫌弃没有成果，曾一度说我毫无能力之类的，外加各种毕业威胁。再后来就直接把我搁置一边，说了句你爱看什么看什么，不再管我，我遇见他打招呼都爱搭不理。偶尔又时不时地发发微信说我没产出等等。}
\BookParagraph{在这种情况下，我很难过，同届学生都做得风生水起，师哥们看到我，都说一句「心疼你」。我压力好大，不想放弃，自己又迷茫，在这种境遇下，心里每天很乱，越乱越学不进去。我不知道，科研如果单纯靠自己，没人指导，是否真的可以做出成果。这种情况下，只靠努力，足够努力，一样会做的好吗？我选择读博，就是希望将来能联培出国，但这样下去，感觉毫无希望。我真的不知道该怎么办，在周围同学的怜悯之中，自己更难过。}
\BookParagraph{您能给点建议吗？我应该怎么面对当下的这一切？在科研上，我该怎么着手，让自己找到前进的方向？导师时不时给我的威胁，让我压力更大。所以，李老师，很期待您的回复，给一个不想放弃却倍感绝望但还想挣扎的我宝贵的建议。您觉得在这种情况下应不应该继续坚持读下去，我一直陷入其中，看不到希望。谢谢老师！（问题来自微博私信）}
\BookSubtitle{一}
\BookParagraph{你想没想过，你一辈子认识的人里的大多数，都仅仅是昙花一现的过客。}
\BookParagraph{小学时两小无猜的发小、中学时一见倾心的初恋情人、大学军训时凶狠无比的排长，时间长了，印象慢慢地淡忘，直到终于想不起来他们的名字，也记不起来他们的样子了。}
\BookParagraph{据说，天底下所有的父亲最黯然神伤、欲哭无泪的时刻，是女儿大婚之时。我猜想，大概是终于意识到，骨肉至亲，二十多年的朝夕相处，终将远去，竟然也无法逃离生命之中「昙花一现」的命数。}
\BookParagraph{学生跟导师读博的那几年时间，你以后会发现，更加是弹指一挥间的事儿。没办法，日子过得太快，青春稍纵即逝，「大江东去，浪淘尽、千古风流人物。」}
\BookParagraph{正因为如此，每一天的日子，归根结底，还要依靠自己。}
\BookSubtitle{二}
\BookParagraph{无论是学习还是科研，都应该静下心来，一点一滴慢慢地积累，这需要很长的时间。在这个漫长的过程中，同学也好，师哥也罢，可能会有人给你提建议，但是没有人可以真正全程陪着你。归根结底，还要依靠自己。}
\BookParagraph{至于这位导师，我赞同毛主席的名言：「人不犯我，我不犯人；人若犯我，我必犯人。」毛主席还说过：「任何方面的横逆如果一定要来，如果欺人太甚，如果实行压迫，那么，共产党就必须用严正的态度对待之。」}
\BookParagraph{当然，咱也可以跟老外学习，无论是Linux内核的创始人Linus Torvalds，还是当年的乔布斯，都是著名的坏脾气，一旦骂起人来，「谈笑间，樯橹灰飞烟灭。」}
\BookParagraph{什么时候可以「中国人民站起来了」，对一个原本只是生命中的过客，有一点儿自己骨气和尊严，不再委屈求全？}
\BookParagraph{归根结底，还是要依靠自己。}
\BookParagraph{我建议：战略上要极端藐视，不要对「成果」太过苛求；战术上则要高度重视，充分利用每天的时间，经常参加各种运动，慢慢放宽自己的心情，每天做一点儿让自己觉得有收获的事儿。}
\BookParagraph{不要忘了，你正值大好青春，不但有的是时间，而且早晚能毕业。}
\BookParagraph{实在心情不好，给微博上的一个叫李老师的陌生人发条私信，也未尝不是一剂良方。}
\BookSubtitle{三}
\BookParagraph{据传说，昙花原来是一位花神，她一年四季，每天都会开花。她爱上了每天给她浇水的帅哥，被玉皇大帝知道了，大发雷霆要拆散她们。玉皇大帝把花神抓了起来，贬为每年只能开一瞬间的昙花，不让她再和情郎相见。}
\BookParagraph{小朋友听到这里一定会问：那后来呢？}
\BookParagraph{这个故事，还是由你继续给小朋友讲吧。}
\BookDate{2018 年 9 月，多伦多}
\BookArticle{不　懂}{不懂}
\BookParagraph{李老师，华为高管被拘留的事儿，又激起了千层浪，对于这类非常不看好中美关系的言论，请问你是怎么看待的？对与即将去美国留学的学生，需要担心什么吗？（问题来自微博私信）}
\BookParagraph{李老师，关于孟晚舟事件您怎么看？（问题来自微博私信）}
\BookSubtitle{一}
\BookParagraph{两位同学，谢谢你们的来信。可惜你们这个问题我回答不出来。}
\BookParagraph{原因很简单：我不懂啊。}
\BookParagraph{这个事件无论是从国际政治、国际关系、中美关系、还是美国的法律来看，我都一窍不通。}
\BookParagraph{这就好像你要是问我二十多岁的女生口红最适合什么品牌和色号，量子通信到底是怎么回事儿，或者正宗的小笼汤包怎么做，我也同样回答不出来一样。}
\BookParagraph{要想逻辑清晰、有理有据、言之有物地写一篇文章，我觉得蜻蜓点水略知一二不行，应该有本事把一个艰深的道理简单而又准确地讲清楚。}
\BookParagraph{而这个本事看似简单，其实并不容易。}
\BookParagraph{不懂就是不懂，我不懂的东西实在太多了。}
\BookSubtitle{二}
\BookParagraph{那不懂的东西就不能研究一下吗？}
\BookParagraph{当然可以啊。不过这也得看你愿意花多少时间。}
\BookParagraph{研究女生口红的品牌和色号，也许有一两个星期就可以粗通；量子通信则不然，怎么也得从高中物理开始，量子力学、量子计算都需要有基本的涉猎，这怎么也得要一两年时间吧。}
\BookParagraph{那孟晚舟事件所涉及到的中美关系呢？}
\BookParagraph{看过我以前文字的朋友可能记得，我有个舅舅叫王缉思，是北大教授，以前担任过国际关系学院的院长，现在是北大国际战略研究院院长，办公地点在北京大学北阁。}
\BookParagraph{他的专业正好就是中美关系。}
\BookParagraph{我记得十几岁的时候去舅舅家玩儿，那时他还在中国社会科学院美国研究所工作。到现在我印象还非常深刻的是一个细节：他的书房里一整面墙全都是书。}
\BookParagraph{我当时想，这么多书，看到什么时候才能看完啊。}
\BookParagraph{想粗通中美关系？没问题啊，先至少看完半墙的书再说。实在不行，基辛格的著作——比如《On China》——至少熟读过吧？不能光看中文书，怎么也得看英文的原版吧？}
\BookParagraph{光看书还不行，还要有能举一反三、触类旁通的能力。这样写出来的文章，才能言之有物，令人信服。}
\BookParagraph{王缉思穷尽多年的精力才在中美关系这个领域学有所成，有所建树，成为国际知名的学者。}
\BookParagraph{哈哈，我还是省事儿一点儿，直接说我不懂吧。}
\BookSubtitle{三}
\BookParagraph{那不懂就不能发表自己的观点了吗？}
\BookParagraph{当然可以。不过这样的观点多半儿只能蒙住不懂行的。真正遇到懂行的人，和他们讨论的时候你大概连插嘴都插不进去。}
\BookParagraph{我这个人有点儿完美主义——或者说是强迫症吧——与其冒着让懂行的人发现我的观点漏洞百出的风险，我宁可没有观点。}
\BookParagraph{所以，即使我花了点时间在网络上研究过女生的口红色号，我也不会去写一篇有关这个话题的微信公众号文章，去争取「十万加」。}
\BookParagraph{所以说啊，我这辈子压根儿就没有当「网红」当「大V」的\BookLineBreak{}「命」。}
\BookParagraph{因为与其什么都懂一点儿，我更愿意用无穷多的「不懂」去换来一生的时间，把一两样东西真正研究到令人刮目相看，拍案叫绝。}
\BookDate{2018 年 12 月，多伦多}
\BookArticle{失　败}{失败}
\BookParagraph{亲爱的李教授你好，关注你的微博很久了，特别欣赏你的人生态度和思想。作为一个人生阅历还不够丰富的人，想请教你，如何拥有一颗坚定的心去对待人生的挫折和失败呢？最近遭受了学习生涯中最大的失败，分数和排名下降很多。心态上的打击很大，也不停地在否定自己。虽然也明白「知耻而后勇」的道理，但在这一个努力的过程，心态很难坚定下来。想听听你的建议，谢谢！（问题来自微博私信）}
\BookSubtitle{一}
\BookParagraph{我觉得咱们首先得琢磨一个基本问题：什么才算是人生中的失败？}
\BookParagraph{「失败」这个词儿让人感觉挺严重的，所以我觉得「人生中的失败」一定和长年累月的错误生活方式有关系，有点儿「在错误的路上越走越远」的感觉。}
\BookParagraph{比如，明知抽烟对自己和家人都不利，却没有下定决心迷途知返，可以算是一种「人生中的失败」。}
\BookParagraph{或者明知成天玩儿游戏影响学习和工作，却深陷其中不能自拔，大概也可以算是一种「人生中的失败」。}
\BookParagraph{那考试没考好呢？}
\BookParagraph{我觉得这不能算，因为考试有很大程度上的偶然因素。比如考前正好身体不适，复习准备和考题相差甚远，或者一步推导出了差错，以致整道题功亏一篑，都可能是考试没考好的原因。}
\BookParagraph{我就有幸认识一位很特别的同学，特别之处正好就和考试有关。}
\BookSubtitle{二}
\BookParagraph{我认识这位同学那一年，和你一样也是刚上大学。}
\BookParagraph{这位同学有一个与众不同的姓氏，他姓「移」。不知为什么，老是让我想起毛主席的文章《愚公移山》。大概也正是因为这个奇怪的姓氏，没有人给他起过外号，大家就尊称他为「老移」。}
\BookParagraph{我刚上大学那年流行穿军装，至少也得斜背个军挎，那种一无所有的感觉，和校园里大喇叭放的崔健黑豹什么的很「搭」。}
\BookParagraph{而老移两鬓略有些斑白，成天穿的是一件西服。西服穿得时间长了，显得有些破旧，风尘仆仆。}
\BookParagraph{老移的爱好也和别的同学不一样。别人的业余爱好至少有好几种，有的人谈恋爱，有的人爱打牌，有的人爱踢球——而老移只看武侠小说。别人白天学习时他在看小说；晚上十点四十五熄灯了，别人开始聊天，他则点起蜡烛，继续看他的小说。}
\BookParagraph{据他自己介绍，五道口租小说的那家店里的所有武侠小说，他都看过好几遍了，有点儿「熟读唐诗三百首，不会吟诗也会吟」的境界。}
\BookParagraph{直到很久之后，才听说了老移当年的轶事。}
\BookParagraph{据说老移的高中时代是在江苏省读的书，88年就考上了清华大学计算机系。可惜的是，上了大学以后因为沉迷于看武侠小说，他好几门课不及格，被学校勒令退学了。}
\BookParagraph{于是老移回到了他江苏的老家，又准备了一年高考。1990年——也就是我入学的那一年——再次考入了清华大学计算机系。}
\BookSubtitle{三}
\BookParagraph{后来的故事估计你都能猜到了。老移因为看小说，又是多门功课不及格。不过这回估计连系领导都被他「我想要回到老地方，我想要走在老路上」的精神所感动了，终于没有再次决定让他退学。}
\BookParagraph{大学五年上完毕业的时候，我们同学里出国的出国，推研的推研，闪婚的闪婚，而「成功」、「失败」、再次「成功」了的老移，终于因为看了五年武侠小说，不及格的课目太多，没有拿到学位，只拿到了一纸结业证书。}
\BookParagraph{说他看了五年小说，其实不是特别准确。准确地说，有一年物理课期中考试前，老移突然大彻大悟，开始学习。结果几天之后的期中考试，他一举拿下全班最高分98分。}
\BookParagraph{考完期中考试之后，他又开始了他的武侠之旅。}
\BookSubtitle{四}
\BookParagraph{老移从清华毕业后，去了建行系统工作了多年，又辗转移民了加拿大。}
\BookParagraph{我十年前最后一次见到他时，他还是那副两鬓斑白风尘仆仆的样子，还是永远穿着一身略显破旧的西装。}
\BookParagraph{他还是和多年以前一样，不怎么爱说话，笑得憨厚而从容。按这几年时髦的词儿来形容，老移一直是一「佛系」青年——六十年代那会儿叫「逍遥派」，正是金庸武侠小说里的常用词。}
\BookParagraph{我向老移的媳妇「表白」，说老移当年一直是我青春无悔的偶像。因为老移太擅长考试了，以至于考好考坏都是他一个人说了算。}
\BookParagraph{有点儿像当年清华校园大喇叭里常放的那首崔健的歌《假行僧》：}
\BookParagraph{我要从南走到北，我还要从白走到黑}
\BookParagraph{我要人们都看到我，但不知道我是谁}
\BookParagraph{新年快乐。}
\BookDate{2019 年 2 月，多伦多}
\BookArticle{盖棺定论}{盖棺定论}
\BookParagraph{和大女儿聊起来她们的历史课。她说她们已经学完美国历史，现在正在学习以希特勒和斯大林为核心的欧洲近代史，下学期开始学毛泽东。}
\BookParagraph{我听着悠然神往，忽然有一种冲动，我也想去上这样的历史课。记得中学时好像就没怎么好好学过历史，而大学一年级学过的「中国革命史」早就忘得差不多了。很多现在还记得的历史，大多是从以前看电视或者看《明朝那些事儿》这类的闲书而来。而对美国或欧洲的历史的了解更是少得可怜，甚至连书都没看过。}
\BookParagraph{女儿还说，下学期即将开始写一篇大作文（extended essay），五千字左右的篇幅，选题不限，科学、数学、英语文学、历史、地理都可以选，自己确定研究的题目和范围。她的一位好朋友准备选《The Nanjing Massacre》作为大作文的题目。}
\BookParagraph{我想，如果我是一个高中生，我不会有胆量写这个作文题目的。}
\BookParagraph{聊着聊着，我想起现在有关「996」的讨论，其实如果懂一点儿美国大萧条的历史和罗斯福制定八小时工作制的起因，绝对开卷有益 —— 其实历史早已盖棺定论。}
\BookParagraph{就像那句名言：「Those who cannot remember the past are condemned to repeat it \footnote{Those who cannot remember the past are condemned to repeat it. — Jorge Agustín Nicolás Ruiz de Santayana y Borrás, The Life of Reason, Book I: Introduction and Reason in Common Sense (1905)}.」}
\BookDate{2019 年 4 月，多伦多}
\BookArticle{读博的情怀}{读博的情怀}
\BookQuote{尊敬的李老师，您好！我是您的书《幸福》的读者，目前在读电子信息专业的硕士研究生，明年就要毕业了，在未来道路选择的问题，想请教一下您的意见。目前我面临毕业后工作还是读博的选择，读博的话就在本校了，工作的话选项多一些，民营企业、国企这些打算都试试，主要是没考虑好要不要读博。我之所以读博，是周围亲戚和自己认为读书就一口气拿到最高学历，这样以后的晋升空间更大，前途选择也更多一些，但自己似乎也没什么学术理想，用不好听的话来讲就是拿个文凭。但是随着年龄的增长，觉得自己迟迟没有挣钱，加上得考虑更多现实问题，现在一直在犹豫了。周围人的观点都站在各自的立场上，所以希望老师能够给我一些意见，帮助我打开视野，谢谢您了！（问题来自微博私信）}
\BookSubtitle{一}
\BookParagraph{「读博」这个词儿，很容易让人误解。}
\BookParagraph{首先，博士虽然是一个学位，但是它不像本科或者硕士学位，不是仅仅靠好好读书好好学习就可以读完的。精确的说，「读」虽然很重要，但「写」比「读」重要得多。而写文章，考验的不是学习的能力，而是创造和发现以前没人想过的新问题和新解法的能力。而创造和发现的过程，则比学习的过程难一些，也苦一些。}
\BookParagraph{其次，博士的英文是「Doctor of Philosophy」，并没有「博古通今」或者「博采众长」等「广博」的含义。有一个成语叫「博大精深」，可是一个人的学问要想既做到「博大」又做到「精深」实在是很难的一件事儿。}
\BookParagraph{能做到「博大精深」的我仅能想到两个例子。一个是我外公的导师赵元任先生，兼通语言学、哲学、物理学、数学和音乐，可以说是近代学者中「博大精深」的范例。还有一个是当今的风云人物 Elon Musk，当年他开创 PayPal 玩儿电子支付那会儿，支付宝什么的还没出生呢。开到一半儿大概他觉得不够好玩儿，靠看书自学火箭技术，继而开创尽人皆知的 SpaceX。英文有句俗语叫「It’s not rocket science」，用来表达「不是很难」的意思 —— 可见火箭技术很难。可想而知，把一个公司玩儿到火箭技术独步天下就更难了。}
\BookParagraph{赵元任和 Elon Musk 之所以能做到「博大精深」，那是因为他们是天才。而如果「博士」只能「博大」和「精深」取其一，我们自然会偏爱「精深」。因为如果只「博大」而不「精深」，英文里叫做「jack of all trades, and master of none」，不是什么夸人的好话。而「精深」的好处在于，假以时日，无需天才，大多数人都可以做到。}
\BookParagraph{所以，哪天你终于博士毕了业去相亲，一位好奇的女生问你研究了这么多年研究的到底是什么东西时，你当然可以泛泛地回答「人工智能」，或者「利用人工智能方法的语音识别」；但是如果你希望能精确一些，其实应该回答「基于五层以上深度神经元网络框架的低信噪比场景下高精确度短语语音识别」—— 因为你的研究以「精深」见长，在这个极其狭窄的领域中水平极高，但是出了这个领域你一窍不通。}
\BookParagraph{当然，如果你善于聊天，可以这么给她介绍：「你玩过 Hey Siri 吗？」}
\BookSubtitle{二}
\BookParagraph{在你的担忧之中，提到了「挣钱」和「现实问题」。}
\BookParagraph{这个简单：如果你觉得收入是一件比较重要的事儿，那应该选择不读博士。}
\BookParagraph{这是因为除了少数大公司的研究部门，需要博士人才的工作 —— 比如大学教授 —— 工资普遍偏低，也没有股票和分红等其他附加收入来源，同时每年的工资涨幅普遍低于企业。}
\BookParagraph{在企业里工作，提高收入的最好办法是跳槽到更加赏识你的地方。而这一条在学术界基本不灵：跳到其他学校涨不了什么工资，要想涨工资得离开学术研究，跳槽去企业任职。}
\BookParagraph{当然了，这还没算为了读博士学位而少工作的这几年的沉没成本。}
\BookParagraph{咱们可以粗略地估算一下。博士毕业假设你去高校工作，工资十万（人民币或美元），每年工资涨幅百分之二，工作三十年的全部工资加起来是 10 (1 + 1.02 + 1.02\textsuperscript{2} + ... + 1.02\textsuperscript{29})，按照等比数列求和，这相当于 10 (1.02\textsuperscript{30} - 1) / (1.02 - 1) = 405 万元。而如果你没有读博直接去工作，起薪二十万，每年工资因为跳槽等因素平均涨幅百分之五，工作三十五年的全部工资加起来是 20 (1.05\textsuperscript{35} - 1) / (1.05 - 1) = 1806 万元，远远超过读博的总收入。}
\BookParagraph{当然话又说回来，去企业工作一般没有「铁饭碗」或者「终身教授」一说，有一定失业的风险。如果失业时正好大环境就业也不景气，一直失业下去也不好说。}
\BookParagraph{不过，我们可以计算一下在企业工作多少年收入才能和去高校工作持平：20 (1.05\textsuperscript{n} - 1) / (1.05 - 1) = 405，解此方程得到 n = 15。也就是说，如果你 25 岁开始工作，只需要工作到 40 岁，总收入就可以超过去高校任职到退休 —— 你瞧，40 岁以后的都是净赚！}
\BookSubtitle{三}
\BookParagraph{当然，这些都是非常粗略的计算，因为收入和很多因素有关，无法一概而论。比如说，我们的计算里没有考虑到现在高校教职等从事研究的工作越来越难找，如果你那「精深」的研究领域正好没有空缺，就连找到第一个工作都可能成为一个问题。}
\BookParagraph{比如十年前的一期《经济学人》杂志就曾经刊登了一篇文章，题目叫《为什么读博多半儿是浪费时间》（Why doing a PhD is often a waste of time）。}
\BookParagraph{其主要的论点是：博士毕业生数量大增，而对博士生的需求相对疲软，很多非常优秀的学生与其去读博士，不如做点儿别的更有意义的事情。}
\BookParagraph{你可能会问：说来说去，那到底什么样的人应该读博呢？}
\BookParagraph{如果你看了我的这篇文字，清晰地意识到了所有读博的坏处和风险，对钱多钱少毫不在乎，反倒是对「基于五层以上深度神经元网络框架的低信噪比场景下高精确度短语语音识别」有一种偏执般的热爱，那你应该读博。}
\BookParagraph{没错儿，读博士需要一种「情怀」，一种「明知山有虎、偏向虎山行」的执着，一种「西出阳关无故人」的悲壮。}
\BookParagraph{因为对于有的人来说，旷日持久地去琢磨一件「精深」的问题，是一种让人开心的事儿，觉得一辈子没白活。}
\BookParagraph{宁可飞蛾扑火，也在所不惜。}
\BookDate{2019 年 9 月，多伦多}
\BookArticle{一览众山小}{一览众山小}
\BookParagraph{为了更有效地管理家里的投资，最近迷上了 Google Sheets。}
\BookParagraph{记得八年前曾经读到过一篇博文，题目叫做《\href{http://www.antipope.org/charlie/blog-static/2013/10/why-microsoft-word-must-die.html}{Why Microsoft Word must Die}》。读完不禁拍案称绝，写了这么一小段微博：}
\BookParagraph{————}
\BookParagraph{读完 Charlie Stross 的博文「\href{http://www.antipope.org/charlie/blog-static/2013/10/why-microsoft-word-must-die.html}{Why Microsoft Word must Die}」，字字珠玑，感慨文字处理软件这一领域二十年来不但没有突破性的创新，还比当年 Mac OS 下的 Word 5.1 退步了很多。速度慢、功能复杂、排版困难、难用得一塌糊涂。觉得学术界的一大好处是不需要被迫用 Microsoft Word 写文章。}
\BookParagraph{————}
\BookParagraph{不知道为什么，这么多年来，我不但对 Microsoft Word 没有什么好感，更是死活学不会 Microsoft Excel。我觉得好用的工具给人的感觉是用的特别趁手，心想事成；而 Microsoft Excel 则恰恰相反，我想做的每一件超级简单的事儿 —— 比如按照某一列排序 —— 都有一种「心想而事不成」的感觉，特别别扭。每次用的时候学会的几招三脚猫的招数，到了下一次再想用又忘了。当然了，全世界其他人都会用的工具我不会用，只能说明我动手能力太差。}
\BookParagraph{这些年来需要处理表格的时候，我多半儿用的是 macOS 自带的 Numbers。Numbers 不仅界面漂亮，也比 Excel 好用多了，那种「心想事成」的感觉又回来了。我虽然觉得它好用，但是因为没什么其他人用，做出来的表格还是需要转换成 Excel 格式以后才能发给别人，还是不甚完美。}
\BookParagraph{管理家里的投资，多年来一直没有找到特别好用的工具。这几天下定决心做了一些调研，发现 Google Sheets 支持一个叫 \texttt{GOOGLEFINANCE ()} 的函数，可以直接在表格里实时更新股票市场的价格。不仅如此，简单学习了一下发现 Google Sheets 跟 Numbers 一样，用起来顺风顺水，省时省力。}
\BookParagraph{一旦爱上了则一发而不可收拾 —— 我花了几个小时，从头设计了一个题为「Live Portfolio Tracker」的表格，专门根据实时市场数据跟踪我所有的股票和债券投资，便于分析调整，颇有一种「会当凌绝顶，一览众山小」之感。}
\BookParagraph{用惯了 Google Sheets 还有个好处：下回再给别人发表格，我一定会洋洋得意地发个 Google Sheets 的链接 —— 想要 Excel? 哈哈自己转吧。}
\BookImage{images/google-sheets.png}{1.0}{}
\BookParagraph{\href{https://baochun.ca/images/google-sheets.png}{原图}}
\BookParagraph{\href{https://docs.google.com/spreadsheets/d/1bzKmkM6HUixKAA48pCFVEQRh645vXejmP3ToPGSdj6o/edit?usp=sharing}{Live Portfolio Tracker 示例}}
\BookDate{2021 年 3 月，香港}
\BookArticle{套瓷信}{套瓷信}
\BookParagraph{一位朋友说起，国内有专业的公司可以帮忙起草润色发给国外教授的「套瓷信」。正好这两天我也收到了一封「套瓷信」，从英文语法到行文的思路，信中有多处一眼就可以看出来的问题。比如信里抬头写的是「Dear Professor Baochun Li」，但是其实应该写「Dear Professor Li」或者「Dear Baochun」—— 前者是常规表示尊重，而后者则剑走偏锋，更为亲切，用「Dear」不用「Hi」来表示正式，而用名字来拉近距离。}
\BookParagraph{记得二十多年前我申请学校，也写过一百多封这样的「套瓷信」。那时清华计算机系只有网络教研组的一台 Compaq 原装 386 机可以收国际电子邮件，我的一百多封信都是用一台当时最为先进的彩色喷墨打印机打印，然后从邮局寄出的。如果教授回复了电子邮件，我再托网络组的同学把邮件用三点五寸的万胜牌软盘拷回来，在宿舍里我自己的计算机上读。}
\BookParagraph{我还清晰地记得，寄信时国际邮资刚刚涨价，一封信涨到了两块九，相当于两份食堂的小炒。这一百多封信，当时真的是下了血本呢。}
\BookParagraph{如今这个抖音微信的时代，自然没人再寄「套瓷信」了。不过我想，如果将来有一天我在办公室里拆看信件的时候，真的发现一封漂洋过海而来的「套瓷信」的话，一定会认真回复的。因为这种特异独行不按常理出牌的路数，说明寄信的人有创新的精神；而明知国际信件的邮资会血本无归还「偏向虎山行」，尤为令人感动。}
\BookParagraph{哈哈，我真的该开个皮包公司，凭着我多年的经验，专门给人起草和润色「套瓷信」。}
\BookDate{2021 年 5 月，香港}
\BookArticle{元宇宙}{元宇宙}
\BookQuote{李老师，最近「元宇宙」这个概念特别火，有人说人类即将进入一个新世界，《黑客帝国》将成为现实，有人说这只是一场雷声大雨点小的骗局。您怎么看「元宇宙」这个概念以及各大科技公司对它的重视？（问题来自微博私信）}
\BookSubtitle{一}
\BookParagraph{「元宇宙」确实有点儿忽悠。}
\BookParagraph{不过不是因为它太玄乎太科幻，而是因为它早就已经差不多实现了 —— 只是为了宣传给戴了一顶冠冕堂皇的帽子而已。}
\BookParagraph{你经常玩儿的各种在虚拟世界中游荡的游戏 —— 比如原神 —— 其实就是「元宇宙」。}
\BookParagraph{前一阵子股票大涨的 Roblox，当然也是「元宇宙」。}
\BookParagraph{在微软的飞行模拟游戏中，你可以单独在晴空万里的清晨中起飞，也可以在电闪雷鸣的夜晚中降落；在经典游戏「The Sims」中，你可以和朋友火热开聊，尽情 party—— 这些也都是「元宇宙」。}
\BookParagraph{如同斯皮尔伯格的电影「Ready Player One」中一样，在这样的「元宇宙」里，你可以以任何自己最钟爱的装扮出现，可以变成一辆自己心仪的跑车或者二战时的飞机，甚至可以在孤岛上建一座古城堡，在塔尖眺望远处的灯塔。}
\BookParagraph{在这样的「元宇宙」里，你可以尽情地爱上一个人，像「Ready Player One」里那样爱上她的迷彩装扮，像电影「Her」里那样迷恋她略带磁性的声音，你们可以一起购物，一起游览世界。}
\BookParagraph{注意，我的「元宇宙」和 Facebook 或者 Meta 的那个南辕北辙，大相迳庭。我的「元宇宙」里没有工作，只有玩儿；不需要什么高级眼镜，只需要一台普通的电脑；没有漂浮在半空中的上半身假人，只有各式装扮的各色人等在争奇斗艳，觥筹交错，人声鼎沸。}
\BookParagraph{我想我会爱上这样完美的「元宇宙」，无以自拔。}
\BookSubtitle{二}
\BookParagraph{虽然说「元宇宙」已经差不多实现了，但其实从计算机技术上来说，离真正实现还有一段距离。}
\BookParagraph{最难实现的其实跟虚拟现实没什么关系，而是一个既可以大家一起上线，又可以超低延迟的计算机网络。}
\BookParagraph{其实计算机技术发展到了今天，保证大家一起上线并不太难，保证低延迟也不是很难，可难就难在两个要求需要同时满足。}
\BookParagraph{比如直播技术可以几百万人一起上线，但是从现场大喊「新年快乐」到大家听到，大概需要十几秒到几十秒的延迟。}
\BookParagraph{如果换成了 Zoom 这样的实时在线会议，延迟可以保证只有一两百毫秒，但是只能最多一两千人同时在线。}
\BookParagraph{你一定会问，现在不是有了 5G 网络，可以支持超低延迟了吗？}
\BookParagraph{其实这个问题远远没有你想象的这么简单。计算机网络就像一个错综复杂的公路系统，有省道、国道，还有纵横交错的高速公路。数据包每走一段就会遇到一个路由器，就像汽车路过高速上的收费站，如果路上没什么车还好，否则需要排队等候 —— 车越多就越拥堵，数据包的延迟就越长。}
\BookParagraph{而 5G 网络的低延迟，只不过让你早点儿开出小区上路而已。}
\BookParagraph{可以说，如果将来「元宇宙」最终没火起来，多半儿是因为你的虚拟现实由于上线的人太多，堵在了高速上。}
\BookSubtitle{三}
\BookParagraph{我有个学生叫王菲。}
\BookParagraph{我一直认为，王菲同学的这个名字起得水平实在太高了，简直是神来之笔 —— 不但中文名极其好记，而且英文「Fei」一样简单易读，毫无歧义，其背后的故事则是不熟的人初次见面自然而然的「ice breaker」。}
\BookParagraph{王菲同学有一句特有哲理的口头禅：「想多了。」}
\BookParagraph{是啊，与其天天琢磨「元宇宙」什么的，不如痛快地踢一场球，坐一回过山车，跳一次伞，喝一顿酒，吃一顿铜锅涮羊肉，再谈一场一辈子无以忘怀的恋爱。}
\BookParagraph{就像 Ernest Cline 的小说「Ready Player One」最后写的：「我们坐在那里，手握着手，沉浸在那种奇妙的快乐中。」}
\BookDate{2022 年 1 月，多伦多}
\BookArticle{虎口遐想}{虎口遐想}
\BookParagraph{周日加班和学生一起写论文时偶然发现，一个已有的算法在我们的实验中运行，性能有点儿过于优异了。}
\BookParagraph{断断续续调试了一整天，终于在凌晨发现，在我一年前写的几行程序中有一个致命的问题，致使需要这几行程序的算法的性能远远超过其他类似算法。}
\BookParagraph{我想，为什么学生从来没有怀疑过这个算法的实验结果呢？}
\BookParagraph{英文中有一句话：「If it sounds too good to be true, it probably is. 」}
\BookParagraph{我一直觉得，无论是搞理论研究还是做工程项目，都要有一点儿科学精神。科学精神中最重要的，就是好奇和怀疑。对未知的事物永远保持好奇，对「too good to be true」的结果持怀疑态度，才能经常纠正错误，推陈出新。}
\BookParagraph{这「怀疑」的态度，我最早是跟当年 UIUC 上学时一个天才教授学的。这位教授是印度人，他的学生每次给他讲新的「idea」，他的回答永远是「This doesn’t work, right? The reasons are the following...」}
\BookParagraph{哈哈，幸亏我不是他的学生，否则肯定得急：这也不行那也不行，您倒是给我想一个「work」的啊！}
\BookParagraph{结果这位教授的学生，后来遍布美国名校。学生的学生，竟然可以回到 UIUC 任教职。}
\BookParagraph{至于「好奇」，我想起了刘雪峰。我近十年前在香港访问了一年，刘雪峰是因为经常跟我一起去学校食堂打饭而熟识的一个朋友。刘雪峰的才情，就表现在他什么东西都特别喜欢深入地琢磨一会儿「为什么」。}
\BookParagraph{比如我当年倡导，写论文时无论想到什么，要先简略迅速地写出来。然后再多给朋友和同学传看，慢慢在讨论中细化论文中算法设计的方案，经过多轮迭代完成。}
\BookParagraph{本来挺直觉的一个小建议，刘雪峰则分析得入木三分，觉得这是「步步为营」和「精益求精」的区别。我倡导的方式是「精益求精」，而传统的先看论文、再做实验，最后才开始写论文的方式，则是「步步为营」。由于研究不像工程项目，风险高而且跳跃性大，更适合「精益求精」。}
\BookParagraph{刘雪峰「好奇」了好几年，到北航任教以后，把他平常产生的各种千奇百怪地想法写成了一本书，把计算机科学和数学思维应用到实际生活中。书名本来叫「算法人生」，后来怕没人看，就换汤不换药地改了一个更为通俗的名字，叫「心中有数」。}
\BookParagraph{现在轮到我好奇了 —— 不知道这本书里有没有那么一章，把数学知识和科学精神发扬光大地用于谈恋爱，追心仪的女生成功概率百分之百，情场永远立于不败之地。}
\BookParagraph{要我是刘雪峰，书名多半儿叫「九阴真经」。}
\BookDate{2022 年 5 月，多伦多}
\BookArticle{Think Different}{Think Different}
\BookQuote{葆春老师，冒昧打扰您了，您接受跨专业学生的申请吗？如果接受的话，对于跨专业学生会更看重哪方面的能力呢？（问题来自微博私信）}
\BookSubtitle{一}
\BookParagraph{什么才是老师这个职业的最高境界？}
\BookParagraph{我一直认为，老师的最高境界不是把一个资质甚高的学生培养得更加优秀，而是把一个完全零基础的学生培养成合格的人才。}
\BookParagraph{比如我推崇备至的姜昆老师，他年轻大红大紫那会儿，收了一个比他更年轻的学生，连中文都说不利落，更别提说相声了。}
\BookParagraph{没错儿，这个学生就是加拿大人、我们多伦多大学的杰出校友，时年二十四岁的 Mark Rowswell 同学。}
\BookParagraph{简称大山。}
\BookParagraph{大山后来果真成为了一名合格的相声和脱口秀艺术家。有多合格呢？我多年以前去看过大山的现场脱口秀，精彩纷呈，令人捧腹。今年春晚里有一段姜昆的相声，很多人说相声里学广东话的那段是抄袭大山几年前的一段脱口秀。}
\BookParagraph{一位原来不会说中文的学生，牛到了竟然有人觉得是老师去抄袭学生 —— 这位老师的水平，可见一斑。}
\BookSubtitle{二}
\BookParagraph{但其实即使是姜昆老师，我猜也不是什么样儿的学生都收的。}
\BookParagraph{这多半儿不是因为他挑剔，而是因为实在是「带宽」不够。}
\BookParagraph{计算机网络中，带宽的单位是比特每秒，延迟的单位是毫秒。老师培养学生，就像一个低带宽高延迟的卫星网络。一个学生读一个博士学位，「延迟」大致是四五年，可以算是「高延迟」了。但是请注意，他的老师的「带宽」也是极其有限的：比如我的「带宽」大致是平均每年收两个学生，假设收到六十岁，那我的职业生涯也仅仅还能再培养二十个学生左右。}
\BookParagraph{不仅如此，一般来说水平越高的老师，相应的「带宽」也越低。斯坦福大学「十年磨一剑」终于把他的 TeX 排版系统「磨」得完美无缺的 Donald Knuth，据说一辈子只带了二十八个学生。}
\BookParagraph{正是由于「带宽」有限，像姜昆这样的优异师资其实是一种极其稀缺的资源。}
\BookParagraph{这就是为什么我从来不是很理解为了招博士生去高校做广告这种做法的逻辑。「酒香不怕巷子深」，吸引学生我认为应该依靠口口相传，如果主动去做广告，岂不有一点儿「酒」其实也不怎么香的嫌疑？}
\BookParagraph{十年前我曾经在微博中提到过，我的理想有三条：「每礼拜有功夫跑二十公里；赚回家的钱全给媳妇儿血拼刷卡给娃上学；哥已不在江湖，江湖还在传说他。」}
\BookSubtitle{三}
\BookParagraph{那大山作为一名「跨专业」的学生，需要什么样儿的素质，才能「天降大任于斯人也」、终将学以致用呢？}
\BookParagraph{我觉得大概也有三条。}
\BookParagraph{第一是愿意下苦工夫完整地做好一件事儿。}
\BookParagraph{二十世纪最著名的数学家之一 Paul Erdős 说过一句名言：「数学家的工作，就是把咖啡变成定理。」}
\BookParagraph{我在美国上学时的第一位导师是当时的系主任 Daniel Reed 教授。他跟现在的时尚明星马斯克差不多，都是每周一百小时的工作狂。我每周去他办公室开会，空气中永远有一丝淡淡的咖啡味儿。}
\BookParagraph{注意，Erdős、Reed 还有马斯克的一百小时跟我们反对的「996」不一样，他们是为了自己喜欢的事儿工作。他们的工作，就像是舞蹈家练舞、钢琴家练琴一样，是为了把一件自己热爱的事儿 —— 比如证明一个定理或者开创一家公司 —— 完整地做好。}
\BookParagraph{第二是有很强的交流沟通能力。}
\BookParagraph{这有点儿像传说中的北京出租车司机，个个都是「侃爷」。你还别小瞧这「侃大山」的能力，据我观察，大多数人其实并不具备。「谈笑有鸿儒，往来无白丁」，无论是几位好友小聚，还是机场咖啡厅闲聊，都可以侃侃而谈，快意人生，这才是优秀的「侃爷」。}
\BookParagraph{第三是喜爱写作。无论中文英文，中心思想明确，言简意赅，逻辑清晰，标点符号正确。}
\BookParagraph{在微博里，我时常可以收到一些陌生人的留言和问题，需要公开回答。每次回答之前，我都会修改一遍问题本身的语言问题。在这个过程中我发现，大多数的问题都需要修改，只有极少数的问题是可以直接引用的。}
\BookParagraph{还有，如果你写的文字不需要修改，说明你自己已经仔细地修改过几遍了。说不定，你属于那种愿意下苦工夫完整地做好一件事儿的人。}
\BookSubtitle{四}
\BookParagraph{十几年来，线上线下，我给国内高校上过不少课，本科生研究生的都有。}
\BookParagraph{在课堂上，我不太敢按照常规中途停下来问学生：「有什么问题吗？」}
\BookParagraph{因为，每次问这样的问题 —— 甚至是提出问题让大家解答 —— 结果都是令人有点儿尴尬的沉寂，直到我开始自问自答为止。}
\BookParagraph{这样一件小事儿，见微知著，也许那些愿意提问或者回答我的问题的学生，已经非常出色了。因为他们打破了「人云亦云」——「人非云亦非云」—— 的成规，因为他们跟常人不一样。}
\BookParagraph{如同苹果九十年代那个扭转乾坤的广告「Think Different」的开场白：「Here’s to the crazy ones.」}
\BookParagraph{与其罗列三条，不如就这么一条。}
\BookParagraph{就像我的学生王菲爱说的：想多了。}
\BookDate{2022 年 5 月，多伦多}
\BookArticle{谢谢你}{谢谢你}
\BookParagraph{英文真的是可以活到老学到老的。}
\BookParagraph{我今天请教了一下英文专业的女儿，「Thank you」现在年轻人一般如何回答。她说一般最常见的回答是「No problem」，有时候也说「My pleasure」，但是很少会说「Sure」或者「You are welcome」。}
\BookParagraph{想起小学二年级刚开始跟姨姥姥学英文的时候，她郑重其事地操着浓重的苏州口音，按照当年的课本，一句一句教我英文的场景。}
\BookParagraph{Long live Chairman Mao. My name is Xiao Ming. This is a desk. That’s a chair.}
\BookParagraph{那年夏天，在北大燕南园的一片蝉鸣声中，跟姨姥姥学到的「Thank you」的标准回答是「Not at all.」}
\BookParagraph{跟女儿学英文的另一发现是「Thank you」的另外一种回答也十分常见：「No worries.」神奇的是，别人说「I am sorry」也可以说这句来回答 —— 反正别人说什么你说「No worries」总是没错儿的。}
\BookParagraph{想起来前两天看到一幅抽象画，画中写着「Live more, worry less.」}
\BookParagraph{是啊，无论何年何月，何时何地，都要「live more」，都要「no worries」。}
\BookParagraph{哈哈，将来如果哪个学生跟我说「no worries」，一定是看过我的文字，说不定可以成为朋友。}
\BookParagraph{如果姨姥姥还健在，我想对她说一句：谢谢你。}
\BookDate{2022 年 9 月，多伦多}
\BookArticle{论文}{论文}
\BookParagraph{刚接到通知，我和王菲合作的论文得到了 IEEE INFOCOM 2023 的最佳论文奖（Best Paper Award）。}
\BookParagraph{二十多年来，我一共在 IEEE INFOCOM 这个计算机网络方向的国际会议上发表了八十篇论文。根据上海交通大学王新兵教授独创的 \href{https://www.acemap.info/ranking}{AceRanking} 网站，以 2000 年以来 INFOCOM 论文的引用 H-index 计算，我排行第一位。}
\BookParagraph{可我从来没有得到过 INFOCOM 的最佳论文奖，这次是头一回。}
\BookParagraph{不过这丝毫不稀奇：会议每年录取两百多篇论文，只有一两篇可以得到最佳论文。按照概率计算，绝大多数学者一辈子也得不到一次。}
\BookParagraph{我一直觉得，一篇理想中的论文应该像一件精心设计打磨的手工作品，需要从选题到实验，从行文到插图，不慌不忙，慢工出细活。论文的作者则各司其职，通力合作。}
\BookParagraph{王菲这篇题为「More than Enough is Too Much: Adaptive Defenses against Gradient Leakage in Production Federated Learning」的得奖论文，历时四个月，三个作者完美合作，缺一不可，基于我这之前一年多来自己写的联邦学习开源框架，实现了我关于论文写作的理想。}
\BookParagraph{论文的选题其实是因为我们实现已有的隐私保护方面论文时，发现这些论文的结论都是错误的。}
\BookParagraph{说它们错了其实也不是特别准确。准确地说，它们把联邦学习理解错了 —— 它们的假设在实际中并不成立。由于这些错误的假设，它们的隐私保护机制过于浓墨重彩，就像开着保时捷 911 去自由市场买菜，远不如骑辆永久二八自然。}
\BookParagraph{我记得当时看到我们的初步试验结果，王菲同学大笔一挥，为她的 \href{https://iqua.ece.toronto.edu/papers/feiwang-infocom23.pdf}{第一篇 INFOCOM 论文} 题名：}
\BookParagraph{More than Enough is Too Much.}
\BookParagraph{就如同她爱说的名言：想多了。}
\BookDate{2023 年 4 月，多伦多}
\BookArticle{成大名}{成大名}
\BookParagraph{今天加拿大工程院宣布了五十五位新晋院士的名单，我有幸成为其中之一。}
\BookParagraph{外公生前曾提及他年轻时有过三个人生理想：成大名、致大才、享大年。按照如今时髦的语言来说，「成大名」是粉丝很多，「致大才」是有真才实学，而「享大年」就更简单了：活得越长越好。}
\BookParagraph{现在想来，大概后两个理想都是为「成大名」准备的。有了真才实学，「上得了厅堂，下得了厨房」，既能写论文又能写代码，一直「为祖国健康工作五十年」，才会有更多的「粉丝」记住你。}
\BookParagraph{外公自己与世纪同龄，一生坚持著书立说，著作等身，桃李满天下。至八十六岁高龄，成为了一位中国近代著名的语言学家，完美地实现了自己的理想。}
\BookParagraph{我想，也许「成大名」的意义在于，百年之后会有人为你写一篇悼词。悼词之中所讲述的一个个与众不同的故事之中，一个鲜活的生命跃然纸上，令人既悠然神往、又黯然神伤。}
\BookDate{2023 年 6 月，多伦多}
\BookArticle{图灵大会}{图灵大会}
\BookParagraph{今年的 ACM 图灵大会提供了一个契机，刘云浩、李向阳和我三人终于又有机会在武汉相聚。}
\BookParagraph{云浩、向阳和我极有缘分。}
\BookParagraph{我们同一年上的清华。唯一不同的是，向阳是以数学奥林匹克竞赛选拔队成员的身份从清华附中提前一年就保送了的，云浩是从京工附中提前录取保送的，而我则参加了当年的高考。}
\BookParagraph{忘了谁曾经说过：向阳是数学家，葆春是作家，而云浩则是演说家。}
\BookParagraph{这之后的三十多年，数学家、作家和演说家走过的轨迹迥然不同，却好像游乐场中的过山车，能奇迹般地交织。数学家当了院长，作家写了本儿《幸福》，而演说家则开创了中国图灵大会。}
\BookParagraph{八年前在香港，我们同时评上了 IEEE Fellow，参加了 INFOCOM 开幕式上的颁奖典礼，并合影留念。}
\BookParagraph{那一年，网上还流传着一篇叫做《不是所有的男神都叫刘云浩》的长文。我跟瑶瑶说，这篇文章夸人的水平太有限了，应该让作家重新写一篇。}
\BookParagraph{回到多伦多一个阳光灿烂的日子，我略加斟酌，挥毫而就了一篇短文，题目就叫 \href{https://baochun.ca/post/长期共存互相吹捧/刘爷/}{《刘爷》}。}
\BookDate{2023 年 7 月，多伦多}
\BookArticle{《生活不是掷骰子：理性决策的贝叶斯思维》序言}{《生活不是掷骰子：理性决策的贝叶斯思维》序言}
\BookParagraph{刘雪峰是我多年的一个朋友，他太能写了。继几年前写了一本《心中有数》，很快又要出版两本书：一本专门讲贝叶斯定理，还有一本讲科技论文写作。他邀请我给两本书各写一篇序言，我欣然答应，在香港去北京的高铁上完成了其中一篇。同时希望可以向他学习，将来多写一些文字。}
\BookDate{2023 年 8 月 8 日}
\BookParagraph{————}
\BookSubtitle{一}
\BookParagraph{关于读书，近年来有两种说法。一种说法是因为网络时代信息越来越碎片化，号召大家多静下心来读书；还有一种说法是，爱拿「读书」说事儿的人，都是一帮不学无术的半文盲，以为声称「爱读书」怎么着能够帮自己装得有文化有追求。}
\BookParagraph{首先，我不是很喜欢「读书」这个词儿。本来「看书」这么一挺普通的事儿，一说成「读书」，就似乎成了一种需要沐浴更衣、正襟危坐才能从事的文化活动。}
\BookParagraph{那在这抖音小红书的时代，看书是不是过时了？}
\BookParagraph{我觉得答案其实挺简单的：这得看是什么书。}
\BookParagraph{有的书吧，写的时候作者就想写成「著作」，看得艰深难懂，确实有点儿看不进去。}
\BookParagraph{比如霍金的《时间简史》，当年风靡大江南北，可是我竟然没看完第一章。}
\BookParagraph{又比如斯坦福大学的著名教授高德纳（Donald Knuth），图灵奖获得者，曾经历经数十载写了一部《计算机程序设计艺术》（The Art of Computer Programming），鸿篇巨制，到现在都没写完。我也曾经看过一章，就没能继续看下去。}
\BookParagraph{但是也有的书，周末午后有个几个小时时间，就可以看个八九不离十，看完还觉得挺有收获，一时半会儿不会忘了。看这种书，我觉得真的可以贴切地用一个四字成语来形容：开卷有益。}
\BookParagraph{我看这本书，就有一种开卷有益的感觉。}
\BookSubtitle{二}
\BookParagraph{「飞机正在剧烈颠簸，会不会出事儿？」}
\BookParagraph{「对方微信回得不怎么积极，我该不该现在表白？」}
\BookParagraph{五花八门各式各样类似这样的选择，我们每天都在做。}
\BookParagraph{我们也经常把那些说话办事儿很靠谱的人归结到「情商高」的一类。确实，情商高的人，所做出的决策往往大多数都是正确的，可以省去很多烦心事儿。}
\BookParagraph{可你有没有想过，其实所谓情商高，有时候就是一种准确的直觉。}
\BookParagraph{而这种准确的直觉，其实是可以慢慢学习和培养的。}
\BookParagraph{这里我有一个坏消息和一个好消息。}
\BookParagraph{坏消息是，这个学习和培养的过程中，需要一点点数学思维。}
\BookParagraph{这门数学就是概率论。}
\BookParagraph{概率论虽然我们在大学都学过，但是学的时候背公式做题去了，学得快忘得快，直到后来全都不记得了。}
\BookParagraph{可不巧的是，想培养这种「情商高」所需要的直觉，还得需要一点儿概率论。}
\BookParagraph{这点儿概率论特别重要。重要到了其他部分的概率论都忘了都没关系，这一点儿概率论需要一辈子记得。}
\BookParagraph{这点儿概率论就是贝叶斯定理。}
\BookParagraph{不过还有个好消息。}
\BookParagraph{贝叶斯定理虽然不容易理解，但是仔细想想，也不是特别难。}
\BookParagraph{这本书的好处，就是可以结合很多生活中的例子 —— 比如点球大战守门员往哪儿扑 —— 让我们可以一个下午就把这点儿概率论想明白。}
\BookParagraph{说不定，想明白了贝叶斯定理，我们就可以慢慢地领悟到这种情商高所需要的直觉呢。}
\BookSubtitle{三}
\BookParagraph{雪峰是我多年来的一个朋友。}
\BookParagraph{他写这本书的初衷，我猜可能是希望我们在待人处事之中，不会是「一帮不学无术的半文盲」，而是真的不自主地用上点儿科学的思维方式。}
\BookParagraph{在这抖音小红书的时代，真的从头到尾认真写本书，其实是挺费力不讨好的一件事儿。}
\BookParagraph{雪峰把写书提高到了「人生意义」的层次 —— 他真心地觉得，写了书、有了读者，人生才真的有了意义。}
\BookParagraph{一个周末的午后也好，地铁通勤的途中也好，这本书或多或少看一小段儿，比起刷抖音小红书，也许更加「开卷有益」。}
\BookParagraph{「飞机正在剧烈颠簸，会不会出事儿？」}
\BookParagraph{答案可以在本书中找到。}
\BookParagraph{「对方微信回得不怎么积极，我该不该现在表白？」}
\BookParagraph{无论回不回微信，都永远不要表白。}
\BookDate{2023 年 8 月，多伦多}
\BookArticle{大陈儿}{大陈儿}
\BookQuote{昨天我看褚明宇老师的微博发了一个跟葆春老师聊天的视频《如何考上清华大学》。我看了下进度条：59 秒，就知道两位大佬是在调侃呢。}
\BookQuote{其实我对这个命题非常感兴趣。我现在是一所高中的物理老师也是班主任，我所带的班级是高二年级最好的班，肩负着考取清北的任务。我猜测葆春老师的粉丝的受教育水平应该都很好，保不齐有本科毕业于清华北大的人。}
\BookQuote{所以我是很想认认真真地讨论这个问题：清华北大是有种乎？还是真的通过后天努力、科学的方法和应考策略，能够达成这样的目标。}
\BookQuote{哦，我的学生都是河北考生。没错，衡水中学和石家庄二中这样的超级中学所在的省份。本科考清华非常难。}
\BookQuote{（问题来自「葆春和他的朋友们」）}
\BookParagraph{这个问题当时被褚老师抢答了。}
\BookParagraph{如果是我来正面回答这个问题，我会说，其实不应该为了上清华学习。如果恰巧考上了清华或者北大，那当然是喜事儿，再接再厉；如果上了其他学校，无论是所谓的 985，还是所谓的 211，其实没有太大的区别，不见得损失了什么。}
\BookParagraph{我上初三的时候，学校本来有一个直升本校高中的名额，可我一度曾想放弃这个名额，参加中考。}
\BookParagraph{我们班班主任大陈儿问我：参加中考干嘛呀？}
\BookParagraph{我说，我想考北京四中。}
\BookParagraph{当时的四中是北京最好的高中之一。}
\BookParagraph{大陈儿问：考上了四中有什么好处吗？}
\BookParagraph{我说：那里可以竞争的对手多啊。}
\BookParagraph{大陈儿的回答，我到现在还记得清清楚楚：}
\BookParagraph{「竞争对手不需要太多，有一个就足够了。」}
\BookDate{2024 年 3 月，多伦多}
\BookArticle{广告}{广告}
\BookParagraph{最近经常看到各式各样的招生广告。}
\BookParagraph{这类广告的意义自然在于，显出一位教授及其实验室的优秀之处，希望能够尽量吸引水平更高的学生来就读。}
\BookParagraph{但是我自从到多伦多大学工作以来，从来没有发过招生广告。}
\BookParagraph{一来我一直认为优秀的教授如同限量版的保时捷一样，无论从他们的时间和经费来源来看，一辈子总共能招收的学生都极其有限，属于一种稀缺资源。比如斯坦福大学著名教授 Donald Knuth，图灵奖获得者，据说其学术生涯一共仅仅指导了 28 个学生。所谓「好酒不怕巷子深」，当「好酒」的香气口口相传远远地飘过来，自然就不需要在胡同口放一块用于引流的广告牌。}
\BookParagraph{二来我虽然没有发过正经八百的招生广告，却多年一直在写属于我自己风格的「广告」—— 我的文字，我的微博，它们就是我的广告。喜欢我的文字的学生读者，如果恰巧也对我的专业方向感兴趣，我虽然没有明说，却是很希望他们能来占用我这里的「稀缺资源」的。也正因为如此，如果学生在「套瓷信」里提到了我以前的文字，我经常会另眼相看，绝对是一个加分项。}
\BookParagraph{但是最近有些事儿让我觉得，其实正经八百发一则招生广告也是一个不错的主意。}
\BookParagraph{我以前一个学生，现在在香港当教授。她收了一个博士生，刚刚读了半年，直接退学了。}
\BookParagraph{无独有偶，我自己的一位博士生，恰巧在教师节的今天，跟我「官宣」不读博士了，准备转为课程型硕士学位毕业，已经开始找工作了。他退学的原因是想早点儿到工业界找工作，早点儿赚大钱。}
\BookParagraph{如果一个学生在决定退学去工作之前问我的意见的话，我会跟他说，人的一生很长，其实不用特别着急打卡上班挣钱。博士学位虽然在读期间会损失不少工资收入，将来毕业了也并不能保证挣到更多的钱，但是博士研究是一种在痛苦中创新的过程，而创新则需要用语言文字交流沟通的能力，也同时需要较强的理论基础和动手能力。试图做研究写论文的过程，就是提高自己领导、沟通、总结、动手和理论推导能力的过程。而这些能力无论将来做什么都挺重要的。}
\BookParagraph{所以我想写一则招生广告。}
\BookParagraph{以前有些人问我，招收学生有什么标准，如何择优录取。我的标准一直挺简单的：数学、英文和编程能力这三项，有一强项即可，如果有俩强项，那就相当罕见了。所以我其实并没有真的「择优」录取，甚至有不少随意的成分。主要是因为我觉得吧，一位水平可以算是「稀缺资源」的教授，其真正的「产品」，不是他的荣誉，也不是他的论文，而是他培养的学生。如果一个学生本来就是天才，那即使不怎么培养，他依旧还是天才，并不能体现教授的水平。而如果一个学生本来基础薄弱，什么都不会，毕业时成为了一名可以胜任大学教职的合格人才，那才能体现这位教授培养人才的水平。}
\BookParagraph{再过两个月我就五十三岁了。可能将来退休之前，没有太多的精力和时间培养很多学生了。如果你是一位学生读者，觉得我属于我自己说的「酒香不怕巷子深」的稀缺资源；如果你数学、英文和编程至少有一项还比较拿手；如果你看过我以前写过的文字，跟我的想法挺有共鸣；如果你同意我「人的一生很长，不用着急打卡上班赚钱」的观点的话，欢迎你来申请加入我在多伦多大学的研究小组，攻读博士学位。}
\BookParagraph{是为广告。}
\BookDate{2024 年 9 月，多伦多}
\BookArticle{个人网站}{个人网站}
\BookParagraph{我平时读过不少网站，玩儿过不少产品，但是很多看过玩儿过就渐渐忘了。为了更好的记住当时的想法，我新开了个英文个人网站 （\href{https://baochun.org}{baochun.org}）。}
\BookParagraph{为什么要用英文呢？}
\BookParagraph{主要原因是由于我的中文水平太低，很多新词儿和专业名词完全想不起来，用中文写的话，速度会慢很多。}
\BookParagraph{而写作的速度，直接决定了一个网站是否可能长期更新。比如说吧，我刚看完一篇论文，有几句话想说，我希望能在手机上迅速把它写出来，电脑上自动同步，然后一个命令就可以完成编译和网站更新，跟发布在小红书上所用的时间相差无几。}
\BookParagraph{而个人网站比微博小红书等社交媒体的强项在于，可以用 Markdown 加入大量链接和原文的引用，可以完全没有评论区，连用于 analytics 的 trackers 都没有，干净利落，专注于内容本身。}
\BookParagraph{作为一种产品，一个经常更新的个人网站，本质上来说回归了互联网的初衷：从人之间的互联，回到了文章之间的互联。}
\BookParagraph{我写的内容主要是和新近流行的技术相关，也时而会关注一下近期的研究和论文。一些我觉得好用的产品，我也计划写一些使用体验 —— 因为没发到微博，不会有人说我「带货」了吧。}
\BookParagraph{比如这几天写的，主要是有关 DeepSeek 的一些链接和想法。}
\BookParagraph{至于中文写作，新的一年还会继续发表在微博和我的中文网站 \href{https://baochun.ca}{baochun.ca}。}
\BookParagraph{祝大家新年快乐。}
\BookDate{2025 年 1 月，多伦多}
\BookArticle{Yes and No}{Yes and No}
\BookQuote{我正在申请读博。看到周围同学十有八九都在做人工智能方向的研究，我是不是也该做人工智能？（问题来自微博留言）}
\BookSubtitle{一}
\BookParagraph{英文中有一个特别滑头的方式回答问题，叫做「Yes and no.」这样滑头的回答在英文语境中很常见，但我发现中文中却很少见。}
\BookParagraph{说它滑头，是因为这么回答似乎相当于什么也没说。甚至逻辑都不对：怎么可能既是「yes」又是「no」呢？这不是自相矛盾吗？}
\BookParagraph{但是其实这是一种相当不错的回答方式：正是因为其逻辑无法自洽，产生了一种突兀的悬念，吸引着你的注意力继续听下去，甚至会听完整个回答 —— 听完「yes」前一半儿还不够，不是还有「no」的那后一半儿吗？}
\BookParagraph{回答你的问题，我想标准答案也应该是：}
\BookParagraph{Yes and no.}
\BookSubtitle{二}
\BookParagraph{咱们先从「no」聊起。}
\BookParagraph{人工智能领域的「顶会」之一 ICML，今年收到了 12107 篇论文，录用了 3260 篇。这个数字有多震撼呢？仅仅是两年前， ICML 2023 的投稿量是今年的一半左右：6538 篇论文，18535 位作者。如果咱们假设作者比例维持不变，那今年 ICML 则有 34323 位作者。按照 ICML 2023 的统计数字，其中三分之二在攻读博士学位，考虑到人工智能领域还有其他「顶会」，我们可以得出一个极其粗略的结论：全世界范围内，至少有两万名学生在人工智能领域攻读博士学位。}
\BookParagraph{那如果咱今天开始攻读博士学位，四年后毕业。在四年期间潜心人工智能核心领域的研究，每年投稿 ICML，毕业时能在海外好一点儿的学校找到教职的概率有多高呢？}
\BookParagraph{我们以北美、香港、新加坡为例。北美地区有一定竞争力的高校只有一百所左右，香港只有六所，新加坡只有三所。假设四年后每一所学校的教职在人工智能领域都有一个空缺 —— 这个假设非常乐观（generous），比如多伦多大学我所在的系，今年就没有人工智能方向的空缺 —— 那相当于四年后的五千多名博士毕业生，竞争一百个空缺职位，成功概率差不多百分之二。}
\BookParagraph{按照项目申请、会议投稿等任何需要「卷」的场合，一旦成功率下降到 15\% 以下，「谋事」虽然「在人」，「成事」就真的只能「在天」了。}
\BookParagraph{如果不考虑研究兴趣，单纯目标驱动，而目标则是去学术界发展，那咱是否应该攻读人工智能核心领域 ——研究结果主要发表于人工智能「顶会」 —— 的博士学位？}
\BookParagraph{你觉得呢？}
\BookSubtitle{三}
\BookParagraph{那如果我们的目标不是去学术界，而是去工业界呢？}
\BookParagraph{去工业界发展，我们需要考虑一个问题：什么技能才是工业界比较稀缺的技能？}
\BookParagraph{我认为，一个人交流和演讲的能力 —— 尤其是用英语交流和演讲的能力 —— 应该是「软实力」中的首选。这是因为，据我观察周围的年轻人，交流和演讲的能力普遍欠缺，所以如果可以在读博期间多交流，多演讲，很可能会在去工业界竞争中占据一席之地。}
\BookParagraph{那「硬实力」呢？}
\BookParagraph{我发现，自从大模型两年前横空出世、各种「copilot」满天飞以来，真正硬核的编程能力，正在迅速地变成一种稀缺资源。}
\BookParagraph{这是因为这么一个怪圈：越是硬核的编程能力，越是需要大量的实战经验。而大量的实战经验，则已经被大模型「卷」成了一种可遇而不可求的奢侈品。}
\BookParagraph{所以，如果目标是去工业界，那首要的是培养硬核编程的「硬实力」和交流演讲的「软实力」。而培养这两种实力都和研究方向没什么关系。甚至可以说，人工智能的核心研究主要是在算法理论领域的研究，对硬核编程能力的要求不高 —— 现在中学生都可以写 Python 了。}
\BookSubtitle{四}
\BookParagraph{王朔的经典之一的题目叫「一半是火焰，一半是海水」。}
\BookParagraph{英文里经常说：「I have some good news and some bad news for you. Let’s have the bad news first. 」}
\BookParagraph{那「好消息」是什么？}
\BookParagraph{即使你不「做」人工智能，你做的其实也是人工智能。或者咱们可以说，你做的不是人工智能的核心，而是其周边领域。}
\BookParagraph{传统的计算机网络、分布式系统、物联网、程序语言设计、数据库、安全隐私保护等领域，自然都需要人工智能技术。}
\BookParagraph{而方兴未艾的新兴领域 —— 比如具身智能，比如多智能体系统 —— 自然更需要人工智能技术来支撑。}
\BookParagraph{我们进入了一个前所未有的新时代：这个时代里，人工智能技术贯穿始终，遍布了每一个角落。大模型、强化学习等技术，无所不在，早已成为任何领域培养人才的必修课。}
\BookParagraph{在这样激动人心的大时代中，我们所有人所「做」的，无论传统还是新潮，无论理论还是实践，其实都是人工智能的周边领域。}
\BookParagraph{人工智能就如同繁星满天的夜空下的大海，深邃、空旷、一望无际，令人既害怕又向往。}
\BookParagraph{而我们的船帆已然升起，等太阳一出来，随时准备启航。}
\BookDate{2025 年 7 月，香港}
\BookArticle{芝麻开门}{芝麻开门}
\BookQuote{冒昧打扰，首先感谢您在微博上分享的内容，我偶然间刷到您的主页，被每一篇文章深深打动，感受到您对学术的热爱与深邃的思考，令我受益匪浅。恰逢今年我也刚刚进入多伦多大学电子与计算机工程专业攻读硕士（授课型），在适应新环境的过程中，我对AI尤其是大模型方向产生了更为浓厚的兴趣。在查阅资料与了解学界动态的过程中，得知您在相关领域有着深厚的研究积累和出色的成果，因此斗胆写信向您咨询，不知是否有机会参与您的课题组或项目，哪怕是以协助的方式，也非常愿意从基础做起，跟随您学习、积累经验。（问题来着微博留言）}
\BookParagraph{前几天，恰巧有机会见到苏黎世联邦理工学院（ETH Zurich）计算机网络领域的著名教授 Adrian Perrig，聊了几次天儿。}
\BookParagraph{Adrian Perrig 教授有多厉害呢？}
\BookParagraph{他在谷歌学术上的 H-index 高达 106，相当于有 106 篇论文引用数在 106 以上。H-index 的特点是，它越高，就越难以继续上升。research.com 显示，根据 H-index 排名，他在计算机科学领域是瑞士第十一位。由于计算机网络领域是一个相对比较冷门的领域，他这样的 H-index，已臻「惊为天人」的水准。}
\BookParagraph{我们「宾主就共同关心的话题，广泛地交换了意见」。比如我们俩都对异步 Rust 编程中的并发机制感兴趣，尤其对于多协程并发时是否应该对数据结构加锁和如何加锁的问题，讨论得还挺深入。}
\BookParagraph{有一次我们吃饭，正好聊到了「套瓷」信如何写的问题。发现我们俩的观察和发现几乎完全独立，却有异曲同工之妙：我们都极少 —— 如果高标准严要求的话，甚至从来没有 —— 收到过令人眼前一亮的套瓷信。}
\BookParagraph{那什么才是令人眼前一亮的套瓷信？}
\BookParagraph{首先，无论是中文还是英文，一封套瓷信可以直接展现作者的文笔。写套瓷信和写论文一样，应该从自己的姓名开始，开门见山，言简意赅，有自己独到的文风。相对而言，大模型写出来的文字，往往显得过于正式而平庸，没有自己的特点。这就是我为什么希望学生尽量不要用大模型 —— 用的时间一长，即使自己亲自写，也可能永远写不出文采飞扬的文字来了。}
\BookParagraph{其次，套瓷信需要有「点睛之笔」。}
\BookParagraph{什么是「点睛之笔」呢？}
\BookParagraph{这个「点睛之笔」的特别之处在于，简单一看，就可以知道这封套瓷信是专门给收信人写的，不可能发给其他任何人。}
\BookParagraph{比如，Perrig 教授十几年如一日，潜心研究和推广下一代全球互联网的设计和实现，名字叫做「SCION」。他为了 SCION 的推广，花费了大量心血和时间，甚至亲自写代码、写书、写用户手册。}
\BookParagraph{然而他发现，收到的套瓷信中，极少有提到 SCION 的，更不用说提到其中设计和实现的细节。甚至可以说，只要套瓷信中提到了这个词儿，他都会另眼相看，至少会花时间回复邮件。如果还能深入浅出地提到其中关键的设计亮点 —— 比如同时采用大量路径的多路径传输 —— 那「套瓷」的成功率就更高了。}
\BookParagraph{原因很简单：连套瓷信都能如此认真对待的学生，将来对研究也多半儿会对细节极其重视。而因为提到「SCION」的学生极其罕见，也突出了作者特立独行、和别人迥然不同的品质。}
\BookParagraph{这就如同《一千零一夜》中那篇家喻户晓的阿拉伯民间故事《阿里巴巴和四十大盗》里，阿里巴巴对着藏宝的山门喊出的暗语：}
\BookParagraph{芝麻开门。}
\BookDate{2025 年 7 月，香港}
\BookArticle{Trivial and Brilliant}{Trivial and Brilliant}
\BookQuote{感谢老师的微博分享，以后我也会更加注重套瓷信的独特细节。}
\BookQuote{我是电气工程专业博士生，曾反复听过您的讲座「Writing Prefect Papers」，受益匪浅，也因此发表过几篇 SCI 论文。部分工作发表在中科院一区 TOP 期刊，但也意识到自己的工作是增量式的（「灌水」），并没有很强的创新性以及理论和工程价值，能发表其实有很大运气成分。}
\BookQuote{我近期去外地参加了一个数学系的优化与运筹的短期培训，发现了一个很小众的领域，在我们专业的文献更是寥寥无几。但是做小众的工作风险也显而易见：即使投入大量的时间和精力，也极有可能是失败的，一将功成万骨枯。与其这样，可能还不如「灌水」，增加论文的数量，或许利于找工作。}
\BookQuote{请问老师，您怎么看待论文的质量与数量的关系？谢谢！（问题来自微信）}
\BookSubtitle{一}
\BookParagraph{我所在的多伦多大学电子和计算机工程系，每年都有一两个计算机工程领域的教职空缺。}
\BookParagraph{我们的面试分为两轮。而第一轮线上面试的问题，一般会提前一两天分享给面试者，可以有针对性的准备。}
\BookParagraph{这里开门见山的第一个问题是：}
\BookParagraph{Can you discuss a recent paper that you are particularly proud of?}
\BookParagraph{「请你介绍一篇自己最引以为豪的论文。」}
\BookParagraph{注意，这里明确说明是一篇论文，而不是几篇。也就是说，即使你有一百篇论文，也请你自行选择其中一篇，给大家讲一讲，这篇论文为什么可以代表你的最高学术水平。}
\BookParagraph{这个问题很明显不是我们的独创，甚至是老生常谈的一个问题。早在 2004 年，我去普林斯顿大学开一个学术会议，开场的基调演讲是信息论领域 MIT 的著名教授 Bob Gallager。他的演讲中就提到，评价一个申请教职的候选人，应该看其最出彩的那一篇论文。}
\BookParagraph{可以说，1960 年时的 MIT，就是因为采用了这个评价标准，才会破格录用了 Bob Gallager 本人当助理教授。他当时的博士论文，提出了信息论领域中石破天惊的低密度奇偶校验码（low-density parity-check codes），但是该论文直到三年后才发表。}
\BookParagraph{你的问题中提到，你有好几篇 SCI 论文已经发表。}
\BookParagraph{这种说法很常见：很多年轻学者开会上见到，自我介绍时也会说自己有多少多少篇「CCF-A」、多少多少篇「顶会」论文。}
\BookParagraph{我想，他们要是真的去找 Gallager 老爷子索要推荐信，他一定会问：}
\BookParagraph{你最出色的一篇论文是哪一篇？}
\BookSubtitle{二}
\BookParagraph{我自己的学术生涯中，也写过好几百篇论文。}
\BookParagraph{其中印象最深的一篇，二十一年前发表在一个名不见经传的学术会议上。这个会议是在我母校 UIUC 附近的一个叫做 Allerton 的小地方开的。我没有查过，但是我觉得大概率连 CCF-C 的会议推荐列表都没上。}
\BookParagraph{而这篇论文呢，其实没有做出任何贡献：它既没有证明任何定理，也没做任何实验。按照现在的高标准，甚至连发表都有点儿难。}
\BookParagraph{这篇论文仅仅是提出了一个猜想：网络编码在一对一的单播网络流（unicast network sessions）的场景下，相比不用网络编码，没有任何优势。}
\BookParagraph{这篇我和我当时的学生、现在任教于清华大学的李宗鹏教授合作完成的论文，后来在编码理论圈内叫做「李和李猜想」（the Li and Li conjecture」。}
\BookParagraph{很多学者研究过这个猜想。但即使是在二十一年之后的现在，也仍然没有被攻克。}
\BookParagraph{乔布斯九七年回归苹果时，把所有的项目负责人叫到他办公室，问了几个问题。其中第一个问题是：「如果资金不是问题，你计划把你的项目发展成什么样儿？」}
\BookParagraph{同样的问题你也可以自己问自己：如果钱不是一个问题，你希望如何利用自己的宝贵时间资源，做什么样的研究？}
\BookParagraph{比如，你会不会也给世界提出一个猜想，简单得出奇，却无人能够解决？}
\BookParagraph{就好像图灵奖得主 Jim Gray 评价同样也是图灵奖得主的 Leslie Lamport 写的一篇题为「Time, Clocks and the Ordering of Events in a Distributed System」的著名高引论文时说：}
\BookParagraph{It’s trivial, and it’s brilliant.}
\BookSubtitle{三}
\BookParagraph{那如果你看着自己的论文，实在找不出一篇自己觉得拿得出手的，怎么办？}
\BookParagraph{答案很简单：现在开始写也完全来得及。}
\BookParagraph{当然了，「钱不是万能的，但没有钱是万万不能的。」生活所迫，我们大概都需要占用大量时间，写一些自己都觉得没什么价值的论文。}
\BookParagraph{可当晚上孩子终于酣睡时，嘈杂的高铁站里等车时，航班上刷剧没有网络时，我们能否利用这些空闲时间，少刷点儿抖音小红书，一鼓作气，励精图治，铺开四百字的稿纸，开始写一篇自己引以为豪的论文呢？}
\BookParagraph{这篇论文，可能需要五年，也可能需要十年，但是它将来终于完成时，你可以自豪的宣布：这是你最精彩的一篇论文，没有之一。}
\BookParagraph{比如我自己，最近五年来，就在写这样一篇论文。}
\BookParagraph{这篇论文提出了一个全新打造的仿真平台，可以将传统离散事件网络仿真的速度，提高二至三个数量级。也就是说，原来几小时才能完成的仿真，现在只需要几秒钟。不仅论文全文是我亲自一个字一个字撰写而成，而且其绝大多数代码，短小精悍，「麻雀虽小，五脏俱全」，也是我几年来在酒店里，在航班上，一点儿点儿亲自设计实现的。可以说，这篇论文对应的系统实现，已经代表了我自从中学学习计算机编程以来的最高水平。}
\BookParagraph{我想，审稿人看完了论文，如果因为「it’s trivial」而拒稿的话，我会回答他：}
\BookParagraph{But it’s brilliant.}
\BookDate{2025 年 7 月，香港}
\BookArticle{高德纳}{高德纳}
\BookParagraph{前两天跟一个北大的学生面试。他问我，为什么研究组里只有六个学生？}
\BookParagraph{我问：「你知道Donald Knuth吗？他整个学术生涯，只毕业了28个博士研究生。」}
\BookParagraph{很可惜，这个学生并不知道 Donald Knuth 何许人也。}
\BookParagraph{闪念之间，我差点儿想给他讲一个长长的故事。讲 Knuth 教授如何花了整整十年时间，自己一个人设计和实现了 TEX 排版系统，仅仅是因为他嫌出版社出版他写的书时，排版太难看；讲他的 TEX 如何三十年前发展到 3.14159 版本之后，每升级一版就在小数点后增加一位；讲他写的算法书里，为什么要引入 MIX 汇编语言；讲他那面值 2.56 美元的支票为什么无人兑现；讲 Knuth 的那句名言：「Beware of bugs in the above code; I have only proved it correct, not tried it.」；讲我们以前的系主任，坚决只用 TEX，只因为嫌同是图灵奖得主的 Leslie Lamport 写的 LaTeX 不够纯粹；甚至想把书架上那本 Knuth 写的《Digital Typography》的书拿下来，给他讲讲里面的 METAFONT。}
\BookParagraph{最后终于想到，现在早已是 AI 的大时代，其实不知道这些陈年旧事儿也完全可以理解。}
\BookParagraph{「Sergey Brin 你知道吧？他是 Donald Knuth 的学生（Jeffrey Ullman，图灵奖得主）的学生（Héctor García-Molina，美国国家工程院院士）的学生。」}
\BookDate{2026 年 2 月，多伦多}
\BookArticle{波澜不惊}{波澜不惊}
\BookParagraph{Andrej Kapathy 最近发布的 Autoresearch 之所以炙手可热，原因很简单：只要有一个benchmark 和一段代码，AI agent 就可以无限循环下去，找针对这个 benchmark 的最优解。}
\BookParagraph{我亲手在 Pi Coding Agent 里试验了一下，一两个小时就把我的 Days 网络离散事件仿真平台的性能提高了26\%，而 Days 现有的实现已经非常不错了 —— 我最近才刚刚花了两天时间，完成过一轮对性能的优化。}
\BookParagraph{Autoresearch 还指明了如何利用 AI 来帮助我们做研究的一个重要方向：要想做好研究，一定要先树立一个 benchmark。无论是运行时间还是安全测试，这个 benchmark 必须是一个可以量化的指标。这其实也是几乎所有工程学科中研究的共性：没有这样的 benchmark，那「公说公有理、婆说婆有理」，研究就无法一步一步在前人的基础上继续开展下去。}
\BookParagraph{一旦有了这样的 benchmark，有了实现一个解决方案的基础代码，理论上来说，AI agent 就可以不断推陈出新：想新的解法，想如何调整现有的算法，去提高这个 benchmark 所代表的性能。「山穷水尽疑无路」对于人来说是一种挫折，但是对于 AI Agent 来说，只是一些 tokens 而已；「金榜题名时」对人来说是一种幸福，而对于 AI agents 来说，也仅仅是一些 tokens 罢了。}
\BookParagraph{我预计，在不久的将来，正是因为这种「波澜不惊」的定力，AI Agent 一定会有能力独立做出有独创性、令人拍案叫绝的研究。}
\BookParagraph{正如 Andrej Kapathy 所说：「This is what post-AGI feels like.」}
\BookImage{images/weibo-5276734135997229-1-2b232b7a4f.jpg}{0.7}{}
\BookDate{2026 年 3 月，多伦多}
\BookArticle{寓言故事}{寓言故事}
\BookParagraph{一篇写了两个月的论文，终于快写完准备提交了。本来想亲自动笔写结语 (Concluding Remarks) 的，但转念一想，说不定让 Fable 5 先写个初稿，再手工修改一下，可以节省一些时间。当然这样做有个风险：如果 Fable 5 —— 是不是可以直译为「寓言故事」 —— 写得水平不高，也许会少了我自己对着一张白纸油然而生出来的那种有感而发的感觉。抱着尝试一下也未尝不可的心态，我写了如下提示词：}
\BookQuote{Now that the experiments section is almost completed, read the entire paper again (the first four sections have finalized and frozen), and REWRITE the concluding remarks with your best possible effort. Do not repeat any previous sentences from the abstract or the introduction. Write in a lucid, convincing way, and emphasize on the originality of this entire project, which spans over 360K LOC including tests.}
\BookParagraph{Fable 5 沉思片刻，洋洋洒洒写了三个自然段。虽然篇幅有点长，文字竟然令人刮目相看，拍案叫绝。无奈篇幅限制，我提示它缩写成一个自然段：}
\BookQuote{Unfortunately, the concluding remarks is too long. It should have only about the same length as the abstract, say fewer than 200 words and just one paragraph (mostly due to space constraints).}
\BookParagraph{Fable 5 一蹴而就，全篇192字。完全没有想到，缩写后的版本，竟然比原来的版本更为惊艳，技惊四座：}
\BookQuote{Self-modifying code and distributed state are, separately, two of the oldest hazards in computing; agent workflows now demand both at once. In this paper, our answer is that one mechanism suffices for both: when the only way a running task can change its workflow is a tagged operation that a transactional store admits once and remembers forever, distribution is survived by durability and self-modification is disciplined by admission. Our theory names the invariants admission preserves, our runtime hides admission behind an ordinary await, and our measurements certify that the hiding is honest, from ten thousand injected faults without a corrupted invariant to iteration priced within one percent of its closed form. Flakes is no veneer over an existing engine: everything from the transactional store to the harness SDK was designed from first principles —— over 360,000 lines of code with its tests —— and a substantial fraction of it was planned, written, and reviewed by harnesses running on Flakes itself. A system that stays safe while its own authors are agents inside it is, we strongly believe, an existence proof for the infrastructure the agent era requires.}
\BookParagraph{如此大家风范的文字，我想就算我亲自动手，还真不一定写得出来。尤其是细节 —— 比如「Flakes is no veneer over an existing engine」 —— 才得以见其写作功力。就像平衡木体操的下法一样，文字的最后一句收尾也极其重要。而这篇文字的最后一句竟然是其不断强调的 agent 写成的，不禁令人扼腕 —— 时代真的不同了。}
\BookDate{2026 年 7 月，香港}
\BookArticle{猜想}{猜想}
\BookParagraph{这几天，王虹和邓煜喜获菲尔兹奖，登上了热门新闻。与其研究他们的经历、个人喜好甚至穿搭，我去仔细研究了一下他们证明的猜想。无论是挂谷猜想还是希尔伯特第六问题，都非常有趣，而且一看就知道难于上青天。不说别的，仅仅是挂谷猜想中的维度定义，我想了半天才搞明白。}
\BookParagraph{说起猜想，我和早年的学生、如今清华大学的教授李宗鹏老师，也曾经于2004年写了一篇论文，提出了一个信息论和网络编码理论中的问题：在无向（undirected）网络中，如果有多个单播信息流（unicast flows），网络编码到底比起不编码有没有什么好处？这篇论文提出的猜想是，网络编码其实没有任何优势。论文发表后，这个猜想被后人命名为李和李猜想（Li and Li conjecture）。后来大家发现，这个猜想也非常难 —— 无论是证明还是举出反例 —— 难到了直到二十二年后的今天，还没有人能够解决。}
\BookParagraph{这几天我忽发奇想，人工智能这么聪明，是否能证明这个猜想呢？说干就干：我的「智能体」们思考了两天两夜，写了两百多页纸的论文，最后也只是证明了一些特例（比如11条边以下的小型网络中猜想成立），离解决问题本身 —— 智能体们把它亲切地称呼为「the global Li and Li」—— 还十分遥远。看来，人工智能也不是万能的：还没到挂谷猜想，仅仅是一个「Li and Li」就能令其望洋兴叹、爱莫能助了。}
\BookDate{2026 年 7 月，香港}
\BookArticle{破折号}{破折号}
\BookParagraph{我的写作风格超级爱用破折号、冒号和分号，现在只有两个办法：一是在论文中加一个脚注，说明破折号都是我自己写的；二是用AI帮我去掉所有破折号和分号。结果我的脚注被AI 审稿人提醒风险太高，建议去掉破折号更保险。}
\BookParagraph{我个人观察预测，大概一两年以后，所有论文都会是AI写的，因为实在写得太快了。如果手写论文，我需要每天写十小时，写差不多两周时间，我还算是写得快的；但是如果我用Fable 5来写，质量相似，可以在三至五小时内写完全文。我为什么不去花两周时间做更好的研究，而是去手写论文呢？}
\BookParagraph{不仅论文将会是全文AI写的，连工作也一定会是AI做的。以我的经验，AI做一天的工作，差不多需要一个小团队做一年以上，还需要是水平很高的团队来做。}
\BookParagraph{所以其实更应该考虑的问题是如何最大程度地利用AI做更好的研究，而不是是否允许用AI —— 时代早已不同，不允许用AI 是螳臂当车，而火车已经离开车站了。}
\BookParagraph{脚注：这篇文字中的破折号是我手工输入的。}
\BookDate{2026 年 7 月，香港}
\BookPart{长期共存 互相吹捧}{长\\期\\共\\存\\[18pt]互\\相\\吹\\捧}
\BookArticle{朋　友}{朋友}
\BookParagraph{褚兄旅居纽约，是我多年的一个朋友。说是朋友，我却没有一点儿比得上他。褚兄德智体全面发展，在众多兴趣爱好均能展现出较高的造诣。他从小弹钢琴，我却是半个乐盲。他朋友众多，三教九流，呼风唤雨，我则枯守灯下，终日伏案，是个书读得并不多的书呆子。他博文强记，引经据典，我则连个名字都背不出来，经常忘事儿。他能操着一口流利的英语，连说半个钟头不犯一个语法错误，我则临阵怯场，发音不准，白白地去了多次清华的英语角。他当学生会主席的年头，我只是个小小的办事员。本着「长期共存，互相吹捧」的原则，他却吹不过我，那只因他的优点实在不难找。如果用八个字来形容他，那就是「中学为体，西学为用」；如果用寥寥四个字，那是「大巧若拙」；如果规定用一个字，那就是一个「牛」。}
\BookParagraph{褚兄和我当年都热爱喝酒和火锅，是典型的「酒肉朋友」。继而以忙为由，天各一方，疏于联络，他走他的阳关道，我则过我的独木桥。他一晃成为纽约的精英一族，出没于华尔街中国精英的聚会中，以其淳厚的文化底蕴，匪夷所思的口才，技惊当场。我则一地鸡毛，买车买房，余墨已尽，沦为小资严重且五十步笑百步的普通中产阶级。褚兄最害怕的是可预知的将来，还有青春不再的心情，在我这儿都成了最好的注脚。}
\BookParagraph{可至少褚兄和我都依然热爱自幼长大的北京，爱北京话，北京人，和北京的小吃。虽然现在的北京有变质的嫌疑，但要是把星吧必胜客什么的搬走，再跟出租汽车司机闲侃几句，那当年的文化和气质没准儿尚在。有点儿不同的是，褚兄经常回根据地北京转悠，买几本书，尝尝街头的小菜，看看咖啡店中活泼年轻的姑娘。我想会有那么一天，我和褚兄在北京城的胡同里偶遇，相约尝尝新开的火锅，辅以两瓶醇香的二锅头。}
\BookDate{2005 年 3 月，多伦多}
\BookArticle{刘　爷}{刘爷}
\BookParagraph{我跟刘爷不熟。一向不是很喜欢各种各样的尊称，比如「大牛」、「大伽」、「大佬」、「男神」等等。在我看来，这些都跟多年前带引号的精英一样，略含贬义。即使是刘教授、刘老师、刘院长等称呼，都尽显敬而远之般地客套。只有京味儿的「刘爷」，干净利落，诚挚而不失跳脱。要是穿越回几百年前，应该叫「刘大侠」，取其豪爽正直之意，行侠仗义之风。}
\BookParagraph{我跟刘爷一共没见过几面，每次都本着「长期共存，互相吹捧」的大政方针，交流共同关心的话题。我们虽然不熟，但是小有缘分：同年同月出生，黑里透白的头发，校友，还都是北京人。可惜除了这几条，其他方面是天差地远。刘爷一眼看去就是标准的\BookLineBreak{}「领军人物」，德智体全面发展，要是早早上个「爸爸去哪儿」，一定红遍大江南北。他专业打过球，研究过清史，写过畅销教材，博闻强记，经常引经据典；我却体育将将及格，没读过几本书更没写过书，记忆力奇差无比，经常记不住人。他朋友学生无数，三教九流，呼风唤雨，觥筹交错之间，谈笑有鸿儒，往来无白丁；我则枯守灯下，终日伏案，既调不了素琴，也阅不来金经。}
\BookParagraph{刘爷著述颇丰，虽然我们研究方向相仿，但坚决算是两类人——有 Mobicom 论文和没有的——而我则隶属于后一类。可惜的是，认识刘爷有年头了，却从来没有机会去人家的基地好好学习天天向上一回，然后合作写几行代码，几段论文，争取哪天能跻身前一类。几年前一次 Mobicom TPC 开会碰到刘爷，说起这事儿，他立马翻出一张 ACM 给他发的信函，找到背面的空白处，大笔一挥，手写了一封邀请函给我，寥寥数语，笔势龙飞凤舞，语气落落大方。信我一直收藏着，仿佛是一张不断升值的猴年邮票，等着套现。}
\BookParagraph{闲下来的时候，有时候也在想，和刘爷的差距是如何炼成的。大概是因为刘爷在全国各大机场等待登机时，我却在超市买今年的新土，准备春天给后院儿的花坛铺上。刘爷著书立说之余正在联系出版时，我却在钢琴前跟孩子们讨价还价，想着那几首还没弹熟的曲子。刘爷盛装在颁奖典礼上对着麦克风发表领奖感言时，我却在暖暖的壁炉前，慵懒地刷着朋友圈。刘爷携家人在迪拜帆船酒店度假时，我却在早春的凌晨，毫无倦意，众里寻他千百度般地苦苦琢磨几行动不动就 crash 的代码。十年二十年一晃而过，刘爷已是中国学术界的精英翘楚，博古通今，著作等身，粉丝遍布大江南北。我则江郎才尽，泯然众人，有点儿像武侠小说中那些着墨不多却玩物丧志无所作为的小人物。}
\BookParagraph{头发虽然白了，至少还有；记性虽然每况愈下，至少还记得住刘爷这样正邪通吃的奇人。如此想来，其实不思进取也挺不错，也用不着套现刘爷的亲笔信去跟他学习了。日子一天天过去，像刘爷这样的男神和少侠们如好人好事般不断涌现，大浪淘沙，在我眼前翩翩起舞，各擅胜场。我乐意当个观众，关注着天涯的八卦，知乎的论战，巨则指点一下江山，细则新写几行代码，让梦想就此止步，渐行渐远。}
\BookDate{2015 年 5 月，多伦多}
\BookArticle{朋友圈}{朋友圈}
\BookParagraph{「向上移动」里有一个问题，问大家最不喜欢什么样的朋友圈现象。我想反其道而行之，聊聊我喜欢的朋友圈风格吧。}
\BookParagraph{我最喜欢的「动态」是原创长文。不是发表在微信公众号或者知乎的原文链接，而是直接可以阅读的长文。回忆如烟的往事也好，描述此时此刻的体验也好，最好分段明确，表达清晰，如身临其境。文字给人的感受，远远超过几张过度压缩的照片。}
\BookParagraph{可惜的是，由于微信朋友圈这个产品自身设计的问题，这样的原创长文不但设有没有任何意义的字数限制，而且无法和照片混排，必须文字在前，照片在后。这就好像一篇议论文无法「夹叙夹议」而必须先叙后议一样，读来索然无味。}
\BookParagraph{即使已经花时间写好了这样的长文，必须分段发表，或者发到广告充斥的微信公众号或微博「头条文章」里。这样发表的文字标题必须引人注目，否则读者会大大减少。而越引人注目的标题，往往是越差的选择。}
\BookParagraph{所以我自己的微信里没有开通朋友圈。}
\BookDate{2017 年 6 月，西安}
\BookArticle{征友启事}{征友启事}
\BookParagraph{每次看到「征友启事」这几个字，总是能联想到刘瑜的小说\BookLineBreak{}《那么，爱呢》开篇的那一段儿：}
\BookParagraph{————}
\BookParagraph{于是，他决定突出一下优势，把原来的那些废话删掉，改写道：I graduated from a first-tier university with a master degree in finance. I work in a Wall Street company right now. My career is very promising. I’m handsome, nice, humorous. I’m looking for a beautiful girl, outside and inside. 翻译成中文，就是：我毕业于一所牛逼的大学，读的是牛逼的专业，在一个牛逼公司工作，职业的前途很牛逼。我长得很牛逼，性格很牛逼，谈吐也很牛逼。我要找一个牛逼的女孩，她要从里到外地牛逼。}
\BookParagraph{————}
\BookParagraph{「征友启事」其实挺难写的。难就难在三言两语之间既要总结自己的经历，又要体现自己与众不同的一面，令人过目不忘。刘瑜笔下的自吹自擂似乎有点儿过于直接了当，不一定奏效；余秋雨笔下酒公张先生「少习西学，长而废弃，颠沛流荡，投靠无门」的自嘲，有凄切之美，但又难免妄自菲薄之嫌。含蓄婉约的文风过于虚无缥缈，无法达成共鸣；而直接写上「有住房有存款有海外亲属」，又人云亦云，无法免俗。}
\BookParagraph{相对而言，我更喜爱平实、普通、简单的文字。与其剑走偏锋，哗众取宠，不如给人一种诚恳、大度、值得信任的感觉：}
\BookParagraph{————}
\BookParagraph{我事业心不强，但是做事认真；学习成绩优秀，但是记忆力差，竞赛没得过奖；平时不够勤俭节约，但是也不太小资；唱歌棋牌球类无一擅长，但是经常运动；不痴迷电脑游戏，不了解飞机舰船，不关心世界大事，但是会带孩子，会洗衣服，会做饭。}
\BookParagraph{世界上有两种人：第一种是一起坐长途飞机都宁可各自玩手机的；第二种是可能多年不联系，但是一聊天儿就能一不小心聊到后半夜的。希望能遇到后一种人，和他们成为朋友。}
\BookDate{2017 年 7 月，多伦多}
\BookArticle{褚明宇}{褚明宇}
\BookParagraph{想听听各位老师如何评价褚明宇老师？（问题来自「向上移动」）}
\BookParagraph{二十一年前的一个晚上，我在美国UIUC校园旁边的一家叫\BookLineBreak{}「新燕京」的小饭馆吃晚饭。跟往常一样，这家中餐馆的快餐颇受欢迎，人声鼎沸，熙熙攘攘。正吃着，冷不丁旁边一个男的跟我说：}
\BookParagraph{「哟，这不李葆春吗？」}
\BookParagraph{褚明宇怎么就认出我的，其实我到现在也没完全明白。虽然我们小时候都住北大燕南园，但他比我大两个年级，路队不跟我一块儿排，我完全不记得还有这么一个人。}
\BookParagraph{褚明宇一个人住一套一室一厅的公寓，就在学校工学院校园旁边儿。而那会儿中国学生喜欢聚居在一个叫「Orchard Downs」的地方，如果几个人合租，房租会更便宜。他开的是一辆六缸发动机的二手美国车，速度表是液晶数字显示，中控甚至是一块触摸屏。在爱买丰田本田车的中国学生里，显得特立独行，与众不同。}
\BookParagraph{不过印象最深的，还是刚认识他时他有多能说。}
\BookParagraph{有个和我一起去 UIUC 的清华同学叫韩牧，北京人，清华计算机系同级里北京市高考录取分第一名，现在都管他这样的叫「学霸」。韩牧不但学习极其出色，作息还非常规律，从来都不晚睡。一次褚明宇邀请我和韩牧到他家聊天儿，三个北京人，竟然聊到早上五点。还记得褚明宇讲的最后一个故事是关于他在北大的女朋友，故事里的女生在我们想象的世界中，长得超凡脱俗，倾国倾城。他最后说：}
\BookParagraph{「她也许随时就会出现在大家面前。」}
\BookParagraph{一个月后，我们果然见到了她，果然明艳不可方物。}
\BookParagraph{除了能说，褚明宇的钢琴弹得实在太好了。我虽然不懂音乐，但无论他是钢琴独奏还是自弹自唱，我都能真实地感受到一种现场音乐会般的氛围，一点儿点儿地蔓延开来，无形之中紧紧抓住所有在场的人。}
\BookParagraph{褚明宇很聪明，而聪明的人其实有两种。前一种是聪明得让人觉得有一种泰山压顶提心吊胆的感觉；而后一种是聪明得让人觉得随和，觉得有很多逗事儿可以开聊，甚至觉得温暖。褚明宇属于后一种。跟他喝酒聊天，会觉得时间过得太快，会害怕筵席终将散场。}
\BookParagraph{但是褚明宇的文字风格，就像小瓶儿的北京红星二锅头，犀利得锋芒毕露，得意得招摇过市。看罢各种「大V」们的「头条文章」后再看他的文字，就好像看惯了山水风光的中国水彩，忽然看到浓墨重彩的西方油画，红则大红，蓝则亮蓝，层层叠叠地冲进视野。他的文字主题鲜明跳脱，从论点到论据，逻辑清晰自然。\BookLineBreak{}「对方辩友」无论是否同意他的观点，都应该承认他的文字确实\BookLineBreak{}「好看」。}
\BookParagraph{而好看的文字越来越罕见了。小时候喜欢王安忆，喜欢王朔，喜欢阿城。后来喜欢刘瑜，喜欢艾米，喜欢辛夷坞，喜欢慕容雪村。可惜的是，他们现在都不怎么写了，而褚明宇还在继续写着他的文字。我有时甚至感到有一丝庆幸，像是在杂草丛生的后院儿中发现还幸存着一株白色的玫瑰花，总觉得是一种错觉。}
\BookParagraph{乔布斯年轻时，除了 Larry Ellison 等少数几个人，其他的芸芸众生一概都瞧不起，管他们叫「bozos」，大概说的就是褚明宇笔下的「下等人」。很多人无法跟乔布斯共事，他的公司众叛亲离，连家具都被低价清空了。而他去世前，Walter Isaacson 问他为什么要写传记，他说：「I wanted my kids to know me. I wasn’t always there for them, and I wanted them to know why and to understand what I did.」}
\BookParagraph{说不定，越是冷清得浓墨重彩棱角分明的文字，越有一种普通人天然的亲和力，它在心里流淌着，会让我觉得温暖。}
\BookParagraph{就像褚明宇自弹自唱的歌，他凌晨五点的故事，还有他开过的那辆二手美国车。}
\BookDate{2017 年 9 月，多伦多}
\BookArticle{百家争鸣}{百家争鸣}
\BookParagraph{进群一段时间意识到自己逼格不够，还需修炼，可以退出吗？}
\BookParagraph{看虎妈那篇文章意识到自己这些年混的不好，原来原因有一部分是不会歧视他人。}
\BookParagraph{看完翟欣欣那篇意识到像我这样谈过几次恋爱的女青年，即使没占到过什么便宜，在别人眼里可能会是什么样子。反正过程没人了解。}
\BookParagraph{阮琦老师不错，但我觉得他对金星的态度确实不太好。}
\BookParagraph{或许我真的不适合这里。（问题来自「向上移动」）}
\BookParagraph{首先，我非常喜欢你的问题。一方面是因为提出这样的问题是需要一点儿胆量的，它多半儿会招致社区里有些人的冷言冷语，甚至人身攻击。绝大多数人大概觉得没有必要做这种「损人不利己」的事儿。另一方面是因为这个问题不仅仅限于一个社区或一个群，同时也可能适用于一个工作单位或一家公司，甚至一个家庭。}
\BookParagraph{这个问题的根本是，你觉得自己在一个地方很不适应，而原因多半是你和这个地方的其他人的理念不同，对事情的看法也不同。}
\BookParagraph{一个社区也许还可以随意退出，但如果是一家公司或一个家庭，大概想「退出」就要付出相对高昂的代价。}
\BookParagraph{你也许认为我这不过是传递「正能量」的「鸡汤」，可是我真的觉得，与其消极退出，不如尝试做点什么。说你想说的话，写你喜欢的文字，表达你鲜明的观点，大张旗鼓地展示你特有的风格。在持续而有建设性的讨论中，去一点一滴地影响和改变这个地方的人。这样的讨论有点儿像新闻联播里大国元首之间的高峰会晤：\BookLineBreak{}「宾主就共同关心的问题广泛地交换了意见，会谈在亲切友好的气氛中进行。」}
\BookParagraph{比如说翟欣欣，褚明宇的观点虽然有的人不一定同意，但是他至少将其观点明确地表达出来了。如果你觉得哪里说得和你的观点不一致，与其因此退出社区，不如把自己的想法写成文字，完整而充分地提出来，给社区里的朋友讨论。}
\BookParagraph{我想，即使是褚明宇自己，也会衷心地希望大家能在社区里踊跃表达不同的观点，形成一种「百花齐放、百家争鸣」的局面。也许，这也是几个月前他在分答开创这个社区的初衷。}
\BookParagraph{乔布斯去世前一年，苹果发布了iPad。有一个用户对这个产品封闭的生态圈非常反感，就给乔布斯写了一封邮件抗议。大概连他自己都没有想到，乔布斯凌晨两点竟然回复了他的邮件，两人在邮件中你来我往，针锋相对地展开了一场争论。最后，乔布斯说：}
\BookParagraph{——}
\BookParagraph{As for us, we’re just doing what we can to try and make (and preserve) the user experience we envision. You can disagree with us, but our motives are pure.}
\BookParagraph{By the way, what have you done that’s so great?  Do you create anything, or just criticize others’ work and belittle their motivations?}
\BookParagraph{——}
\BookParagraph{与其退出，不如自己动手，开创一个新的世界。}
\BookDate{2017 年 9 月，多伦多}
\BookArticle{下等人}{下等人}
\BookParagraph{李老师，我是一名大三的学生，经常觉得身边的人很「下\BookLineBreak{}等」，是我自己的问题还是环境的原因呢？怎样才能交到真心的朋友？（问题来自分答「在行一点」）}
\BookParagraph{正好今天收到一条微博私信。发信人不同意褚明宇最近的一篇微博，于是来问我：「关注您很久了，气质和褚明宇很不一样。不知道您怎么和这个上等人成朋友了？」}
\BookParagraph{这两个问题出发点颇为不同，但是言外之意大致都是：如何尽量「向上移动」，别跟「下等人」成为朋友？}
\BookParagraph{其实我们来想一想，我们知道的人里，哪些人是「上等人」呢？巴菲特算吗？马斯克算吗？布隆伯格算吗？斯皮尔伯格算吗？如果限于国内，那马云算吗？马化腾算吗？陈晓卿算吗？刘瑜算吗？}
\BookParagraph{想好以后我们再来假想一下，如果给你一个和他一起吃午饭的机会，你作何感想——欣喜若狂？还是激动万分？}
\BookParagraph{要是我有这样一个机会，只会觉得紧张。就跟要去谷歌面试一样。}
\BookParagraph{比如有一次我和一位教授吃饭，他一语双关，提到数学里著名的「李群」（Lie group），我一时半会儿没反应过来。从那儿以后每次见到这个教授，我就有点儿紧张。}
\BookParagraph{而真正的朋友，至少你和他吃饭不会觉得紧张。}
\BookParagraph{我有一个朋友是大学同学，因为他们家住得挺近的，有时会带孩子串个门儿。孩子一起玩儿的时候，我们也会一人一瓶啤酒，聊一会儿天。这位同学有两大爱好：分析台湾形势和研究中国军事问题。很可惜由于平时时间有限，我对台湾和军事无知得很，超级不熟，聊不起来。而我喜欢的主题，比如汽车，比如电影，他同样无话可说。}
\BookParagraph{不过后来我们发现一个知识交集：安卓操作系统。他的工作和安卓相关，恰好我也懂一点儿。从那儿以后我们每次见面，聊几句就聊到了安卓，竟然也言之有物，侃侃而谈。}
\BookParagraph{所以我觉得，无论是二十多的小姑娘还是六十多的老爷子，每个人其实都有一些自己喜欢的人和东西，就像一个「频谱」。两个人之间，只需要各自的频谱有一处很窄很窄的频段相同，就能产生共鸣。而我们像在用一个老式的收音机收听遥远的电台，只需要不停地扭动大大的旋钮，找到那个频段，就找到了「真正的朋友」。}
\BookParagraph{毕竟盖棺定论之时，无论你是「上等人」还是「下等人」，总会希望世界上还有人能够时时想起你，觉得你是他的一个真正的朋友。}
\BookDate{2018 年 4 月，多伦多}
\BookArticle{张戎}{张戎}
\BookQuote{张戎老师上微博吗？想关注这位牛逼的老师。（问题来自微博评论）}
\BookParagraph{张戎虽然不怎么玩儿微博，却有一个 \href{https://www.zhihu.com/people/rong\_zhang/}{知乎帐号}，里面有不少精品长文，主题涉猎颇广。他文风独到，很多文字看起来是随心所欲一蹴而就，但其实是精益求精仔细琢磨过的，很值得一读。}
\BookParagraph{他的知乎长文中我最为偏爱的一篇，回答的是这么一个问题：\href{https://www.zhihu.com/question/20894671/answer/376383314}{有哪些十分钟就能学会并终生受用的生活技能}？}
\BookParagraph{这个问题其实挺难回答的，反正我自己想了半天没想出什么比较像样的答案来。张戎的文章开门见山，直奔主题 —— 我直接引用几段吧：}
\BookQuote{当然要有 Alfa-Bravo-Charlie!}
\BookQuote{你在电话里跟老外拼读自己姓名的时候，有没有一种鸡同鸭讲的感觉？}
\BookQuote{扯着脖子喊话好几遍，还是 Z-C-V 不分，B-D-P 混淆。即使对方把字母组成单词来确认，如果信号不好还是分不清「 Bob」还是「Dog」，对吧？}
\BookParagraph{我第一遍看的时候，简直有一种眼前一亮的感觉。每次电话里对方问：「Can I have your name, please?」我就有点儿紧张。「B as in … “Bob”, N as in — “Nancy”?」}
\BookParagraph{再看下去，张戎就开始了他的「营销」：}
\BookQuote{你看一下 Alfa-Bravo-Charlie 这套标准单词，能不能花 10 分钟背下来，将来每次电话拼读英文姓名地址电邮都能省掉 2 分钟「B！B！B！不是 D！」的呐喊时间，5 次电话就回本，划算吗？}
\BookParagraph{这篇文章我看了好几遍，还真花了不少时间 —— 绝对不止十分钟 —— 把这套「国际无线电通讯拼读字母表」背了下来。不久之后果然派上了用场，电话中信心百倍地回答：「B as in Bravo, A as in Alfa, …, and N as in November.」}
\BookParagraph{一篇文章的结尾就像平衡木体操中的下法一样，极其重要。让我们来欣赏一下张戎这篇文章的「下法」：}
\BookQuote{大家知道英文中字母表这个词（Alphabet）实际上是希腊字母表头两个字母 Alpha 和 Beta 组成的。对于国际无线电通讯拼读字母表这个东西，与其取那个容易和国际音标（International Phonetic Alphabet）混淆的简称（NATO Phonetic Alphabet），还不如干脆叫「Alfabra」来得简洁性感。}
\BookDate{2019 年 11 月，多伦多}
\BookArticle{祝福}{祝福}
\BookQuote{李老师您好，想向您请教一个问题。给朋友，长辈，老师或领导发节假日的祝福短信，如何做到不流于形式显得真诚不虚假？感觉祝福信息已经沦为了一个礼节性的过场，双方都认为是个虚假的客套。同时也祝您父亲节快乐。（问题来自微博私信）}
\BookParagraph{哈哈，我最不擅长的就是发祝福短信了。为了更好地回答你的问题，我专门咨询了两位朋友：一位是国内高校的年轻老师，人又善良又漂亮；另一位是跟我差不多年纪的教授，朋友三教九流，多得满世界都是。}
\BookParagraph{咱们假设，你发节假日的祝福短信是希望让对方记住你。但是由于逢年过节发祝福短信的人太多，你的祝福往往很难显得「特立独行」。而这种特立独行跟别人不一样的风格，正符合我一直推崇的做一个「contrarian」的哲学。我三年多前写过一篇《\href{https://baochun.ca/post/\%E5\%A5\%BD\%E5\%A5\%BD\%E5\%AD\%A6\%E4\%B9\%A0\%E5\%A4\%A9\%E5\%A4\%A9\%E5\%90\%91\%E4\%B8\%8A/\%E9\%80\%86\%E6\%B0\%B4\%E8\%A1\%8C\%E8\%88\%9F/}{逆水行舟}》，说的就是这种哲学：}
\BookQuote{而我更钟爱的是一种逆水行舟式（contrarian）的风格：别人都买房我不一定买，别人都看武侠小说我不一定看，别人双十一都抢购我不一定抢，别人都玩儿朋友圈，我则自始至终「金盆洗手，淡出江湖」。}
\BookParagraph{那发祝福短信这种小事儿，怎么「特立独行」，让别人更好地记住你呢？}
\BookParagraph{我「锦囊」中的第一条建议，是「出其不意」。很多人每逢佳节，无论是农历新年还是端午中秋，都要发一轮祝福。这样的祝福短信自然不能群发，而需要给每一位领导或朋友「私人定制」，自然要花不少时间。但其实即使是跟相亲的对象也没必要每天都发信息问候，领导和朋友一年问候祝福一次其实足够了 —— 关键是这一次祝福哪天发。}
\BookParagraph{本着「出其不意」的原则，发祝福最好的时间选择，是在对方的生日那天。生日是一个挺特殊的日子：它既有些私密，一般只有家人和最要好的朋友才能想得起来；又有些公开，即使不太熟悉，你要是打听一位朋友的生日，对方多半儿会很大方地告诉你。而现在大家都很忙，即使是平时叱咤风云的领导，每年的生日其实也不会有几个人能想起来。所以，你如果在生日这天准时发了祝福短信，文字简洁明快、情真意切，一定比农历除夕发同样的短信「特立独行」的多。}
\BookParagraph{我的第二条建议，是「持之以恒」。我向美女教授咨询这个问题时她告诉我，她有一个现在已经毕了业的学生，找她的时候会给她发好长的信息，而没事的时候就音讯全无。我觉得这很常见，但正好是你体现自己「特立独行」的一个机会：如果你可以每年同一个日子都发一封祝福信，长此以往，十几年甚至几十年如一日，给对方的印象一定比大多数人要更加深刻。}
\BookParagraph{但是为什么要局限于发祝福短信呢？我有一位同系的同事，每年的新年前夕都会给我发一封邮件，邮件里是一封精心设计制作、图文并茂的信。每年的信只有一页，但是包括了全家人过去一年中经历过的大事和去过的地方，还有一张全家旅游时的合影。他的「年终大事记」近十年来没有断过，看着孩子们一年年慢慢长大，让人不禁感慨光阴若白驹之过隙，岁月已蹉跎。}
\BookParagraph{其实，如果是几位特别的朋友或者领导，按照我这位同事的思路，完全可以手写一张祝福的卡片，再附上几张刚洗出来的彩色照片，提前邮寄过去。「出其不意」得自然爆表，再能持之以恒，无论如何也不会成为「礼节性的过场」了。}
\BookParagraph{当然，如果「出其不意」和「持之以恒」太麻烦，写情真意切的祝福文案又太强人所难，我的另一位朋友有一个建议：}
\BookParagraph{发条「新年快乐」，然后转 6666 就行了。}
\BookParagraph{哈哈，如果你哪天真的成了「财富自由」的土豪，想让人记住你很简单。}
\BookParagraph{父亲节快乐。}
\BookDate{2021 年 6 月，香港}
\BookPart{为祖国健康工作五十年}{为\\祖\\国\\健\\康\\工\\作\\五\\十\\年}
\BookArticle{长　跑}{长跑}
\BookParagraph{从小到大就害怕长跑。小时候怕四百，中学怕八百测验，大学锻炼时，少不了一千五。后来才知道，其实心里一直最向往的是一种和长跑有关的感觉，可能是最后冲刺时的超脱，也可能是跑完了又歇了半个钟头后的轻松。长跑永远地和人生观世界观的最高境界联系在一起，象「为祖国健康工作五十年」，「走出教室，走出宿舍」，「现在是锻炼时间」。长跑象征着满地的风沙，食堂的广播，青春的冲动，还有代代相传的火种。}
\BookParagraph{很久以后才发现，长跑就象我的生活，想停也停不下来，因为停了就再也跑不动了。自己一个人跑永远是孤单的，总想要有一个伴儿，时不时可以拉个手，聊个天儿，跑得也轻快些。长跑象解数学题，象追女孩儿，「养兵千日，用兵一时」，追寻的是那体操表演漂亮的下法，和那越跑越快喜洋洋一切完美无缺的感觉。一直向往当那众星捧月的焦点，疾速冲线后挥舞着五星的红旗，热泪盈眶地听那半空中的国歌，还有那全体起立经久不息的掌声。这些，可能都从心中对长跑的幻觉而来。}
\BookParagraph{长跑就象喝酒，跑完过一阵子还想跑，但是总有一天发现，青春成了历史，跑不动也喝不动了。曲终人静的时候，总是希望年轻时的长跑留下了些什么，也许是空空的房子里泛黄的照片，火红的壁炉旁尘封的情书，还有遥远的记忆中相扶的一双手。倘徉在圣诞歌声中满眼英文的商场里，还有时端详着那双减价的耐克，想起那年倾囊而出而来的回力牌儿篮球鞋。总想着酒醒明月下之时，还能与当年同跑的女孩儿相拥而眠，生死相许。午夜梦回，又似乎听到了她那轻柔的歌声，看到墨绿色的客厅里高朋满座，觥筹交错，哈哈大笑。天明时合上眼睛，俨然看到了当年朋友婚礼上来宾在留言名册上意气风发的题词：将爱情进行到底。想起当年那同名的青春偶像剧，说不定，长跑，也是可以接力的。}
\BookDate{2005 年 5 月，多伦多}
\BookArticle{教　授}{教授}
\BookParagraph{俗话说三百六十行行行出状元，但我看不说别的，只说教授这行，就并不一定是件什么好事儿。咱们拿「教授」这两个字来说吧，本来还算挺乐观向上的一个好名字，MIT BBS 上偏偏给这个职业按照音译起了个外号叫「发考题」，看了就让人有一种直冒冷汗的感觉——天天给人发考题什么的，再跑到个没窗户的会议室里判个卷子，何乐之有？就连英文名字叫 professor 也无缘让我产生特别美好的联想，每次看到这个词儿就老想到在进入美国边检时浑身紧张实事求是的对答 ——人家问：「What is your occupation?」 我回答：「I am a professor at the University of Toronto.」人家追问：「What do you profess?」 我差点儿砰地蹦出一句：「I have no idea!」}
\BookParagraph{要是说谁谁干教授这个职业，总是会让人联想到他一定很忙。我收到的 email 里，但凡不是太熟的人希望让我帮他个忙的，一般都是以「I understand that you are very busy」开头，再以「Thank you very much in advance for your time!」收尾。按照我的经验，他这话说的那是相当地靠谱：这份工作虽然跟别的工作一样都不过是挣个死工资，却不是一般的忙。文山会海基本占齐了，再砸上各种行政职务，经费申请，项目结题，同事来访，讲课备课，发发考题什么的，基本就差吃饭在Chinatown，睡觉在办公室了。最恐怖的可以说还是各种各样的 deadline，基本三天一个小 deadline，五天一个大 deadline，deadline 之前还有若干个认真负责的 reminders，生怕你把 deadline们给忘了。要是拿不到经费做梦都想，拿到了又要汗流浃背地当什么「老板」，废寝忘食地当什么「大牛」，到了deadline还不能忘了给学生打工。人家古代唐朝时还兴个「因过竹院逢僧话， 偷得浮生半日闲」，至少还有那闲情逸致写两句诗什么的；要搁现在这当教授的，不要说没时间碰见个老僧就闲扯两句，就算是真的到了竹院，估计都辨认不清方向，只因心里还在嘀咕经费申请的截止日期是否就在这几天。}
\BookParagraph{教授忙是忙了，好像也没见忙出了什么成果。据谣传，天天对着带显像管的显示器工作会影响下一代生男孩儿，我当年听到这谣言时正好更换了昂贵的液晶屏，可孩子生了还是女儿之后还接着女儿。正好我更偏爱女儿嘛也就罢了，好不容易周末忙到半夜赶在什么夜里三点的 deadline 之前字斟句酌地跟生孩子似的写了几篇论文，让一些不知哪儿找来的 reviewer 们不费吹灰之力就给拒了，还好意思说「We regret to inform you」云云！要是这时自己还有心情再重新展卷重读那些被拒绝的精品，感觉就像帅才无缘伯乐、日志发到了人人一样，也就只能自斟自饮喝个闷酒，「成果」却是永不见天日了。}
\BookParagraph{教授之忙，那是忙的最高境界。我发现教授里头发掉得快的是大有人在，著作虽然还未曾等身，头发可能已经掉光。其实有没有头发也不是什么头等大事，没头发的教授可能比我这样暂时还有头发的更帅，更加让他的学生们拥戴。教授的忙上加忙，那只能说是不用「鞠躬尽瘁，死而后已」来形容就不够「给力」。当教授的，因为太忙得了癌症还是个晚期的例子俯首皆是。等到自己知道大势已去的时候，有像 Ralf Koetter 那样落叶归根的，有像 Randy Pausch 那样到处去培养粉丝开励志讲座的，更有像 Jennifer Hou 那样判作业带学生直到最后一天的。就像今年上海的高考作文题目一样，有人争相发表文章说「一切都不会过去」，证明他们都没有白忙；我倒更赞同「一切都会过去」的悲观论调，教授这个职业长江后浪推前浪，不用说掉了几根头发，没有了谁会议还不是照开地球还不是照转！依我看 deadline 这个英文词儿应该好好改一改，至少改成一个禁忌词汇，否是天天「死期」「死期」的英文挂在嘴边，我看离那真的「死期」也就不甚遥远了。要不然怎么前有教授当着当着就变节投奔了微软的，后有哈佛留不住撂了挑子去了谷歌的，人家聪明着呢，知道这教授的职业远不如什么微软谷歌脸书推特来得香甜。}
\BookParagraph{就这么个怎么想都不怎么地的职业，近年来倒成了时髦，几百个竞争一个位置的大有人在。它就像一辆早已满员的绿皮火车，虽然车破而又人满为患，无奈适逢春运高峰，一票难求。每逢开会总会有学生像英语角碰到老外似的主动跟我聊天，聊到痛快淋漓之际，不外乎又要感慨一下教授这行工作多难找，谁谁谁多少篇论文都没着落等等。 让我受宠若惊的同时又感慨僧多粥少，不知他们将来能不能挤上那班春运的火车，去他们梦寐以求的地方。}
\BookParagraph{我总在想，这些论文青年勇于投奔教授这个职业，图的是什么呢？是不是以为当教授的就是一年好几个月到处以开会为名义旅游？是不是不知道当了教授就无缘标新立异、注册了人人都没人加你为好友？是不是不知道当了所谓老板没人管还不如当个小兵儿有人惦记？是不是不知道文章发表了是中了彩票儿，拒绝了才是过日子种菜？是不是不知道等到摇滚的青春已尽燃的时候，一切都已过去，却未必能盼来那全场起立经久不息的掌声？说不定，魔力归根结底还是在「教授」这两个神奇的字中间，让大家前仆后继，争发考题。}
\BookDate{2011 年 4 月，多伦多}
\BookArticle{讲　课}{讲课}
\BookParagraph{这学年快过完了才发现，现在终于可以举重若轻地宣布，我已经讲了「十几年」的课了。讲课是个难事儿，难就难在无论怎么讲都很不靠谱。你要是高屋建瓴吧，他说你没有细节的渲染；你要端着放大镜注重细节吧，他说你讲得太深，听不明白；你深入浅出花无数时间精心准备的幻灯片，他说你讲得太快；你自带一盒粉笔在黑板上举一反三慢下来呢，他说你考试的范围没讲到。语言也要细细斟酌：你活泼吧，有没人听明白你笑话的风险；你严肃吧，他听得腻了下次压根儿就不来听你的了。要是在国内讲就更难：用中文讲专业词汇量不够，用英文讲又有在座的听不太明白，又不好意思提问，因为他英文口语可能不够自信。所以说，讲课是一门艺术，只有极少数人才能讲得游刃有余的艺术。}
\BookParagraph{清华上了五年，但上过的课还不如当时的民办新东方印象深刻。那会儿的新东方主讲就是俞敏洪，现在自然是微博粉丝百万，十七年前只不过是个民办英语学校的主打贫嘴而已。但是俞敏洪贫嘴和新东方其他人就是不一样。俞敏洪上的verbal，我认识的无论是清华的什么大牛还是北大的什么精英，全都到中关村一个小破阶梯教室听他贫，无一例外。主要是因为他特别能白活，没什么动静的英文词儿都能上个段子，要是稍稍能引起点联想的，马上上纲上线直接谈交友谈哲学谈社会，当然最多的还是聊出国。那会儿出国相当难，清华北大还没什么国际影响，申请出国要放弃本校免试读研，算个风险很大接近买彩票的工程实践。有了俞敏洪的课，很多人不但学了几个新词儿，更重要的是把本来试试看的心理演变成了全面出国誓死不回头的热情，以致很快俞敏洪的新东方免费课有上千人参加。俞敏洪的课，没有板书，没有教材，没有幻灯，没有提纲，没有作业，也不用记笔记，在场的上千人竟然全听傻了。唯一的期末考试就是GRE，考好了就有可能中彩票去美国。说白了，他讲的不是英文，是他无法掩饰的人格力量，是他声音洪亮口若悬河张牙舞爪的表象下蕴含的激情。}
\BookParagraph{俞敏洪讲东西，行云流水，没有一句听不懂的。这一点，特像果粉的偶像乔布斯。唯一的区别，是乔布斯的境界又高了一层。他的演讲全是精心准备、字斟句酌过的，大到时间的安排，小到每一个笑话，都是事先安排好的。但他的演讲，妙就妙在让人听了觉得他是在临场发挥，在谆谆善诱地跟你拉家常。他用的语言一概简单明了，似乎没有任何修辞，没有任何做作，没有任何掩饰。他凭着如此治大国若烹小鲜的天才，让在场和不在场的粉丝年复一年地物我两忘只盼时间停滞不前。这种一边讲一边和听讲的粉丝达成一种默契和互动，继而左右他们未来的生活方式和轨迹的，实在难到了极点，可俞敏洪和乔布斯都做到了。}
\BookParagraph{下礼拜又要回国讲课了。每次跟人家说，我一上讲台就紧张，人家都说，啊？你还紧张呀。其实我做梦都是讲课讲错了的噩梦。但要是真等清醒了开始上课的时候，总还是盼着第二堂课的人数不要减少太多，盼着有朝一日能像俞敏洪乔布斯一样，「遥想公瑾当年，小乔初嫁了，雄姿英发。羽扇纶巾，谈笑间，樯橹灰飞烟\BookLineBreak{}灭。」等有那么一天，俞敏洪乔布斯都不再上讲台的时候，我想我会怀念他们，还有他们讲过的课。}
\BookDate{2011 年 8 月，多伦多}
\BookArticle{皮包公司}{皮包公司}
\BookParagraph{前几天我跟一个申请读硕士研究生的学生聊天，问他将来毕了业有什么打算。他不假思索：「读博士，读完了争取当个教授。」就好像当个教授是一切优化问题的最优解。}
\BookParagraph{我倒是觉得，当个大学教授和当个小学老师差不多，只是一个普通的工作，并没有那么令人悠然神往。力荐教授这个职业的，大多数自己都是教授。借用电影《肖申克的救赎》里的一个词儿，这叫「institutionalized」。}
\BookParagraph{当个教授确实有机会到处开会，可惜去开会其实行色匆匆，疲惫不已，大多数时间是在工作，而不是在旅游。}
\BookParagraph{当个教授确实有一笔稳定的收入，可惜工作压力太大，时间太长，基本没有周末，平均每小时税后的薪金，大概跟在星巴克卖咖啡差不了太多。}
\BookParagraph{当个教授确实有做学术研究的自由，可惜研究的成果和水平都是别人说了才算。和其他工作差不多，少数人在金字塔尖儿，占有大多数的「声望」和资源。}
\BookParagraph{当个教授相当于开了一家皮包公司：品牌推广、市场营销、产品创收、人力资源，都是一个人管。从经费申请到学生推荐信，从立项到结题，从本科课程到组织会议，没有一件事儿是省油的灯。}
\BookParagraph{我做这个工作，或许是因为我有自知之明：没有足够的胆量去「海归」创业，没有足够的能力去谷歌编程，也没有足够的人气去当「大V」网红。}
\BookParagraph{或许只是为了将来女儿可以跟朋友介绍：「我父亲当年是个大学教授。」}
\BookDate{2018 年 3 月，多伦多}
\BookArticle{专　注}{专注}
\BookParagraph{前一阵子收到一封从学院一个办公室发来的电子邮件，要求所有教师把期末考试的草稿上传到一个学院的网站，以便专人勘误。}
\BookParagraph{一位老教授告诉我，根据他一辈子的经验，这种邮件在收到第五次以前不用处理或回复。「发送邮件的人如果真的想要联系你，自然有办法找到你。」}
\BookParagraph{我想起了著名计算机科学家 Donald Knuth 拒绝使用电子邮件的故事。想跟他联系，只能写信。}
\BookParagraph{果然，我没有上传期末考试，也没有人再来催促我，考试照常如期进行。}
\BookParagraph{其实邮件还算好的，还有不回复的自由。华尔街日报报道，越来越多的公司为了提高员工工作效率，开始启用 Slack 这类的实时通讯工具。从 CEO 到接电话的前台，每人开个帐号，上传自己的头像，在「Slack channel」里实时讨论问题，踊跃发言。}
\BookParagraph{天啊，这不就是没有红包可以抢的微信群嘛。}
\BookParagraph{近来有人撰文希望把「智能手机」改得「傻」一点儿：把iPhone 界面变成黑白的，删掉所有像 Facebook 这样的社交媒体，并关掉绝大多数推送消息。和老教授的经验相比，这些建议异曲同工，殊途同归。}
\BookParagraph{我认为，真正要提高效率，唯一的办法是专注。晚回或不回邮件，不看微信群，不看朋友圈，才有可能去专注于最重要的工作。但凡一时半会儿想不起来的事儿，都不是重要的工作，没有时间去完成。}
\BookParagraph{想提高工作效率，著名投资人 Marc Andreessen 建议：「Don’t keep a schedule.」}
\BookDate{2018 年 4 月，多伦多}
\BookArticle{大学老师}{大学老师}
\BookParagraph{应该以怎样的心态做一个大学老师呢？（问题来自「向上移动」）}
\BookParagraph{你这个问题的范围有点儿宽，很难一一详细作答。我要不然给你讲两个故事吧。}
\BookParagraph{我在多伦多大学有一位同事，我们一起教一门本科生一年级课程。这位同事是前系主任、IEEE Fellow、ACM Fellow，著作等身，曾获多伦多大学教学最高荣誉——校长教学奖。他曾经成功创业，住在多伦多市区黄金地段的一个大宅里，差不多相当于北京二环内的四合院儿吧。无论如何，这样的资历搁国内都可以称作是\BookLineBreak{}「男神」、「大牛」或「人生赢家」吧，应该满世界去做「人间指南」的报告。}
\BookParagraph{再加上国外大学实行的是终身教授制，教学质量高低，基本没什么区别。}
\BookParagraph{可是他在教学方面所达到的境界，如果我没亲身经历过，我大概不会相信。一门课四十多课时，他的讲义是全程手写的。写的时候分蓝色红色两种墨水，红色专门用于眉批、修改、脚注和强调。他为了激发学生的兴趣，专门请来知名公司的校友来给学生做报告。他每天晚上逐条回复论坛上大量学生的问题，考卷的草稿多次修改，力求完美。}
\BookParagraph{这门课，虽然我每年的学生评估成绩都相当高，但没有一次能够超过他的——因为他的教学评估分数，已经接近无懈可击。}
\BookParagraph{小时候在北京大学燕南园听姥姥讲过一个故事。故事说她有一次在北大路过外公上课的讲堂，听到外公在给学生上课。讲课的声音洪亮、饱满、富有激情。她有点儿好奇，就停下来看了一眼。结果发现讲堂里一共只有两个学生。}
\BookParagraph{套用现在时髦的词儿，有点儿当教师的「情怀」，其他的「just let it go」。}
\BookDate{2017 年 6 月，西安}
\BookArticle{人　设}{人设}
\BookParagraph{作为一个情感专家，如果有一天厌倦了这个人设怎么办？（问题来自「向上移动」）}
\BookParagraph{首先，你能提出这个问题，说明你在随时审视自己（be reflective），这种自省能力，需要极高的才华和修养，一般人不容易做到，值得「手动点赞」。}
\BookParagraph{「人设」这个词儿用在名人身上，我认为可以理解为一种「品牌」。比如褚明宇常常提到「人分上等人和下等人」的理论，就具备一种品牌效应，其相应的产品力有助于他在市场竞争中异军突起。和普通品牌不一样的是，你作为「情感专家」，品牌价值中的一大部分直接和你的名字联系在一起，就像 Coco Chanel 的 Chanel 或者乔布斯的苹果，很难剥离开来。}
\BookParagraph{但是如果哪一天自己都不喜欢自己的品牌了，怎么办？我认为无非「渐变」和「跳变」两种策略。}
\BookParagraph{选择「渐变」，就是选择有意地将产品逐步转型，或者慢慢淡出大众的视野。李安导演的电影，从《饮食男女》到《色 · 戒》到《Life of Pi》，是成功的转型；而张扬导演的电影，则从当年《爱情麻辣烫》、《洗澡》的成功，到不为路人所知，是低调淡出了大家的视野。}
\BookParagraph{我开始在 UIUC 读书那年进来一个名叫 Vaduvur Bharghavan的教授。他在 UIUC 任教期间培养出来的学生成绩斐然，后来都去美国名校当了教授，他也毫无悬念地评上了终身教授。传说他升职后突然悄悄地离开了学校去了加州，没有临别赠言，甚至没跟任何人打招呼。这就是我说的「跳变」。}
\BookParagraph{「求变」就像很多人的「中年危机」，无论是买辆跑车还是积极创业还是热衷组织校友会，其实本质原因不外乎是觉得同一件事儿做的时间长了，想重整旗鼓，卷土重来，实现少年时代的梦想，活出别样的精彩。可是，为什么一定要「变」呢？说不定一辈子当一名小有名气的情感专家，一位不知名的教授，一个「有趣的人」，甚至一名空乘或者油漆工，也自有一种稳操胜券的幸福。}
\BookParagraph{以不变而应万变，以无招而胜有招。粗茶淡饭，说不定强于山珍海味。品牌也好，事业也罢，「maintaining the status quo」也许是最聪明的选择。}
\BookDate{2017 年 7 月，西安}
\BookArticle{搞研究}{搞研究}
\BookParagraph{在学术界还是在工业界搞研究？（问题来自「向上移动」）}
\BookParagraph{众所周知，投资有两种主要方式：投资股票（stocks）和投资债券或定期存款（fixed income）。}
\BookParagraph{投资股票的风险相对比较高，但是长期来看回报也高，「在大风大浪中成长」，适合承受风险能力比较强的人。而投资债券则正相反，风险很低，但是相应的回报也低，适合喜欢风平浪静的人。当然通常更好的策略，是年轻的时候尽量多投资股票，进而逐渐加入债券投资，到退休前把风险降到最低。}
\BookParagraph{工业界有点儿像投资股票，而去学术界则更像投资债券。}
\BookParagraph{同样的人才去工业界，收入一般比去学术界高很多，有时候甚至是两倍以上。那为什么每一个学术界的空缺，都通常会有大量具有博士学位的申请者呢？}
\BookParagraph{学术界大多有终身教授制（tenure）。一位教授晋升以后，大学就放弃了解雇他的权利，相当于以前国营企业的「铁饭碗」。去工业界则因为各种原因被解雇的风险较高，尤其是随着年龄增长，竞争力逐年衰退，有可能失业时正是孩子上大学最需要收入的时候。}
\BookParagraph{很多人会说，那些失业的朋友，很快都会找到新的工作，说不定工资比以前还高，没有任何损失啊。这是因为在经济发展快的社会或时代，工业界工作的失业风险自然低一些。但是我们无法准确预测未来，也许到我们自己最需要工作收入来源的时候，经济发展减缓，甚至出现经济危机，这也是风险的一种表现形式。}
\BookParagraph{所以，工业界回报高就必须承担一定的风险；而学术界「铁饭碗」的工资必然会低不少。去工业界还是学术界，要考虑的一个重要因素就是更偏爱「大风大浪」，还是更钟情「风平浪静」。}
\BookParagraph{其实，有关收入的横向比较还没有这么简单。我们不但要考虑退休之前的工资，还要考虑退休后的退休金收入。假如30岁开始工作，65岁退休，活到100岁，那么退休以后的时间和退休前大致相当，退休金收入对于保证退休后的生活水平至关重要。}
\BookParagraph{一些大学和很多政府部门相似，实施一种叫「defined benefit」的退休计划，退休以后每月提供一笔数额丰厚的退休金，一直发到去世；此后再继续给配偶发退休金，直至去世。如果工业界没有类似的计划，保守起见，则需要将退休前收入中的至少一半存起来留给退休后来花。大多数人无法做到这一点，也许只能降低退休后的生活水平。}
\BookParagraph{当然，大学教授如果教学科研之外还有余力，可以同时去公司咨询或创业，享受「大风大浪里成长」的高回报，又没有相应的高风险。就像一个「portfolio」里既投资股票又投资债券，有点儿\BookLineBreak{}「两全其美」（have your cake and eat it）的感觉。很多人还提及诸如「我就喜欢教书育人」或「我就热爱那种没人管我的自由」等\BookLineBreak{}「情怀」方面的因素。可惜这些因素使我们本来就不简单的「优化模型」更加复杂，就不在这篇短文里展开论述了。}
\BookDate{2017 年 8 月，上海}
\BookArticle{减　肥}{减肥}
\BookParagraph{如何不那么痛苦的减肥？（问题来自「向上移动」）}
\BookParagraph{我所了解的有关健康和运动的常识，基本上都是从书里看来的。}
\BookParagraph{——}
\BookParagraph{十年前我在一次多伦多至北京的航班上，一口气读完了一本书，可以说改变了我后来的运动习惯。书名很长，叫做「The Fitness Factor: Every Woman’s Key to a Lifetime of Health and Well-Being」，作者Lisa Callahan是一位医生，专业是运动受伤恢复。其实这个书名起得有点儿误导，书里介绍的常识不仅适用于女性。}
\BookParagraph{书里介绍了大量有关减肥、饮食、健康和运动的常识，但是其中最重要的一条，其实可以用一句话概括：无论男女，无论在什么年龄段，保持健康的最有效的手段是力量训练（weight training）。作为一种无氧运动，力量训练不仅可以增加肌肉力量和比例，肌肉恢复的过程还可以提高新陈代谢的速度，也就是常说的「睡觉时都在减肥」。另外，它可以增强骨骼的密度，大大延缓中老年人骨质疏松的过程。}
\BookParagraph{正是因为这些显而易见的好处，我和国内的朋友聊天时，经常推荐他们开始力量训练。可惜的是，无论是年轻学生还是中年领导，大多数人都不以为然，认为女生没必要练力量，或者认为力量训练只是为了肌肉好看，所谓「六块腹肌」（six-pack）。但其实绝大多数人力量训练的频率和重量极其有限，很难练到「好看」的水平。力量训练其实是减肥和保持健康最有效的捷径，是运动必不可少的一部分，男女老少皆宜，甚至只需要家里有哑铃就行。}
\BookParagraph{就像那本书里说的：力量训练就像一种万灵药，可以延缓衰老，预防重大疾病，提高身体的综合素质，而且任何年龄都可以开始。关键是你想不想得到这些好处。}
\BookParagraph{——}
\BookParagraph{四年前一个偶然的机会，我在美国一个机场的书店里买了另外一本书，书名叫「The 8 Hour Diet」，作者是 Men’s Health 的主编。美国的 Men’s Health 和国内的不同，以健康和运动为主题，极少提及时尚，是一本相当专业的杂志。}
\BookParagraph{这本书的中心思想叫「Intermittent fasting」。每天无论吃多少都可以，但是要保证所有的进食都在8小时内完成，比如十点到六点之间，其他时间只喝茶或水。另外重视保证8种健康食品的摄入，同时坚持每天至少8分钟的有氧和力量训练。这和很多人所说的「少吃多餐」「一定要吃早餐」等常识是截然相反的。}
\BookParagraph{「8-Hour Diet」的好处很容易理解。让身体每天都有一段时间不能得到满足，这样可以「居安思危」，激励细胞更加活跃。书里提到，「8-Hour Diet」若能与健康食品和运动结合起来，不但对减肥很有好处，还可以有效地预防癌症、心血管疾病和糖尿病这三大杀手级疾病。}
\BookParagraph{遗憾的是，我把这本书的理论给很多朋友介绍过，但极少有人认同这些想法。有时候不禁会感慨，一个人的想法即使是误区，一旦形成，确实很难改变。}
\BookParagraph{·  ·  ·}
\BookParagraph{希望这篇短文，可以像十年前那本书改变我一样，改变微博里几位朋友的想法。}
\BookDate{2017 年 8 月，上海}
\BookArticle{完美主义}{完美主义}
\BookParagraph{李老师，你好，突然想问你一个问题，有点唐突。看您的微博以及您写的文章，感觉您是一个追求完美的人。我是一个大学生。我本身也是一个有些完美主义的人，这给我造成了极大的困扰，导致了我严重的强迫症，已经好几年了，现在虽然不会感觉那么痛苦的焦虑感了，但依然是被强迫症困扰，还是会有焦虑感，我想问下李老师会有这种焦虑感觉吗，又有什么方法可以减轻吗，谢谢老师。（问题来自微博私信）}
\BookParagraph{没错儿，我当年刚刚到多伦多大学开始工作的时候，的确是一个「完美主义者」。工作中，每一堂课都会认真地准备；每一篇论文都写到截止日期的那个凌晨，经常睡在办公室；回复每一个邮件都会非常细心地重新读一遍再发出。生活中也一样，新买的房子外墙油漆颜色我不喜欢，就亲自把房子又重刷了一遍。}
\BookParagraph{小时候学「青少年修养」，老师鼓励我们要「有理想」、「有追求」、「积极进取」，而不要「做一天和尚撞一天钟」。我想，我们每个人大概或多或少都有点儿追求完美的情怀吧。}
\BookParagraph{时间一年一年过去，我逐渐发现无论是工作还是生活里的事情，已经多到了忙不过来的程度。邮件需要回复的每天好几百，很难再逐条检查错误；孩子大了需要接送，很难在办公室通宵达旦地工作。可是作为一个「完美主义者」，事情做不完，就会经常处于一种极度焦虑的状态。越是时间紧张，陷得就越深。}
\BookParagraph{直到有一年，我开始失眠了。}
\BookParagraph{我听说过一种理论，说一个人百分之九十七都在被下意识控制着，只有百分之三是主动意识。我自行从这里衍生出来的理论是，追求完美而不得，进而陷入的那种焦虑状态，实际上直接融入了下意识。失眠也就在所难免。}
\BookParagraph{因为长期的失眠，一次讲课我讲到一半儿讲不下去了，竟然休息了一分钟。}
\BookParagraph{系领导跟我聊天，问我为什么会焦虑。我给他举了一个例子。孩子学校放假白天在爷爷奶奶家，我怕她们会感到无聊。领导说：「Tough luck.  They will have a tough life anyway.」}
\BookParagraph{这句话我至今记得。的确，很多事儿不必这么放在心上。没能给孩子一个完美的童年，她们一样会长大；邮件忘了回复，人家自然会再找上门来；办公室的书桌乱七八糟，重要的文件早晚能找到。}
\BookParagraph{说不定，与其去「追求完美」而精益求精，还不如「做一天和尚撞一天钟」的悠然自得。}
\BookParagraph{我后来想明白了这一点，开始用百分之三的主动意识去跟多年来追求「完美主义」的下意识作斗争。我时不时提醒自己：车到山前必有路，船到桥头自然直。}
\BookParagraph{至少，跟失眠比起来，其他的一切都不重要。}
\BookDate{2017 年 8 月，多伦多}
\BookArticle{码　农}{码农}
\BookParagraph{想听褚老师和李老师谈谈「码农」这个职业。（问题来自「向上移动」）}
\BookParagraph{我很少的几个偶像之一， Chris Lattner，就是一个标准的「码农」。}
\BookParagraph{Chris Lattner是我和褚老师在 UIUC 的校友。他在读硕士学位期间就设计和实现了一种叫「LLVM」的系统，将传统编译器（compiler）模块化，大大提高各个模块的可重用性（reusability）, 开源以来彻底颠覆了几十年来的编译系统设计。他博士毕业后加盟苹果，将一种叫「Automatic Reference Counting」的内存管理机制引入了 Objective C 和 iOS 平台，替代了当年最流行的「Garbage Collection」机制，大大提高了内存管理的效率，保证了iOS apps 在内存有限的情况下一样可以流畅运行。}
\BookParagraph{2010年以来，他独自一人开始设计一种名为「Swift」的全新编程语言，汲取了二十多年来现代程序设计语言几乎所有的宝贵经验，真正实现了既能快速开发、又能高速运行的「Holy Grail」，现在已经成为iOS平台上的主流程序设计语言。2013年他获得了UIUC 的特别校友成就奖（Distinguished Alumni Achievement Award），可以说他的工作改变了整个软件产业。}
\BookParagraph{如果说 Chris Lattner 是一个标准的「码农」，那么如果能再年轻一回，我的理想就是当一个像他一样的「码农」。}
\BookParagraph{可惜的是，这几年「码农」这个词儿无论是用来叙事还是用来自嘲，都带有一种天生的贬义。这种贬义甚至影响了大量学生，觉得再怎么好好学习计算机，将来不也就是当一个「码农」或「程序猿」而已。}
\BookParagraph{但是实际上「码农」这个职业虽然显得像个「农民工」，却先天具备两个极为有利的优势。第一，不仅仅是计算机行业，很多其他行业也需要开发和维护越来越复杂的软件系统。第二，一个水平高的「码农」就像一个天生丽质的模特，相当罕见。乔布斯曾经说过：「The difference between the average software developer and the best is 50:1.  Maybe even 100:1.」}
\BookParagraph{这是因为当一个好的「码农」，不仅需要从离散数学开始，掌握一种叫做「计算思维」的能力；还需要从计算机体系结构、操作系统、算法分析入手，掌握几十年来计算机工程的无数成功和失败的设计。就好像学一门乐器，需要多年静下心来苦心孤诣地积累，才能成为一代名家，才能自然而然地具备一种「一夫当关，万夫莫开」的气场。}
\BookParagraph{当然至为登峰造极的功夫，是将程序设计从工程提升到艺术的层面，发现其中的魅力所在。这一点，和当年梁思成研究建筑、杨振宁研究物理相比，颇有异曲同工之妙。Chris Lattner 主导的 Swift语言中每一个版本的演进，都是在以一种近乎工匠般严格的精神，仔细地衡量每一行代码的语言风格如何更加「漂亮」。比如程序界面里需要定义一个表示蓝色的变量，早先写成：}
\BookParagraph{let blue = UIColor.blueColor()}
\BookParagraph{后来演进为：}
\BookParagraph{let blue = UIColor.blue}
\BookParagraph{如果一个所谓的「码农」能够返朴归真，达到如此「以无招而胜有招」的境界，我会像刚刚看完伍迪艾伦的电影一样，产生一种久久挥之不去的感动。}
\BookDate{2017 年 8 月，多伦多}
\BookArticle{工程师}{工程师}
\BookParagraph{如何成为一个合格的工程师？（问题来自微博私信）}
\BookParagraph{有一位工程师、一位牧师和一位医生在高尔夫球场打球。他们发现前方有几个人占着场地，打得又慢又不专业，就问球场的服务员。服务员说：这几个人是消防员，曾经在球场的一次火灾中，因为帮忙救火而失明了。}
\BookParagraph{牧师动容地说：「我今晚一定会替他们祷告，祈祷他们重见光明。」}
\BookParagraph{医生惋惜地说：「我联系一下眼科的同事，看看能不能为他们做点儿什么。」}
\BookParagraph{工程师不解地说：「那他们为什么不能晚上打球啊？」}
\BookParagraph{——}
\BookParagraph{工程师的世界和科学家艺术家不同，不但不是非黑即白，而且没有完美可言。一个工程师的工作，是如何在一些约束条件下完成一个系统的设计和实现。这样的系统和真实的世界一样，永远是不完美的。就像需要在一年内设计一辆造价不超过十万元的汽车，一位合格的工程师，需要在各种取舍之间权衡利弊得失（trade-offs and compromises），最终得到一种最好的解决方案。}
\BookParagraph{在我们这个高尔夫球场的故事里，牧师和医生都在表示同情，只有工程师在想如何解决这个问题。}
\BookParagraph{而约束条件中最常见的，往往是资金和时间成本。一家航空公司的电话订票系统，如果只有一个电话服务中心，很有可能在最繁忙时无法应付需求。但是如果在全国建立多个服务中心，虽然可以提高系统的可扩展性（scalability），但由于各个中心之间的数据需要实时同步，提高了系统设计的复杂度，增加了系统实现的资金和时间成本。我们常说的「KISS （Keep It Simple Stupid）」，就是希望通过简化设计而减少资金和时间的投入。}
\BookParagraph{其实，日常生活和工作里的很多事情，与其一味追求尽善尽美的「工匠」精神，不如提倡一点儿工程师式的「实用主义」，在有限的时间里做完。做饭与其追求做得多好吃，不如做得快；写论文与其追求「沉鱼落雁之容」，不如「在规定时间里」把问题和解法写清楚。一个合格的工程师需要大量工程实践和训练，而这种「多快好省」的生活和工作，同样需要一种「工程师」式的理念和智慧。}
\BookParagraph{一个玻璃杯里有半杯水。悲观主义者想：「玻璃杯已经空了一半。」乐观主义者想：「玻璃杯还有一半水。」}
\BookParagraph{而工程师会想：「玻璃杯的设计有失误，容积可以减半。」}
\BookParagraph{我喜欢这种「工程师」式的哲学。}
\BookDate{2017 年 9 月，多伦多}
\BookArticle{Gambit}{Gambit}
\BookParagraph{李老师好！您曾经说过，找工作就跟找对象一样：「没有的时候有无穷多种选择，非常灵活机动；有了就成了唯一选择，不满意的话就只能听之任之，或者付出昂贵的“跳槽”成本。」我本科快毕业了，想请教您一下，您是怎么看待人生第一份工作的选择的呢？（问题来自微博私信）}
\BookParagraph{你听说过一个叫「gambit」的英文词儿吗？}
\BookParagraph{这个词儿最初是从一种叫「the Queen’s Gambit」的国际象棋开局来的。这种开局有点儿匪夷所思：第二步棋就放弃一个卒，从而让皇后有充分发挥其优势的机会。这种古老的开局是典型的以退为进，韬光养晦，采取的是「先天下大乱而后天下大治」的方针。然而它却极其有效，曾经在国际象棋大师赛中风行一时。}
\BookParagraph{找第一个工作就像下国际象棋。有时候吃一点儿亏、冒一点儿风险，往往如同「the Queen’s Gambit」一样，是一种行之有效的开局。}
\BookParagraph{我有一个舅舅叫王缉思，是我从小到大的偶像。他在国际关系尤其是中美关系领域是世界知名的学者，现任北京大学国际战略研究院院长，曾被美《外交政策》杂志评为世界风云人物一百人之一。去年我回北大时他曾给了我一张名片，地址是「北京大学北阁」。我从小在北大长大，这个地址在我印象里可以和「宾夕法尼亚大道1600号」媲美。不过我最佩服的还是他一口流利的英文，略带英国口音，以至于我长大以后每次听到英国口音就会想到他。}
\BookParagraph{而他的第一个工作，是在内蒙古放羊。而且一去就住了十年。据说，他那会儿为了学习英文，只能用短波收音机收听美国之音。}
\BookParagraph{也许，如今第一个工作再不尽如意，也不会比他的经历更糟糕了吧。但是我想，也许他后来所有的成就，都离不开那个时代给他带来的阅历；也许他一口标准而流利的英文，也离不开那个短波收音机。}
\BookParagraph{所以，如同那句「天将降大任于斯人也」的老话里说的，第一个工作不一定非要踌躇满志，锋芒毕露。放弃一些优越的条件，冒一点儿有意义的风险，静下心来，从长计议，踏踏实实地一点一滴提高自己真实的水平，提高自己在市场里的价值。说不定将来可以充分发挥「皇后」的优势所在，迅速崭露头角。}
\BookParagraph{最近有一个新闻：谷歌的 AlphaZero 同时开动了五千多个TPU，不但靠「左右互搏」的战法四小时学会国际象棋，而且棋风有人的「灵性」，已臻「独孤求败」的境界。}
\BookParagraph{而 AlphaZero 最喜欢的两种开局之一，就是古老而凌厉的「the Queen’s Gambit」。}
\BookDate{2018 年 1 月，多伦多}
\BookArticle{找工作}{找工作}
\BookParagraph{李老师您好，经常默默地看您的微博，包括对于青年人问题的解答，觉得您的文风特别的实在和中肯。所以我也想请教您一个问题，不知道能不能有幸得到解答。}
\BookParagraph{我是学内燃机的工科生，得到了一个车企台架工程师和一个知名运营商企业客户经理的offer。我倾向于前者，因为我觉得在车企好歹是一门技术，不管将来怎么样都能找到工作，尤其是年纪大了之后有技术更稳。但是就职于运营商的一个表哥却强烈建议我去他这家公司，说待遇好而且也锻炼人。但我想做这种销售类的工作谁都可以干，尤其是年龄大了以后出来能好找工作吗？}
\BookParagraph{我想这也是一个中国相当部分工科学生面临的问题。当毕业季找到的工作在性质上有巨大的差别时，到底是坚持一技在身，还是转去做一些销售服务类或者其他跨行业的工作？（问题来自微博私信）}
\BookParagraph{如果你去问李开复或者去看TED的演讲，他们会告诉你两个字：「激情」（passion）。这种「follow your heart」听起来不错，可是它的假设里有一个致命的缺陷：很多人对什么工作都没有什么热情。}
\BookParagraph{这就好像你去相亲，有两个对象。一个长得好看但有可能跟别人跑了，另一个有学识有文化什么都好就是相貌一般。可哪一个你都不是「一日不见如隔三秋」，就想赶紧把这婚结了。你该选哪一个？}
\BookParagraph{工作和谈恋爱一样，不是每个人的运气都好到曾经爱一个人爱到茶不思饭不想，可以写进电影剧本。}
\BookParagraph{那如果对什么都没什么热情，工作这么重要一件事儿，应该如何选择？}
\BookParagraph{首先，千万别全盘接受别人的建议。表哥也好，李老师也好，给出的建议都是基于他自己的工作背景和偏爱，往往不一定中肯。你可以仔细研究他这篇「议论文」的论据和论证方法，但不要直接接受他的论点。}
\BookParagraph{其次，对工作没什么热情，正好可以理性地把它们放在一个多维度的空间里，分析一下它们的特性。跟投资差不多，一个工作可以简单地用三维空间表述：风险（risk)，流动性（liquidity），还有回报（return）。风险是指这个工作是否很容易失去（投资产品是否会大跌）；流动性是指你是否可以随时换个工作（投资产品是否可以随时卖出变现）；而回报呢，则代表这个工作的总收入。}
\BookParagraph{和投资一样，风险高的工作潜在的未来回报自然也高一些；而你很容易换其他工作，多半儿公司把你换成别人也不难。}
\BookParagraph{如何选择一个最好的工作？一个最简单的办法，你可以把一个工作在这三个维度各转化为一个 [0, 1] 区间里的量，零代表最差，一代表最好，再计算一下这个工作在三维空间里的位置到零点的距离。距离最长者为胜。}
\BookParagraph{比如「创业」这个选择，回报比普通工作高百分之五十，风险则可能高百分之一百，在三维空间里，也许是一个很糟糕的选择。}
\BookParagraph{当然，如果你有一个女朋友或男朋友，最好的方式是进行「优化投资组合」：一个人的工作风险较小，流动性和长期回报较低；另一个人去选择风险高、潜在回报高、流动性较强的工作。这样可以保证一定回报，又尽量降低风险。}
\BookParagraph{而这个理论的前提是两个人需要相互信任。翟欣欣之所以找苏享茂，苏享茂之所以找翟欣欣，大概都是「优化投资组合」的结果。}
\BookDate{2018 年 2 月，多伦多}
\BookArticle{自我介绍}{自我介绍}
\BookParagraph{李老师，您好！欣赏您看待问题又中肯又不落俗套，常常让人耳目一新。再过两个月我就要离开学校去公司上班了，想请教您一个问题：进入一家新的公司，如何做一个好的自我介绍呢？（问题来自微博私信）}
\BookParagraph{自我介绍与其追求标新立异，不如突出重点。而介绍自己时最关键的信息，是自己的名字。我一向觉得，中文名字和英文单词一样，是最不好记的。绝大多数的父母给孩子起名字的时候，偏爱一些美好或者励志的寓意，完全不考虑别人是否容易记住。我去开会的时候见到很多陌生人，上来的自我介绍是「咱俩加个微信吧」，导致最后我甚至忘了微信里哪个陌生人对应的是他。}
\BookParagraph{这一点英文就好很多，无论姓氏多么难记，至少名字千篇一律，非常好记。比如系里有一位教授叫Frank Kschischang，他幽默地把自己的难记的姓氏画成一幅编码示意图；我的导师叫 Klara Nahrstedt，她干脆建议大家不要记她的姓算了。而无论是Frank还是Klara，都可以第一次见面就轻易记一辈子。}
\BookParagraph{我自己的名字算是中文里相对容易发音的，但是对于外国人来说并不好记。在俱乐部打球时的对方队员，有时候上来就知道他自己记不住这个复杂的名字，于是问「Can I call you Bo?」简直就是当下要给我起个好记的外号。}
\BookParagraph{对于中国人来说，我的名字则算是比较好记的。即使如此，我自我介绍时，多半还要补充一句这个名字的来历：「“葆”是草字头的“葆”。毛主席说：永葆革命青春。意思是人虽然自然生命会衰老，但是政治生命却可以永葆青春。」}
\BookParagraph{我发现一提到毛主席，很多人都可以加深印象，更容易记住这个名字。}
\BookParagraph{所以啊，你介绍自己的名字时可以介绍一下这个名字的由来和一条尽人皆知的典故，尽量让大多数人一次记住。不过也不要高兴得太早，以后见面时需要再把自己的名字重复几次，加深别人的印象，就基本大功告成了。}
\BookDate{2018 年 4 月，多伦多}
\BookArticle{智　能}{智能}
\BookParagraph{您好。我这里有篇文章，题为《目前的计算机还没有实现真正的智能》。希望您能给些意见。（问题来自在行一点/分答）}
\BookParagraph{这个问题你问错人了，我不是研究人工智能领域的，所以不懂人工智能。}
\BookParagraph{不过，不仅仅是我不懂，你提到的文章作者，也照样不懂人工智能。咱们就来分析分析他的立论，到底目前的计算机是否已经实现了真正的智能。}
\BookParagraph{关键的问题是，到底什么才是「真正的智能」呢？}
\BookParagraph{我想起来二十年前莱温斯基事件中，克林顿总统对大陪审团说的一句名言：「It depends on how you define “alone”.」}
\BookParagraph{借用克林顿的句式：「It depends on how you define intelligence.」}
\BookParagraph{这个看上去挺形而上学的哲学问题，其实著名计算机科学家阿兰·图灵（Alan Turing）早在1950年就已经深入地研究过，还把他的研究成果发表在一篇题为「Computing Machinery and Intelligence」的论文里。论文中他提到，直接回答「机器是否能思考？」（Can machines think?）这个问题比较难，因为「思考」这个词儿太抽象。他提议把问题改为：「计算机是否可以玩好模仿游戏？」（Are there computers which would do well in the imitation game?）后人管这篇论文的中心思想叫「图灵测验」（the Turing test）。}
\BookParagraph{模仿游戏其实挺好玩儿的：假设一个男生、一个女生和一个提问者在不同的房间里，只能用文字交流。男生的目的是让提问者觉得他才是女生，而女生则尽力帮助提问者做出正确的判断。图灵认为，如果用计算机来代替这个捣乱的男生，提问者判断错误的频率和以前一样，那计算机就具备了「真正的智能」。}
\BookParagraph{这就是为什么几年前的那部有关图灵生平的电影的名字就叫\BookLineBreak{}《模仿游戏》。}
\BookParagraph{您提到的这篇文章所讨论的问题，六十多年前的图灵早就想过了。图灵还写了论文，他的生平后来还拍了电影，这篇文章的作者竟然一无所知。所以说该作者比我也高明不到哪儿去，也不懂人工智能。}
\BookParagraph{那计算机现在是否已经通过了图灵测验呢？}
\BookParagraph{您看过前几天 Google I/O 2018 里有关 Google Duplex 的演示吗？计算机可以打电话给餐厅预约时间订座位，对话清晰自然，连英文都像常人一样略带犹豫和停顿。这个演示极其震撼，集谷歌多年的研究之大成，只有一个英文词儿可以形容：jaw-dropping。斯坦福大学前校长 John Hennessy 说，如果仅限于打电话预约时间这件事儿，计算机其实已经通过了图灵测验。（In the domain of making appointments, it passes the Turing test.）}
\BookParagraph{当然，约个时间还不能算是真正的智能。几年前我非常喜欢的一部叫做《她》（Her）的电影里，男主人公 Theodore 深深地爱上了电影里管它叫「操作系统」的计算机。这个超级智能的「操作系统」既可以旁征博引，又可以打情骂俏；既可以风趣幽默，又可以风情万种。绝对称得上是「此中有真意，欲辨已忘言。」}
\BookParagraph{所以我认为，真正的智能，是在谈恋爱这个领域通过了图灵测验。就好像一台计算机通读过魔鬼咨询师阮琦老师的三部曲，却不是生搬硬套其中的好词好句，而是活学活用其思想，进而举一反三，充分地发挥其主观能动性，搭讪也好，约会也罢，顽皮时令人莞尔，情到深处时则令人动容。}
\BookParagraph{到那时候，世界上会有两种人：一种是谈恋爱不知道怎么接话时咨询一下计算机，我管这叫「增强智能」（augmented intelligence）；还有一种直接和计算机谈恋爱，谈得乐不思蜀。}
\BookParagraph{也许到那时候，我们真的要好好反思一下克林顿的名言：「It depends on how you define “alone”.」}
\BookDate{2018 年 5 月，西安}
\BookArticle{胆大心细}{胆大心细}
\BookParagraph{李老师，你好。想再麻烦你一次。请问在工作中，如何让自己变得细心，不出差错？是一遍又一遍地提醒自己「要小心，慢慢来」吗？前几天工作出了差错，很内疚，对自己很失望。谢谢。\BookLineBreak{}（问题来自微博私信）}
\BookParagraph{我们小学中学参加各种考试时，都可能遇到考试结束时间还没到，题已经做完了的情况。这时你是会提前起身交卷离开呢，还是反复检查答案，直到最后一分钟？}
\BookParagraph{相信你的老师一定曾经强调过，做完了题一定要反复检查，不要提前交卷。}
\BookParagraph{老师说得很有道理啊：题目大多有陷阱，自己也可能粗心大意，好好检查一下，可以提高成绩。提前交卷，你早几分钟离开教室还能干什么去？}
\BookParagraph{可我倒是觉得，这也得看这个考试是否重要。如果是高考或者GRE，那当然要全力以赴；如果只是个普通的期中考试或数学竞赛，提前交卷有诸多好处：不仅自己可以更加自信，心情愉快的同时，说不定还能给同学们带来一点儿心理压力。}
\BookParagraph{所以说，做任何事情，不出差错的概率和你投入的时间是相关的。你要是投入了足够的时间，自然可以把出差错的概率降到最低。可是这样做到底是否有必要，那要看这件事儿的重要程度，是否值得你的时间。不要忘了，把一件事情迅速、高效、果断地完成，既能在领导面前体现你的能力，自己也能从中得到不少乐趣。}
\BookParagraph{我在清华上学时就曾经参加过一次全校的英语竞赛。那场考试我不但第一个交卷，还得了第一名，并获得了一笔奖学金。如果仅仅是第一名而不是第一个交卷，我想自己也不会到现在还记忆犹新，到这个回答里来津津乐道。}
\BookParagraph{不去斤斤计较地避免出错还有一个好处，可以减轻工作给自己带来的心理压力。我是一个细心的人，事必躬亲，力求完美，直到后来连做梦都梦到上课的时候把一道习题讲错了。凌晨从梦中惊醒的时候，我意识到其实一个人对心理压力的承受能力比自己想象的要差很多。}
\BookParagraph{这样说来，不那么重要的事儿就不要花太多时间，尽量拖延到最后，然后高效粗线条地完成。即使出了错，说不定结果也没有你想象的那么糟糕。}
\BookParagraph{当然，如果这件事情确实非常重要，甚至重要到可以决定今后的事业发展，那绝对要投入大量时间，认真做得尽善尽美。但是有很多事儿，比如分析财务报表或者开发程序，需要非常专心才能保证万无一失。对付这种重要而又容易出错的事儿，我认为「不能一口吃个胖子」，「仗要一个一个打」，要「论持久战」。一次不要工作太长时间，前一天的工作后一天检查，分而治之，每一段时间尽量高效地完成一个小目标。如果需要和别人合作，就更需要提前开始，预留足够的时间检查合作者的工作中是否有问题。}
\BookParagraph{「胆大心细」，重要的事儿投入大量时间尽善尽美，不那么重要的事儿意识到「离了你地球照转不误」，是我对付工作的哲学。}
\BookDate{2018 年 6 月，西安}
\BookArticle{华山论剑}{华山论剑}
\BookParagraph{李老师好！我是通过褚老知道的您，特别想问问，您是怎样保持身材和身体素质的（褚老说您的身体保持的极好）。据说男的三十以后身体开始走下坡路，我快29岁了，平时也有时间并且也在锻炼，很想和您请教一个科学的健身和饮食计划，感谢！（问题来自微博私信）}
\BookSubtitle{一}
\BookParagraph{运动和投资一样，其实我们最需要的是一个长期的策略。也就是说，和科学的健身饮食计划比起来，最关键的问题其实是：如何能在四十、五十、六十、七十岁时一直坚持运动。}
\BookParagraph{很多人年轻时时间比较充裕，运动健身非常积极；而一旦到了四五十岁，体能直线下降的趋势一旦显现，就失去了运动的热情。}
\BookParagraph{而绝大多数有关运动和健身的建议，都在强调「三个月见效」，如同投资建议强调短期收益一样。}
\BookParagraph{别告诉我「运动产生多巴胺，会带来快乐的感觉」：相比躺在沙发上看电影，长跑或者举重绝对是一种痛苦的体验，毫无快乐可言。}
\BookParagraph{长期坚持一种给自己带来痛苦的东西，太不容易了，依靠的动力不可能一成不变。所以最重要的是如何想尽办法，找到适合自己的动力（motivation）。}
\BookSubtitle{二}
\BookParagraph{很多人喜欢在微信朋友圈晒自己的成绩，点赞数一多，动力也相应而来。}
\BookParagraph{但是靠朋友圈提供动力有两个问题。一是每周三四次运动如果全都晒出来，很快朋友们就看得没劲儿了，长期来看点赞和评论只有可能单调递减。}
\BookParagraph{比如这篇微博是我刚跑完步在星巴克一边喝咖啡一边写的，正打算晒一下刚才的成绩。你会发现，点赞数多半儿不如平时。}
\BookParagraph{二是像力量训练这样必须长期坚持的运动无法截屏晒图片，不容易分享出来。}
\BookParagraph{当然，孙俪可以让邓超给她录一小时的健身视频，再专业剪辑加速发到微博。她一定会收到无数赞美，「身材太好了」、「这什么器械我都不知道啊」。}
\BookParagraph{而你要是这么费劲儿折腾视频来分享，一定回应寥寥。}
\BookParagraph{所以，朋友圈微博晒成绩，不是一个长久之计。}
\BookSubtitle{三}
\BookParagraph{那报名跑马拉松呢？有了目标是不是就可以经常运动了？}
\BookParagraph{我同意有一个近期的目标确实非常重要。但是马拉松是极限运动，对体能要求很高，多半儿跑了几次就懒得跑了，能常年坚持的是极少数。而且，马拉松训练要求每月跑量至少两百公里以上，相当于每周需要五六个小时为跑马拉松「工作」。大多数人没有这么多时间，所以我认为跑马拉松是一种时间上的「奢侈品」。}
\BookParagraph{太重视跑步还有一个问题：运动的项目过于单一。长期运动最重要的是以力量训练为「主食」，辅以花样繁多的有氧训练当「配菜」。如果每周只有五六个小时，正确的做法是拿出一半时间做力量训练，余下的时间尽量尝试各种运动项目的结合：跑步、动感单车、游泳、拳击、跳绳、瑜伽、普拉提、羽毛球，不一而足。}
\BookSubtitle{四}
\BookParagraph{我在美国上学时有一个好朋友叫张戎，当年是清华十杰，绝对是一等一的「大牛」学霸。现在戎爷年龄摆在那儿，已跻身「大佬」行列。张戎是运动健将，当年跑铁人三项的，现在经常参加马拉松比赛，竟然跑过很多人梦寐以求的Big Sur International Marathon。}
\BookParagraph{几年前张戎来大瀑布玩儿，我们早上一起沿着尼亚加拉河约跑十五公里。他一路闲庭信步，谈笑风生，我则不得不全力以赴。但我跑完了觉得心情好极了，平时多运动，竟然可以和当年可望而不可及的运动「大佬」比肩。}
\BookParagraph{所以说，更靠谱的目标应该是自己生活中的朋友。无论什么年龄段，都可以找到三两好友，或遥相「PK」，或亲临其境，成为一种永恒的动力。}
\BookParagraph{我前一阵儿在西安一个小馆子新认识了两个「酒肉」朋友。大家酒过三巡，相约有机会时一起轻装上路，以最快的速度挑战华山。}
\BookParagraph{今天早上，我想着这个约定，顶着炙热的朝阳，以均速六分四十五跑完十一公里。水壶里弹尽粮绝时，我跑在一条几乎没有人的自行车道上。远处的高压线告诉我，这一切的痛苦没有什么意义。}
\BookParagraph{想着的只有那个和以前素不相识的两个人的一个约定：一起以最快速度奔赴华山绝顶，和「华山论剑」那块大石头合影留念。}
\BookDate{2018 年 7 月，多伦多}
\BookArticle{领导艺术}{领导艺术}
\BookParagraph{我来问李老师个问题。push 学生 push 不动怎么办？}
\BookParagraph{组里有几个学生，进度有快有慢。叫他们发周表报告进展，有的人一眼就能看出来一周都没干什么。一件事情布置下去，问一次没啥进展，再问还是没啥进展，我都懒得问了，显得老师反而急吼吼的沉不住气。你说批评吧，怕批评重了想不开。不批评吧，那就有点进展表扬一下，其实也很违心，因为实在是太差，找不到亮点，硬要表扬，他们也没法进步。灌鸡汤吧，也没有什么效果。所以就处在一种明明知道学生很差，却又怕告诉他们事实会导致工作更不好开展的局面。咋办呢？（问题来自微信）}
\BookParagraph{「Push」这个词儿老是让我联想到马克思说的「资本家剥削工人的剩余价值」，似乎有点儿贬义，要不然咱换一种说法。}
\BookParagraph{咱们就来聊聊「领导艺术」吧。}
\BookSubtitle{一}
\BookParagraph{有一特好看的电影叫「When Harry Met Sally」你看过吗？电影中男主角 Harry 有一句名言：「There are two kinds of women: high maintenance and low maintenance.」}
\BookParagraph{其实员工也有两种：「high maintenance」和 「low maintenance」。管理这两种员工的风格截然不同。}
\BookParagraph{「High maintenance」的员工，需要一种细致入微、甚至有点儿封建家长式独断专行的管理风格。比如敏捷软件开发（agile development）的核心叫「Scrum」，一度非常流行，就是这种管理风格很好的一个例子。「Scrum Master」事无巨细事必躬亲地去跟踪员工的进度，每天都要站着开一个十五分钟的会。每个人都必须发言，并且要向所有成员当面汇报你昨天完成了什么，承诺你今天要完成什么。}
\BookParagraph{作为领导，你自然就要多花很多时间和心血，严格控制和跟踪项目的进度，制作和更新 Gantt chart。所以说，这绝对是「high maintenance」。}
\BookParagraph{你也许会问，科研往往需要的是一点一滴的积累和井喷的灵感，怎么能靠这种天天检查进度的办法来管理呢？这个问题其实2010年马里兰大学的两位教授就研究过，发现其实确实是行之有效的，还专门写了一篇论文来总结他们的成功案例 {\BlissSymbols ①}。}
\BookParagraph{但是另外还有一种员工，是典型的「low maintenance」。他们的求知欲和创新能力都很强，更乐于独立思考，独立管理自己的时间，直至独立完成一个引以为傲的项目。管理这种员工最好的方式是充分信任他们，「无为而治」，英文叫「laissez-faire」。唯一需要做的就是充分提供所需的资源，他们有问题时也可以随时向你咨询。}
\BookParagraph{当然最可怕的一种员工是眼高手低，自己没有独当一面的能力，却好高骛远。就像「When Harry Met Sally」里 Harry 对 Sally 说的：「You are the worst kind: you are high maintenance, but you think you are low maintenance.」}
\BookSubtitle{二}
\BookParagraph{一个微博上的问题，如果没有任何回报，大概不会有人去回答这个问题。}
\BookParagraph{但是如果回答这个问题有一千块钱收入，一定会有不少人愿意去回答。}
\BookParagraph{这符合博弈论里的一个基本假设：每个玩「游戏」的玩家都是理性的（rational），他们不仅有愿望而且有能力可以将自己的收益最大化。}
\BookParagraph{诺贝尔经济学奖得主、著名经济学家 Milton Friedman 说过：「The world runs on individuals pursuing their self interests.」}
\BookParagraph{如果干得好坏一个样儿，不会有人愿意努力工作。作为领导，你可以设计一种合理而有效的激励机制，在「比学赶帮超」的气氛中，让多快好省地完成项目的员工脱颖而出，得到他们应得的「分红」。这就是为什么公司里每年有绩效评估的原因：绩效评估的结果直接和奖惩机制挂钩。}
\BookParagraph{当然，科研工作的「分红」不一定是奖金，也可能是去夏威夷开个国际会议的机会，是某个全校优秀论文的奖项，是提前一两年毕业，甚至仅仅是因为在截止日期之前写完了一篇重要论文所得到的通报表扬。}
\BookParagraph{关键之处在于，你设计的激励机制里这些潜在的回报都绝不是你临时想出来的，而是早就白纸黑字有章可循，公平合理。绩效评估的设计绝对不能变相鼓动员工之间你死我活相互竞争，而是尽量鼓励合作双赢。在最好的激励机制下，当每一个员工都竭尽所能将自己的收益最大化时，实现你自己既定目标的可能性趋于最大。}
\BookSubtitle{三}
\BookParagraph{古老的「鸡汤」，或者说好听点儿叫「情怀」，也不一定就一点儿用都没有。}
\BookParagraph{比如有的公司——我不说估计你都能想起来是哪家——依靠的正是其领军人物的领袖魅力（charisma）。}
\BookParagraph{而且往往最优秀的员工所看重的，正好就是这些散发着领袖魅力的「情怀」。}
\BookParagraph{你一定会问，可这样的领袖魅力可遇而不可求啊，跟咱这种\BookLineBreak{}「衣带渐宽终不悔」的普通中青年小领导有关系嘛？}
\BookParagraph{我倒是觉得，这样的领袖魅力其实就是一种自然而然的亲和力。就像夏天里家里的姥爷坐在胡同口的藤椅上，一边扇着扇子，一边讲着一个引人入胜的三国故事。}
\BookParagraph{而你给大家的，不是什么「鸡汤」，也不是什么「情怀」，仅仅就是茶余饭后润物细无声地讲了几个故事而已。}
\BookDate{2018 年 7 月，多伦多}
\BookParagraph{{\BlissSymbols ①} M. Hicks, J. S. Foster, “Adapting Scrum to Managing a Research Group,” http://www.cs.umd.edu/projects/PL/score/.}
\BookArticle{工作和生活}{工作和生活}
\BookParagraph{李老师，您好，关注您微博有一段时间了，想咨询您一个问题。}
\BookParagraph{在您看来应该如何平衡科研与生活？}
\BookParagraph{现在我是一名工科硕士生，现在感觉每天都呆在实验室，一周有时超过八十个小时，几乎没有私人时间，感觉这样的日子并不是一个正常人的生活。在我看来一个正常的生活是应该去做一些工作之外的事情的，但是现在如果一有懈怠就会达不到老师的要求。}
\BookParagraph{对此您怎么看？ 谢谢。（问题来自微博私信）}
\BookParagraph{正好前一阵子和一新认识的研究生吃饭聊天儿，聊到她的老板老是半夜三更地发微信找她办事儿。}
\BookParagraph{我说：「不回老板微信不就得了？等到第二天上午再回也不晚啊。」}
\BookParagraph{她反问：「你还想不想毕业了？」}
\BookParagraph{她给我看她的微信记录，不但工作到凌晨第一时间把事儿办好发回给老板，还友好地加上一句：「您早点儿休息。」}
\BookParagraph{她跟我解释，没办法啊，都是为了毕业。师兄有的就不常回微信，结果老板连工资都发得少的可怜。}
\BookParagraph{看样子你的情形和这个差不多：每周科研时间少了，就达不到「老师」的要求。}
\BookParagraph{如果这个「少了」的定义是少于八十小时，那我只能说，这位老师不配称作「老师」，应该按照国内的行规，称作「老板」。}
\BookParagraph{以前在清华上学的时候，每天下午四点半广播里都会响起一个嘹亮的声音：「同学们，让我们走出教室，走出宿舍，去参加体育锻炼，争取为祖国健康地工作五十年。」}
\BookParagraph{一周工作八十小时，估计你多半儿没法儿为祖国健康工作五十年。就像如果你按四百米短跑的速度去跑马拉松，多半儿跑不完全程一样。}
\BookParagraph{工作和生活哪个占第一位？当然是生活第一。}
\BookParagraph{疲惫不堪时，就应该休息；浑身乏力时，就应该运动；该吃就吃，该玩儿就玩儿。无论何时何地，你所爱的人，你的父母，你的孩子，都比你的工作更重要。}
\BookParagraph{那我老板是个「workaholic」，动不动半夜发微信怎么办？}
\BookParagraph{哈哈，建议你早睡早起，跟老板的作息时间完全「错峰出行」。实在不行就跟苹果的 Tim Cook 学习，每天晚上九点睡，早上四点起床，起床后运动一小时，五点开始工作。}
\BookParagraph{这样老板的微信不是装作没看见，是真的没看见。}
\BookParagraph{当然最理想的状态是不把科研当成一个为了老板而不得不完成的工作，而是当成为了自己的将来而「投资」的一种学习过程。这样无论每周时间多还是少，「好好学习，天天向上」，始终是一种因为学到了新东西而让人感到开心的体验。}
\BookParagraph{可惜这种开心学习的体验需要慢慢地一点一滴培养，却很容易被一条微信粉碎，一切退回原点。}
\BookParagraph{就像一个容易失眠的人，好不容易数着绵羊快要睡着了，忽然手机震动了一下，瞬间又异常清醒，各种烦心的事儿纷沓而来，终究无法功德圆满地进入梦乡。}
\BookDate{2018 年 8 月，多伦多}
\BookArticle{前　途}{前途}
\BookParagraph{李老师您好！从褚老师那认识您，一路看您的问答过来，感觉您就像身边的一位长者朋友，如沐春风，沁人心田。不知道您是否还记得，1988年版本的《新华字典》中有这么一句话：「张华考上了北京大学；李萍进了中等技术学校；我在百货公司当售货员：我们都有光明的前途。」}
\BookParagraph{这句话最近被网友改成了「张华考上了985大学；李萍进入了某垄断国企的干部培训班；我在世界五百强商超当供应链主管：我们都有光明的前途。」以此来笑谈现在的阶级分化和贫富差距。我之前翻阅最新版的《新华字典》也看到了这句话，但我的第一反应是，当时的人们感情真挚，不会因环境变化而自卑，情谊疏远。}
\BookParagraph{我现在是一名正在准备考研的学生。周围的同学也都想通过考研改变一些现状，对未来充满了期待。但我也知道和以前优秀的同学终究是有很大的差距。如何面对这种差距，正确地放弃幻想，轻装前进呢？（问题来自微博私信）}
\BookSubtitle{一}
\BookParagraph{俗话说得好：人比人，气死人。任何事儿都千万别跟别人比，别给自己树立标杆儿和榜样，否则一定跟小时候家长拿自己和「别人家的孩子」比如出一辙，最后是自个儿给自个儿找不痛快。}
\BookParagraph{尤其不能比的，是好久以前的同学。时间这个东西，是一种奇妙的催化剂：时间长了，一切早已面目全非。当年和同学分道扬镳，你走你的阳关道，我过我的独木桥，这么多年来每一次微不足道的决定，都会影响后来的「前途」。}
\BookParagraph{大多数同学聚会，除了特别要好的几个之外，与其说是去重温旧情，不如说是好奇心理使然，去看看其他同学都混得如何。同学聚会是一种典型的有偏采样：自己觉得混得不好都没去参加，参加了的都有点儿「来头儿」，有可以炫耀的资本。所以啊，这样的同学聚会，将来还是尽量少参加为妙。}
\BookSubtitle{二}
\BookParagraph{那些「前途光明」的同学，当年是怎么出人头地，和大家拉开差距的？}
\BookParagraph{这需要三样东西，缺一不可：资质、兴趣、运气。}
\BookParagraph{资质就不用说了。即使是在微信公众号靠写文章赚钱的，没有一点儿天赋，也很难写出全国人民纷纷转发的「十万加」来。}
\BookParagraph{以前有一种说法叫「干一行，爱一行」。「兴趣」就是把因果关系反过来，改成「爱一行，干一行」。你决定干的事儿一定是你不吃饭不睡觉地热爱着的事儿。「前途光明」的数学家不是为了毕业而学的数学，而是因为热爱得不可开交，天天浸淫其中，乐不思蜀。「混得好」的企业家也一样，完全是因为一种执着的、持之以恒的、发自内心的热爱。}
\BookParagraph{资质和兴趣都具备了，还要有足够好的运气。这一点 Gladwell在他那本著名的《Outliers》里已经说得很明确了：乔布斯和盖茨等人能成功的一大因素是他们赶上了好时候：他们都是1954至1956年出生的。}
\BookParagraph{你看，想有个「光明的前途」真不是那么容易的一件事儿。}
\BookSubtitle{三}
\BookParagraph{我小学时妈妈跟我说，她年轻时特别喜欢骑自行车，所以对未来一直有一个美好的理想。}
\BookParagraph{我问：是当个自行车运动员？}
\BookParagraph{她笑着跟我说，你猜错啦，是当个邮递员。}
\BookParagraph{现在想来，那个年代当个邮递员和新华字典里说的在百货公司当售货员大概异曲同工：一辈子有稳定的收入，有劳保，有退休金，绝对是一份非常令人羡慕的工作。}
\BookParagraph{这有点儿像在加拿大啊挪威啊这样的高福利国家当个空乘，当个小学老师，或者当个公共汽车司机，有工会制度保护，有稳定的收入，有劳保，有退休金。}
\BookParagraph{这样的生活波澜不惊，既没有「过独木桥」的风险，又没有和风险相应的回报，是一种「一眼能看到火葬场」的「小日子」。你可以精打细算，量入为出，有时间尽心尽力地照顾家人和孩子，尽享天伦之乐。}
\BookParagraph{那将来如果是这样红红火火的「小日子」，算是「光明的前途」吗？二十年后能有面子跟有点儿「来头儿」的同学们觥筹交错、共叙旧情吗？}
\BookParagraph{如果你觉得可以，那么我有个好消息要告诉你：这样的「光明前途」既不需要太高的资质，又不需要「爱一行，干一行」的兴趣，连运气都不需要太好，兢兢业业地把每天的事儿做完做好就可以了。}
\BookParagraph{妈妈后来清华毕业，但终究没有当上邮递员。}
\BookDate{2018 年 8 月，多伦多}
\BookArticle{经典}{经典}
\BookParagraph{不知道为什么，这么多年来我一直在思想上无法与时俱进，接受「PPT」这个日常用语。但凡有学生找我要「PPT」，我都会礼貌地告诉她：我没有「PPT」，因为我的幻灯片是用 Keynote 制作的。}
\BookParagraph{但其实我心里也明白，她所谓的「PPT」其实就是「幻灯片」的意思，不是特指「用 PowerPoint 制作的幻灯片文件」，没有必要较真。不过下回换了一个学生来问，尤其是那些竟然写成「ppt」的学生，我又会像一个回零的秒表，忘了「不要斤斤计较」的初衷，再次郑重宣布：我没有「PPT」。}
\BookParagraph{「幻灯片」这个词儿，就像「黑胶唱片」、「上海手表」或者「TDK 磁带」一样，给我一种很亲切的感觉。我想大概是因为我在美国读书时经常用到投影幻灯机，需要事先打印好一摞那种透明的幻灯片。每张幻灯片左边都有一张长长的纸条需要撕掉，超级有仪式感。印象最深的是博士生资格考试那天，我紧张得凌晨就醒了，睡不着就慢慢地撕幻灯片的纸条，等着天亮去考试。}
\BookParagraph{不过后来我发现，幻灯片还不够「经典」。刚到多伦多大学头一两年，讲课的学生反馈总是在系平均水平上下浮动，没什么进步。直到后来我才明白，最好的教学方式不是用透明幻灯片和投影幻灯机，而是带上一盒粉笔直接写黑板。如果能够「更上一层楼」完全脱稿，从算法描述到数学证明，粉笔在里外两三层的黑板上下翻飞，游刃有余，那才是讲课的最高境界。}
\BookParagraph{理查德·费曼当年物理课上穿着笔挺的西装在黑板上写字时的那种令人着迷的范儿，就像大礼堂草坪上弹吉他的校园歌手，像学生节台上美丽的主持人，无论如何设计如今的「PPT」，我想都无法匹敌吧。}
\BookDate{2019 年 5 月，多伦多}
\BookArticle{百度}{百度}
\BookParagraph{看到微博里有国外的老师回国工作，抱怨百度不好用。恰巧我前一阵子也回国教了一门短期课程，惊喜地发现百度其实并不是一无是处。}
\BookParagraph{比如上周因为备课需要，我想安装 PyTorch + Python 3.7 的最新版本。正常情况下，应该是命令行一个命令，再等一杯咖啡工夫的事儿。结果我花了一整个晚上的时间一次一次地尝试，为了下载一个几十 MB 的 PyTorch 安装包。不但进度极其缓慢，而且下到一半儿就断了，下次必须重新下载。}
\BookParagraph{我有点儿不服气：作为一名学计算机专业的博士，修不了 Windows 也就罢了，连一个安装包都下载不了，水平未免也太低了吧。}
\BookParagraph{这时我竟然莫名其妙地想到了百度。我用百度直接搜索安装包的名字，第一个链接就是清华大学的镜像网站。从这个网站下载，十几秒钟就顺利完成了，解决了一大疑难问题。}
\BookParagraph{不过也别高兴得太早了。课程结束后我到北京转机回加拿大，转机时间较长，计划进城看望一下父亲。我抱着试试看的态度，用百度地图搜索从机场去马甸桥的最佳公交线路。百度地图告诉我，最好的办法是机场快线去三元桥，然后出站转 847 路。我听说机场快线需要现金才能买票，还专门找了一个取款机取了现金，做好了上路前的充分准备。}
\BookParagraph{到了三元桥站，我发现情况有点儿不妙。三元桥站有 A、B、C1、C2、C3 五个出站口，没有任何标识说明哪一个出站口出来可以坐公交车。跟着百度导航出站，发现是一个错误，只好再重新回到地铁站里，换一个出站口出站。运气不太好，换到第三个出站口才找到 847 路公交车站。等了好一会儿没来车，只好放弃了这个「小试验」，后一半儿行程改成打车。}
\BookParagraph{打车时发现百度地图里也可以叫车 —— 不过这时我已经没有「探索未知世界」的心情了。我打开「滴滴」，车一会儿就到了，仿佛又回到了自己熟悉的「地盘」。}
\BookParagraph{回到加拿大，又想起了这件小事儿。「完美主义」使然，我又做了一点儿调研。结果发现百度开始的建议就是错误的：最好的办法是坐去公主坟方向的机场大巴，可以直接在马甸桥站下车。看来即使在这个「大数据」和「人工智能」的时代，还是不能完全依赖「临阵磨枪」，还是要事先做好调研，取好现金，才能「学好数理化，走遍天下都不怕」。}
\BookDate{2019 年 7 月，多伦多}
\BookArticle{极致的简单}{极致的简单}
\BookParagraph{前一阵子回国讲课，预定的行程比较紧张，只讲了一周时间，几乎每天上午都要上课。由于这次的课是第一次开设，需要大量时间准备。结果是每天从下午开始工作，晚饭后准备到晚上十一点，早上三点多起来继续，往往需要准备到七点钟离开酒店去上课之前。}
\BookParagraph{即使如此认真地抓紧时间，要不是用上了「手写幻灯片」这一大法宝，其实备课时间还是远远不够。}
\BookParagraph{比如说，如果幻灯片需要数学公式，最快的传统方法是用 LaTeX 排版，完成以后再以矢量图的形式拖拽插入幻灯片。准备五六张幻灯片排满数学公式，我尝试了一下，再快也得一个多小时。如果中间还要插入文字，或者把排版数学公式换成画示意图，时间大概需要成倍增加。}
\BookParagraph{而一节课四十五分钟，需要至少二十至三十张幻灯片。「换算」成时间，要是按照传统方式设计幻灯片，五个小时估计都打不住。}
\BookParagraph{手写绝大多数幻灯片，一个大问题是选择 iPad 里哪个 画图工具效率最高。尝试了 Procreate 等好几种可能，发现还是这几年常用的一款最好用：Paper。}
\BookParagraph{Paper 好用不是因为它的功能多强大，恰恰相反，是因为它功能不多，但是每一个功能在设计上都用了心思。我所有的幻灯片一共就只用到了三种笔：粗笔、细笔、水彩笔。粗笔能不能再细一点儿？水彩能不能再透一点儿？不行。就连颜色我都全程只选红色、蓝色和黑色。}
\BookParagraph{这就好像一家有点儿「小资」的馆子，把菜单简化到只有几个菜，点菜自然就不用千辛万苦地斟酌再三了 —— 反正每个菜都可以体现这家馆子的最高水平。}
\BookParagraph{又好像苹果手表虽然功能多得目不暇接，但是真正能用到的却很少。倒是「看时间」这个最为重要的功能，与其戴苹果手表，不如戴一块卡西欧电子表。}
\BookParagraph{极致的简单才能带来极致的高效。}
\BookParagraph{一周时间我准备了十六节课，几百张幻灯片。一旦内容想好了，每张幻灯片只需要一两分钟就可以一气呵成，再用 AirDrop 迅速传到电脑上，心情愉快。即使在每小时三百公里疾驰的高铁上，因为心思全都集中在内容上，工作效率和在酒店几乎完全一样。}
\BookParagraph{当然了，手绘幻灯片还有一大好处：图文并茂、色彩鲜明、与众不同，充分体现了一个完全不会画画儿的理工科中年教师对可望而不可及的艺术世界长期以来孜孜不倦的向往。}
\BookDate{2019 年 7 月，多伦多}
\BookArticle{跟着感觉走}{跟着感觉走}
\BookQuote{李老师，好久没看您发微博啦。最近看一个剧里有句台词触动很大：人生一世，选条路，不退让，不更改，一直走下去，是件幸事。我感觉，随着认知面的拓宽，自己越来越会辩证地去思考问题，且随着知识情感等不断地丰富，这个辩证的过程也会变得越发的复杂甚至最终被困于其中。就好像我有两个我在打一场少有结局的辩论赛。所以有些时候看到一些「一根筋」的人，会羡慕他们的单纯可爱并为其事迹感动。所以，是不是我「想太多」呢？又怎么让自己不想太多呢？（问题来自微博私信）}
\BookParagraph{想得多一点儿，其实是好事儿啊。如果「一根筋」的结果是「一条道走到黑」，而非「浪子回头金不换」，其实也不是什么特别让人羡慕的事儿。}
\BookParagraph{当然了，咱也不能想得太多 —— 天天想，月月想，想过来再想回去，结果没真正做成什么事儿，自然也不可取。}
\BookParagraph{尽量在两个极端中找到一个「sweet spot」—— 中文直译过来我觉得可以叫 「甜点」—— 想得不能太多以致一事无成，也不能太少以免「一招不慎，满盘皆输」。}
\BookParagraph{那这个「甜点」在哪儿呢？}
\BookParagraph{我个人认为，可以沿用类似法律体系中的那套路子，给自己确立几条「公理」—— 有点像法律体系中的宪法 —— 只要不违背这些「公理」，其余的不用想太多，「跟着感觉走」就可以了。}
\BookParagraph{那这些「公理」都是什么呢？}
\BookParagraph{我今天就抛砖引玉，说两个例子给你参考吧。}
\BookSubtitle{一}
\BookParagraph{第一条「公理」是：不要人云亦云，要货比三家。}
\BookParagraph{我考大学的那会儿，很多人都认为学生物工程是最好的专业选择。因为当时离二十一世纪还很遥远，人工智能还处于寒冬，大家都觉得等中国人民的生活实现了小康、国家实现了四个现代化之后，生物工程是未来的发展方向。结果很多当年最杰出的人才都去学了生物，后来去了美国留学，读完博士发现找不到工作，只好继续读博士后。读了博士后还找不到工作，只能转行学计算机，因为学计算机好找工作。}
\BookParagraph{时间过得很快，现在「三十年河东，三十年河西」，很多人都认为应该去学计算机，甚至再细分一点儿去学人工智能，量子计算，或者区块链。可是，现在热门的专业，不等同于十年学成以后还热门啊。}
\BookParagraph{这条「公理」就派上了用场：不要人云亦云，要货比三家。}
\BookParagraph{可是怎样才能做到「货比三家」呢？}
\BookParagraph{我的建议是，不能只听信一些「专家」、「大神」的意见，更不能直接相信满屏的公众号文章，而是正反双方立场各找出一两条高质量的意见，仔细研读一下，然后再通过独立分析得出自己的结论。}
\BookParagraph{选读别人的意见时，还要考虑到作者是否有利益相关。比如一位研究人工智能的学者，多半儿会说一些有利于人工智能研究的观点。我们可以了解他的观点，但是也要观察其观点是否有失公允。}
\BookParagraph{「货比三家」说得是，参考两三种不同意见就可以当机立断地拍板了，不需要太多。选好了就开始执行下去，就算后来发现选错了，也已经尽力而为，不用后悔。}
\BookSubtitle{二}
\BookParagraph{第二条「公理」是：生活第一，工作第二，做一天和尚撞一天钟。}
\BookParagraph{每个人的时间都是有限的资源，需要在工作和生活之间分配。很多人喜欢把自己的工作放在第一位，生活放在工作之后。我的这条「公理」是说，生活才永远是应该第一位的，工作远远没有我们想象的那么重要。}
\BookParagraph{你想象一下，如果知道很快就要离开这个世界，今天晚上躺在床上会想些什么？}
\BookParagraph{我也许会想我的家人，我的孩子，我的宝贝汽车和音响，甚至我的梦中情人。}
\BookParagraph{但我想我一定不会去想我的公司，我的单位，我的论文，和我得到过的荣誉。}
\BookParagraph{大多数幸福和快乐，都来自于「业余时间」的生活，而非工作。比如我现在给你写这篇回答，是我生活的一部分：高兴了就写点儿，没什么可写的就懒上一阵子 —— 这也是为什么我可以很长时间不发微博的原因 —— 写完了有人喜欢看有人点赞，让人心情愉快。但是如果让我把写公众号文章当做工作来做，每天一定要完成多少写作的工作量，我想我一定不会喜欢，幸福感会大打折扣。}
\BookParagraph{所以我觉得，在分配自己的时间和做各种取舍选择时，应该永远把生活放在第一位，因为生活才是幸福的源泉。工作可以认真地「做一天和尚撞一天钟」，一切顺风顺水自然很好，如果正好不尽如意，那就听天由命好了。}
\BookParagraph{想起了我在美国第一位导师 Dan Reed 的一句名言：}
\BookQuote{Too many good books, and too little time.}
\BookSubtitle{三}
\BookParagraph{只要把这几条「公理」记住，余下的就交给运气和时间吧。}
\BookParagraph{别人如果「一条道走到黑」，不一定是一件「幸事」，咱们也不必羡慕他们，也不必为他们的事迹而感动。}
\BookParagraph{咱就一普通人，踏踏实实地过着一普通人的日子罢了。只要自己开心，既可以「从一而终」，也可以「朝令夕改」。}
\BookParagraph{生活就像一艘帆船 —— 咱们要想要做的，是如何把帆迎着海风展开，让它鼓得高高的，全速前进。}
\BookDate{2020 年 1 月，多伦多}
\BookArticle{萝卜白菜}{萝卜白菜}
\BookQuote{李老师好，有个非常纠结的选择需要咨询您。我刚毕业留京，到了很忙的基层当公务员，正在读在职博士。工作以后被离家远的漂泊感、北京的买房等生活压力压迫得忍无可忍，准备取消录用回家乡那边的省会新一线考公或者事业，或者等博士毕业进学校。可是父母反对，他们觉得省会生活压力也不小而且发展弱、考公竞争太大容易直接失业心态更不好。我自己是很不愿意继续在京基层公务员了，因为生活压力太大实在没有幸福感可言，可是父母的观点也很有道理的样子，所以导致我非常犹豫。您能帮忙指点一下么？}
\BookQuote{我是独生女也没有对象，自己离家太远这个漂泊感和生活压力太可怕了。而且放眼望去也觉得可改善性很小，本人不是喜欢打拼的性格，希望能过亲朋好友围绕的温馨生活，所以一度想冒险裸辞回去省会重新开始。但是又不确定自己考虑的对不对。}
\BookQuote{我最近几天也在反省自己。可能是刚出校园，独立应对生活的能力还不够，一个人没有援助力所以受不了压力了。而且因为我身处基层，单位平台低人员素质低，直接领导是个人品和作风很差的军转中年油腻男，天天被当成情绪垃圾桶和新兵蛋子挖坑欺负，工作特别忙还总受气，痛苦又没价值的工作，一个独生女远离家的孤独，还有未来的压力和恐惧。一来我就在各方面遇到加难题型，我真的都怕自己转正进入服务期走不掉就抑郁了。可是另一方面又觉得自己是不是在逃避压力，去了省会只是房价和家远的压力稍微缓解，还得面临可能短时间找不到满意工作的综合压力，找到了是不是又是另外一个坑。我真的挺纠结，又患得患失的。}
\BookQuote{这段时间思考去留的问题，让我精疲力竭十分痛苦，还请李老师指点一二。谢谢老师！（问题来自微博私信）}
\BookSubtitle{一}
\BookParagraph{我们先来做一个假设：工作无论做什么，都会有其难处，世上不可能有十全十美的工作。}
\BookParagraph{你不觉得吗 —— 有的工作挣钱多但是太忙，有的虽然闲但是挣钱又太少；有的岗位受别人排挤，有的又责任重大，风险太高；有「痛苦又毫无价值」但是却四平八稳不容易丢的工作，就自然有挺有责任感但随时会失业的工作。}
\BookParagraph{如果我们的假设成立，那你现在的工作虽然离完美有一段距离，可你将来回到省会城市重新找到的工作，很可能比现在也好不了多少。}
\BookParagraph{如果不考虑工作好坏这个问题，那我们的这道选择题就简单多了：如果咱今后打算工作三十年，是住在北京呢，还是住在省会城市？}
\BookParagraph{回答这个问题，我绝对不是一个合格的人选：我北京出生北京长大，肯定带有北京人的偏见。}
\BookParagraph{我的答案是：我不但会选择住在北京，还会找机会到世界上其他的著名大城市都住上一阵子 —— 比方说上海、香港、首尔、东京、纽约、还有伦敦。}
\BookParagraph{也许正是因为我选择的是北京，将来才会有更多的机会去更多的地方。}
\BookSubtitle{二}
\BookParagraph{我有一位教授同事，曾经成功创立过一家外汇交易的公司。因为这家公司创立至今发展得非常成功，他早已衣食无忧 —— 按照现在时髦的词汇，叫「财务自由」了。特斯拉刚刚推出 Model S 的那一年，他就买了一辆来玩儿。他戴的智能手表是当年大名鼎鼎的 Pebble，论资历，苹果手表得管它叫「爷爷」。}
\BookParagraph{一次闲聊时，这位同事说起他每个周末把自家的豪华特斯拉开出来，在 Uber 上接单。我觉得很神奇：他也不缺这点儿外快啊，为什么要临时开出租？}
\BookParagraph{他说：就是为了好玩儿啊。}
\BookParagraph{就跟他多年前买了点儿比特币，无论涨啊跌啊从来没卖过 —— 也是为了好玩儿。}
\BookParagraph{工作二三十年，退休了还有二三十年，这么长的时间，我觉得「好玩儿」特别重要。}
\BookParagraph{事业成功家庭和睦自然很重要，但「好玩儿」其实也是一种挺重要的体验，因为还有很多其他可以「玩儿」的。比如周末没事儿去开开出租，是一种「玩儿」法；在单位跟领导斗智斗勇，是一种「玩儿」法；认识一大帮子新朋友一起撸串喝酒侃大山，也是一种「玩儿」法。}
\BookParagraph{而北京这样的大城市，玩儿法多半儿会比回到省会城市多一些。玩儿法一多，日子就变得更加有滋有味。}
\BookParagraph{就像计算机科学里有一种叫「图灵测验」的概念一样，我一直笃信一种测验，我管它叫做「deathbed test」—— 这里的「deathbed」可以说是一个中性词，没有伤感的含义。虽然意思猜都猜出来了，还真没有对应的中文词。}
\BookParagraph{这个测验是在问，临终前对自己玩儿过的这几十年，是像酒足饭饱的那种心满意足的感觉，还是后悔莫及？}
\BookParagraph{我想，如果既跟领导缠斗过，又跟朋友快活过；既开过出租、又当过房奴；玩儿过文艺，弹过吉他，又飙过跑车，到「deathbed test」的那一刻，就不太容易后悔了吧？}
\BookParagraph{否则，如果换了我，我总是会好奇：如果当年选择留在北京住下去，后来会是什么样？}
\BookSubtitle{三}
\BookParagraph{当然了，「萝卜白菜，各有所爱」，很多人就跟我截然不同，更喜欢小一些的城镇。这就好像既有喜欢儿孙满堂、也有喜欢丁克的人一样。}
\BookParagraph{我今天读到了一则在环球邮报网站上的留言，说一个加拿大人自己一个人住在阿尔伯塔省的一个小镇附近自己的农场里，离他最近的一户人家在一公里开外，每天只有一条狗陪伴着他。他只需每一两周到小镇上买一点儿水果，其他的吃的喝的农场里都有大量储备。他每天上网逛逛，在家做做木工，日子过得也挺充实。最近的世事纷争和他没什么关系，唯一的影响是小镇上给他安排的义工取消了。}
\BookParagraph{也许是因为我当年是个城里长大的孩子 —— 就算把这个农场送给我，木工玩儿腻了的时候，我也注定会感到孤单。}
\BookParagraph{如果真有世界末日的那一天，我想我会去上海徐家汇或者纽约曼哈顿找一家小酒吧，安静地等待它的到来。}
\BookDate{2020 年 3 月，多伦多}
\BookArticle{看不见摸不着}{看不见摸不着}
\BookQuote{葆春老师您好，我是一个不知名大学的本科的文科生。我父母文化不高，高考填志愿时不会帮我选专业，结果调剂到了我不喜欢的文科专业 —— 酒店管理。}
\BookQuote{我挺羡慕理科生的。理科生可以自然而然地毕业后从事自己专业的工作。}
\BookQuote{我一直都是个没有目标的人。我好像从未对什么提起过很大兴趣，我也不是个勤奋努力的人。}
\BookQuote{今年就要毕业了。我想去上海，但我不知道自己到底可以做什么。我想过考研，但也没有明确的考研目标，我真的很迷茫。}
\BookQuote{我很后悔自己以前的选择，这两年来都沉浸在痛苦中，非常希望葆春老师能指点一二，谢谢葆春老师，盼复。（问题来自微博私信）}
\BookParagraph{我以前在引用一个打算回答的问题时，问题中所用的语言和标点符号，一般都会粗略地修改一遍。比如有的人喜欢用很多逗号，经常一句话特别长。我会把问题中的一部分逗号改成句号，再多分几个自然段，让问题更加清晰易读。}
\BookParagraph{看了你的问题，我发现无论是语言还是标点，通篇竟然没有一处需要修改。我喜欢你的问题中每个自然段里第一句话不用逗号，而是直接用句号结束。这样表达更加简洁明了，颇有古人「环滁皆山也。」的开门见山，我一看到就挺有好感的。可以说，我喜欢你这个「不知名大学本科的文科生」的文风。}
\BookParagraph{那咱们就步入正题吧。}
\BookSubtitle{一}
\BookParagraph{我小的时候有一种说法，叫「学好数理化，走边天下都不怕」。高中时，很多学习成绩不错的同学都对理工科趋之若鹜。大家高考考的都是「理工农医类」的考卷，班里学习最好的同学报的是北大生物系，而前几名的同学也清一色地都报了理工科。我上大学的时候，清华还是一所以工科为主的大学，我学的虽然是计算机，但学过金工实习，学过机械制图，感觉毕业了随时可以分配到大中型企业当一名工程师。}
\BookParagraph{然而不知你发现没有，国外很多大学的惯例是将文科和理科归入同一个学院，工科则归入另外一个学院。与其分为「理工科」和「文科」，不如分为「文理科」和「工科」更为贴切。加州大学伯克利分校更加别出心裁，把文理学院叫做「College of Letters and Science」，一个学院涵盖了从文学到音乐到历史到经济到生物的方方面面，涉猎之广，令人目不暇接。}
\BookParagraph{所以可以这么说，文科 —— 也就是伯克利的「Letters」—— 至少和理科是可以相提并论、分庭抗礼的。那这个「Letters」到底毕业后有没有用、是不是白学了呢？}
\BookParagraph{我很晚才领悟到，其实这个「Letters」说不定比金工实习机械制图比当年热门的金融和生物 —— 甚至比从二十多年前到现在一直热门的计算机 —— 都更加有用。}
\BookParagraph{原因很简单：跟金工实习或者计算机都不同，文科学的很多东西，属于「看不见摸不着」的技能。正是这种飘忽不定的感觉，让很多人觉得时间白白浪费全都白学了，毕业了也不知道要干什么好。可如果从长远来看，也正是这种「看不见摸不着」的技能，才能真正左右一个人长期的职业发展。}
\BookParagraph{我喜爱的电影「Jerry Maguire」和「Chicago」中的女主角 Renée Zellweger，刚上大学时学的是英语文学（English Literature）。上了一阵子她开始选修戏剧表演方面的课程，后来毕业后参演了一些小成本制作的电影。共和党著名参议员 Mitt Romney 在去哈佛大学学法律之前，本科学的也是英语文学。}
\BookParagraph{那他们的英语文学是不是白学了呢？从「看不见摸不着」的角度来看，我觉得非但没有白学，还造就了他们后来的事业。如果他们本科时学的是工科，那说不定现在也只不过是在一家大公司的研发部门供职，而不会成为好莱坞星光大道中的一员，不会成为政坛的一代枭雄了。}
\BookParagraph{乔布斯零六年在斯坦福的演讲中说：「You can’t connect the dots looking forward; you can only connect them looking backwards. So you have to trust that the dots will somehow connect in your future.」}
\BookParagraph{英语文学也好，酒店管理也罢，或许越是没什么用的文科，就越没有白学。}
\BookSubtitle{二}
\BookParagraph{所以说，学什么专业干什么工作是「有用」的，还真的不好说。}
\BookParagraph{1873 年，美国纽约有一个二十出头的小伙子，叫 Frank W. Woolsworth。一天，他衣衫褴褛地去应聘一家商店的营业员职位。老板跟他聊了一会儿，对他说：「行了，你下周一来上班吧。」 他非常开心，连忙问：「那我能拿多少工资？」 老板大叫道：「工资？你还想拿工资？我教你怎么做生意，你应该反过来给我发工资！」}
\BookParagraph{那时候美国的劳动法还没有最低工资保护有关的条款，Frank 得以顺利得到了这个头三个月不拿工资的工作。头几年他确实很辛苦，每天只能打扫卫生，早上七点干到晚上九点，大家都瞧不起他。后来，他逐渐开始揣摩商品摆放位置和定价的门道，终于投资了几百美金，跟几个合伙人成立了一家以自己名字命名的「初创」公司，开了自己的第一家商店。他的生意越做越大，连锁店开始遍布全美，继而又遍布全世界。1913 年，他的公司在纽约曼哈顿建成了当时全世界最高的摩天大厦。}
\BookParagraph{试想，如果当年的 Frank 看不上这个既不发工资又要打扫卫生的工作，大概就悟不出来后来得以大获全胜但却转瞬即逝的商机了。公司成功之后，他果然很感激当年的老板给了他这个不拿工资的工作，还请这位老板到他自己的公司来当合伙人，「反过来」给他发了工资。}
\BookParagraph{这样「励志鸡汤」的例子自然俯拾皆是。年逾九十的斯坦福大学著名经济学家 Thomas Sowell 在他的著作《Applied Economics》中认为，无论是上学、看书还是工作，都是一个获取新的知识和掌握新的技能的过程。像在给自己「投资」一样，是一个积累「human capital」的过程。Sowell 提到：}
\BookQuote{The growing importance of human capital tends to create greater income inequality between those people who have been assiduous in acquiring knowledge and mastering skills and those who have not. It tends to create greater inequality between those who work and those who do not.}
\BookParagraph{他的意思是，只要学点儿什么干点儿什么，永远比什么也不学什么也不干天天打「王者荣耀」强很多。}
\BookParagraph{如果你还是觉得自己的学业和工作一时半会儿没有什么用处，那就想想 Frank Woolsworth —— 他当年「初创」了公司，结果开的第一家商店很快就倒闭了。}
\BookParagraph{放心吧，只要多运动增强免疫力，人一辈子有的是时间。}
\BookSubtitle{三}
\BookParagraph{当然了，要想「学点儿什么干点儿什么」将来有用，你学的东西一定不能是别人强加给你的，得是自己感兴趣的。如果是别人强加的，换了谁都很难做到长期「勤奋努力」—— 你瞧，连李开复都明白要「追随你的心」。}
\BookParagraph{如果你想来想去，除了打「王者荣耀」和看「乘风破浪的姐姐」，怎么也想不出还对其他什么东西感兴趣，那怎么办？}
\BookParagraph{我建议，竭尽心力，自己设计制作一样需要一两个月时间才能做好的玩意儿。}
\BookParagraph{一两个月时间不算太长，但也不太可能一蹴而就。}
\BookParagraph{你提到正在考虑读研的可能性。读研的过程虽然千差万别，但是都有一个共同的特点：要毕业就得写几篇论文。}
\BookParagraph{虽然写出来的论文本身多半儿没什么用处，然而一篇论文就像一座雕塑，每一笔每一划都是自己设计制作出来的，都是自己的呕心沥血之作。在这个过程之中得到的「看不见摸不着」的技能，意义远远超过论文本身。}
\BookParagraph{如果你对写论文不感兴趣，那就拍一部微电影，上线一家淘宝店，写一本书，开一个线上直播，修好一辆没人要的旧车，或者实现一个网站。}
\BookParagraph{注意，去云南大理旅行或者种一株玫瑰送给男朋友不能算 —— 花点儿钱或者浇点儿水就行，太简单了。}
\BookParagraph{比如我这两个月，就竭尽心力，从头学起，花了所有的业余时间和所有周末，实现了一个全新的网站。}
\BookParagraph{其实这跟跑马拉松差不多 —— 虽然没什么用，过程也很艰苦，但大功告成之后，我看着网站背后的每一行代码，感觉每一个字节都在我的指挥下跳动，觉得非常过瘾，感到了一种「更上一层楼」般的快乐。}
\BookDate{2020 年 6 月，香港}
\BookArticle{寄语}{寄语}
\BookParagraph{清华计算机系校友会前一阵子跟我联系，说「目前系友会准备做一期招生季的宣传推送，将以 “酒井伉俪” 的形式推出」，希望我和郭芳可以提供一小段各自取得的成就，一张照片，还有一段对即将报考计算机系新生的祝福和期望。}
\BookParagraph{我大致想了想，还真的没有什么「成就」可以拿来吹的。MIT 教授 Bob Gallager 曾经说过，学术界真正的成就，不是靠论文的数目，而是看最突出的一两篇论文是什么水准。没办法，只得不要脸地再次把外公搬出来「捧场」:}
\BookParagraph{「李葆春教授家学渊源，外祖父是中国著名语言学家、曾任清华大学和北京大学教授的王力先生。李葆春教授 28 岁出任加拿大多伦多大学助理教授，33 岁晋升为终身教授，36 岁晋升为正教授，43 岁评为 IEEE Fellow。他在网络、云计算、安全和隐私保护、多媒体系统等领域发表了 380 多篇论文，H-index 达到 80，在世界华人学者中名列前茅。」}
\BookParagraph{郭芳帮忙找出了一张九五年五月大学毕业时在北京香山拍的照片。到了要写「祝福和期望」这段话时我颇踌躇了一会儿，因为很容易写成口号式的空话。虽然我觉得不会真有「即将报考计算机系新生」看到 —— 但是万一真有的话，他们看到的将是以下这段「寄语」：}
\BookQuote{当代著名作家王蒙曾经说过：「不怕暴露自己的缺点，乃至敢于自嘲，意味着清醒更意味着自信，意味着活泼更意味着真诚。」有了这种自嘲式的自信，认识到自己的长处和短处，就可以不人云亦云，不随波逐流，基于事实形成自己独到的见解。英文里有一个单词叫「contrarian」，大致意思是「逆流而上」。我想做一个「contrarian」，一个逆流而上的人。}
\BookDate{2020 年 7 月，香港}
\BookArticle{忙}{忙}
\BookQuote{李老师，昨天跟一个小学同学见面，很有感触，想跟您聊聊。}
\BookQuote{她家境不错，学业认真，人品也很好，一路顺利进入一个国内 985，想着毕业就出国，但毕业前夕因为家庭变故，她「忍痛」留在家乡一个国企工作至今。毕业这两年，每次春节回家同学聚会，她不是在学习就是在考证，都不来。我觉得她太不轻松了。很多人轻而易举、云淡风轻的事情，在这位同学这里，似乎是不得了的高山，而她会特别努力，去攀登越过，再追求一座更高的山，更难的问题。}
\BookQuote{半年过去，这位同学从地方国企总部，跳到了一个分部。现在她每天的生活很忙，虽然每天焦头烂额一度抑郁，但因为现在的工作是随时有突如其来的麻烦、突如其来的新问题要解决。所以虽然很紧张，但也很刺激。之前在国企总部，属于每天做一模一样的工作，毫无变动。所以这种忙碌感，让她觉得很充实，并认为充实的日子让人内心很健康，感觉每天都在释放自己。}
\BookQuote{不知道李老师怎么看待「忙碌是无所事事者的避难所」这句话。对比之前有大把时间但是很焦虑，现在的她很忙但是很充实。作为朋友我为她高兴，但还是不能理解这种状况。不知道李老师怎么看。（问题来自微博私信）}
\BookSubtitle{一}
\BookParagraph{说起「忙」这个话题，我想到了外公的杂文集《龙虫并雕斋琐语》中的一篇文章，题目就叫《忙》。}
\BookParagraph{外公将「忙」分为三种。青年时代除了读书之外，体验的是恋爱之「忙」：}
\BookQuote{恋爱的青年有闲中之忙，有忙中之闲。所谓闲中之忙，是因为游山玩水，步月赏花成为一种功课，一种手段，你闲也要你闲，你不闲也要你闲。这样的情形，我们可以叫做「忙于装闲」。}
\BookParagraph{而人到了中年，却又该为事业而忙了：}
\BookQuote{恋爱的忙，虽忙不苦；事业的忙，有时候既忙且苦。当然，以一身系天下安危的人，多忙一分，则民众多受一分的德泽；就是为自己而忙，只要忙得有意思，忙得有新花样，忙得顺利，也就高兴去忙。不过，世界上高兴忙的人实在太少了，苦忙的人也实在太多了。}
\BookParagraph{那你可能会问，等到事业有成，是否就可以不忙了呢？在外公的笔下，「忙」中的最高境界是为了应酬而忙：}
\BookQuote{一个人得到了社会所知以后，似乎他的世界就应该被社会所糟蹋。这一种忙，忙得最苦。既不为食，又不为色，只为的是怕得罪人。我们家乡有一句俗话说：「三十又忧名不出，四十又忧名不收。」古人入山唯恐不深，就是「收名」之道；如果你「自嗟名利客，扰扰在人间」，随便怎样苦忙，也只算是自作自受了。}
\BookParagraph{哈哈，如果你现在正在为应酬而忙，那恭喜你事业有成了。}
\BookSubtitle{二}
\BookParagraph{「忙」本身不是问题，但是忙得上了瘾就不妙了。}
\BookParagraph{什么叫忙得上了瘾呢？}
\BookParagraph{有的时候一个人实在太忙了，忙得稍微闲下来一点儿就坐立不安；忙得而没有时间把任何一件小事儿做得完美、做得精益求精。}
\BookParagraph{而如果做点儿事儿一直捉襟见肘，无法静下心来精心完成，最后很可能功亏一篑，一事无成。}
\BookParagraph{比如，一位教授接受了邀请，需要在一次会议上做一个报告。而因为他太忙，没有时间准备报告的幻灯片和演讲稿，只得责成其学生帮他准备，他上台念一下就可以了。一次两次这么做虽然不雅，也无可厚非；但是如果常年如此，这位教授就会渐渐失去了独立思考和写作的能力，脱离了真实的环境。以致很可能演讲时经常望文生义，无法应对现场针锋相对的提问，「知其然不知其所以然」。那这位教授有没有可能自己准备演讲稿、甚至脱稿演讲呢？其实如果他可以回绝一些为了应酬的「忙」，是完全可以做到「自力更生」的。}
\BookParagraph{所以说啊，稍微忙点儿虽然挺好，但忙得上了瘾不提倡。有的人把自己折腾得太忙了，是因为他们对将来的期望有点儿太高。绝大多数人没有马斯克的天分和马云的机遇，即使跟他们一样忙，天天想着跟他们一样创业，也很难「改变世界」。也有些人太忙，是因为过于重视别人对他的看法，怕得罪人。别人都忙着加班，他虽然觉得加班很累人，也跟着加班；别人都忙着玩创业，他虽然没有创业的灵感，也跟着玩创业。当他们忙着「释放自己」、忙得上了瘾时，日子疾驰而过，会错过很多美丽得令人心动的瞬间，错过了就再也回不来了。}
\BookParagraph{「忙」很像看电视。看得时候很开心很充实，但是看得太多看上了瘾，就逐渐忘了还可以在街边杀一盘象棋，还可以抱一把吉他自弹自唱，还可以什么都不干，任由光阴从眼前流逝。}
\BookSubtitle{三}
\BookParagraph{一个多月前，我正收拾好行装准备赶往机场飞香港的途中，不意间收到了一条微信。}
\BookParagraph{是一位很少联系的朋友发来的一句简短的问候：「你们还好吗？」}
\BookParagraph{我秒回：「还挺好的，正好今天早上要飞香港。」}
\BookParagraph{对方没有再回信。}
\BookParagraph{我简单想了想，这一天好像不是端午节，也不是双十一不是五二〇不是六一八不是七夕不是任何网红或传统的节日。}
\BookParagraph{也许是我的结交面儿太窄，我已经不记得上次有朋友还没到节假日就给我发来简单的问候是什么时候了。}
\BookParagraph{也许是因为我们都太忙，忙得没有时间也没有心情，忙得很难忽然没来由地就想起远方的朋友来。}
\BookParagraph{那天早晨的这样一句简单的问候，让我怦然心动。}
\BookDate{2020 年 8 月，香港}
\BookArticle{开关}{开关}
\BookParagraph{今天一大早起来，走路去一家牙医专科诊所拔牙。}
\BookParagraph{正常的拔牙其实比较简单，普通牙医就可以做。我这颗牙算是个反常的例外，我自己的牙医对着 X 光片看了半响，说离鼻腔实在太近，还是专科医生比较保险。我于是上周约了门诊去专科医生咨询。他看了看，说只要答应他拔牙以前八小时不吃东西，前两小时不喝水，「then it’s going to be a walk in the park.」我问他会疼吗？他笑着说：「Is a walk in the park painful?」}
\BookParagraph{后来跟护士聊了一会儿才明白，医生说的是要静脉点滴全身麻醉。护士说，这是一种介于局部麻醉和全身麻醉之间的麻醉方式，麻醉过程中还可以下意识地听医生和护士的指示行动。比如护士要是说「把头向左转一点儿」，人就会下意识地转头。}
\BookParagraph{我听着像天方夜谭，觉得好神奇啊。不过万一英文听力不过关，下意识里听不懂英文，就不知道会怎样了。}
\BookParagraph{今早在诊所帮我准备静脉点滴的是另外一位护士。她一边给我胸前加入几个测心率的传感器，一边夸我心率低，问我是不是经常运动。我开心地说是啊，每天早上跑步。然后跟她聊了几句，夸她的仪器真先进。她说这些仪器先进啥，你还没去过 ICU 呢 —— 接着给我讲她在 ICU 帮医生装心率调节器（PaceMaker）的经历。讲了几句她说：「我一会儿开始给你加一点儿药量。」}
\BookParagraph{这是我现在记得的最后一句话。}
\BookParagraph{醒来时已经在休息室了，医生没有见到，护士也不见了。我回家吃了止疼药想，与其说是「walk in the park」，不如说是「sleep in the park」，醒来时没有梦。}
\BookParagraph{不过如果后来一直没有醒来呢？}
\BookParagraph{乔布斯的传记作者 Walter Issacson 为了写这本传记，多次采访了乔布斯。Issacson 还记得最后的一段采访：}
\BookQuote{I remember sitting in his backyard in his garden one day and he started talking about God. He said, “Sometimes I believe in God, sometimes I don't. I think it's 50-50 maybe. But ever since I've had cancer, I've been thinking about it more. And I find myself believing a bit more. I kind of—maybe it's ’cause I want to believe in an afterlife. That when you die, it doesn't just all disappear. The wisdom you've accumulated. Somehow it lives on, but sometimes I think it's just like an on-off switch. Click and you're gone.” He said—paused again, and he said, “And that's why I don't like putting on-off switches on Apple devices.”}
\BookParagraph{Click—and you are gone.}
\BookParagraph{还好，我今天只是重启了一下，不是关机。}
\BookDate{2021 年 10 月，多伦多}
\BookArticle{知识青年}{知识青年}
\BookQuote{你好李老师，我非常喜欢你。}
\BookQuote{今天北大一教授说，青年失业率已经逼近百分之二十。想向你请教一下，如何看待青年的就业问题？希望你的这篇文章，能从历史或者宏观的角度，讲讲在这样的环境下，青年人如何平复找不到工作的心态。谢谢李老师，希望能得到你的回复。（问题来自微博私信）}
\BookParagraph{「能从历史或者宏观的角度」—— 感觉你这问题怎么有点儿像高考作文题啊。}
\BookSubtitle{一}
\BookParagraph{1933 年，美国正处于大萧条时期，失业率居高不下。就职仅仅两个月的罗斯福总统，向美国国会提出了一项全新的动议：由政府出资，雇佣未婚年轻男子参加一个叫「Civilian Conservation Corps」的组织，简称「CCC」。这个组织的活动保证不会与任何已有的成熟行业竞争，而将从事「林业、防止水土流失、防洪等项目」。}
\BookParagraph{虽然罗斯福希望 CCC 的工作对美国有用，但这其实不是重点。就在一年前，罗斯福的前任胡佛总统（Herbert Hoover）被迫部署军队来驱散第一次世界大战中失业的退伍军人在华盛顿的游行。罗斯福心里很明白，数以百万计的失业者是多么危险。}
\BookParagraph{国会立马就批准了罗斯福的这项动议。在三个月内，有 25 万年轻男子参加了 CCC，住在全国各地营地的帐篷里。只有军方才有足够的能力来管理这么多人，所以罗斯福让麦克阿瑟将军来全权负责。}
\BookParagraph{在这以后的九年里，CCC 彻底改变了美国美丽富饶的农村。大约三百万年轻男子参加过 CCC。他们在野外坚持不懈地与森林大火作斗争，种植了三十多亿棵树，修建了十二万英里的公路和一万三千英里的小道，全新铺设了八万九千英里的电话线，在数百个州立公园修建了各种小房子、鱼类孵化场和野生动物保护区。}
\BookParagraph{这三百万人每周六天早起锻炼，过着有点儿像童子军也有点儿像军营的规律生活。每位员工每月收入大约 30 美元，并一律必须将其中 25 美元寄回家。虽然大多数人刚开始工作时营养不良，但由于这样军营一般的规律生活，第一年平均体重增加了十二磅。}
\BookParagraph{CCC 被证明是大萧条时期最受欢迎的政府计划。}
\BookSubtitle{二}
\BookParagraph{我第一次读到这段「人定胜天」的历史时，竟然有点儿悠然神往。}
\BookParagraph{罗斯福的这项计划简直就是教科书般的神来之笔：几百万青年人不仅有了工作，还在军营般的规律生活中提高了体力和技能，培养了组织性纪律性，将来的一生受用无穷。}
\BookParagraph{这简直就是所谓「躺平」的镜像反面啊。}
\BookParagraph{我无端地认为，在当年那经济大萧条的年代，每一个人可以做到的事儿，相对大环境而言，是极其渺小的。但是人多了铁锹多了，要铺的路要修的桥，早晚有一天可以通车。}
\BookParagraph{所以，与其天天关心时局，不如先从自己做起，从小事做起，从六点起床锻炼做起。与其「躺平」—— 这个流行词的具体含义其实我到现在也没有完全明白 —— 不如问问自己到底具备什么别人不具备的学识、才干和能力，可以让市场的天平早晚有一天会向自己倾斜。}
\BookParagraph{说来说去，找工作其实和找对象差不多，都是一个概率问题。如果你是个阳光帅气健谈的小伙子，早睡早起、经常运动，虽然无法保证有喜欢你的姑娘，但是概率自然会高不少。如果你喜爱而且擅长你的专业，不断提高自己的专业水准，你能找到工作的概率自然也会高于常人。}
\BookParagraph{「要创造我们的幸福，全靠我们自己。」有首歌脍炙人口，其歌词这样写道。}
\BookSubtitle{三}
\BookParagraph{我二十多岁的日子，大部分是在美国一个美丽的农村度过的。}
\BookParagraph{有一年我参加了学校一个叫「中国学生学者联谊会」—— 简称「CSSA」—— 的组织，连夜和当时 CSSA 的主席褚明宇同学一起撰写排版一本叫「新生指南」的小册子。}
\BookParagraph{小册子里面的内容不太记得了，不外乎是罗列一些办理入学手续还有租房买菜什么的注意事项。}
\BookParagraph{只记得凌晨时分，褚明宇为了这本《新生指南》的序言怎么写，沉吟了许久。}
\BookParagraph{天快亮了，他终于大笔一挥，开门见山，一语石破天惊：}
\BookParagraph{「伟大领袖毛主席教导我们：知识青年到农村去。」}
\BookDate{2022 年 5 月，多伦多}
\BookArticle{导师}{导师}
\BookParagraph{我 22 年前在伊利诺伊大学的研究生导师 Klara 最近准备参与评选一个 Women in Science 的国际奖项，想起了我，邀请我写一封推荐信。}
\BookParagraph{写这封信的时候，我又想起她修改论文时龙飞凤舞的字迹，想起她晚上 11 点多发的那些「come up to my office」的邮件。一个优秀的导师，应该既有大局观，又能在细节上下功夫；既能自己从早到晚努力工作，又能对待学生「像春天般温暖」。}
\BookParagraph{Klara 就是这样的导师。}
\BookDate{2022 年 5 月，多伦多}
\BookParagraph{——}
\BookParagraph{May 29, 2022}
\BookParagraph{To whom it may concern:}
\BookParagraph{It is my great pleasure to write this letter in strong support for Professor Klara Nahrstedt’s nomination for the L’Oreal-UNESCO For Women in Science International Award. Klara has been my thesis advisor at Illinois throughout my Master’s and Ph.D. studies, for a period of four years from 1996 to 2000. Since 2000 and immediately following my graduation, I have been teaching at the Department of Electrical and Computer Engineering, University of Toronto, where I am currently a Professor, IEEE Fellow, and the Bell Canada Endowed Chair in Computer Engineering. In my professional career, I have published over 450 co-authored research papers, won several awards, and achieved an H-index of 84 according to Google Scholar. With H-index as the metric, I am currently ranked \#12 in Canada in the field of computer science, according to \href{https://research.com/u/baochun-li}{research.com}.}
\BookParagraph{At Illinois, I have benefited tremendously from Klara’s supervision, advice, and stewardship. Klara’s influence on my student experience and my ensuing professional career was so significant that it simply could not be overstated. With her mentoring and leadership, I have grown from a new graduate student with very little experience on research to a mature researcher winning a best paper award, all within a short time span of four years. Suffice it to say, without Klara’s supervision at Illinois, it is impossible for me to start my academic career at Toronto in the first place, let alone remaining productive throughout my career. She has had a life-changing impact on my academic and professional career.}
\BookParagraph{I believe that an excellent researcher should be both visionary and paying attention to detail. Klara taught me both with her mentorship. With respect to research vision, she consistently demonstrated an excellent grasp of the “big picture,” especially when it comes to new, emerging, and interdisciplinary areas of research. I have benefited tremendously from Klara’s vision as I worked on my PhD thesis work, and from our numerous meetings and discussions in her office. Ultimately, one of our papers won the Leonard G. Abraham paper award from the IEEE Communications Society in 2000, and was cited more than 500 times since its publication in an IEEE journal.}
\BookParagraph{Klara also paid a lot of attention to detail when we were working together on research. She enjoyed revising my research papers by handwriting. As she discovered problems from minute typographical mistakes to structural and clarity concerns, her handwritten communication has always inspired me and refined my research skills effectively. She was particularly demanding with respect to timely and effective communication, both in more formal talks and over more informal email messages. She cared about interpersonal and communication skills when her students, including myself, attended international research conferences. But more importantly, she was as nice as anyone to her students in the research group, caring passionately and deeply about their career choices when their graduation drew close.}
\BookParagraph{I have read a number of bitter stories after I started teaching at Toronto, told by graduate students and their supervisors in other universities. Some students developed negative feelings towards research and towards academia in general, some even poked fun at the harshness towards getting a Ph.D. degree. I can honestly, and proudly, say that I experienced none of these in my own graduate studies. My own experience felt more like sailing smoothly in the Caribbean, rather than stalling in the midst of a winter snowstorm. I was always fully charged, upbeat, and energetic to attain my research goals. In hindsight, though there were turbulences along the way as some of my research papers were rejected even after putting in the hard work, the supporting role that Klara played was so reaffirming and so warm, I never felt discouraged. This continued well into my academic career after I started at Toronto, as I inherited the same philosophies when mentoring my own students, over sixteen of which started their own academic careers at universities around the world. I guess I was lucky from the very beginning, and I hope that my experiences at Illinois are convincing enough to offer my strongest support for Klara’s nomination for the L’Oreal-UNESCO For Women in Science International Award.}
\BookParagraph{Sincerely yours,}
\BookLeftParagraph{Baochun Li, Ph.D. (Illinois, 2000)}
\BookLeftParagraph{University of Toronto}
\BookArticle{欢庆六一}{欢庆六一}
\BookParagraph{好几个朋友给我转发了一篇题为《\href{https://mp.weixin.qq.com/s/ZgyXRMuNeFQrDWsB6gnxSg}{他被邓小平摸了后脑勺，无数人的命运由此改变}》的公众号文章。}
\BookParagraph{文章开篇讲了一个著名的故事。邓小平当年摸着李劲同学的头，说了一句名言：「计算机的普及要从娃娃抓起。」继而提到景山学校有「两个戴红领巾的男孩，一个对着屏幕做出指点江山状，一个操作着电脑打字。屏幕上是两行字：我们爱科学、欢庆六一。」这两个孩子，后来一个上了北大，还有一个上了清华。}
\BookImage{images/childrens-day.jpg}{0.7}{}
\BookParagraph{哈哈，那个操作着电脑打字的孩子就是我。那是 1983 年的六一儿童节，当时我十一岁。}
\BookParagraph{没想到的是，四十年之后，我竟然还在编程序，只不过是用快一些的计算机和不一样的编程语言而已。由此看来，我大概是所有该文提到的人里「混」得最差的一个。}
\BookParagraph{不过这篇文章有一个重要的细节没有提到。}
\BookParagraph{同样是 1983 年，邓小平十一那天，曾经给北京景山学校题词：「教育要面向现代化，面向世界，面向未来。」}
\BookParagraph{这句早已耳熟能详的话，现在回想起来，竟然还一点儿也不过时 ——}
\BookParagraph{教育要面向现代化，面向世界，面向未来。}
\BookDate{2022 年 5 月，多伦多}
\BookArticle{时间}{时间}
\BookParagraph{几个月前，有个朋友曾经提到回头会在微博上问我一个有关「时间管理」的问题。直到现在这个问题也没有等来，我想要不然先聊起来，等问题来了再行补充。}
\BookParagraph{我一直觉得「时间管理」这个词儿不是很贴切：时间不是财产，无论咱是花还是不花，都会自然而然地流逝，实在是没办法「管理」。无论是「只争朝夕，力争上游」（也就是所谓「开卷」），还是「终日昏昏醉梦间」（也就是所谓「躺平」），都仅仅是不同人的选择和生活方式，很难说孰优孰劣，也很难评价「时间管理」的水平。}
\BookParagraph{不过话又说回来，我当了二十多年的教授，确实经常有忙得时间不够用的时候。从这个角度来看，如何更有效地利用时间，还真是个亟待解决的难题。其解法众说纷纭，有人甚至设计了极其纷繁复杂的所谓「get things done」的系统，甚至看个电影吃个饭都纳入了「时间管理」的范畴。出于好奇，我还真尝试过一两次，结果发现需要花时间和精力去「管理」时间，还带来了不少不必要的焦虑，有些得不偿失。}
\BookParagraph{我的博士是在伊利诺伊大学香槟分校读的。当年我读书的时候，有一位德高望重的老教授，名字叫 C. L. Liu。这位刘教授的重要贡献之一是一篇 1973 年的论文（Liu and Layland），首次提出 rate monotonic scheduling，算法简单而有效，成为实时操作系统中调度算法设计的里程碑。}
\BookParagraph{我发现，rate monotonic scheduling 中有一些设计理念，同样可以用来「管理」时间。}
\BookParagraph{这个算法有两个基本假设和一个核心设计。}
\BookParagraph{一个假设是同一个时间只能做一件事儿，比如写论文或者编程的时候不玩儿手机。这个假设我觉得挺重要的，一段时间只做一件事儿，才能更加专注，质量更高。}
\BookParagraph{第二个假设是优先级高的任务来了立即开始，优先级低的任务先暂停执行，比如男（女）朋友要叫个外卖，论文写作立即暂停，无论是否马上可以写完。}
\BookParagraph{这个核心设计，就是采用固定的调度优先级（static priority）。这是什么意思呢？我们平常决定先做什么，一般是看这件事儿是否紧急，截止时间早的先完成。Rate monotonic scheduling 也是用的这一原则，但是它是根据不同任务的紧急程度先确定优先级，然后优先级高的先做，不再看截止时间。比如上级领导交代的任务，优先级相对更高，无论什么情况下，任务来了马上开始。}
\BookParagraph{Liu and Layland 这篇著名论文证明了，用这样的算法设计，对时间的利用率无法达到百分之百。即使任务非常多，最多可以到 ln 2 (69\%)。这就是说，有百分之三十的时间浪费了。看着挺差劲儿的，但是因为浪费的时间可以用来刷抖音小红书，其实可以算是生活工作三七开吧。}
\BookParagraph{以「浪费」百分之三十时间为代价，优先级固定，高优先级任务来了立即执行，老板的邮件和微信秒回，这就是我的「时间管理」。}
\BookParagraph{刘教授一辈子待人以诚，著作等身，培养了二十六位博士研究生。离开伊利诺伊大学以后，曾任新竹清华大学校长，两年前去世，享年八十六岁。}
\BookDate{2022 年 8 月，多伦多}
\BookArticle{科技论文之美}{科技论文之美}
\BookParagraph{据说，有一个简单的判据可以知道自己什么时候博士快毕业了 —— 刚开始做研究的时候，论文看不懂，会觉得自己怎么这么无知；等快毕业时再看不懂，会觉得这篇论文怎么写得这么晦涩难懂，难以理解。这说明对待晦涩难懂的论文，自己的心态已经有了飞跃，开始有了一种「一览众山小」的感觉。}
\BookParagraph{乔布斯曾经说过，他理想中的产品，是科技和艺术的交融。我们写一篇科技论文，有没有这样一种可能，不仅通俗易懂，还能给读者一种艺术之美感呢？}
\BookParagraph{您一定要说，那肯定是要多看「顶会」论文啊，摘抄一些好词好句，多学习，多模仿，不就可以越写越有感觉了吗？}
\BookParagraph{我觉得吧，「顶会」发表的论文不一定是写得漂亮的优秀论文；而那些给人一种艺术之美感的论文，却不一定会在「顶会」发表。学习和模仿什么论文，其实大有讲究。}
\BookParagraph{我们来看一篇六十年前 Gale and Shapley 的经典论文 “\href{https://baochun.ca/images/Gale-Shapley.pdf}{College Admissions and the Stability of Marriage}.” 短短七页的论文，刚开场就极其直白，颇有一种「环滁皆山也」的气势：}
\BookQuote{Introduction. The problem with which we shall be concerned relates to the following typical situation: A college is considering a set of n applicants of which it can admit a quota of only q.}
\BookParagraph{这么简单的一句话，为什么我竟然极为赏识呢？}
\BookParagraph{首先，这句话开门见山，目的是要开始阐述这篇论文准备解决的是什么问题。我看过很多近来的论文，往往第一页看完，一直在介绍背景、文献和算法，却并没有真正说清楚这篇论文到底想解决什么问题。其次，这句话语言极其简单，而简单明了的语言，有助于读者轻而易举地看懂论文的思路，看完有一种豁然开朗的感觉。}
\BookParagraph{其实不仅仅是开场，这篇论文通篇都非常易读易懂。即使读者开始对这个研究方向不感兴趣，都会被一种简单的好奇心驱动，一点儿点儿通读全文。}
\BookParagraph{我们再来看看论文的结尾：}
\BookQuote{What, then, to raise the question once more, is mathematics? The answer, it appears, is that any argument which is carried out with sufficient precision is mathematical, and the reason that your friends and ours cannot understand mathematics is not because they have no head for figures, but because they are unable to achieve the degree of concentration required to follow a moderately involved sequence of inferences.}
\BookParagraph{您瞧，是不是有一种气势磅礴、更上一层楼的感觉？作者 —— Shapley 2012 年获得了诺贝尔经济学奖 —— 当年就告诉我们，一篇论文如同一盘围棋，只要每一步推理计算精确，表达清晰，就可以说是数学；而这样的数学很多人看不懂，并非由于数学本身多么艰深，而只不过是因为他们没有集中注意力罢了。}
\BookParagraph{这段似乎是与主题完全无关的结尾，提到的问题却是现在很多论文都普遍存在的。这些论文似乎包罗万象，不仅有大量的公式和证明，还有通篇的各种算法。但是如果仔细读下来，不但没有体会到艺术之美，而且还不知所云，由于算法描述忽略了一部分细节而无法复现（reproducibility）。究其原因，在于作者并没有把一步一步的推理过程和算法步骤通过数学和个例说清楚，让读者觉得云山雾罩，知其然而不知其所以然。}
\BookParagraph{说到这里，我们应该已经可以发现，其实科技论文之美，不是因为其中用到了什么好词好句，而是因为逻辑清晰，条理鲜明，甚至不需要作者提供源程序，仅仅依靠集中注意力通读论文本身的内容，就可以完美复现其提出的算法。这有点儿像黑白照片，虽平淡无奇，同样令人怦然心动。我们甚至可以说，一篇论文的语言越是简单，甚至简单到一位不同领域的学者都可以看懂，就越有一种令人拍案叫绝的美感，离科技与艺术的交叉路口就越近。}
\BookParagraph{看到这里您可能会问，写得太简单、太容易懂，那审稿人难道不会觉得论文水平太小儿科，没有内涵和新意吗？}
\BookParagraph{当然有这种可能。不过，如果因为怕语言过于简单而把论文写得艰深难懂，不但有点儿本末倒置，审稿人也不一定会因为看不懂而给出很高的评价，是吧？而且，我一直坚信，一篇论文真正的内在价值，并不是一两位审稿人的负面评价就可以改变的。我们写论文，不是为审稿人而写，而是为将来的读者写的。}
\BookParagraph{What, then, to raise the question once more, is a beautiful paper?}
\BookParagraph{如果一篇科技论文设计新颖，巧夺天工，算法步骤清晰，可以直接复现，开场和结论简洁明快，就是一篇「美」的论文。但凡美的事物，都让人有一种流连忘返的感觉，像一场精彩绝伦的电影，看完一遍，还想再看第二遍。只不过与电影艺术不同，论文的美在于它开门见山式的简单，在于它步步为营式的精密，洗尽铅华，是一种超凡脱俗的美。}
\BookParagraph{这样的论文，即使没有被「顶会」录用，那又何妨？}
\BookParagraph{如果博士快毕业了，咱们还可以说，这些注意力无法集中的审稿人，水平也未免太低啦。}
\BookParagraph{就如同 Gale and Shapley 大作中的最后一句话：}
\BookParagraph{For them the foregoing may serve as a useful illustration.}
\BookDate{2023 年 12 月，多伦多飞往东京的航班上}
\BookArticle{INFOCOM Achievement Award}{INFOCOM Achievement Award}
\BookParagraph{新年刚过就收到了一份特大喜讯：我荣获今年的 INFOCOM Achievement Award。}
\BookParagraph{这一殊荣有点儿不太容易翻译成中文。理论上来说，这是一项终身成就奖，奖励在 IEEE INFOCOM 这一计算机网络领域的「顶会」中取得了突出成就的学者。但是它的英文却是「Achievement Award」—— 没有「Lifetime」，简单而低调。}
\BookParagraph{和年迈的父亲分享了这一喜讯。电话线那边，父亲声音洪亮，显得非常开心：「嗯，那当然是很厉害的啦！」}
\BookParagraph{二十多年前的导师 Klara 也发来了贺电。邮件的落款是「Your proud academic mother.」}
\BookParagraph{会议请我提供一小段研究简历，便于在官网发布。我慢慢地磨洋工时，忽而想起二十年前。那一年的 INFOCOM 一如既往是七月底截止，我以「洪荒之力」连续奋战了三个月，投了六篇论文。发榜那天晚上，邮件一一打开，结果全军覆没，黯然神伤。}
\BookParagraph{那年的会议在香港会展中心隆重召开，我虽然没有论文，还是去参加了。第二天晚上的酒会上，有一位刚刚认识的参会者礼节性地问我：「你的论文是什么时候讲？」我回答：「我没有论文，是来跟大家学习的。」}
\BookParagraph{觥筹交错，谈笑之间，二十年一晃而过。我想，将来封笔江湖、故国神游之时，江湖上的传说里，至少不会觉得我写不出论文了吧。}
\BookDate{2024 年 2 月，多伦多}
\BookArticle{随遇而安}{随遇而安}
\BookQuote{职业倦怠怎么办呀？工作太繁重，快两年没发奖金，自我察觉有些职业倦怠，下班了之后想去做一些提升自己的事，但我感觉休息时间不够我补充工作造成的消耗，已经没精力和时间去提升自己了。想问问大家有没有啥方法可以参考一下？}
\BookParagraph{（问题来自「葆春和他的朋友们」）}
\BookParagraph{我觉得吧，无论是学习还是工作，都不可能一直顺风顺水，你说呢？}
\BookParagraph{比如工作吧，那肯定经常会既忙碌又无聊，这是没有办法的事儿。比如我现在已经是凌晨时分，才刚刚做完一天的工作。今天的工作中，大多数时间不是在开会，就是在处理各种和人事招聘有关的数据、文档和表格。事情虽然繁杂 —— 对涨工资也没什么帮助 —— 但是总得有人去做这些工作，对吧？领导把这些任务交给我，我觉得也是对我的一种信任。即使涨不了工资，如果我可以圆满地完成这些工作，也是对这种无形的信任的一种报答。}
\BookParagraph{我愿意相信，报应 —— 也就是英文中的「karma」—— 是真切地存在的。今天我为领导为单位不计得失地努力工作，也许明天就真地涨了工资，或者天上掉馅饼儿，收获了一份重大喜讯。所以，为了积攒「karma」，当别人在玩儿或者在休息时，我愿意起早贪黑不分节假日地工作，甚至可能会忙到没有时间去「提升」自己。}
\BookParagraph{可为什么一定要「提升」自己呢？比起「好好学习，天天向上」，我更喜欢一种随遇而安的生活。相比现在的时尚词汇「躺平」，我觉得「随遇而安」体现的心态中，少了一份选择，多了一点儿平和。空中飞人也好，宅在家里也罢，有兴趣的事儿可以多做一点儿，没意思的事呢，则可以少做一点儿，「做一天和尚撞一天钟」，「船到桥头自然直」嘛。说不定哪天发现自己不自觉地「提升」了，那也未可知呢。}
\BookParagraph{比如刚才，我看到你的问题，小有感触，就随波逐流地写了以上的文字。就这样多回答几个问题，说不定还能提升我的写作水平呢。}
\BookDate{2024 年 3 月，多伦多}
\BookArticle{EIC Fellow}{EIC Fellow}
\BookParagraph{周末参加了加拿大工程研究院（Engineering Institute of Canada）的院士颁奖典礼。多伦多大学工程学院的院长 Christopher Yip 亲自参加了典礼，工程研究院的主席亲自颁奖，十分隆重。}
\BookParagraph{颁奖过程中的幻灯片需要一张小时候的照片，我选的是一张五六岁时在北京新大北照相馆的留影。还依稀记得当年我住在北大外公家，妈妈出差来北京，带着我坐公共汽车到新街口附近去照的，当时牙还没有长齐。}
\BookParagraph{系里一位老教授曾经跟我说，荣誉和奖励其实没有什么意义。我有点儿好奇，问他为什么。他说原因很简单，无论得的是什么荣誉，几百年之后不会再有人记得了。我一想其实也有点儿道理 —— 几百年后，可能只有牛顿爱因斯坦这样的科学家还能青史留名吧。}
\BookParagraph{十九世纪英国作家和艺术家 John Ruskin 曾经说：「The highest reward for a person's work is not what they get for it, but what they become because of it.」或许正如他所说，对辛勤工作最有意义的回报，是在追求完美的过程中收获的那种「山重水复疑无路，柳暗花明又一村」的开心劲儿吧。}
\BookDate{2024 年 4 月，多伦多}
\BookArticle{智力测验}{智力测验}
\BookParagraph{今天和一位办公室就在左近的老教授在走廊中偶遇。}
\BookParagraph{老教授不但著作等身，是其研究领域世界级的专家，近年来还一直不断地在发表新的论文，非常活跃。}
\BookParagraph{他说：哎，下午两三点钟看到一个学生在走廊这头等人。我问他等谁，他说在等李教授。我跟他说，李教授的办公室在走廊那头儿啊。}
\BookParagraph{我想了一下，跟他说还真是，有个学生第一次来找我，找到我学生的实验室去了，还给我发了一封邮件。后来到我网站才找到了我的办公室的门牌号。没想到还等错了地方。}
\BookParagraph{我自己办公室和实验室的门牌号码，都在我网站中可以找到。}
\BookParagraph{他笑着开玩笑说，也许这是一个智力测验 —— 一次找不到办公室的学生不能收。}
\BookParagraph{我想，其实这可能和智力无关，而和是否重视生活中的细节有关。俗话说「魔鬼都在细节之中」—— 往往一件事儿是否可以高效率高质量的完成，也和是否重视各种细节相关。}
\BookParagraph{想起来前两天。另外一个学生来办公室找我，也是等了一整天都没看到人。原因很简单：我的办公室门外的走廊还有一个门，这个门平时有可能是锁着的。}
\BookParagraph{但是这个学生其实可以给我发短信或者拨打我办公室的电话询问一下 —— 如果我正好在办公室，我完全可以帮她开门。}
\BookParagraph{我办公室的电话也可以在我网站找到。}
\BookParagraph{看来，不仅仅需要注重细节，还要有一定的临场应变和敢于交流的能力。}
\BookParagraph{注重细节、临场应变、敢于主动交流 —— 这些也许就是著作等身的老教授说的「智力测验」吧。}
\BookDate{2024 年 9 月，多伦多}
\BookArticle{向上移动}{向上移动}
\BookSubtitle{一}
\BookParagraph{我有一个学生，平常爱用大模型。前一阵子写了个给几位教授的邮件，内容也是用大模型先生成，再自己详加修改。想法虽然挺好，可惜大模型犯了一个小错误：它把 DOE 自作主张展开写成「Department of Education (DOE)」 — 其实也不是完全没道理，学校的邮件嘛 — 但是在多伦多大学「DOE」代表「Departmental Oral Examination」。由于这是学校内部的简称，大模型不太可能知道。}
\BookParagraph{学生没检查出来，等我收到抄送才看到，已经晚了。}
\BookParagraph{所以啊，人工智能和大模型虽然发展迅速，日新月异，怎么利用，还是挺有学问的。}
\BookParagraph{7 月 3 日，我将做客褚明宇的「向上移动」第三季，聊聊如今的人工智能和大模型。}
\BookSubtitle{二}
\BookParagraph{我一向不怎么喜欢特别忽悠的那种所谓「卖课」的广告。}
\BookParagraph{自从 2017 年的第一季，褚明宇的「向上移动」一直不是一个产品，一门课程，而是一个社区。}
\BookParagraph{怎么说呢，无论是第一季的文字，还是后来的直播，这个社区从来没打算过要教点儿什么内容，而是以交流为主。}
\BookParagraph{在现在微信群漫天飞舞的时代，其实真正有点意义的在线社区，我感觉倒是越来越少了。大多数微信群反映的是已有的社会关系 — 同学、同事、校友 — 而非随机生成的全新社区。在人工智能领域中，大多数微信群是「exploitation」，而非「exploration」。}
\BookParagraph{你这时一定会说：这不就是粉丝群嘛，见过见过。}
\BookParagraph{我觉得也不尽然。「粉丝」这个词儿有点儿被滥用了，但其实不可能所有人都完全认同褚明宇的观点。我感觉这个「向上移动」社区更像是由对生活对工作比较爱较真儿，比较乐于问个为什么的朋友们组成的。}
\BookParagraph{对，「朋友」这个词儿可能比「粉丝」更贴切。刚开始是陌生人，时间长了，就慢慢儿成了朋友。}
\BookSubtitle{三}
\BookParagraph{褚明宇的这个社区，还请到了我们上学时的共同好友张戎跟大家交流一场。}
\BookParagraph{张戎是个神人。清华上学时就是十杰候选人，后来学了机械工程，玩过汽车，还去斯坦福读了 MBA。知乎上的文字既有工程师的严谨，又掺杂一点儿北京人的幽默。}
\BookParagraph{他去年写了本书，插图都是他亲手画的。我看完一直想写个书评，由于工作繁忙，到现在都还没动笔。}
\BookParagraph{在这人工智能的大时代，我们更需要褚明宇和张戎这样的朋友：他们一直立场鲜明、论据靠谱、逻辑清晰。}
\BookParagraph{不是所有人都是「博主」、「主播」。有极少数人，你不一定同意，但会经常想起他的观点。只要想得到，你就会去想他们说的是不是正确的，而听过想过，就比完全没听过没想过更有意义。}
\BookParagraph{如果「粉丝」是这么定义的话，那我也算褚明宇和张戎的「粉丝」之一。}
\BookParagraph{下个月线上见。}
\BookDate{2025 年 5 月，多伦多}
\BookArticle{卷}{卷}
\BookParagraph{今天朋友给我发过来一篇王朔的专访，其中有一句话印象颇为深刻：人生的归宿是在墙角刷抖音。}
\BookParagraph{换句话说就是，终于不用再「卷」了。}
\BookParagraph{连王朔都刷抖音了 —— 想起八十年代他的处女作《空中小姐》中的那份张狂、那份柔情，不禁有一种「the end of an era」的伤感。}
\BookParagraph{「卷」这个词儿挺新的。大概比「颜值」还晚几年才出来，可一出来就红遍了大江南北。比如说吧，前一阵子端午节我初次访问南京大学，和好几个新老朋友吃饭聊天，听到最多的词儿就是「卷」。}
\BookParagraph{也许是我一厢情愿吧 —— 我愿意相信，「卷」有时也是个褒义词。}
\BookParagraph{为期一周的访问一晃而过。在校园里骑车、在食堂吃饭、在办公室写代码之余，印象尤为深刻的是有幸结识了几位计算机专业的青年学者，个个著作等身、思路清晰、看问题的视角独特，令人眼前一亮。按照现在流行的说法，都是什么什么「神」。长江后浪推前浪，青出于蓝而胜于蓝 —— 确实太「卷」了。}
\BookParagraph{当然了，江苏省人杰地灵，南京大学又是全国少有的顶级高校之一，这里青年学者的学术水平高，那也是意料之中的事儿。但他们的「卷」，还不仅仅体现在善于发表论文方面，更体现在其演讲、提问、和交流的能力，体现在对人工智能和物联网等新技术发展趋势的细致观察。说来也巧，这里的国家实验室就叫做「软件新技术」，有的教授主持着工业界广泛应用的云计算开源社区，有的则发明了简单到一听就懂、又复杂到需要好几年时间积累而得的巧思奇想，「卷」到了令人匪夷所思的地步。}
\BookParagraph{有趣的是，南京还有一个科举博物馆。我忙里偷闲去参观，发现当代之「卷」和古人之「卷」比起来，完全不可同日而语 —— 古人穿越时空到如今，是绝对的「卷王」。博物馆中收藏了不少当时科举考试中考生的试卷，在八股文的条条框框中照样文采飞扬，蝇头小楷则美不胜收，简直就是文学和艺术集于一体的巅峰之作啊。}
\BookParagraph{遥想几十年前的王朔，那也自然是意气风发，大街小巷的书摊上，随处可见他的文集。「五十步笑百步造句 —— 前五十步笑，百步以后哭」，幽默而不失深意。他当年的「卷」，卷成了家喻户晓的风云人物，卷出了可以传世的作品，卷到了在墙角刷抖音都能刷到自己的时候，人生真的无憾了吧。}
\BookParagraph{南京的一次接风席间，一位青年学者，熟读历朝历代诗词绝句，西瓜汁过了三巡，他站起身来，说：「今天没喝酒，本来当场不一定会有灵感，但是既然恰逢端午，李老师又远道而来，我吟诗一首给大家助兴如何？}
\BookQuote{端午黉门迎李公}
\BookQuote{南雍端阳驻鸿儒，\BookLineBreak{}菰叶葆香润春珠。\BookLineBreak{}千舟竞发深求索，\BookLineBreak{}犹记沧浪清与浊。」}
\BookParagraph{「卷」到如此境界的人生，也是一种可遇而不可求的幸福。}
\BookDate{2025 年 6 月，多伦多}
\BookArticle{代理}{代理}
\BookParagraph{受邀在褚明宇开创的「向上移动」第三季中做有关人工智能的直播，没几天时间了，这几天在构思讲点儿什么。}
\BookParagraph{我一直认为，一场为时一小时左右的演讲，要想质量高，即使是不计时间成本地准备，其实也挺难的。}
\BookParagraph{一小时的时长，讲什么主题，这个主题对应用什么思路展开，围绕这个思路如何组织素材，原始素材之间如何起承转合，都是需要首先解决的问题。}
\BookParagraph{比如，主题和素材之间如何才能动态保持平衡？既不能虽然「干货满满」但却是「只见树木、不见森林」，又不能虽然「高屋建瓴」但却显得太空泛，没有实例佐证。}
\BookParagraph{人工智能这个话题，在这全社会都在「做」人工智能的大时代，我觉得尤其难讲。虽然线上线下，讲人工智能的人多如牛毛，但是既能见到「森林」、又能见到「树木」的演讲，我发现其实凤毛麟角。大多数演讲，我听了觉得没听懂 —— 当然这很可能是我脑子反应不行中间没跟上。}
\BookParagraph{大框架有了，也就仅仅相当于一部电影的编剧完成了剧本，还需导演对具体细节的发挥。比如如何设计演讲，才能既有核心技术内容，又有一点点儿轻松愉悦的调侃？乔布斯当年现场给星巴克打电话的那句「4000 lattes to go, please」，至今记忆犹新。}
\BookParagraph{七月三日晚，如果您正好是褚明宇「向上移动」社区的一员，敬请关注我正在准备的演讲。演讲题目暂定为：「没事儿别瞎玩儿代理！」}
\BookDate{2025 年 6 月，多伦多}
\BookArticle{迷茫}{迷茫}
\BookQuote{李老师好，您是我非常崇敬的老师，同时我也是您的笔迷。我刚刚博士毕业，现在的状态总结为可用「迷茫」来概括，若老师有空，是否能指点迷津，不胜感激。}
\BookQuote{我的家境一般甚至偏下，有两三篇还算拿得出手的论文成果。老师们说我年龄大了，应该找个有钱人结婚生子解决经济问题。也说我应该去海外博后，拿了海优再回国。我在毕业前夕找了互联网薪资还算不错的工作，后来觉得还是想做教职，就退掉了。但其实海优其实难度较大，不一定能申请得上，而且近些年毕业的博士越来越多，在海外呆几年再回国，可能现在有的 offer 也找不到了。}
\BookQuote{而且本来做科研就是件辛苦的事，我有时看到学术圈里的人搞圈子，真的觉得心灰意冷。我不知道应该早早去上班赚钱，还是继续在染缸里坚持下去，做博后做科研找教职，又或者是早点退出这里放弃理想结婚生子赡养父母... 对不起老师，可能我的观察过于片面，思考也有偏见，如果老师有空，能否讲讲您的看法？我将受益终生。（问题来自微博私信）}
\BookSubtitle{一}
\BookParagraph{很多时候，做一个关键的决策，都像是在解决一个优化问题。}
\BookParagraph{简单起见，咱这个优化问题只有一个优化目标，而把其他所有目标都转化为约束条件。}
\BookParagraph{比如，你的问题中，是希望既能收入越多越好，又能有一份和自己的兴趣爱好相近的稳定工作，还能找到称心如意的伴侣结婚生孩子。而你的选择有三种：去工业界工作，去学术界工作，或者找一个优秀的 —— 最好是有钱的 —— 对象结婚。}
\BookParagraph{这里首先有一个问题。最后一个选择，和你前两个选择其实并不矛盾。理论上来说，找对象结婚需要花一些时间，但是即使想找特别优秀或者特别有钱的，也不太可能需要全职。既然如此，咱们可以假设，你的选择有两种：}
\BookParagraph{是去工业界还是去学术界工作？}
\BookSubtitle{二}
\BookParagraph{如果我们根据第一性原理摒除一些海优啊圈子啊关系不大的细节，两者之间的区别在于：工业界工作收入高不少，但是不太稳定，工作本身你可能不太感兴趣。学术界工作呢，收入虽然低，但是相对稳定一些，工作本身也更有可能和自己的兴趣爱好相近。}
\BookParagraph{我们粗线条地大致分析一下：如果你的优化目标很明确地定义为最大化收入，那似乎工业界是首选。而如果你的优化目标是兴趣和学术理想，学术界则成为首选。}
\BookParagraph{当然，这个优化问题的分析还需要加一两个脚注。如果咱们假设学术界你可以顺利拿到编制或终身教职，可以工作三十年。而工业界风险甚高，十年后就会丢了工作。那工业界的税后工资需要是学术界的三倍或以上，才会在优化收入时首选工业界。同理，如果学术界也有因为「非升即走」而「走」的风险，而这个风险还不小，那可能工业界工资只需要是学术界的两倍左右，就可以首选工业界。}
\BookParagraph{所以，咱们化繁为简，如果你想将收入最大化，这个优化问题的答案，完全可以靠计算工业界和学术界每年收入的比例得出。比如，如果某一所大学给了你每年税后收入四十万元的 offer，等同于工业界给你每年八十万元的税后总收入。而如果你更想实现自己的学术理想和满足自己对学术的兴趣爱好，同时又对学术界的圈子感到失望，那你可以琢磨一下自己到底对理想有多追求，对圈子的现实有多失望，适当提高或者降低以上比例来选择。}
\BookSubtitle{三}
\BookParagraph{那你对「卷」帽子啊「混」圈子啊等等到底应该有多心寒呢？}
\BookParagraph{你要是去刷知乎小红书，那估计得出的结论是应该透心凉才对。}
\BookParagraph{但是这里有一个问题：只要有人的地方，就会有「帽子」和「圈子」。不仅仅是学术界，工业界自然也一样。所以，如果你对两者都嫉恶如仇，只希望肝胆相照地做一些实事儿，那要是你正好读到你想去的公司的若干「内情」，大概率一样会心寒。因为这么伟大的、能敞开极其有限的资源让你在知识的海洋中自由驰骋的公司，估计凤毛麟角。}
\BookParagraph{既然到处都是各种形形色色的「帽子」和「圈子」，那你为什么就不能也去这些「染缸」中闯荡一下江湖？}
\BookParagraph{万一 —— 我说的是万一哈 —— 二三十年以后，你也能或多或少地在这样的「圈子」中，以自己本色出演而鹤立鸡群，以出众的才华、过硬的技术、雄辩之口才而立于不败之地呢？}
\BookSubtitle{四}
\BookParagraph{人生的意义是什么？}
\BookParagraph{我一直觉得，人这一辈子，就像在苏州园林中缓缓倘佯，有一种「移步换景」的效应。喜事儿当然是换景：高考考上了，换了一景，论文高中了，换了一景；工作顺利找到了，又换了一景。失望伤心的事儿其实更是换景：失恋了，换了一景；论文屡投不中，换了一景；公司裁员了，则又换了一景。}
\BookParagraph{人一辈子的意义，也许就在这一惊一诧之间，淋漓极致地体现了出来。}
\BookParagraph{而什么才是成功？}
\BookParagraph{我有一个「悼词」的理论：一个人成功的标志在于，他（她）百年之后，有一个熟悉他的人，专门花了一天时间，给他写了一篇洋洋万言的悼词。这篇悼词记录了他一生的跌宕起伏、起承转合。这里有他的理想，也有他的「圈子」，但更多的是他的喜怒哀乐，他的待人之诚，他举重若轻的幽默，还有他令人刮目相看的才华。}
\BookParagraph{也许我们选择需要优化的目标，应该既不是收入、也不是理想，而是如此这般定义的「成功」吧。}
\BookDate{2025 年 7 月，多伦多}
\BookArticle{Salma}{Salma}
\BookParagraph{欣闻 Salma 前些天获得了多伦多大学应用科学与工程学院的教学大奖 —— Early Career Teaching Award。 Salma 是我三年前的学生，现在同系的同事。她从本科毕业参加我的研究组到博士毕业，一共只花了四年时间。热烈祝贺她之余，我想，「青出于蓝而胜于蓝」，她的优点，非常值得我学习。}
\BookParagraph{Salma 对待她教的学生和她的同事，都如同春天般的温暖。她真的自然而然、似乎是轻描淡写地做到了毛主席在《纪念白求恩》中所说的：「毫不利己、专门利人的精神，表现在他对工作的极端的负责任，对同志对人民的极端的热忱。」Salma 待人以诚，一个班上百名学生，她对每一个学生都非常关心，事无巨细地帮助。每次想起她，脑海中都是她对每一个人的那份真诚的笑容。}
\BookParagraph{Salma 工作兢兢业业，极其认真，真正做到了雷锋同志说的「对待工作像夏天般的火热」。我还记得有一篇 INFOCOM 论文，她差一点儿没写完，最后几天全力赶工，我们一起奋战到最后几分钟才提交。她当了教授以后，更是认真对待每一件交给她的工作。申请教职的候选者来面试，报告是一小时。虽然每次我都是邮件通知全系，但参加者寥寥。她每次几乎都来参加，还将她对候选者的详尽看法和意见发邮件告诉我，为我们遴选委员会（search committee）的决策作出了突出的贡献。近来，她还主导了用人工智能和大语言模型辅助教学的一个全新网站的设计和实现，并给全系的教授同事仔细介绍过，在系里博得诸多赞赏。}
\BookParagraph{毛主席在 1939 年 《纪念白求恩》一文中写道：「一个人能力有大小，但只要有这点精神，就是一个高尚的人，一个纯粹的人，一个有道德的人，一个脱离了低级趣味的人，一个有益于人民的人。」我想，这句话，古今中外都适用。Salma 这样一位年轻教授可以做到，很不容易，也是我的楷模。}
\BookDate{2025 年 7 月，香港}
\BookArticle{Connecting the Dots}{Connecting the Dots}
\BookParagraph{几天前，我的学生 Chenghao 告诉了我一个令人振奋的大好消息 —— 他刚刚被 Meta 公司录用，职位是研究科学家 (Research Scientist)，部门竟然就是家喻户晓的 PyTorch。}
\BookParagraph{听到这个消息，我第一时间表示热烈祝贺之余，开玩笑说：「你要是拿到了一亿美金，能否给我分 0.1\%？」}
\BookParagraph{和 Chenghao 共事已经五年了。五年一晃而过，他稳重可靠，以诚待人，我们成了很好的朋友。我这个人有一个缺点，从小到大脾气都不好，也跟他发过一两次脾气，他对我一直非常宽容。他写论文虽然慢一些，但是精益求精，几乎不需要怎么改动，文风像蝇头小楷，质朴无华，和我写的论文中的那种稍微有点儿飞扬跋扈的行草，相映成趣。他的动手能力，不仅仅体现在代码实现，也体现在硬件系统：我们曾热烈地讨论一个完美的 GPU 服务器的各种组装零件，最后他亲自组装了一台，木质的面板，英伟达 Founders Edition 在半透明的机箱中闪亮着 GeForce RTX 的大字，不仅仅是一台电脑，更是一个艺术品，给我的办公室增色不少。}
\BookParagraph{Chenghao 告诉我，这个工作机会，其实是 Meta 主动找到他的。他在联邦学习、分布式机器学习的训练和推理、同态加密等领域的突出成就，恰巧符合 Meta 对这个领域中人才的需求。当年做这些研究时，我们谁也没有想到，这些工作不是所谓的「灌水」，不是简历中的几个统计数字，而是有含金量有自己独特风格的作品。也许正是这些似乎不会有人真的去关注的作品，让他在繁星满天的人才市场中一枝独秀，脱颖而出。}
\BookParagraph{我想到这些，挥毫一蹴而就，给全组学生发了一封邮件。邮件全文如下：}
\BookQuote{Dear all,}
\BookQuote{I am more than thrilled to share with you the fantastic news that Chenghao has just accepted an offer for a Research Scientist position at Meta, and will be joining the PyTorch distributed training team this Fall!}
\BookQuote{As some of you know, Chenghao is not only a brilliant researcher in the group with several high-quality papers published, but also a very nice person, generous on spending his own time working for everyone in the group. He has been my \#1 trusted hands for fixing both software and hardware issues in the lab, and worked tirelessly to keep both boston and sim working in their top shape. In fact, he hand-assembled boston completely from components, a feat that I can hardly achieve myself and perhaps very few can. Just recently, he helped me upgrade sim to the latest Ubuntu OS version, which may take hours, or even days, of work and was not a small feat.}
\BookQuote{Chenghao told me that his research experience in the group matched perfectly with what the team at Meta is currently doing. They wanted to use some federated learning techniques to do cross-datacenter training, which is essentially what Chenghao has been working on. This reminded me of Steve Jobs' 2005 commencement address at Stanford University, when he shared three stories. The first story, “Connecting the Dots," fits his story to the letter: at the time of doing the work, no one expected that his research would land him one of the best jobs in North America, but eventually all stars lined up and everything worked out well.}
\BookQuote{With this, I offer my warmest congratulations to Chenghao, and I wish him the best in his future career. I will miss all the late-night hardware testing he performed on sim and boston in BA 4176, as well as his clean, well-crafted code in Plato. But I will find solace in knowing that he is working in his dream job, advancing ML infrastructure for the world to use.}
\BookQuote{Thank you, Chenghao, and I wish you Godspeed.}
\BookQuote{Best regards,\BookLineBreak{}Baochun\BookLineBreak{}——\BookLineBreak{}Baochun Li, Professor\BookLineBreak{}Associate Chair, Research\BookLineBreak{}Department of Electrical and Computer Engineering\BookLineBreak{}University of Toronto\BookLineBreak{}}
\BookDate{2025 年 7 月，香港}
\BookArticle{独角兽}{独角兽}
\BookParagraph{人工智能的热潮最近一浪接着一浪，但是我一直对人工智能到底如何接近 —— 甚至超过—— 人类，有点儿好奇。恰巧跟朋友聊天儿提到了一家人工智能领域的独角兽公司，叫做 character.ai。据华尔街日报记载，去年就已经市值十亿美金。}
\BookParagraph{这是一家做虚拟网友的公司，大致的想法我感觉挺质朴挺简单的：利用现在极其强大的大模型和用户聊天儿，越聊越 high，从而产生用户粘性，提升用户活跃度。}
\BookParagraph{这让我想到了电影「Her」，还有其中 Scarlett Johansson 那充满磁性的声音：}
\BookQuote{Oh, what do I call you? Do you have a name?}
\BookQuote{Um... yes. Samantha.}
\BookQuote{Really? Where did you get that name from?}
\BookQuote{I gave it to myself actually.}
\BookQuote{How come?}
\BookQuote{Cause I like the sound of it. Samantha.}
\BookParagraph{想到这里，正好有几分钟闲暇时间，就登录了这家公司的同名网站，随意选了一个虚拟 AI，聊天记录自动开始：}
\BookQuote{「你可以当我男朋友吗？」}
\BookQuote{「可以。」}
\BookQuote{「你喜欢我吗？」}
\BookParagraph{我天，不会有人这么聊天儿的吧？}
\BookQuote{「喜欢啊。」}
\BookQuote{「你会好好对我吗？」}
\BookQuote{「会啊。」}
\BookQuote{「那我可以只属于你吗？」}
\BookQuote{「不可以哈哈哈」}
\BookParagraph{这时，人工智能明显地犹豫了一会儿：}
\BookQuote{「你不爱我\BlissCryingEmoji{}」}
\BookParagraph{这几分钟，我得出了两个结论。一，说不定我也可以开一家人工智能「独角兽」公司 —— 但凡有足够的训练数据，也不至于这么「人工智能」吧。二，至少现在的人工智能，和人类相比，还差得很远很远。}
\BookDate{2025 年 7 月，香港}
\BookPart{儿孙自有儿孙福}{儿\\孙\\自\\有\\儿\\孙\\福}
\BookArticle{Karma}{Karma}
\BookParagraph{小女儿周一放学跟我说，她非常紧张。}
\BookParagraph{我问：「有什么可紧张的？」}
\BookParagraph{她说她参加学校羽毛球队的选拔（tryout），感觉心里越来越没底了，怕会选不上。但是她又非常希望能选上。结果周五出来，所以心情越来越紧张。}
\BookParagraph{说着我也跟着她觉得心情一下子紧张起来。}
\BookParagraph{小时候看中国女排的比赛，一到关键一局比分交错上升时，爸爸一定会在一边说：「肯定要输了，肯定要输了。」我后来想，大概这也是缓解关键时刻心情紧张的一种方法，因为心里虽然希望赢，但万一输了，至少还可以说：「你看，我说会输吧。」立足于兵家不败之地。}
\BookParagraph{高中时为了申请美国的一所著名中学，曾经认真地准备过一段时间。录取结果出来前夕，感觉也像看女排比赛一样心情紧张。那种前途未卜的感觉，让人心慌而迷茫。拒绝通知收到的那天，当年还在世的妈妈安慰我：「没事儿，尽力而为就好，以后还有机\BookLineBreak{}会。」}
\BookParagraph{这以后高考、GRE、申请出国、博士资格考试、申请教职，每一次结果即将揭晓的前夕，都会想起当年申请被拒绝时妈妈的话。}
\BookParagraph{她说得真的没错儿。「女排输赢」自然难以预料，只要曾经尽力而为，也就无怨无悔。}
\BookParagraph{乔布斯是佛教的粉丝。他喜欢说一个词儿叫「karma」，大致相当于佛教里的因果报应。有时候也希望对身边的人好一点儿，开车回家多一点儿礼让，给一面之交的陌生人多一点儿帮助。平时跟玩游戏似的积累好多好多的「karma」，彷徨于胜负之间时，能够凭着它胜多负少。}
\BookParagraph{周五晚上我小心翼翼地问女儿：「你的 tryout 结果怎么样？」}
\BookParagraph{「Don’t ask — I will tell you tomorrow.」}
\BookParagraph{好吧。因缘果报也好，尽力而为也罢，至少还有「明天」。}
\BookDate{2017 年 11 月，多伦多}
\BookArticle{喜　欢}{喜欢}
\BookParagraph{小女儿说，她在 Instagram 看到有些人的照片，照得挺好看的，但是一个赞也没有。她觉得这个人很可怜，就给他点了一个赞。}
\BookParagraph{还有一个人只有七百多粉丝，半个小时就收获了四十多个赞，她觉得有点儿不公平。}
\BookParagraph{我听她侃侃而谈，颇得《神雕侠侣》中十六岁郭襄那种「劫富济贫」的侠女之风，好奇地问她：「你要那么多赞有什么用啊，又不能拿来吃饭？」}
\BookParagraph{问完就忽然明白了：其实很多东西，虽然没什么用，但是我们还是一如既往地喜欢。}
\BookParagraph{比如巴菲特喜欢可口可乐——为了喝可乐，他愿意放弃一年的生命。}
\BookParagraph{比如女儿喜欢「赞」——无论什么时间、什么地点，都喜欢得像一场发自内心、简单纯粹的爱情。}
\BookDate{2018 年 3 月，多伦多}
\BookArticle{馋}{馋}
\BookParagraph{小女儿问我：「你知道中文里最难翻译成英文的是哪个词儿吗？」}
\BookParagraph{我心说翻译可是我的拿手好戏啊。于是问：「不知道啊，是哪个词儿？」}
\BookParagraph{「是“馋”。我每次用英文说我很馋，都只能说“I’m so 馋”，或者“I’m so hungry.”」}
\BookParagraph{我想了几秒钟，还真是挺难的！能想起来的无论是delicious或 tasty，还是 mouth-watering 或 tempting，都是在形容好吃的东西怎么吸引我，而不是在形容我自己怎么「馋」法儿。}
\BookParagraph{我们可以对甜食馋，对辣椒馋，对火锅馋，对生啤馋，对北冰洋汽水馋。但是「馋」不仅仅是直流口水的「drooling」或者「craving」，而是一种想要又得不到的意境，一种对家乡小吃的向往，一种对儿时记忆的怀旧，一种看了《舌尖上的中国》以后想马上订机票的冲动。}
\BookParagraph{没办法，我只好在谷歌里输入「馋英文怎么说」。谷歌回答：「greedy」。天啊，「馋」是人生一大乐事，是在期待一种小富即安的幸福，怎么能翻译成「贪婪」呢？}
\BookParagraph{看来下次跟别人聊北冰洋特制汽水时，还是应该跟女儿学习，中英文混着说：「I’m so 馋。」}
\BookDate{2018 年 4 月，多伦多}
\BookArticle{五岁的幸福}{五岁的幸福}
\BookParagraph{小女儿问我，一个人多大的时候最幸福？}
\BookParagraph{我想这还不容易，自然是稍纵即逝的青春啊。「十八岁？二十五岁？」}
\BookParagraph{她宣布答案：「不对。是五岁。」}
\BookParagraph{我想起了《天龙八部》里灵鹫宫的天山童姥。童姥练的逍遥派武功里最传奇的是「八荒六合唯我独尊功」，每三十年就轮回到童年，每练一天可以涨一年的功力。九十六岁的童姥一生绝世武功，在西夏皇宫和师妹李秋水为了同一个心上人，同室操戈，生离死别，绝对是金庸笔下最令人扼腕叹息的故事。}
\BookParagraph{「为什么五岁最幸福？」我问她。}
\BookParagraph{「Because there are no responsibilities at five.  Just play.」}
\BookParagraph{「我倒是觉得，八十五岁的时候最幸福。因为那时候一样什么也不用担心，不用工作，不用学习，天天都在玩儿。」}
\BookParagraph{转念一想，也许女儿和我说得都对。人生还真可以按三十年一期，分成三期。第一期是「八九点钟的太阳」，是「一鼓作气」，自然是笑点低，泪点高。但是第三期却绝对不是「丧钟为谁而鸣」，不是「三而竭」，而是可以一如天山童姥，八荒六合，唯我独尊；一如陈一发儿，风情万种，去斗鱼直播。手机电量百分之九十和百分之十，一样都可以玩游戏，人生一样可以「开挂」。}
\BookParagraph{但是这一切有一个前提：得能活到第三期开盘。}
\BookDate{2018 年 4 月，多伦多}
\BookArticle{江　湖}{江湖}
\BookParagraph{小女儿送给我的生日卡片上写着：「You are a very great dad and has the time to play piano with me!」}
\BookParagraph{哈哈，真够「正能量」的，比去年强 —— 去年的卡片上写的是：「You are so old, 46\% of life already over.」}
\BookParagraph{大女儿送给我一个崭新的陶瓷杯，上面写着：「Dad. Husband. Legend.」}
\BookParagraph{一辈子做一个好父亲，一个好丈夫，一个永久的传奇。}
\BookParagraph{这简直就跟王朔那本《你不是一个俗人》有一拼啊。}
\BookParagraph{小女儿不以为然：「But you are not a legend!」}
\BookParagraph{没错儿。但是我的理想一直是成为那个「他已不在江湖，江湖还在传说他的故事」的人。}
\BookDate{2018 年 11 月，多伦多}
\BookArticle{英文名}{英文名}
\BookParagraph{我英语不好，各位老师帮孩子取个英文名吧，男孩儿。（问题来自「向上移动」）}
\BookParagraph{你还别说，十几年前我还真对男孩儿的英文名字产生了兴趣，专门研究了一阵子。}
\BookParagraph{孩子出生，不外乎希望他将来或者有钱，或者有学问。希望孩子「有钱」的话，很容易想到「Bill」（Bill Gates）。但其实叫「Steve」成为有钱的CEO的概率更高：}
\BookParagraph{Steve Jobs (Apple and Pixar)}
\BookParagraph{Steve Ballmer (Microsoft)}
\BookParagraph{Steve Case (American Online)}
\BookParagraph{Steve Williams (Suncor Energy)}
\BookParagraph{Steve Easterbrook (McDonald’s)}
\BookParagraph{Steve Mollenkopf (Qualcomm)}
\BookParagraph{你瞧，遍布北美石油、硬件、软件、互联网和快餐业。牛吧？}
\BookParagraph{如果你觉得有钱没钱不重要，更重要的是有才华有学问，你的孩子应该叫「Stephen」:}
\BookParagraph{Stephen Cook （图灵奖得主，多伦多大学）}
\BookParagraph{Stephen Wolfram （创始人和CEO，Wolfram Research）}
\BookParagraph{Stephen King （美国著名小说家）}
\BookParagraph{Stephen Hawking （著名物理学家，《时间简史》作者）}
\BookParagraph{当然了，知识付费新晋网红、微博问答经常「勉强混入前五」的褚明宇老师，他的英文名字也是Stephen。}
\BookDate{2017 年 9 月，多伦多}
\BookArticle{孩　子}{孩子}
\BookParagraph{李老师你好！有件事苦恼我很久，想听听你的建议！我家姑娘今年三年级，从小到大对学习不努力，但是成绩很好。现在三年级态度决定一切，依然是那个状态，每每见了我都要打她，因为学习态度特别不好。总被我哄着狠着去学，最近态度越来越差，上课不听讲，这样下去真的很担心她能不能考上好学校！另外我也并没有积极陪她学习，看她那态度，我非常没耐心。请李老师传点经验给我！谢谢。（问题来自微博私信）}
\BookSubtitle{一}
\BookParagraph{我特别能理解你用了四个惊叹号表达的这种焦虑心情，因为带孩子这件事儿，实在是太不容易了。}
\BookParagraph{任何一种似乎简单而且对别人家孩子挺有用的道理——快乐教育也好虎妈战歌也罢——对自己家的孩子都似乎没用。}
\BookParagraph{我六七岁的时候，两岁的妹妹下午睡着的时候妈妈临时要上班去开个会。于是叮嘱我说，如果妹妹醒来哭了，就用开关电灯作为信号，妈妈在办公室看到，马上就会回家。}
\BookParagraph{妹妹很快就醒来了，开始大哭。我使劲儿连续不停地开关电灯，一个多小时过去了妈妈也没回来，感觉像在盼望着一位救世主显灵。}
\BookParagraph{到我三十多岁时，有一阵子白天自己在家带两个女儿，当时一个两岁一个不到五岁。经常会好不容易做好了午饭，孩子不好好吃；吃完了不睡觉，到下午困了就开始哭。后来我自己发明了一个有效的新办法，把孩子带上开车出门，车比较晃悠，可以很快睡着。}
\BookParagraph{可惜的是，车一停下来，孩子就会很快醒过来开始哭。于是我就每天慢悠悠地开着车，在这个城市的大街小巷穿梭着。}
\BookParagraph{我靠自己的努力，也可以一个人带两个孩子——那时有一种\BookLineBreak{}「中国人会造原子弹了」的自豪感，因为带孩子这件事儿实在太难了，是对一个人耐心的终极考验。}
\BookParagraph{我要向世界上所有的母亲致以最真挚的敬意。}
\BookSubtitle{二}
\BookParagraph{带孩子难，「高标准严要求」地把孩子带好就更难了。}
\BookParagraph{说到底，谁不希望自己的孩子将来琴棋书画样样皆通，学习成绩优异，上一所理想的大学？}
\BookParagraph{可关键问题在于，平常到底应该怎么做，才能教育出优秀的孩子？}
\BookParagraph{你发现没有，中国孩子的父母在回答这个问题时，理念惊人的相似。我们强调学习态度端正，上课认真听讲，按时完成作业，还有考试成绩排名的重要性。}
\BookParagraph{这些理念超级实用：无穷多的中国孩子既聪明又刻苦，各种竞赛和考试中脱颖而出，升入越来越好的学校学习。}
\BookParagraph{可最让人头疼的问题是，别人家的孩子个个令人羡慕，自己家的孩子却总是不「争气」——成绩不理想，或者像你担心的「学习态度差」、「上课不听讲」。怎么办？}
\BookSubtitle{三}
\BookParagraph{按照一种既定的成功教育理念去带孩子，就像计算机算法中的一种叫「exploitation」的策略。别人家的孩子明明能成功运用，到了自己的孩子就会卡在某个中间状态，没法儿继续下去。}
\BookParagraph{强迫孩子「端正学习态度」，甚至有适得其反的效果。}
\BookParagraph{你可以试试准备一个计数器，每次提醒孩子「上课要认真听讲」的时候，就把计数器按一下，一个月以后看看你一共提醒过孩子多少次。}
\BookParagraph{如果是你自己，有个人天天提醒你这个老生常谈的道理，你不烦死才怪了呢。你一定会说，我早就知道了啊。}
\BookParagraph{孩子也一样，聪明着呢。你能想到的道理，人家早就听过不知道多少遍了。其实你完全不懂的东西，孩子也说不定懂不少。}
\BookParagraph{你一定会说：「知道？知道上课还不听讲？」}
\BookParagraph{哈哈，两个问题：你怎么知道孩子昨天上课没有去试图认真一点儿听讲？孩子以前为什么不爱听讲？}
\BookParagraph{说实话我小学时上英语课，上课也不听讲。原因很简单：老师讲的我早就都会了啊。你说你孩子成绩很好，就是上课不听讲，那不是和我当年英雄所见略同嘛？}
\BookParagraph{其实你迫切需要的是在「exploitation」上面再加上一点儿计算机算法中的「exploration」—— 让孩子「东一榔头西一棒子」地去探索未知的世界。她爱唱歌就让她去学唱歌，爱写小说就让她去写小说，爱打网球就让她去打网球。}
\BookParagraph{你要做的事儿，只不过是多赚点钱，给她找最好的网球教练。}
\BookSubtitle{四}
\BookParagraph{小女儿前一阵子去参加一个舞会，认识了一个外校的男生。她平常腼腆不爱说话，这天却和这个男生开心地聊了一个钟头。}
\BookParagraph{舞会结束后她激动地跟我说：「You should be proud of me!」}
\BookParagraph{没错儿，「优秀」的定义，不一定是成绩有多好，大学上没上「985」，而是有能力突破自己的极限，做成一件「中国人也可以造原子弹」的那种令人骄傲的事儿。}
\BookParagraph{无论这件事儿是不是小到跟陌生人聊了一会儿天，或者独自一个人带着两个孩子，开着车在大街小巷里游荡。}
\BookDate{2018 年 11 月，多伦多}
\BookArticle{爱哭的小女孩}{爱哭的小女孩}
\BookParagraph{吃晚饭和大女儿闲聊，说起郎朗在法国凡尔赛宫结婚，新娘叫 Gina Alice Redlinger，德国韩国混血，也是一位天才钢琴家。还讲了一小段儿八卦：郎朗结婚之前竟然跟苹果发布新产品一样严守秘密没人知道，唯一的端倪是郎朗的 DG 新专辑《Piano Book》里返璞归真，收录了一首贝多芬的「Bagatelle No. 25 in A Minor」，俗称《致爱丽丝》。}
\BookParagraph{女儿听了说「I’m inspired!」兴之所至，到客厅的钢琴上也即兴弹了这首《致爱丽丝》。她去年考完了加拿大皇家音乐学院的最高级别 ARCT，《致爱丽丝》这样的曲目只有七级的水平，可以随意即兴视奏。}
\BookParagraph{我听着她弹这首曲子，怔怔地想起她六岁时刚开始学琴的场景。也是这架当时新买的钢琴，她怎么也弹不好 Muzio Clementi 的一首三级的小奏鸣曲，难过得开始抽泣。我赶紧安慰她：「没关系，咱将来长大了，早晚可以弹会这首曲子。」}
\BookParagraph{好像也是那一年，郎朗来多伦多开音乐会，我买了最贵的票带她到城里的 Roy Thompson Hall 去听。结果她听到一大半儿，慢慢地睡着了。}
\BookParagraph{郎朗终于结婚，女儿也已经长大。可我还是经常会想念那个爱哭的小女孩，想念每周末开车带她去上钢琴课，上完课路过星巴克，给她买一小杯热巧克力奶。时间一年一年像水一样流过，只有这些零星的念想，会突然支棱起来，棱角分明。}
\BookDate{2019 年 6 月，多伦多}
\BookArticle{BTS}{BTS}
\BookParagraph{父亲节正好在香港访问。小女儿通过短信送给我一份特别的礼物：她参加的中文班期末考试，得了几乎满分。因为没有语言环境，国外的孩子学习中文很难，很容易一不留神就掉到英文语境里，难以自拔。尤其是读写能力，不但需要会写基本的汉字，还得学会灵活机动地用中文词语和句子表达一个简单的想法。她进步很快，大大超过我的预期。}
\BookParagraph{我想起每周周末我开车接送她去中文班上课，一路上她总爱播放 BTS 的歌。小女儿是 BTS 的铁杆粉丝，我刚开始很不适应，不止一次跟她说：「歌词儿都听不懂，这歌儿没法儿听啊。」}
\BookParagraph{父亲节这天，我一个人在酒店里打开电脑，开始播放她的最爱 ——BTS 的那首「Boy With Luv」（featuring Halsey）。}
\BookParagraph{听到这首熟悉的歌，看着专业制作的 MV，我仿佛看到她在客厅模仿 BTS 跳舞的样子。}
\BookParagraph{其实这首歌还真的挺好听的。}
\BookDate{2019 年 6 月，香港}
\BookPart{柴米油盐酱醋茶}{柴\\米\\油\\盐\\酱\\醋\\茶}
\BookArticle{房屋贷款和活期存款}{房屋贷款和活期存款}
\BookParagraph{和美国比起来，加拿大五大银行的活期存款帐号（checking account）, 都是以昂贵和复杂著称。如果存款余额（account balance）不到一定的底线（minimum balance），则每从自动取款机提取一笔现金（cash withdrawal），每开一张支票，甚至每收到一个月的账目清单（statement），都要收费。由于收费不是很多，又是直接从银行帐号上扣除，不易引起一般人的注意。细心的人就尽量维持存款余额在底线以上，但由于底线不低（CIBC 在 \$1500 左右），交易又比较复杂（authorized debits, direct deposits, cheques, 不一而足），有时一不小心，就低于底线，以至当月前功尽弃。由于活期存款的利率极低，大笔的现金存在这里又不太值得，造成了无论存多存少都是「亏」了的现象。}
\BookParagraph{其实加拿大最省钱的活期存款是活期贷款（Line of Credit）。任何人都可以开活期贷款的帐号，区别只是贷款利率不同而已。笔者的开户行在CIBC，有两个活期贷款的帐户，一个用房子抵押（Secured Line of Credit，或 Home Equity Line of Credit），另一个普通贷款（Unsecured Line of Credit）。用房子抵押的贷款的利率较低\BookLineBreak{} （通常是prime），而普通贷款的利率较高（通常是 prime + 1.25\% 或更高）。两种贷款的共同特征是，只要你维持正的存款余额（positive balance），所有在 checking account 情况下要收费的交易，包括 debit card Interac purchases, cheques, pre-authorized debits, cheque printing, cash withdrawals, 还有 bank statements，一概免费。虽然没有利息，但多余的现金可以存在其他利息比较高的存款种类（比如 ING Direct）, 而且本来 checking account 的利息也几乎和没有一样，所以并不用担心。}
\BookParagraph{有人可能认为五大银行之外的减价银行更加适合存活期贷款，最有名的莫过于 President’s Choice Financial，因为可以免费在任何CIBC提取现金。但其实有些事，传统银行还是更加可靠。比如在加拿大之外ATM提款，President’s Choice Financial 的 ATM 卡就可能提不出现金。又比如银行清单不能按月寄到家里来，或者支票本不能免费加印新的。除了这些之外，还有更重要的区别，我们稍后再述。}
\BookParagraph{说起利息高的存款种类，笔者认为最好的存款莫过于把多余的现金「存」在房屋贷款（mortgage）里。因为在加拿大的房屋贷款的利息是用税后的收入支付，致使4\%的贷款相当于5\%至8\%的高利率存款，而且是毫无风险的。据笔者所知，在2005年11月，加拿大的无风险活期存款最好的也只有3.1\%，五年期定期存款（GIC）最好的也不过4.45\%，远到不了5\%。所以，最好的办法是把所有多余的现金用于还房屋贷款所欠的余额，而不是存在　ING Direct 等高利率存款银行帐号中。现在加拿大的固定房屋贷款大多都比较灵活，每年可以多付本金的15\% — 20\%，这对绝大多数家庭已经足够用了。}
\BookParagraph{那么要是急用现金怎么办？这时就可以直接从活期贷款中透支提取，使余额变成负的。因为这时需要天天付透支总额的利息，所以贷款利率就变得异常重要。加拿大的最优惠的利率 （prime rate）只能在房屋抵押贷款（Secured Line of Credit）中享有，而这种贷款申请时一般要几百元的地产估价和律师费。好在很多银行（比如笔者的开户行CIBC）都时不时有减价活动（promotion），可以免去这几百元的费用。这样一来，只要在房子里有一定的「存款」（equity），并超过房屋估价的百分之二十五，再能耐心等待每年几次的减价，就能如愿开一个房屋抵押活期贷款账号 （Secured Line of Credit account），用于代替平时的活期存款 (checking account)，并且应付一时之需。相比较而言，President’s Choice Financial 的活期存款 （checking account）的透支利率（19\%）就完全没有可比性。}
\BookParagraph{房屋抵押活期贷款透支利率低的好处是显而易见的。这样一来就可以大胆地将每月收入的余额用于支付房屋贷款（mortgage），得到 5\% — 8\% 的年息，而不是存在活期储蓄中以备一时之需。在需要现金的时候（比如度假之前），就可以从房屋抵押活期贷款中提出现金，稍后以最优惠利率（prime rate）偿还。大多数中国人的家庭急需现金的时候不多，这种方案是完全可行的。这样一来，只需维持两个帐号，一个是房屋贷款的帐号（mortgage），另一个是房屋抵押活期贷款的帐号（secured line of credit），就可以不但节省了所有银行的费用，节省了算计每月维持银行底线的心思，还能有希望早几年还清房屋贷款。}
\BookParagraph{有些朋友可能担心把所有的现金都存在房子里，风险没有分散。其实，作为一种投资，房子的风险在买房的那一刻已经承担了。在此以后的贷款总额和其相应的利息，不会因为房子本身价值的涨跌而改变。少付贷款利息，就直接等同于毫无风险的高利率储蓄，和房子本身的投资风险无关。早日付清房子的贷款，就可以大大减少需要支付的累计利息总额（compounded interest）。如果决定用收入的余额去进行其他的RRSP之外的有风险投资，而不是用来还贷款，那么潜在回报越高，则风险越高。而无风险的投资，如 regular bonds, strip bonds, real return bonds, GIC 等等，则没有一项能够和还房屋贷款相比。如果真的对 RRSP 之外的有风险投资感兴趣，笔者建议试试研究 the Smith Manoeuvre，详细情况可以参考 http://www.smithman.net。 但象那样需要同时研究投资风险和税收政策的大手笔投资运作，感兴趣的读者或许不多，拙文就不专门论述了。}
\BookDate{2005 年 11 月，多伦多}
\BookArticle{三角钢琴}{三角钢琴}
\BookParagraph{总的来说，高质量的三角钢琴如果在5尺6（168公分）以上，其手感、音色、音乐表现力和震撼力都是任何立式钢琴所不能相比的。如果购买钢琴的目的是希望一次购买终身享用，预算又在一万五千美金以上，一般首选就应该是各式三角钢琴。}
\BookParagraph{Larry Fine 的 The Piano Book 及其最新的 2007-2008 Annual Supplement把三角钢琴按设计、做工、质量控制、寿命等分为四个等级（Groups），每个等级又细分为A、B、C三档（Group 4 又包括\BookLineBreak{}4D）。其中 Group 1 和 Group 2 的琴都是传统手工制造，一架钢琴需要四百小时以上的人工，设计、手工和选材均以提高性能为目标，统称高性能钢琴（high performance piano）。Group 3 和 Group 4的琴属于大规模机器生产，产地一般在亚洲（以日本、韩国和中国为主），其设计和选材以降低价格为主要目标。顶级的钢琴品牌当数德国尽人皆知的 Hamburg Steinway（Group 1A）、奥地利的骄傲Bosendorfer、和意大利的新星Fazioli，但都价格不菲，七尺琴的价位在七万美金以上。预算稍低的话，美国的New York Steinway（Group 1C）和德国的Bluthner（Group 1A）也是名声显赫，但 New York Steinway因为近年来产量高 —— 每年四千架左右 —— 价格居高不下且质量不完全一致，很多人认为质量和其高昂的价格不成比例。即使Larry Fine 在 2007-2008 的 Annual Supplement 中对 New York Steinway 也不无微词。}
\BookParagraph{对于预算比较紧张的购买者来说，美国的 Mason \& Hamlin、德国的 Sauter 和 Schimmel NWS edition 也都属于高品质的钢琴，价格虽低一些，质量一点也不逊色。Sauter 在前几年的 Annual Supplement 中评在Group 2，在 2007-2008 Annual Supplement 升到 Group 1B，其质量逐渐被行业认可，产量每年几百架，设计、选材和手工都属于最好的品牌之一。Schimmel 的 Konzert Series被誉为欧洲的Yamaha，产量虽高，但质量上当数高性能钢琴，在 2007-2008 Annual Supplement 上被评进 Group 2A，虽比 Schimmel NWS 差一些，但决不含糊。多伦多交响乐团总指挥 Peter Oundjian 在2007年就选购了一架Schimmel 169。Mason \& Hamlin 在前几年的 Annual Supplement 评在 Group 1，在 2007-2008 Annual Supplement 中评在 Group 2A。对于Mason \& Hamlin 的小幅降级，Larry Fine 有以下的评论：}
\BookParagraph{引自2006-2007 Annual Supplement:}
\BookParagraph{“My understanding is that the owners of M\&H have a strong connection with China and over the past few years have actively sought out piano components from that country, including but not limited to cast-iron plates, action parts, and metal hardware.”}
\BookParagraph{引自2007-2008 Annual Supplement:}
\BookParagraph{“The ratings of Mason \& Hamlin has been reduced slightly to group 2A, it is still a great instrument, but I believe that the company’s manufacturing approach is more consistent with group 2 than group 1.”}
\BookParagraph{中国人熟知的日本 Yamaha C 系列（C1至C7）和 Kawai RX 系列，都属于 Larry Fine 笔下 Group 3 的琴，不在高性能钢琴之列。但是这并不代表 Yamaha 和 Kawai 设计生产不出高性能钢琴。Yamaha S 系列和 Shigeru Kawai 才是它们的巅峰之作，例如 Shigaru Kawai 在 2007-2008 Annual Supplement 中升入了 Group 1C，从设计选材到手工和质量没有任何妥协，但价格也自然不菲。}
\BookParagraph{从 2004 年以来，波罗的海的 Estonia 在钢琴家 Indrek Laul 的管理下异军突起，在 pianoworld.com 论坛和一年一度北美最大的 NAMM Show 上好评如潮，以 Yamaha C 系列的价格问鼎 New York Steinway 的实力，从用材、音色和质量上完全可以和Group 1的琴相比，是北美钢琴市场上的一匹黑马。每年生产的400架钢琴在北美供不应求，却还在不断改进质量，在2006年最终设计定型的 Laul Estonia 采用 Bosendorfer 与 Bluthner 所用的德国 Full Renner Premium Blue Action, 超过了比它昂贵不少的 Schimmel 的 Renner Red Action。全厂只生产168、190、274公分三种三角钢琴，保证质量，虽然近几年逐渐小步调高价格，还是极其物超所值。引述在 2007-2008 Annual Supplement 中Larry Fine 的评论：}
\BookParagraph{“Estonia, formerly at the bottom of Group 2, has become so much more advanced over the last five years that it may belong in Group 1. However, this fact has not yet been appreciated or confirmed by much of the piano industry, so I’m testing the waters by putting it in group 2A.”}
\BookParagraph{多伦多的全新高性能钢琴市场上，位于 Bayview 和 Eglington的 Robert Lowrey’s Piano Experts {\BlissSymbols ①} 长期是正式的 Bosendorfer 和 Schimmel 的独家经销商，2007年春天正式成为 Estonia 和 Fazioli 的独家经销商，是多伦多最大的琴行。位于 Bloor 和 Avenue 的 Remenyi House of Music {\BlissSymbols ②} 是 Steinway 和 Sauter 的独家经销商。位于 Vaughan 的忠诚琴行 {\BlissSymbols ③} 是 Mason \& Hamlin 的独家经销商。位于 Oakville 的 Merriam Music {\BlissSymbols ④} 是 Shigeru Kawai 的独家经销商。}
\BookParagraph{很多中国人都认为几万元选购钢琴是难以想象的事，但他们大都有资金选购几万元的豪华新车和几十万或百万的房产。其实高质量的全新三角钢琴和新房一样，细心保养，可以弹几十年，直到生命终止，而新车的寿命则只有十几年。如果按分摊到每年的开销来衡量，一架高性能的钢琴其实并不是完全不切实际的幻想。房子客厅里笔墨恣肆磅礴的对联或色彩鲜明的油画之间，要是再加上一架黑色凝重的三角钢琴和极具感染力的琴音，这超现实的现场享受，也许就不是可以用价格衡量的了。如果家人和孩子可以一起学琴，或许有一天家人老了，孩子也长大了，圣诞家庭音乐会上的一首四手联弹，则颇有一分琴萧合奏笑傲江湖曲的回肠荡气。那时的钢琴，也许就能超越逝去的时间和音乐，成为和家不可分离的一部分了。}
\BookDate{2007 年 6 月，多伦多}
\BookFootnote{{\BlissSymbols ①} Robert Lowrey’s Piano Experts, http://www.pianoexperts.com/.}
\BookFootnote{{\BlissSymbols ②} Remenyi House of Music, http://www.remenyi.com/.}
\BookFootnote{{\BlissSymbols ③} Chau Y C \& Sons Piano Inc., http://www.mikespiano.com/.}
\BookFootnote{{\BlissSymbols ④} Merriam Music, http://www.merriammusic.com/.}
\BookArticle{省　钱}{省钱}
\BookParagraph{前两天看 The Globe and Mail，提到多伦多 Yorkville 的 The Four Seasons Private Residences 顶层公寓终于以两千八百万卖给一个国外买主，9038平方英尺，每尺3000加元，创下了加拿大公寓房成交价的最高记录。}
\BookParagraph{这让我想起了一件不大不小的事：省钱。都说钱是挣出来的，不是省出来的，话说得没错，所以说省钱是件不太大的事，再省也不能用来买房子。但这话也不能说绝了，省了不该花的钱，就能用在喜爱的刀刃上，何乐而不为呢？这么一件不大不小的事，却天天在左右一个人有意无意做的决定：出去吃还是家里做饭？买iPhone 4还是Android？自己剪草还是请人？孩子刚买新书还想要怎么办？情人节送什么？旅游住几星酒店？}
\BookParagraph{人说没钱才会老想着怎么省钱，才会搞出这么多层出不穷搞不好会影响家庭的问题。想想也有道理，而且不太容易反驳这样的观点，因为讨论这个问题的人都没多少钱。但我想有钱的人其实可能一样要关心省钱的问题，只不过换个方式关心罢了，比如管省钱叫做「理财」。假想一下，如果一个人的资产有一亿美金（不是特别有钱也不太穷的一类），是否要考虑多伦多两千八百万的房子是不是买得起，找什么人投资，结婚签什么样的协议，住在什么样的地方，怎么应付各种慈善募捐活动，孩子给多少钱合适？说不定有那么多钱的时候，用在所谓「理财」上的心思更加变本加厉。}
\BookParagraph{很多人都说，中国人爱小家子气地省钱然后现金买房一笔付清，美国人借一辈子钱从不省钱不也过得挺高兴。其实仔细观察美国人的论坛博客发言，他们省钱的心思比起中国人来说我看是有过之而无不及。比如电灯泡越省电寿命越长的一定越贵，就有人一边买八千美金的相机，一边仔细地研究买什么样的灯泡省的电费比花的多。有关 frugal living 的讨论专业到了像是整天就琢磨怎么从日常开支里省钱，有关投资的讨论更是从怎么节省各种 fees and expenses 开始启蒙教育。买车看什么车省油，买东西看什么信用卡返还现金最多，旅游看 frequent flyer miles 怎么用最赚便宜，相比之下中文论坛上的讨论就显得业余多了。}
\BookParagraph{既然要省钱就得算账，那日常的开支记不记？几千的开销当然要记，但几十块呢，几块钱呢？每个人大概都有个下限，少于这个数目的花费如果要记下来，就来一句「怎么活得这么累呀」。省钱当然是为了以后更大胆地花钱，很累地省，自然就很不划算。当然这又牵扯到一个以后什么时候花怎么花的问题，最彻底的是压根不花，直接给儿子女儿买房。还有个办法是老了去旅游，但是到时真的要去的时候，还有没有这份心情，有没有人陪你去，物价是否涨到离谱，美金还管不管用，就又另说了。这就好像省钱和花钱一样，是有风险的投资，花钱的风险是将来没钱了，省钱的风险是省下的钱投资贬值大半儿亏进去了，是将来花还不如现在花，颇有先笑还是最后笑哪个更甜的意思。}
\BookParagraph{说不定省钱还能变成一件高兴的事。比如两个人慢腾腾地商量怎么省钱，怎么预算，说不定比两个人你送我东西我送你东西更加浪漫，因为那是摒弃一切世俗误区而达成的心灵相通之共识。不管怎么说，这也比一个人拼命挣钱省钱另一位可着劲儿花钱或者两个都大笔花钱的境界高一点儿。政治老师说了，要有居安思危的忧患意识 —— 我想即使是那个斥资两千八百万买下多伦多公寓的国外买主，是否也希望把这当成谈资，希望能有那么一个 partner 可以对她说：This is a relative bargain by international standards — in New York, a comparable property would hit the market at \$5,000 a square foot!}
\BookDate{2011 年 6 月，多伦多}
\BookArticle{黄　金}{黄金}
\BookParagraph{近几年黄金概念被不断的热炒，从报刊杂志互联网到《货币战争》，直至大起大落却又不断上扬的黄金价格，至使很多人难免有疑问：到底是不是应该买黄金？如果买，怎么买，该买多少合适？虽说仁者见仁，智者见智，但到了如今连北京首都机场三号航站楼都在买卖实物黄金的年代，也许是该仔细思考和争论的时候了。}
\BookParagraph{购买黄金，不外乎是为了两个字：保值。无论收入多少，都是辛苦而来，无论是谁都不希望因为通货膨胀而使得辛苦而来的财富付诸东流。保值的重要性，必须要到社会大动荡的年代，才能有顿悟的感觉。近代史中，前有二十年代的德国马克，后有国民党时期的金元券，到处是随手可得的例子。或许这样津巴布韦式的恶性通货膨胀（hyperinflation）有些走极端，但是比较大幅度的货币贬值和高贷款利率，即使在美国加拿大，七八十年代也曾经出现过。多伦多 Bloor 街以北的 Annex 社区现在价值百万的普通房子可以在一九七八年以四万加元成交，如果说美元加元没有大幅度地贬值，我想谁也不会相信。}
\BookParagraph{}
\BookParagraph{那么将来货币是否就一定会贬值呢？固然历史不一定能代表未来，但凡是仔细研究过美国奥巴马政府的经济刺激方案和国会预算赤字的，都会对每年一万亿美金 {\BlissSymbols ①} 左右的赤字感慨万分。这样惊人的数字所代表的问题是深层次的，是不可能以提高税收、压缩开支或增发国债等正常手段解决的。它不是一个短期的问题，甚至不是美国一个国家的问题。唯一的出路就是让美金不断贬值，能够控制的也只不过是幅度的大小而已。近几年美联储实施的所谓数量宽松政策（Quantitative Easing, 包括QE1和QE2），正是它靠贬值来帮助购买国债的举措。这样再清晰简单不过的逻辑，使美金和其他世界主要货币将来的逐步贬值的过程变成了一个时间的问题，而不是「是否会」的问题（ when, not if ）。}
\BookParagraph{在一个货币贬值的年代，历史证明了持有黄金或白银的重要性。国民党时期用小金条和美元现钞{\BlissSymbols ②}做为流通的硬通货，就是市场的产物。如果目的不是「投资」（investment）而是「保值」（preservation of wealth）的话，我认为实物黄金是必然的选择。首先，实物黄金可以在大城市任何大银行轻易买卖，不存在房屋买卖时常见的流动性（liquidity）问题。其次，黄金体积很小，极其便于携带和储藏，十万美金的财富也就是两三公斤的金币而已。再次，黄金不会随时间变质，不需要管理，卖出时得到的增值部分（capital gain）不需要报税，留给孩子作为遗产不需要报遗产税。作为货币贬值环境中的保值手段（inflation hedge），实在是没有更好的选择。}
\BookParagraph{当然，任何事物都有两面性，买黄金有什么坏处呢？我想，首当其冲的问题就是黄金是不能「升值」的。其次，因为绝大多数黄金在各国的中央银行，政府买卖可以轻易操控市场，短期和中期内，黄金价值的风险偏大，一定要长期才能有保值的效果。再次，实物黄金的储藏不是很方便，保管不得当有丢失的风险。这些，确实都是值得仔细考虑的问题。}
\BookParagraph{}
\BookParagraph{很多人可能会想，我们家没钱买黄金，想都不用想这个有钱没处花的人考虑的问题。其实我认为，但凡有一定收入和储蓄，都应该持有一部分黄金，以备万一。一块一盎司的金币（比如 Royal Canadian Mint 发布的 Maple Leaf coins）一千二百加元 {\BlissSymbols ③}，绝大多数人能够负担得起。如果还是觉得太贵的话，不正是因为货币已经贬值了么？}
\BookParagraph{有一部分人已经大概意识到了黄金的重要性，但是他们买的是纽约股市上的黄金ETF（比如GLD），而非实物黄金。这样所谓\BookLineBreak{}「投资黄金」的做法，我认为不是很明智。黄金是不会「升值」的，所以所谓的「投资黄金」，不外乎是「投机黄金」，希望能短期\BookLineBreak{}（一两年内）黄金成交价大幅「升值」，而这恰恰是和黄金的优势背道而驰的。无论是 ETF、互助基金、pool account、还是 gold certificates，都有所谓的 counterparty risk，如果为其背书的公司出现问题，很难兑换成实物黄金。而持有黄金的作用，正是要在经济或社会出现重大问题的年代，保护个人和家庭的资产。如果连是否持有黄金本身这个问题本身都存在背书公司倒闭的风险，我认为是不值得的。}
\BookParagraph{我虽然主张购买一定数量的黄金，但我更希望有生之年永远用不着它。如果社会经济大动荡的年代真的出现了，历史证明黄金可以成为棋局中最后的一步棋，充当暴风雨中少有的一处避风港。\BookLineBreak{}电视剧《潜伏》中把小金条藏在鸡笼里，总有它的用武之地。}
\BookDate{2009 年 11 月，多伦多（初稿）}
\BookDate{2011 年 6 月，西雅图（再稿）}
\BookFootnote{{\BlissSymbols ①} 美国一万亿美元的财政赤字是 2009 年的统计数字，2011年大致是一万五千亿左右。}
\BookFootnote{{\BlissSymbols ②} 当时 35 美元可以自由兑换一盎司黄金。}
\BookFootnote{{\BlissSymbols ③} 1200 加元是 2009 年的价格，2011 年的价格大致是 1600 加元。}
\BookArticle{名　字}{名字}
\BookParagraph{也许是因为老了，近年来经常记不住别人的名字。这对于我这种到处开会认识各色人等的职业来说，基本属于致命的问题。特别羡慕克林顿那样的 photographic memory：见到别人毫不犹豫地直呼其名，那该有多亲切热络啊。更令人费解的是，我越费尽心机地想记住名字，甚至问人家是哪几个字，那名字竟忘得越快。}
\BookParagraph{不过反过来想，也许是因为中国人的名字善于用典，变化繁多，确实难记。单字的名字虽然好记些，可惜重名太多，只好尽量用生僻字。一个大学同学单名一个「翀」字，结果没人认识，早期的电脑中文字库又不支持，只好输入「羽中」或「冲」代替。无独有偶，同班还有个叫「喆」的同学，无法想象她吃了多少苦！}
\BookParagraph{僻字不敢擅用，只好剑走偏锋，用虽然好认好记但不常用在名字里的僻名。外公在《龙虫并雕斋琐语》里的《姓名》一文里说，他给自己起的单名一个「力」字，就属于当时的僻名，因为古人认为「王道以德，霸道以力」，无人用于名字。可惜他不知道解放后文化摇身一变，变成了时髦的好词儿，经常用于起名字。}
\BookParagraph{十几岁时看王安忆的小说，主人公单名一个「一」字，当时有眼前一亮的感觉，觉得这名字起得已臻化境，又简单又少见。没想到还是难逃厄运，姓氏笔画排序里，经常能见到叫「丁一」的高居榜首。看来，父母如果不肯狠下心来给自己的千金起名叫「李二」「张三」的话，就无法实现既好记好认又没有重名的境界。}
\BookDate{2013 年 11 月，多伦多}
\BookArticle{Apple Watch}{Apple Watch}
\BookParagraph{Apple Watch发布后众说纷纭的讨论之中，我最关心的问题是如果几年以后再回头看看，它到底是不是一个如同1984年的Macintosh、2001年的 iPod、2007年的 iPhone 那样划时代的产品。虽然谁也无法准确预测未来，我认为很有希望。iPhone 6 Plus的热销，证明了喜欢超大屏手机大有人在；而手机越大，手表的机会就越多。}
\BookParagraph{原因很简单：手机越大，有些场合中拿出来使用就越不方便。例子之一是苹果的 Apple Pay。它虽然支持用 iPhone 上的 Touch ID完成付款，但其实它对 Apple Watch 的支持才是真正的亮点。大屏手机一般需要放在包里，拿出来付款其实和拿出钱包一样不方便；而手腕上的 Apple Watch 就简洁明快地彻底解决了这个问题。}
\BookParagraph{例子之二是错过重要推送信息的问题。iPhone放在包里，自然会经常错过一些重要电话短信和邮件。而 Apple Watch 只需轻微震动，就完美地解决了这个问题。我认为 Apple Watch 之所以有用，极有可能和不宜使用大屏手机的各种工作、生活或社交场合息息相关。贴身的手表因为其天生的私密性，甚至比手机更加重要。}
\BookParagraph{因为 Apple Watch 在一千美金以下的定价比较低廉，很难成为那种因「越贵越要买」而趋之若鹜的奢侈品。长期来看，买它的人多半是因为有大屏手机无法满足的需求。2010年推出的 iPad，其实即使到现在，生活中真正的刚性需求并不明朗，所以销量开始下降。而手表不同，包括运动健身在内的几种刚性需求非常明确。}
\BookParagraph{手表设计最大的问题是，一个人最多只能同时戴一块手表，而戴什么表直接体现了他的品味。这也正是 Patek Philippe 和 A. \BookLineBreak{}Lange \& Söhne 等一线品牌的机械表两万美金起价还能热销的原因。苹果不算一家传统的制表公司，但在这一点上表现出了对制表业细节惊人的重视和驾驭能力，能基本满足最挑剔的机械表玩家。}
\BookParagraph{可以清晰地看出，无论软件硬件，苹果在 Apple Watch 的设计上已经竭尽全力。比如 Home screen 上的圆形图标，是直接效仿传统瑞士机械表中 complications 对应的圆形表盘设计。苹果还为 Apple Watch 的 UI 设计了基于 Helvetica Neue 的全新字体，这在其历史上极其罕见­——苹果上次亲自设计全新的字体还是二十多年前。}
\BookParagraph{这样看来，说不定多年以后，Apple Watch 可以成为一个比iPad更为重要的产品。人们会从 13" 笔记本电脑 ＋10" iPad + 4" iPhone 的传统组合逐渐演变成电脑 + 5–6" iPhone + Apple Watch的组合，因为后者更轻盈，更省事儿，也更私密。同时，由于手表是一种时尚，变化无常，更新周期短，利润可能会超过 iPad。}
\BookDate{2014 年 9 月，多伦多}
\BookArticle{买　房}{买房}
\BookParagraph{分答上回答了一个问题，关于年轻人是否应该买房。可是回答时间限制太严格，有点儿意犹未尽，想写一篇关于是否买房的说明文。中学时学说明文时，强调的是只陈述客观事实，而不发表主观论述和观点。说明文的影响力在于，「摆事实」以正视听之后，每个读者自然会形成自己的观点。}
\BookParagraph{是买房还是租房住，自然和有多少钱直接相关。如果没有足够的资金，当然应该租房住。可是钱有两种：一种是一笔能拿得出来的存款，一种是每月的收入。举个常见的例子，如果我能一笔拿得出一百万，每月收入一万五，是否应该买房？}
\BookParagraph{这个问题关系到买房有什么好处和坏处。房子对人有两种价值：一是提升生活品质，包括让心情更愉快。二是作为一种投资产品，将来可能会升值。}
\BookParagraph{只考虑生活品质的话，理论上说房子越大、地段越好，生活品质越高，心情越愉快。按国内一二线城市当前的租售比分析，每月同样的开销，很明显租房可以租到更大更好的房子。另外，租房的灵活性很好，随时可以把家搬到更好的地段，更大的房子，生活品质相应提升。可是心情是个难以捉摸的东西，很多人没买房之前心情很差，买了房心情非常愉快，感觉生活一下子好起来了。虽然没什么根据，但是对这些人来说，显然买一个小房比租一个大房生活品质更高，心情更好。}
\BookParagraph{但如果只考虑投资价值的话，租房虽然没有任何投资价值，但是因为租同样的房子的每月开销比买房更低，省下来的钱可以投资股票等其他投资产品。而跟其他投资产品相比，房地产属于不动产 （illiquid assets），买卖成本偏高，而且不是想卖就能立即出手。从投资成本来说，不适合短期投资，也不适合「集腋成裘」式的投资。另外，因为买房一般都需要贷款，相当于借钱投资，有杠杆效应，投资风险更高。比如如果首付一百万，贷款四百万，那么房价要是下跌百分之十，就相当于投资损失百分之五十。因为风险更高，房地产投资只适于三四十年的长期。这样长的时间里，大多数的房产可以至少做到保值，也就是说抵消了货币贬值的损失。五年甚至十年都属于短期投资，因为杠杆效应的存在，风险很高。相比而言，长期投资股票等产品是可以升值的，因为企业可以把资金用来提高生产力。}
\BookParagraph{但是不买房转而把资金投资其他产品也有一个很重要的坏处：这样的投资方式需要很强的自制力，每月主动把剩余的资金用于投资，以保证将来退休以后的生活品质。可惜这很不符合有钱就想花的人之常情，属于逆水行舟。买房相当于强制性的储蓄，对自制力稍差的人来说是一件好事。可即使强制储蓄，也要有个上限，大致应该是每月收入的三分之一。也就是说，如果每月收入一万五，应该最多拿出五千去还供房贷款。否则不仅降低了平时的生活品质，还因为杠杆效应，增加了风险。}
\BookParagraph{有人说，想这么多干什么，别人买我就买。这样的人适合买房，因为他们认为和别人不一样（contrarian）风险很高；有人认为不买房就找不到他心目中的女神或者男神，那么他应该买房，因为找到终生伴侣很重要；有人认为不买房就结不了婚，那么他该买房，因为婚不能不结。但是无论如何，至少应该明白买房带来的风险和相应失去的其他机会，明白历史不一定等同于未来，这样才能知己知彼，立于不败之地。}
\BookDate{2017 年 3 月，多伦多}
\BookArticle{咪　蒙}{咪蒙}
\BookParagraph{看到一则新闻：「咪蒙付费音频课程将上线，教你月薪五万，听课三年后加薪不超 50\% 可申请退款。」}
\BookParagraph{从咪蒙及其公司的角度出发，这又是一场策划周密而专业的市场营销。其志在必得的亮点在于「如果没成功可以退款」，非常吸引人。}
\BookParagraph{这让我想起了网络普及以前美国有线电视里全天播放的一个叫Home Shopping Network 的频道。思路是找两个能说会道的主持人推销各种生活用品，保证价廉物美，可以分期付款，而且至关重要的是没有风险，不喜欢的话可以全额退款。}
\BookParagraph{这个频道在美国一度非常成功，充分说明了美国人多么喜欢买可以退款又便宜的东西。}
\BookParagraph{有趣的是，咪蒙许诺的退款是一种「期货」：您要是不喜欢，不能马上退款，必须等到三年以后。}
\BookParagraph{很明显，如果每年投资回报率是3\%，即使所有学员都申请退款，咪蒙这个产品的营收都可以达到本金的 9.3\%。}
\BookParagraph{当然更好的办法是，三年以后公司申请破产倒闭。即使破产之后公司还有可以清算的资产，也是债权人和有投票权的股东优先，绝对不可能退给当年注册付费音频课程的这些学员。这属于违约风险（default risk），一个提供「知识付费产品」的新媒体公司，其违约风险可想而知。}
\BookParagraph{一边是一场策划周密而专业的营销，另一边是大量有理想、有情怀、希望为「月薪五万」而付费的有志青年。游戏的结局可想而知。}
\BookParagraph{明白违约风险的人，根本就不是该「知识付费」产品的目标客户群。}
\BookParagraph{就好像数学家不会去买彩票一样。}
\BookDate{2017 年 11 月，多伦多}
\BookArticle{遗　嘱}{遗嘱}
\BookParagraph{打完一场冰壶比赛，和几个老外一起喝酒聊天，聊起了「遗嘱」这个话题。}
\BookParagraph{「遗嘱」这个词儿最能体现中英文之间的区别了。中文给人一种沉重而不吉利的感觉，随时可以联想到小时候经常听到的「我国卓越的无产阶级革命家、思想家、教育家、诗人」、「因病医治无效」等等。而英文里相应的「will」则是中性词，没有任何悲伤的成分，也挺符合自然规律——无论早晚，总有这么一天。「It will happen.」}
\BookParagraph{所以如果在国内一帮同事一起去打完保龄球或唱完歌再奔街边烤肉店喝两扎啤酒聊一会儿，绝对不会提到「遗嘱」这个话题。}
\BookParagraph{一起聊天儿里一个老外是个会计师，讲了一个真实案例。一哥们儿立好的遗嘱锁在银行保险箱里，连律师授权书（power of attorney）都设好了。如果监护人（trustee）恰好也去世了，可以由律师代为执行遗嘱。可惜忘了把开保险箱的权力授予这个律师，结果银行不允许开保险箱。}
\BookParagraph{这个话题过于艰深，正好我不熟接不上话，只好说一句：「I will need to do some research.」}
\BookParagraph{二十世纪著名的数学家 Paul Erdos 说过，最完美的死法儿（the perfect death）应该是这样的：课堂上正在讲一个定理的证明，学生举手提问：「What about the general case?」教授回答道：「I will leave that to the next generation.」然后立即去世。}
\BookParagraph{我聊着聊着走神儿了想，最完美的死法儿应该是去世的时候没有财产，一辈子自己挣的所有财富都被自己花光了吧。}
\BookParagraph{当律师办完所有的手续，终于打开了那个保险箱时，里面是一封很多年前写好却一直未曾寄出的情书。}
\BookDate{2017 年 12 月，多伦多}
\BookArticle{勤俭节约}{勤俭节约}
\BookParagraph{从小到大受到的教育里，「勤俭节约」一向是一种美德。}
\BookParagraph{比如《小学生守则》里就有一条叫「生活俭朴，爱惜粮食，不挑吃穿，不乱花钱」。它虽然排在「热爱祖国，热爱人民，热爱中国共产党」之后，但是那种「一分钱掰成两半花」和「新三年旧三年缝缝补补又三年」的哲学，却自始自终刻骨铭心。}
\BookParagraph{现在一想，怎么就没有人举手提问：「为什么要存钱呢？」}
\BookParagraph{老师一定会回答：「存钱是为了等以后买房、投资、买比特币啊。」}
\BookParagraph{按理说，这个理儿讲到这儿就差不多了。大多数人也是这么身体力行的：省吃俭用，买车买房，如果还有闲钱去「理财」，投资点儿黄金美股比特币什么的。}
\BookParagraph{但关键的问题在于，拿省下来的钱去「理财」干嘛？}
\BookParagraph{老师一定会说：「理财是为了将来花啊。」}
\BookParagraph{这个回答里暗含了一个假设：钱将来花比现在花要强。}
\BookParagraph{可惜这个假设是错误的，反过来倒是多半成立。首先，钱是会贬值的，现在一百块可以买三杯咖啡，三十年以后也许只能买一杯了。其次，人是会逐渐衰老的，年轻时花钱玩儿过山车是享受，年纪大了说不定就是折磨。年轻时一定要喝上好的葡萄酒，年纪大了喝什么酒已经没什么区别了。}
\BookParagraph{而且无论贫富，永远有这样一种可能：快要离开这个世界的时候，发现钱还远远没有花完。}
\BookParagraph{所以正确的做法是，每年的收入全部花掉，还完房贷之后，银行的存款永远是负数，就像很多外国人。到了退休的时候要是缺钱花，再把已经基本付清的房产用「反向房贷」（reverse mortgage）抵押给银行提现。}
\BookParagraph{这种「做一天和尚撞一天钟」的哲学有诸多好处。没有机会积累大量的积蓄，就没有那种跟中了彩票似的挥霍无度的风险；钱花掉了，自然不用担心它将来贬值，也不用担心比特币大跌是否需要「清仓」；孩子没有继承财产的可能，更容易形成独立自主、自力更生的性格；在年轻时充分享受生活，爱过恨过，吃香的喝辣的，走遍大江南北。}
\BookParagraph{我有一个理想：将来有一天回首浮生往事，如果能再重来一次，该怎么玩儿怎么折腾，一切的选择，一如当年。}
\BookDate{2018 年 1 月，多伦多}
\BookArticle{敬　酒}{敬酒}
\BookParagraph{周末去开会，见识了一回上千人的晚宴上的中国特色：敬酒。}
\BookParagraph{敬酒和给客人夹菜一样，是酒席上历史悠久的优良传统。我甚至怀疑，粤菜酒席上的公筷，最初是专门为了给客人夹菜而设。可惜夹菜这个传统随着能旋转的桌子这一发明的普及，日渐式微：一个能坐得下二十人的大桌，就算您想费心给客人夹菜，也力不从心了。}
\BookParagraph{而敬酒则不然。夹菜够不着的贵宾和领导，我们敬酒时大可放心，只需离席而立，再走上几步，就可以满面春风地出现在领导的面前。不仅如此，我们还可以离开自己的酒桌，长途跋涉，不辞辛劳地到其他酒桌上去敬酒，机动性不逊一支机械化步兵师。}
\BookParagraph{若是上百桌的酒席，大型吊灯的金碧辉煌之下，无数天南海北的嘉宾互相敬酒，无论初次相识还是忘年之交，棋逢对手，捉对交锋。那你来我往的场面，不要说是老外，即便是我，震撼之程度都堪比参加了一场人民广场的大型相亲活动。}
\BookParagraph{只是可惜了刚上桌的一桌好菜。据我观察，开始上菜之后不会超过五分钟，敬酒的环节就一定会蹒跚登场。刚开始是涓涓细流，然后逐渐湍急起来，很快就如瀑布般震耳欲聋。无论是敬酒的还是被敬的，是学生还是老师，是初出茅庐还是久经沙场，都要在极其有限的时间内敬遍全场，哪儿还有时间品尝那道刚刚出笼的清蒸鲈鱼？}
\BookParagraph{敬酒的环节就像一道博弈论的习题：即使您想继续吃您的鲈鱼，只要别人都在敬酒，您的最优策略一定是立刻加入战团。按照中国人尊老爱幼的规矩，博士生要给优博获奖者敬，优博要给优青敬，优青要给杰青敬，杰青要给长江敬，以此类推，举一反三。}
\BookParagraph{不过敬酒的辞令似乎有所退步。缺了一点儿「问世间，情为何物？直教生死相许」的柔情，「李白斗酒诗百篇」的气势，或是\BookLineBreak{}「高谈雄辩惊四筵」的豪迈，取而代之的是「加个微信吧」，还有\BookLineBreak{}「你扫我还是我扫你？」}
\BookParagraph{料想得到，若是《饮食男女》中的老朱看到他那道绝活儿「龙凤呈祥」遇到这般冷落，一定会在一边暗暗地摇头叹息吧。}
\BookDate{2018 年 5 月，西安}
\BookArticle{比特币 （一）}{比特币 （一）}
\BookParagraph{大家如何看待「比特币」等虚拟货币？（问题来自「向上移动」）}
\BookParagraph{买股票也好，债券也好，房子也好，虚拟货币也好，花钱买一样儿东西如果目的是期望将来它升值，不一定都是一种投资（investment），也可能是一种投机（speculation）。}
\BookParagraph{投资和投机的一个重要区别在于，投资看重的是长期，所谓\BookLineBreak{}「放长线钓大鱼」；而投机看重的是短期，所谓「捞一把赶紧撤」。}
\BookParagraph{比如如果您觉得某场维密秀的门票或下一代的iPhone会火，于是去买来准备转手卖了，就是一种「投机」。}
\BookParagraph{大多数投机都是有风险的。下一代iPhone很可能供货充足，转手卖的价格比入手价格还要低。为什么说是大多数呢？因为确实有的投机基本上是稳赚不赔，比如法拉利或保时捷的某几款跑车。可惜的是，法拉利的那几款跑车您肯定买不到，因为很多人都排队等着呢，稳赚不赔正是因为它的稀缺性。听说那场维密秀的门票也并不好买。}
\BookParagraph{买比特币属于一种投机。而且比特币一点儿也不稀缺，随时可以买到。}
\BookParagraph{哎，不是都说比特币的总量有上限，这不是正说明了它的稀缺性吗？}
\BookParagraph{按照前几天的交易价格，比特币的总量大概折合一千八百亿美元，按照一辆法拉利跑车五十万美元计算，相当于三十六万辆法拉利跑车——这好像也不是特别稀缺：特斯拉自从十几年前创立以来，一共卖出了大约二十三万辆，可绝对不会有人认为倒卖特斯拉能赚钱。}
\BookParagraph{那为什么今年一年比特币交易价格大涨了十几倍呢？}
\BookParagraph{这是因为人都有一个弱点，就是喜欢「跟风儿」，别人买什么我也要买。所谓「传销」、「庞氏骗局」，都是利用这个弱点投机。如果您今年上半年买了比特币，那么恭喜您深刻地理解了人性的这一弱点，投机圆满成功。}
\BookParagraph{可是问题在于，同样的弱点还会有另外一种表现形式：别人都卖什么我也要赶紧卖掉。}
\BookParagraph{十七世纪的荷兰炒郁金香，十九世纪英国炒铁路公司的股票，二十世纪初美国佛罗里达炒房，1929年的黑色星期二，1987年的黑色星期一，八十年代末的日本，九十年代末的纳斯达克。历史像是在机械地重复，泡沫永远以「这次不一样了」（It’s different this time）开场，而却是完全一样的结局。}
\BookParagraph{不过您还别说，比特币这次还真有点儿跟以前不一样。投机比特币有一种投机其他东西没有的风险：比特币不容易保管。}
\BookParagraph{首先绝对不能让某个「云」里的比特币交易公司托管。如果因为该公司的安全漏洞，您的比特币让小偷给偷走了，公司一定早就跟您签署了「概不负责」的协议。}
\BookParagraph{wired.com前一阵子发表了一篇长文，讲的是一个计算机高手买比特币的故事。他2016年1月专门研究了很长时间，认为保存比特币最稳妥的办法是买一个 Trezor。Trezor 是一个像U盾一样的硬件加密钱包（hardware wallet），可以把比特币存在里面，用密码保护起来。他把密码写在一张纸条上，存在办公室的抽屉里。一次出门度假收拾东西时，他把那张纸条带回家放在女儿卧室的枕头底下。万一飞机失事，女儿还可以用这个密码取出比特币。}
\BookParagraph{度假回来，他惊恐地发现枕头底下的纸条不翼而飞，原来是清洁工打扫卫生时把纸条给扔了。他试图回忆密码，可是想起来的所有可能的密码都是错的，而且每输错一次密码，Trezor都会等加倍的时间才能允许下一次输入。他为了回忆正确的密码，一年多来茶不思饭不想，晚上甚至睡不着觉。直到今年8月，他终于以0.85比特币的代价，雇了英国一位15岁的黑客，利用 Trezor 1.4 版固件的安全漏洞，成功地破解了他自己的钱包，找回了他的密码和比特币。}
\BookParagraph{所以啊，就算您是一位杰出的计算机科学家，熟读比特币交易所依赖的「区块链」和「拜占庭一致性」理论，要是您想投机比特币，最好先想好一个万全之策，记住您的密码。}
\BookParagraph{如果是我，大概还是去排长队买法拉利的跑车或者维密秀的门票更为靠谱。}
\BookDate{2017 年 12 月，多伦多}
\BookArticle{比特币 （二）}{比特币 （二）}
\BookParagraph{李老师，前两天您在社区回答了关于比特币的看法，认为它是完全的投机不建议参与。想问下您怎么看待区块链技术和基于区块链技术开发的应用而产生的代币，如以太坊，eos和一些山寨币，您觉得他们都没有实际的价值么？（问题来自「向上移动」）}
\BookParagraph{显然，无论是比特币还是以太坊（Ethereum），都是有实际价值的。如果你有数量足够的比特币或其他数字货币，理论上来说你可以用它去买法拉利的跑车。}
\BookParagraph{但我想你更关心的问题其实是，比特币和以太坊等其他数字货币，到底是否有投资价值，也就是说它们长期来看（五至十年）是否会升值。}
\BookParagraph{我在前一篇回答里盖棺定论，把买比特币定性为「投机」（speculation），你也许认为过于武断。}
\BookParagraph{2009年，比特币最初的设计目标是作为一种数字货币用于电子支付，也就是说当钱花用来买东西。它的创始人 Satoshi Nakamoto设计的区块链技术从学术角度看是一种优美而杰出的设计，重点是完全不需要信任任何监管体系，而是依赖公开记录并验证所有交易。}
\BookParagraph{但是比特币最初实现区块链技术用到的算法中出现了一个问题。它处理交易的吞吐量（transaction processing capacity）不高，形成了性能上的瓶颈。因为大量拥堵，造成了交易费用居高不下。从这一点来看，比特币作为货币用于支付的目标基本已经失败，很难进一步推广。}
\BookParagraph{这也是包括以太坊在内的其他数字货币如雨后春笋般地不断涌现的原因。这些新型的数字货币的交易处理吞吐量都比比特币高得多，那是否可以将来取代比特币作为主流数字货币实现电子支付呢？}
\BookParagraph{我认为希望十分渺茫。}
\BookParagraph{首先，如果其他数字货币开始流行，对持有比特币的人不利，这些人一定会反对推广其他货币。他们持有比特币的目的早就不是为了要拿来买东西，而是用来保值（store of value），甚至升值。其实，这个目的和最开始设计比特币的初衷是背道而驰的：如果作为货币用来交易，需要的是流通甚至通过增发略微贬值，而不是升值。黄金可以用来保值，但却很难用来作为货币买东西。}
\BookParagraph{其次，传统货币用于电子支付或转账，现在早已司空见惯。别说改用比特币或以太坊，就算是让您从微信支付宝改成其他支付平台，多半儿您都不乐意。}
\BookParagraph{好吧，用来电子支付买东西是没希望了。那是否可以跟买股票一样用来长期投资期望它升值呢？}
\BookParagraph{这相当于期望越来越多的人会买入比特币或其他数字货币。而买入的原因不是为了电子支付买东西，而是因为大家都期望它长期可以升值，是一种优秀的投资方式。}
\BookParagraph{您多半儿听说了，比特币历史上曾经创造过一年上涨50倍的辉煌，也曾经一年跌过60\%。就像货币历史专家、加州大学伯克利分校的经济学教授Barry Eichengreen说的：「Expectations can change.」}
\BookParagraph{而「区块链」并不能增加比特币长期升值的概率，因为绝大多数人不明白它是怎么回事儿。}
\BookParagraph{如果您听说的都是「投资」比特币成功的故事，那是因为这是一种「有偏采样」：投资失败了人家不会告诉你。}
\BookParagraph{不过投机是一种娱乐，只要您充分意识到了其风险，是非常刺激好玩儿的。}
\BookParagraph{祝您过山车玩儿得开心。}
\BookDate{2017 年 12 月，多伦多}
\BookArticle{奢侈品}{奢侈品}
\BookParagraph{你买过什么普通价格、但是能很大程度提升幸福感的东西？\BookLineBreak{}（问题来自「向上移动」）}
\BookParagraph{就像七十年代的手表、自行车和缝纫机一样，我也想到了三样东西可以提升幸福感：Tylenol、洗碗机、笔记本。}
\BookParagraph{Tylenol 中文好像叫泰诺，主要的西药成分是对乙酰氨基酚，俗称扑热息痛。它价格低廉、老少咸宜，可以在吃药半小时后开始有效缓解病毒性感冒所伴随的高烧头疼症状。那种迷迷糊糊之间浑身出汗高烧迅速消退的感觉，无异于「久旱逢甘露，他乡遇故知」，幸福感远超登机时发现升到了公务舱。扑热息痛对孩子的疗效尤其显著，副作用甚微，非常安全。可惜竟然有很多人还不了解它的好处。}
\BookParagraph{洗碗机和烘干机一样，是国外家庭的标配，可惜国内并不常见。洗碗机十分彻底地把「洗碗」这项家务劳动一笔勾销，也同时根除了家里类似「我忙一天了，你怎么不洗碗」的争执。洗碗机洗碗可以高温消毒，保证每次洗完碗筷都像新的一样。同时它效率很高，用水量大概是手洗的五分之一。就像 Tylenol 一样，每次用洗碗机，我都感慨于它的发明者对人类做出的贡献之伟大。}
\BookParagraph{在这电子产品至上的年代，一个小小的笔记本竟也成了一种精神上提升幸福感的奢侈品。时间长了我发现，无论是朋友圈还是电子书，无论是虚拟现实还是无人机，很多东西如同过眼云烟，看过玩儿过就忘记了。而一支笔、一个笔记本给予我的，是一种难得的慢节奏，像一支慢三拍的交谊舞，简单、快乐、自然。在笔记本上写字是一种冥想，它有一种魔力，可以让心情安静下来，把瞬时的火花变成永恒的经典。它会让我感到幸福。}
\BookDate{2017 年 12 月，多伦多}
\BookArticle{保　险}{保险}
\BookParagraph{最近听朋友聊到买香港保险的事，打算买重疾和理财保险产品，想问问社区里有没有懂这行的，香港保险靠不靠谱？（问题来自「向上移动」）}
\BookParagraph{前一阵子去香港开会，一个大学的教授请我去一个挺高档的馆子吃饭。他妻子也参加了，她正好在一家保险公司供职，大概是推销保险产品的吧，简称「理财顾问」。据说她收入颇丰，远远超过大学里的教职。}
\BookParagraph{三句话不离本行，我们就聊了一会儿保险产品。她给我介绍说，她们公司的保险产品，可以保证每年百分之六的回报率。}
\BookParagraph{「那要是出现像零八年那会儿的经济危机呢？」}
\BookParagraph{「那也能保证百分之六的回报。」}
\BookParagraph{我礼节性地笑着，没有再继续聊这个话题。据说她的客户大多是国内来的。}
\BookParagraph{任何投资产品，但凡提到「保证」什么超过定期存款的回报率，都会十分可疑。这是因为市场是无法保证短期投资收益的，潜在的回报越高风险越大。既然是投资，就要清晰地认识并接受其风险。这就是所谓的「风险和机遇并存」。}
\BookParagraph{香港的保险业，往往是用这种「保证回报」的金字招牌作为其最大的亮点来推销其产品。而买这类保险产品的客户，则常常正是被如此「物美价廉」的「投资」产品所吸引，竟然忘记了购买一份保险产品的初衷。}
\BookParagraph{保险最基本的职能，其实是在投保人出现意外死亡（或者重大疾病）时，保险公司给家属可以赔付一笔免税的赔偿金。}
\BookParagraph{而实现这个职能最简单的险种叫定期人寿保险（term life insurance）。举个例子，如果一个定期保险一年交一千美元，在六十岁过期。投保人六十岁前去世赔付五十万美元，但是如果六十岁前没有去世，则所有的保费就白交给保险公司了。这种保险产品因为比较简单易懂，在当今十分成熟的保险业，保费会相对低廉，甚至会包括头十年或终身保费不会上涨的条款。相对而言，保险公司的利润空间也不会太大。}
\BookParagraph{而往往需要「理财顾问」推销的保险产品，基本上都属于「终身保险」（universal life insurance）。表面上看，这种保险产品兼具定期保险和投资两种职能：如果意外去世可以赔付一定金额，但没有去世则可以拿回大部分「投资」的本金及其分红——「best of both worlds」。}
\BookParagraph{但是不要忘了，但凡什么东西听起来非常诱人，往往一定暗藏玄机，英文叫「too good to be true」。那具体来说这「玄机」在哪里呢？对不起，这可是「可口可乐秘方」，保险公司是不会明确告诉你的。比如如果你提前终止保险合同，有可能大部分的「投资」无法收回，损失惨重。而这种「只能买入不能变现」直接影响到投资的「流动性」（liquidity），是投资的大忌。}
\BookParagraph{我个人认为，要是想买保险就买简单明了的定期保险，要是想投资就专门去投资，两者最好不要混为一谈。实在要买，也建议花几天时间仔细研读合同里的条款（fine print），尽量找到其「玄机」所在。而无论保险还是投资，一个好的产品往往即使「待字闺中」依然炙手可热，不需要推销，更不需要给推销员分红。}
\BookParagraph{一个最好的例子是美国著名的基金产品之一，西蒙斯创立的Renaissance Technologies 旗下的 Medallion Fund。该基金自从1993年以来只对本公司员工开放，二十多年来每年平均回报70\%以上。}
\BookParagraph{这不仅仅是「待字闺中」，简直就是「闭门谢客」了吧。}
\BookDate{2017 年 12 月，多伦多}
\BookArticle{投　资}{投资}
\BookParagraph{30岁左右，买完了房子和车子之后，还有一点儿小钱应该怎么理财？（问题来自「向上移动」）}
\BookParagraph{六个月的「向上移动」社区快要结束了。我最后写一篇「收官之作」，聊聊关于投资的一些常识性的建议，以此祝大家新年快乐。}
\BookParagraph{首先，选择投资方式和产品之前，需要充分想明白自己到底能承受多大的风险。要是拿自己的退休金去投资，可能三十年以后才需要变现，对风险的承受能力肯定比拿几年后就要花的钱去投资要强。本质上来说，投资其实就是一个在风险和回报之间寻求一种动态平衡的游戏。}
\BookParagraph{那么让我们来看看这个「游戏」普通人应该怎么来玩儿。}
\BookParagraph{首先，「不能把鸡蛋都放在一个篮子里」。这就需要考虑资产配置（asset allocation）的问题。如果你的资金是一块蛋糕，资产配置就是如何切这块蛋糕。投资分两大类：股票（equity）和债券等固定收益产品（fixed income）。股票风险高，固定收益投资则风险较低。所以，如果风险承受能力较低（比如临近退休），这块蛋糕应该在固定收益投资上多分一些。}
\BookParagraph{短期债券（short-term bonds）等固定收益投资产品的表现一般与股票市场是「负相关」的，可以有效地降低风险，也不会损失太多回报。这就是为什么无论风险承受能力高低，都建议分一部分「蛋糕」给债券投资。}
\BookParagraph{具体到股票部分的投资，应该进一步细分到不同的地理区域\BookLineBreak{}（甚至产业）。这样可以进一步平衡不同地域之间的风险，所谓\BookLineBreak{}「东方不亮西方亮」。}
\BookParagraph{比如Canada Pension Plan（CPP）是世界十大退休储蓄计划之一，管理着三千多亿加元的资产，其中55\%投资股票。这55\%中，只有3\%投资加拿大的公司，28\%投资了美国和世界其他发达国家——所谓EAFE（欧洲、亚太、远东），6\%投资中国、俄罗斯、巴西、印度等新兴市场（emerging markets），还有18\%投资了私募股权（private equity）。}
\BookParagraph{这也太复杂了吧？普通人到哪里去买这么多地方的股票啊？}
\BookParagraph{ETF登上舞台的时机到了。}
\BookParagraph{ETF（exchange traded funds）的中文好像叫交易所交易基金，它有三大好处。其一是可以直接在股票交易市场中随时买卖，流动性（liquidity）非常好。其二是它的管理费用（Management Expense Ratio）比一般的基金低很多。比如美国有一种叫Vanguard Total Stock Market的ETF，管理费用只有0.04\%。也就是说，十万美元的投资，每年只收取40美元的费用用于管理和交易。因为投资可能是几十年的长期，费用高低明显可以影响投资回报率。}
\BookParagraph{在美国加拿大等股票市场里，ETF的最后一个好处是种类繁多，琳琅满目，选择余地非常大。仅仅美国Vanguard公司，就有七十多种。想买世界上哪部分市场（或产业）的股票或者债券，买相应的ETF即可。在加拿大买ETF，甚至连交易费用都是免费的。中国市场的 ETF 产品虽然还不成熟，但至少可以买比如美国标准普尔指数等常见的 ETF 种类。}
\BookParagraph{那蛋糕切好了，股票和债券按季度分红，是不是可以「躺着数钱」了呢？}
\BookParagraph{差不离儿了，但是还有一个小常识不要忘了。股票和债券的分红（dividend），一定要尽快投资，这样才能有「利滚利」（compound interest）的效果，极其重要。很多股票交易经纪公司（brokerage）支持所谓的「Dividend Reinvestment Plan」，简称 DRIP，可以自动将分红用来买同一种股票，不需要交易费用。即使没有 DRIP，也可以定时手工用红利买入股票或债券的 ETF。}
\BookParagraph{只要保证分红再投资，再定时根据风险承受能力重新调整一下「蛋糕」的分法儿，普通人其实也可以长期管理自己的投资，完全不用花时间和精力去研究全球经济问题。}
\BookParagraph{不是都说该买对冲基金比特币或者什么最新最时髦的投资产品吗？}
\BookParagraph{巴菲特在 2008 年 1 月跟一位叫 Ted Seides 的对冲基金经理打了一个一百万美金的豪赌。巴菲特认为，截止 2017 年 12 月的十年时间内，投资五种对冲基金产品再扣除管理费用的收益，不如直接投资美国标准普尔指数。}
\BookParagraph{他的立场是，对冲基金的算法即使再天才，相比其高昂的管理费用，长期来看还是不合算。}
\BookParagraph{结果最近揭晓：巴菲特赢了。}
\BookDate{2017 年 12 月，多伦多}
\BookArticle{雷克萨斯}{雷克萨斯}
\BookParagraph{李老师，您好。刚看到您的最新微博，说喜欢汽车。可以说说雷克萨斯吗？褚老师曾经说过，所谓「高端」汽车品牌里他最没兴趣买的是雷克萨斯。（问题来自微博私信）}
\BookParagraph{挑选汽车就像挑选葡萄酒一样，理论上说不能以品牌一概而论，而需要考虑具体的型号和出产年份。但实际上大多数人买车主要看的是品牌，所以品牌价值是汽车厂商的兵家必争之地。}
\BookParagraph{丰田公司八九年创立雷克萨斯这个品牌，是为了推销LS400。当年的LS400可不是一般的豪华轿车：它是专门为美国市场设计的，采用4.0升8缸发动机，还有豪华车标准的纵置后驱结构，无论是做工的精益求精还是驾驶舱的安静程度，都可以「秒杀」当年的奔驰和宝马。当年连比尔 · 盖茨都买了一辆LS400，并且连续开了五年。}
\BookParagraph{丰田也是很早进入中国市场的品牌之一，而且当年在中国市场对品牌价值的树立非常重视。它八十年代初「车到山前必有路，有路必有丰田车」的广告词，可以说和九十年代初「摩托罗拉寻呼机，随时随地传信息」一样，几乎代表了一个时代。国内汽车广告里，大概只有《蜗居》里宋思明的那句「开车的男人，有血性的，都希望拥有一辆路虎」可以和它媲美，不过那已经是零九年左右的事儿了。}
\BookParagraph{可惜的是，丰田对雷克萨斯这个品牌的核心价值一直举棋不定，让人无可适从。最大的错误，就是当年把好好的「凌志」改成了\BookLineBreak{}「雷克萨斯」这个名字。「凌志」这个名字好就好在它就像「奔驰」、「宝马」这些名字一样，有一种灵动的感觉（soul），让人浮想联翩，令人心旷神怡，耳目一新。而「雷克萨斯」这个名字则普通、平庸，没有令人悠然神往的那种风骨。}
\BookParagraph{雷克萨斯的轿车虽然还是和当年一样，纵置后驱、做工精良、质量稳定、行驶安静，但这些和它的品牌形象一样，过于朴实无华，没有任何让人拍案叫绝之处。}
\BookParagraph{但是雷克萨斯最大的败笔，还是它的纺锤形格栅。这个设计的初衷是希望一改以前老气横秋的气质，能够用一种激进的设计，让年轻人喜欢。可惜这就像一个平时生活朴素的女生，忽然开始浓妆艳抹，像换了一个人，只能让人觉得莫名其妙，避而远之。更何况汽车的设计最重要的是设计元素几十年如一日的传承，比如二战期间吉普车的造型，一直沿用至今。设计风格改来改去，反而有一种「暴发户」的范儿。它多了一番算计，却少了一份优雅。}
\BookParagraph{这就是为什么《蜗居》里宋思明开的车绝对不可能是雷克萨斯的原因。}
\BookDate{2018 年 4 月，多伦多}
\BookArticle{城　市}{城市}
\BookParagraph{李老师好！请问您如何看待生活城市的选择呢？对于年轻人来说，北上广深这样的一线城市发展机会仿佛更多，但生活成本往往特别高，买房买车对于普通人来说几乎是不可能的事儿，但回到二三四线城市又有些不甘。您怎么看？（问题来自微博私信）}
\BookParagraph{我前几天去重庆的一所大学访问。傍晚时分，一行三人驱车前往南山吃晚饭。一路上全程高速，高速公路双向八车道，是崭新的。路过横跨长江的寸滩长江大桥，也是崭新的，去年年底才建成通车。在南山上的一个路口，我们走错了路，看到了路边一排排建筑设计风格很有品味的联排房屋，还有很多新建的咖啡店和特色餐厅。}
\BookParagraph{重庆我以前只去过两次，对我来说是一个完全陌生的城市。可是在南山偶遇这么一片颇具特色的小区，我竟然会想，其实在这里定居也挺美的。}
\BookParagraph{西安也是这样一个城市。我十年前第一次去西安的时候，是因为一个偶然的机会去讲课。那时的西安和前几天的重庆一样，整个城市我一个人也不认识。我那时曾经从一个叫小寨的地方沿着长安中路步行北上，通过城墙，直达钟楼，一路看到老人们围在一起打牌，还有时尚男女出入各种名牌商店。也曾经路过这个城市南郊规划的新区，感慨于那里市容环境的整齐划一。这些年来多次去西安，虽然还是不认识很多人，但是却一直觉得这个城市很熟悉，很亲切。它和世界上任何城市一样，会有其长处和短处，但如果在这里定居，我绝对不会有什么遗憾。}
\BookParagraph{当然，北京上海这样的城市，无论从工作还是生活的哪一方面来说，都有其显而易见的优势。但是，它们的问题也同样明显：人口密集，交通拥堵，房价居高不下，生活压力太大。就好像纽约曼哈顿当然很好，但是如果收入不够高，住在那里捉襟见肘，生活品质不但没有提高，反而会大幅度降低，得不偿失。而和生活品质直接相关的因素，比如子女教育和居住环境，国内很多城市和北京上海深圳相比，完全可以分庭抗礼。}
\BookParagraph{我选择城市的哲学是，与其倾其所有去一线城市贷款买房后过节俭的生活，不如在二线城市好好找一个环境宜人、格调高雅的地方定居。最幸福的日子是这么定义的：房价虽然贵但是还没到贵得离谱的程度，自己上班不需要太长时间，孩子上学也在附近，手头虽不阔绰但是也有点儿富余，生活丰俭由人。}
\BookParagraph{要是想念一线了怎么办？随时订个房间，到高铁站买张票，北京国贸也好，上海法租界也罢，小憩几天。}
\BookParagraph{放心吧，时间长了，无论你在哪里，都一定会适应、习惯、进而爱上这座城市的。}
\BookDate{2018 年 6 月，西安}
\BookArticle{吃}{吃}
\BookParagraph{李老师好！想起您以前说过您想念姥姥做的泡菜，用坛子泡的切段黄瓜，让人印象深刻。能聊聊吃的吗，您心目中最完美的一顿晚饭是什么样子？该有哪些菜？（问题来自微博私信）}
\BookParagraph{我经常坐长途飞机，飞机上看过无数的电影。但是我发现，这些电影看的时候真切而投入，有的惊心动魄，有的唏嘘不已，可一下飞机，竟然很多连名字都记不起来了。而记忆犹新的电影，很多都是年轻时专程去看的，有的甚至连台词儿都能背下来几句。}
\BookParagraph{吃对于我来说，也是如此。最好吃的菜里排名第一的，那绝对是清华东门外那个小饭馆里的酸辣土豆丝了。上大学时一个月的伙食费也就差不多一百块，食堂一块五的小炒，称得上是奢侈品。有一次和郭芳去清华东门下馆子，酸辣土豆丝六块八一盘儿。吃完了她意犹未尽，我当场拍板又点了一盘儿，因此一战成名。}
\BookParagraph{所以说，吃得不过瘾到能再点一盘儿一模一样的，才是吃里的真爱。无独有偶，十几年前到武汉游玩，郭芳的一个同学在一个小馆子请客。他看到一盘儿爆炒野鳝我特别爱吃，当机立断，叫服务员再上了一盘儿一模一样的。大快朵颐之时，这位同学的形象立马高大起来。从那以后我每次想起他，就有一种莫名的好感，像个多年的朋友。}
\BookParagraph{吃有两种。一种是饭局，不怎么认识的人通过吃来结交，在觥筹交错之间，慢慢熟络起来。饭局上往往是高朋满座，包间儿里灯火通明地开好几桌儿，席间有舞狮子的，有变脸的，有弹古筝的，不一而足。像钱钟书笔下所说：「吃饭时要有音乐，还不够，就有“佳人”、“丽人”之类来劝酒；文雅点就开什么销寒会、销夏会，在席上传观法书名画；甚至赏花游山，把自然名胜来下饭。」}
\BookParagraph{饭局里从饭前的开胃小菜到饭后的主食，各有特色，百花齐放，百家争鸣。好吃的菜太多，用情不专，反而哪个菜都没有什么印象。而且圆桌大得令人生畏，和左右两人有一搭没一搭地「尬聊」，虽然湖光山色，佳人伴唱，酒过三巡之时，盛情难却之下，还是犹如长途飞机上的电影，少了一分真情实感。}
\BookParagraph{而我更喜欢的一种吃是三两老友平常见不着，好不容易有机会聚到一起，找个小馆子叙旧。西餐早就领悟了这种「小馆子」的哲学，越是正宗的西餐厅桌子越少，每个桌子也小得可怜——人家压根儿就没打算让你一大帮子人一起来吃饭。}
\BookParagraph{「小馆子哲学」的中心思想是菜单上没几道菜，甚至不用点菜，厨师想做点儿什么就吃什么，所谓「today’s special」。菜吃得用情专一，吃鱼就不能吃牛排，吃色拉就不能喝汤。而桌子是长方的，保证无论坐在旁边儿还是对面儿，都能随时华山论剑，火热开聊。聊到high了，服务员正好不失时机地礼貌打断，奉上甜点的菜单。}
\BookParagraph{当然，相比西餐厅，我更常去的是街头巷尾俯首皆是的小馆子。长桌上铺着塑料的桌布，三两好友一起，点了爽口木耳就点不起夫妻肺片，菜单上没有西湖醋鱼，也没有宋嫂鱼羹，但是点上一碗炸酱面，也照样其乐融融。华灯初上，熙熙攘攘，如果这时再有六块八的酸辣土豆丝，再有姥姥亲手用长筷从坛子里夹出来的泡菜，就此生无憾了。}
\BookDate{2018 年 7 月，多伦多}
\BookArticle{穷爸爸}{穷爸爸}
\BookParagraph{李老师您好，无意中看到您给别人的解答，蛮有感触。我有个工作抉择的问题想请教您：一个是工作轻松，朝九晚五，包吃，五险，月薪偏低，有年终奖；另一个是主办，忙一点儿，工资高两千，有年终奖，五险一金，包吃，双休。相对来说第二个可以学习锻炼的更多，但第一个更稳定轻松。今年我二十八了，年纪也很尴尬，我不知道如何抉择。（问题来自微博私信）}
\BookParagraph{为了全面地回答你的问题，我又重读了 Robert Kiyosaki 写的一本书，叫《富爸爸，穷爸爸》（Rich Dad, Poor Dad）。}
\BookParagraph{这本书分了八章，但是其实最重要的结论只有三个。}
\BookParagraph{1. Your house is a liability, not an asset.}
\BookParagraph{你的房子为什么不是你最重要的资产（asset）呢？因为资产的定义非常明确，是不需要工作就可以给你带来现金收入的。绝大多数人的房子包括抵押贷款，不但无法通过出租的方式带来现金收入，每月还需要偿还贷款，是一种现金支出。如果是在国外，即使你的房子全款付清，但仅仅是放在那儿空着（比如cottage \BookLineBreak{}homes），因为需要交纳房产税，也一样不是一种资产，而是一种负债（liability）。只有买来就是为了出租的房产，才是一种资产。}
\BookParagraph{你大概会问，买房子和我选择哪个工作有什么关系？}
\BookParagraph{太有关系了。因为大多数人工作的目的就是为了挣钱，挣到了一笔钱就用来付房产的首付。将来工作每月的大部分收入，都用来偿还贷款和日常生活开销。一辈子从来没有拥有过给他带来现金收入的资产，所以需要一直工作，直到退休。}
\BookParagraph{他们从来没有想过，如何正确地投资，让钱更好的「为自己工作」（let money work hard）。}
\BookParagraph{从这个角度分析你的问题，我们假设你认为第一个工作的工资收入足够维持衣食住行的日常开销，那么应该选第二个工作，但同时多余的两千元不是像「月光族」一样花掉，也不是存在银行里，而是拿来投资。这时你年轻的优势就体现出来了：和一位三十八岁的投资者比起来，同样的钱，你的投资带来的长期收益会高很多，这是由投资是「利滚利」（compound interest）决定的。}
\BookParagraph{当然了，这里的「投资」不包括拿钱去买房子，因为房子不是你的资产，而是你的负债。}
\BookParagraph{2. Work to learn — don’t work for money.}
\BookParagraph{你听说过无线通信理论里有一个常用的方法叫「spread spectrum」吗？}
\BookParagraph{这个方法的核心要义是，我们不用很窄的频段去发送信号，而是把要发送的信号分到更宽的频段。这样可以有效地减少各种干扰，对抗各种偷听，还可以支持很多人同时一起通信（multiple access）。}
\BookParagraph{我们耳熟能详的WiFi，就用到了「spread spectrum」这个技术。近来物联网常用的超长距离通信（long range communication），也用到了这个技术。}
\BookParagraph{我们去投资，不能把鸡蛋都放在一个篮子里（diversify the portfolio），和「 spread spectrum 」的想法相似。我们在工作中，也应该尽量什么都会一点儿，而不是「术业有专攻」。}
\BookParagraph{所以说，工作的目的不是为了赚钱，而是为了拓展自己的技能，尤其在自己薄弱的环节。很多人一辈子的工作都在用同样的技能，虽然对长期工作中的业务非常熟悉，但是无法克服自己的薄弱环节。恰恰相反，我们应该尽量什么都学一点儿，什么工作都有一点儿经验和涉猎（learn a little about a lot）。}
\BookParagraph{这样说来，你的第二个工作选择「可以学习锻炼的更多」，自然是更好的选择。}
\BookParagraph{你甚至应该去再找一个第二职业，在第二个工作里去学一个完全不同的技能。比如如果你平时的工作从事软件开发，那你的第二职业可以去做市场销售。就像去健身房运动一样，运动的过程中确实很累，但是对自己的发展卓有成效。}
\BookParagraph{3. Take (measured) risks.}
\BookParagraph{机遇永远和风险并存，没有亏过钱的人，说明他从来没有去投过资。}
\BookParagraph{工作也一样有风险，因为未来无法预测。和投资一样，我们不能去找一个毫无风险的工作，因为这样也同样可以保证将来毫无进一步发展的可能。你想啊，如果你二十八岁时我告诉你有这么一个工作，将来从事这个工作需要一辈子节衣缩食，而且没有任何致富的可能，你会去选择这个工作吗？}
\BookParagraph{这么说是不是哪个工作都不该选，而是应该大学毕业直接去创业？}
\BookParagraph{哈哈别忘了，「 take risks 」这个短语里还包含了一个脚注：「measured」。虽然是「高风险，高回报」，但是也不能盲目地去追捧风险。创业虽然非常有意义，但是需要大量的经验和技能，在不具备这些经验的情况下，很多时候甚至是一种「高风险低回报」的工作。}
\BookParagraph{这就像投资比特币或者特斯拉一样，风险很高，但不一定有与其风险相应的回报，不如去投资亚马逊或者苹果。}
\BookParagraph{当然如果你把现在的钱都去买了房子，将来的钱都去给孩子买了房子，那任何投资都无从谈起了。}
\BookDate{2018 年 8 月，香港}
\BookArticle{赌　博}{赌博}
\BookParagraph{李老师，最近深受赌博困扰。缘起于世界杯期间朋友介绍的一个赌球软件，从赌球到后来在网上赌百家乐，从几十块到几千块，最后输掉四万多。家里想办法还了这笔钱，我自己却闹下了病根。昨天去了天主教堂，希望能借助神的力量渡过难关，去之后心里发愿，悬崖勒马，走正当渠道好好挣钱。}
\BookParagraph{赌博实在是一种心理上的疾病。这个感觉太可怕了，一旦进入赌博，没日没夜，夜晚也如白天。自己失去了的时间和不断的苦闷，甚至差点儿觉得人生无趣，一个平时看起来很正常的人结果竟然可以变成这样。不知为何想和您唠叨这些，我是从褚明宇的推荐关注的您。愿能借主的力量，让我早日克服心魔，走出昏暗。（问题来自微博私信）}
\BookParagraph{你听说过有一种说法叫「house edge」吗？}
\BookParagraph{有了「house edge」，一个赌场——或者赌球网站——为你设计的赌博玩法，虽然看似公平，赌一次两次甚至你还有可能赚钱，但赌的次数多了，赌场一定会有一定比例的「提成」。}
\BookParagraph{如果去赌场赌博想尽量输得少一点儿，你应该去玩儿「house edge」最低的赌法儿。比如玩一副牌的Blackjack，如果你用最优策略去玩儿，house edge 大概在百分之二左右。}
\BookParagraph{当然了，拉斯维加斯或者澳门的赌场里， Blackjack 很难找到只有一副牌的，也很难见到那种滴酒不沾，不苟言笑，最优策略早已烂熟于心的玩家。大多数人玩儿的老虎机，house edge 竟然高达25\%以上。}
\BookParagraph{概率论里的大数法则（the law of large numbers）表明，当你下赌注的次数趋近无穷大时，你的平均回报接近随机变量的期望值。正因为有了「house edge」，这个「期望值」是赌场一定会赚，你一定会亏。}
\BookParagraph{那如果是一场不存在「house edge」、完全公平的赌局呢？}
\BookParagraph{数学家早在十七世纪就讨论过这类问题。因为人天生的赌徒心理——赢了还想再赢更多——结果还是一样会输，概率论里叫做「gambler’s ruin」。这个结论非常斩钉截铁：无论你用什么方式去赌，即使赌局的期望值对你有利，你也一样会输个精光。}
\BookParagraph{所以说，如果你确实想靠赌博赢钱，你只能赌一次。赌一次还有赢的希望，次数一多必输无疑。可惜的是，一旦第一次赢了，绝大多数人会误以为赢钱很容易，还会像你一样继续赌下去，学不来刘正风「金盆洗手，退隐江湖」之风。}
\BookParagraph{其实，即使你真的可以一直赢下去，赢来的钱又能有什么用呢？}
\BookParagraph{很多沉迷于赌博——或者炒股——的人，会自然而然的想，有了钱还怕花不出去吗？我可以拿钱来买豪宅、玩儿跑车、周游世界。}
\BookParagraph{其实到了那时候，所有这些用钱能买来的玩法，刺激程度都远远不如再玩儿一把更大的赌注。}
\BookParagraph{这让我想起一部叫「Click」的美国电影。故事里男主人公觉得生活很无聊，于是就用一个神奇的遥控器把时间向前「快进」。直到终于有一天发现，其实生命的所有意义，都在于那些「飘风骤雨惊飒飒，落花飞雪何茫茫」的细节之中。}
\BookParagraph{还是李白说得真切：「人生得意须尽欢，莫使金樽空对月。天生我材必有用，千金散尽还复来。」}
\BookDate{2018 年 9 月，多伦多}
\BookArticle{命中注定}{命中注定}
\BookParagraph{李老师，想请教您一个问题：我一直以为自己是真正在为别人着想，但结果好像是处在我自己的环境里为别人着想。因此，我说的话好像会时不时地伤害到别人。}
\BookParagraph{例如，别人送我一件礼物，我首先表示很高兴，然后说，太贵了，退了吧，等你赚钱了再送我。别人就会觉得我是在说他没有钱，在伤他自尊。说者无心，听者有意，这个现象好像会经常出现在我身上，并且我不知道自己说的话伤害到别人了，直到别人告诉我。}
\BookParagraph{我试着站在别人的角度考虑，但是这好像并不是一件容易的事情，毕竟别人所处的环境可能很难模拟出来，我还是会被我所处的环境或多或少地影响着。}
\BookParagraph{我该如何改变？如果出现这种情况我又该如何弥补呢？（问题来自微博问答）}
\BookParagraph{哈哈，这个问题问我绝对是问错人了。}
\BookParagraph{世上有两种人。一种极具亲和力和领导才能，人见人爱，平时聊天随和、自然，让人感到如沐春风，心情愉快。}
\BookParagraph{还有一种人性格内向，「情商」不高，容易说错话得罪别人，说完了呢一想又会后悔，但是早已覆水难收。}
\BookParagraph{我母亲年轻时是前一种，而我则恰巧属于后一种。}
\BookParagraph{不仅如此，而且多年来我发现，性格以先天的基因传承为主，就好像长得不好看很难改一样，性格也是很难后天改变的——改的了一时，改不了一世。}
\BookParagraph{有很多书里，会列举各种各样生活中的例子，教育像我这样的读者，在某一个场景中的每一句话，应该这样说，而不应该那样说。}
\BookParagraph{可是问题在于，对于前一种人，这种书简直就是小儿科，他的话本来就是正确的说法啊。而对于我这样的人，就像你的故事里说的，无论初衷多好，总不可能每句话都想半天再说吧，那也太费劲儿了。而靠直觉说的话，性格使然，早晚还是要说错。}
\BookParagraph{别说看书了，就算摆一个活生生的标兵模范在你面前，天天耳濡目染，也照样是要不然东施效颦，要不然画蛇添足。}
\BookParagraph{那怎么办？}
\BookParagraph{这个问题，一个相貌平平的女生看到别人漂亮，或者一个脑子不够聪明的男生看到同年级的学霸，一样都会问。}
\BookParagraph{既然改已经没戏了，我建议：别想这么多，想怎么说话，就大大方方地敞开了说好了。}
\BookParagraph{那些「有意」的「听者」，无论如何，都无法让他回心转意，让他喜欢上你。}
\BookParagraph{说他「没钱」怎么了——我们大家不是也都没钱吗？}
\BookParagraph{你将来命中注定的朋友，即使你长得不够漂亮、脑子不够聪明、说话不够完美，他们也早晚会喜欢你的。}
\BookParagraph{就像二十多年前黑豹的那首《无地自容》里唱的：}
\BookParagraph{人潮人海中有你有我}
\BookParagraph{相遇相识相互琢磨}
\BookParagraph{人潮人海中是你是我}
\BookParagraph{装作正派面带笑容}
\BookParagraph{不必过份多说自己清楚}
\BookParagraph{你我到底想要些什么}
\BookParagraph{不必在乎许多更不必难过}
\BookParagraph{终究有一天你会明白我}
\BookDate{2018 年 9 月，多伦多}
\BookArticle{消费观念}{消费观念}
\BookParagraph{李老师你好，想跟您聊聊「消费观念」的事。}
\BookParagraph{我的「消费观念」算是比较落后。我的家境并不贫穷，但我在买东西的时候却非常在意商品的价格，如果价格有点贵，即便我很喜欢也不买，因为我觉着省钱是种本事。而我的女朋友、父母和周围的很多朋友都很会享受生活，觉着喜欢就买喽。他们也常吐槽我的这种苦行僧式的生活，说这也太屌丝了。李老师对此有什么看法？（问题来自微博私信）}
\BookParagraph{哈哈，这明明是个优化问题——你问一个学理工科的算是问对人了。}
\BookParagraph{银行里一共就这么多钱，怎么花才能最快乐最幸福？}
\BookSubtitle{一}
\BookParagraph{我前一阵子看到一篇文章，讲述在美国小镇里生活的一个大学教授，年薪五万，无论如何精打细算、精简开支，每月都所剩无几，基本只能当一个月光族，英文管这叫「live from paycheck to paycheck」。刘震云的成名小说《一地鸡毛》，说的也是这样一个故事。}
\BookParagraph{这会让人陷入一种思维上的怪圈：越精打细算、量入为出，平时就越把精力用在如何花钱这件事儿上，越难以打破这种举步维艰的局面，去提高自己的收入。}
\BookParagraph{所以说啊，长期来看其实有一个爱花钱的女朋友是好事儿：钱花得越多，越能激励你们俩去思考如何多挣钱。}
\BookParagraph{「开源节流」——「开源」永远比「节流」更重要。}
\BookParagraph{你看美国借的国债一年比一年多，经济不是照样红红火火、蓬勃发展嘛。}
\BookSubtitle{二}
\BookParagraph{精打细算其实不难，但是如果像你说的那样只关注价格，长期来看，不一定是最省钱的。}
\BookParagraph{原因很简单：贵的东西因为质量更高，说不定更耐用。比如你买一辆自行车，与其买辆最便宜的骑了几天坏了，不如买一辆贵一些的，说不定可以骑很多年。所以说，如果计算平均每年在自行车上的开销，买贵的比买便宜的更划算。}
\BookParagraph{当然，我们到现在为止做了一个假设：你买自行车完全不在乎一辆好车给你带来的快乐。}
\BookParagraph{这个假设几乎永远是错的。如果你的女朋友买一个Louis Vuitton 的包比买一个 Coach 快乐二十倍——虽然她不一定会拼Louis Vuitton 这两个单词——但价格仅仅贵了十倍，那她买 LV 的决定就绝对明智。}
\BookParagraph{如果你买的东西可以用一辈子，一身西装、一件实木家具、一架三角钢琴，那也绝对不应该买最便宜的。因为一旦买了，穿着用着一点儿也不开心，早晚你收入增加了还会换更好的，长期来看，花的钱更多。}
\BookSubtitle{三}
\BookParagraph{当然了，很多奢侈品其实卖给你的只是一个品牌的价值，比起普通商品，不一定能带来更多的快乐。}
\BookParagraph{一块几万的机械表，每隔几年就要维修，走时还不如一块几百块的石英表。一架几万的单反相机，却很少带在身边，平时照相还是要用手机。}
\BookParagraph{那很多人为什么还要买机械表买单反相机呢？}
\BookParagraph{不是因为需要，只是因为得到一样梦寐以求的东西的那一瞬间，有一种快乐的感觉。}
\BookParagraph{如果你不在乎这些，甚至越省钱越快乐，那少买奢侈品确实是「节流」的一个好办法。}
\BookSubtitle{四}
\BookParagraph{那女朋友喜欢奢侈品怎么办？}
\BookParagraph{你当然可以晓之以理、动之以情，提倡少买又贵又不合算的东西。}
\BookParagraph{不过我更赞同另一种思路：你私下把她最想要又一直不舍得买的东西悄悄买下来，在一个街边的小馆子正吃饭的时候，不经意地拿出来送给她。}
\BookParagraph{「Happy wife, happy life」—— 「你快乐所以我快乐」是花钱的最高境界。}
\BookDate{2018 年 10 月，多伦多}
\BookArticle{预测未来}{预测未来}
\BookParagraph{李老师，您好，新年快乐。很抱歉新年第一天就用这样的方式打扰您。我有一个困惑，刚才看到有人转发2019年属相蛇犯太岁，我不太明白，就去查了一下。是描述属蛇的人各种运气不好，甚至会有血光之灾，我不知道该不该相信。李老师，您怎么理解这些说法？想听听您的看法。（问题来自微博私信）}
\BookSubtitle{一}
\BookParagraph{先给你一个结论吧：任何试图预测未来的观点，都别盲目相信。}
\BookParagraph{比如如果你看了一篇朋友圈转发的文章，说美国股市已经见底，2019年一定会稳中有升；或者看了一篇公众号里的「十万加」，说新的一年国内二三线城市房价必然会大跌，你千万别盲目信以为真。}
\BookParagraph{原因很简单。如果文章的作者确实有这种预测未来的能力，他完全可以靠这种「超人」的能力自己在股票或房地产市场里大赚一笔，没有必要把他的预测结果费尽千辛万苦写成文章免费提供给你阅读。实际上，即使作者的预测只是「有很大的把握」，他也一样可以依靠期货等各种金融衍生产品对冲风险，有多大的把握就可以赚多少。}
\BookParagraph{金融和房地产市场的本质特征是，虽然几十年的长期可以升值或者保值，但一年甚至三五年的短期未来是无法准确预测的。}
\BookParagraph{你也许要问，那像巴菲特这样的「股神」呢？}
\BookParagraph{像巴菲特这样在股市中可以依靠准确的投资选择长期跑赢市场的人，其实也不是只有他一个。但是如果你希望他们替你投资，你得支付不菲的佣金费用。况且即使是巴菲特，从几年的短期来看，也不是每一次预测和投资选择都是正确的。}
\BookParagraph{天下没有免费的午餐。If it sounds too good to be true, it probably is.}
\BookSubtitle{二}
\BookParagraph{本质上来说，炒股有关的建议和看属相看星座运势差不多：它们都是利用你「想赢怕输」「不怕一万就怕万一」的心理，吸引你去关注。}
\BookParagraph{预测你新的一年有「血光之灾」，后面多半儿会给你一些「化解之法」。要是预测你新年会「时来运转」「福星高照」呢，那后面多半儿会话锋一转，说几条「注意事项」，给你出几条高招。}
\BookParagraph{这些运数预测文章的作者，多半儿希望能利用你愉悦或害怕的心理，吸引你去传播、打赏，去为这些「知识」花钱。说到底，咱们现在不是正身处「大数据」和「知识付费」的时代洪流中嘛。}
\BookParagraph{可是我一直不太理解为什么很多读者会真的相信跟属相或者星座有关的分析。你想啊，同一年生的上千万人，或者同一星期生的几十万人，无论遗传基因、家庭背景、政治面貌、文化程度有多大差异，运数和性格都大致相同，这怎么可能呢？}
\BookParagraph{有人可能会反驳，不对啊，我认识一个发小，用属相和星座的那一套去分析她的性格，特准！}
\BookParagraph{别忘了，a broken clock is right twice a day.}
\BookSubtitle{三}
\BookParagraph{微博的粉丝分析中，甚至都能找到粉丝星座的分布图。在我的微博粉丝里，魔蝎座的人最多。}
\BookParagraph{我对星座性格分析不熟，就上网搜索了一下。对待感情，魔蝎座的人「很难相处，有种高处不胜寒的感觉。碰到喜欢的人，不会主动的表白，只是做一些事情，引起对方的关注，当对方注意到的时候，摩羯座的人又会显得特别扭捏，尽管这就是自己想要的。」}
\BookParagraph{哈哈，我前一段不是一直在说千万不要表白嘛——看来我的文字最能和魔蝎座的人产生共鸣。}
\BookParagraph{但是其实真正的原因多半儿是，很多粉丝在开通微博帐号时不想透露自己的生日，就随便填了一个一月一日。}
\BookParagraph{而一月一日正好是魔蝎座。}
\BookParagraph{本文此后还有百分之九十有关魔蝎座的精彩内容，晓之以理，动之以情，有理有节，不容错过。您可以猛戳以下二维码加入本文专属的收费微信群，每月酌情收取一元人民币费用。名额有限，先到先得，欲购从速。}
\BookDate{2018 年 12 月，多伦多}
\BookArticle{Too Good to be True}{Too Good to be True}
\BookParagraph{李老师您好，我的问题是有关赌博的问题。我是一个北方人，几年前来到广东，发现这边的人对于赌有着无限的热爱。有段时间我也参与其中，后面就收手了，明白了久赌必输的道理，但是有一样东西我还是没有搞清楚（已经戒赌了，只是还想知道答案）。之前听朋友说关于时时彩有很多软件做了时时彩计划，就是每期这个软件告诉你怎么买，准的时候可以做到二三十期全中，他们说只要能按计划下注就肯定赢。李老师，这东西靠谱吗？这种软件的逻辑是啥？它每期出的计划是有什么测算公式还是怎么得出来的？？\BookLineBreak{}（问题来自微博私信）}
\BookParagraph{这类的问题我一般都不会回答，主要原因是苦口婆心说了半天多半儿没什么用 —— 和中医有点儿像，「愿者上钩」，该花钱买的人你再怎么说，他永远都是要花钱的。不过呢，你最后这两个大问号让我动了恻隐之心，觉得你还有救。}
\BookSubtitle{一}
\BookParagraph{我想了半天，如果把我的回答浓缩成一句话，中文里还真不容易想出和以下这句英文一样绝妙的句子：}
\BookParagraph{If it’s too good to be true, it probably is.}
\BookParagraph{中文可以说「天下没有免费的午餐」，但是不是特别贴切。这句英文想说的意思是，你要是哪天觉得赚了便宜了，多半儿是要亏的。}
\BookParagraph{比如说你去一家馆子吃饭，结账时服务员给你介绍说，如果你「免费」加入他们的「金卡会员」，买一张五百块的充值卡就可以打八折。你多半儿会想，再吃一两顿就可以把充值的钱花完，打八折的便宜不赚白不赚啊。}
\BookParagraph{可是你的计算中一定隐含着一个这家馆子会一直开下去的假设——你没有考虑「跑路风险」。但实际上，很有可能过几天馆子关张了，你的「金卡」变得一文不值。要是馆子号称充五千块钱可以打五折，表面上便宜似乎赚得更多，但实际上他的「跑路风险」却大大增加。}
\BookParagraph{这是一个「too good to be true」的典型案例。你的正确做法是，按照贵一些的正常价格结账走人，别充值什么「金卡会员」。}
\BookSubtitle{二}
\BookParagraph{这个案例中，店家虽然利用了我们天生爱占小便宜的心理，但到底钱的数目不多，无关大局。甚至可以说，在这种小事儿上摔一两跤是一件好事儿，就当交了学费，下次碰到大「便宜」，说不定就不会去占。}
\BookParagraph{比如去年挺火爆的「P2P理财」这种动辄号称每年保证百分之十回报的「理财」产品，对于稍微懂一点儿理财的人来说，简直就是「司马昭之心，路人皆知」啊。百分之十的回报当然是有可能的，但问题就在这「保证」二字。}
\BookParagraph{一年多前我在香港和一个教授朋友吃饭，他妻子是卖保险产品的，据说做得风生水起，收入比教授还高。她大概把我当成了潜在的客户，三句话不离本行，说如果买她的保险，可以保证百分之六的回报。看她的样子，是真的相信有这样的「保证」。我一笑置之，没有和她继续聊这个话题。}
\BookParagraph{为什么一旦「保证」百分之十的回报，就千万不能买呢？咱们可以先假设这个「保证」真金白银，童叟无欺，那这个利率和普通债券市场的利率之差，就是一个「arbitrage opportunity」，中文大概叫做「套利机会」。这种机会利用的是市场暂时的低效，是稳赚不赔的。比如，我们大可到银行以百分之八的利率贷一笔钱，然后「投资」到这个「P2P理财」产品，赚取那百分之二的差价。}
\BookParagraph{而市场的一大特性，正好是聪明的人到处都是。一个稳赚不赔的套利机会，发现的人一多就没有了，所以这样的机会一定是稍纵即逝的。既然这些保证百分之十的理财产品一直没有降低利率，就肯定说明它们的保证是有风险的，而且保证的利率越高，风险越大。}
\BookSubtitle{三}
\BookParagraph{那有没有基本毫无风险的投资呢？有啊。如果你去买美国政府或者加拿大政府发行的债券，又可以一直持有，直到它们十年甚至三十年以后到期，那只要政府不要倒闭，你这笔投资的本金一定可以收回，利息也一定可以收到。}
\BookParagraph{但其实即使是买这样的政府债券，因为货币贬值的因素，债券到期后你拿到的钱的实际购买力说不定还会有所下降。要想保证实际购买力不会下降，你需要去买一种美国政府发行的叫做TIPS（Treasury Inflation-Protected Securities）的产品，是一种real return bonds，可以保证的利息是基于物价涨幅之上浮动的。}
\BookParagraph{而投资这样的产品可以比物价涨幅高多少呢？百分之零点六。}
\BookParagraph{对，你要是希望你的投资有一个真正没有风险的「保证」，你的投资每年的把通货膨胀排除后的实际购买力，只有百分之零点六的涨幅。}
\BookParagraph{「可是我有一个很靠谱的朋友，说香港的保险保证百分之六的回报啊？」}
\BookParagraph{「微信里看到过一家理财产品，保证百分之十，这么多人都买了，他们不可能都不懂吧？」}
\BookParagraph{别忘了 ——「If it’s too good to be true, it probably is.」}
\BookSubtitle{四}
\BookParagraph{与其去买「P2P理财」、买「时时彩」的彩票、或者买一个教你如何稳赚不赔的软件，你其实可以去买几本好书，或者注册几门课程，把有限的资源去为自己的未来「投资」。}
\BookParagraph{如果你还有闲钱，我郑重地建议你去尝试一下低空跳伞。}
\BookParagraph{一套带预备伞的降落伞装备，大约二十多公斤。我建议你背上经典的 T-11 降落伞，看看自己是否有胆量从飞机上跳下去。}
\BookParagraph{你跳之前一定听说过，跳伞是一项非常安全的运动 —— 主副降落伞都打不开的概率，大概只有十万分之一。和无痛分娩所用到的局部麻醉造成瘫痪的概率差不多，甚至比美国政府无力偿还债券的概率还低一些。}
\BookParagraph{如果你不敢冒险跳下去，那其实也是人之常情。我也不比你强多少 —— 连过山车我都不敢坐。}
\BookParagraph{但是，你为什么就敢冒险买赌博软件或者「P2P理财」呢？}
\BookDate{2019 年 4 月，多伦多}
\BookArticle{信任}{信任}
\BookQuote{李老师，您好，元旦快乐。辞旧迎新的年底，是我比较纠结的时候，想问问您的意见。}
\BookQuote{我和男朋友是做教育培训的，不是正式的公司，两人一起合作招生。学生的学费都是转到他的支付宝账号里，一年结束，我真的很想让男朋友理一下把钱分了，这是我付出一年应该得到的。我之前试探过他，他的意思是在一起了，钱也应该在一起。}
\BookQuote{作为女生我比较希望自己经济独立，买车买房可以一起还，但是退一万步讲，离开你我依旧可以过得很好，不想依赖谁。我不知道该如何沟通关于钱的事情？让他明白又不伤害感情，我是爱他的，但我觉得这是两回事，独立与爱他并不冲突。}
\BookQuote{辛苦您了，新年快乐！（问题来自微博私信）}
\BookParagraph{我的缺点之一是，不够大度，小家子气。}
\BookParagraph{所以我建议：斩钉截铁地立刻摊牌表明立场，把你自己的那部分收入从他的支付宝账号要回来。}
\BookParagraph{他越是找理由不给你，你就越要仔细留意他不给你的理由。如果他的理由像你说的是「既然在一起了，钱也应该在一起」，怎么办？}
\BookParagraph{哈哈那你可以回应：「那就把所有的钱都转到我这里好了。」}
\BookParagraph{他的逻辑中有一个明显的问题。把两个人的钱放在一起当然很好，但是必须是共同拥有的账号。在加拿大，共同拥有的私人账号（joint account）非常常见，很容易开，把两个人的名字都写进去就可以了。当然还有一个办法是注册一个公司，两个人都是公司的股东，共同拥有公司的财产。}
\BookParagraph{如果不是共同拥有的账号，那可不是「在一起」的钱，而是他一个人的财产。按照这个逻辑，把所有钱都转给你，岂不是也符合「在一起」的定义？如果他找出种种理由不愿意全都转给你，那他的「在一起」的逻辑则完全站不住脚，不攻自破。}
\BookParagraph{那如果将来你们确实生活在一起了，钱应该如何保管才是最优策略呢？}
\BookParagraph{不知道你是否听说过，计算机上运行的操作系统中有一个概念，叫做「线程」。在程序运行的过程中，两个线程既共同拥有一些全局变量（global variables），又各自拥有各自的私有变量（local variables）。}
\BookParagraph{两个人在一起其实也可以这么过日子。既各自有各自的账号，又同时拥有共同的财产，这样有百利而无一害。其实对大多数家庭来说，与其说是「共同财产」，不如说是共同的负债，因为两个人需要一起还房子的贷款。贷款越早还清，支付的利息总额就越低，所以每月开销的盈余，大部分其实都拿去还贷款了，两人各自的账号其实也仅仅只是一点儿「私房钱」，也就够时不时给对方买个礼物罢了。}
\BookParagraph{如果是在国外，两人各自的账号还包括退休储蓄。因为退休储蓄账号（比如美国的 IRA 或加拿大的 RRSP）可以少交税，理论上就又复杂了一个层次。比如，收入多的一方可以把钱存入收入少的一方的账号中，叫做「income splitting」，可以最大程度省税。}
\BookParagraph{钱在各自的账号中，既可以互送小礼物，又可以一起还贷款，一起研究如何储蓄、如何投资、如何省税，两个人小日子不是也挺红火？}
\BookParagraph{电影《肖申克的救赎》中 Andy 在狱中的转折点，源于他对狱警 Mr. Hadley 说的一句话：}
\BookQuote{Mr. Hadley, do you trust your wife?}
\BookDate{2020 年 1 月，多伦多}
\BookArticle{量力而行}{量力而行}
\BookQuote{李老师新年好。看过您的一些回答，觉得回答中肯，都很有参考价值。我最近的问题是这样的：和女朋友已经相处一年多了，近日她又对我讲买房子的事情，大意是婚前我要买一套属于自己的房子，哪怕小户型也好。}
\BookQuote{因为她自己的同事们要么结婚要么在深圳周边的城市如东莞都买房供房了，她的意思是如果现在不咬牙买下今后就更加买不起了。我们都是工薪族，我比她大好几岁，职业发展一般。之前谈到买房的话题时两个人还闹过矛盾，因为我对买房有点抵触，不想在婚前就背负房贷做房奴，可是买房对她来讲就是安全感的所在。我也不懂什么投资理财，房子用她的话讲就是一个可以升值的商品，比把钱存到银行要强很多。}
\BookQuote{我不想跟她就买房的问题再有矛盾，请问李老师，我怎么在大方向去规划买房这件事，不是不买，也不是跟风似的买。谢谢。（问题来自微博私信）}
\BookParagraph{房子当然可以买，不过你们最好能先大致了解一些买房的常识再买。这样无论将来出现什么问题，不至于因为房子的事儿影响两个人的感情。}
\BookParagraph{从几十年的长期来看，买不动产虽然不能像你女朋友说的那样保证可以升值，但是确实是有可能可以保值的。也就是说，随着货币的贬值，房产的纸面价值不断上升，而其购买力则大致保持不变。这是因为随着时间的推移，房子本身是贬值的，但房子所在的地段和社区逐渐成熟，地价逐渐升值，两者相抵，房产基本保值。}
\BookParagraph{要是想长期投资可以升值，其实应该投资股市。这是因为如果资本市场非常成熟，从长期来看，投资股市的一部分钱是用来扩大再生产的，确实创造了新的价值。}
\BookParagraph{当然，几十年的长期不确定因素太多，没有人可以预测。比如如果二十年前在美国买房，其投资收益远远不如把同样的钱投资美国股市。但如果换到国内，结果则恰恰相反：二十年前投资股市远远不如买房。}
\BookParagraph{但是投资有一个常识：以往的投资表现不代表未来。说不定下一个二十年应该反过来在美国买房、在国内投资股市呢。世界上房价不升反降的例子比比皆是。比如日本东京的房价自从九零年崩溃之后，由于人口逐渐老化，到现在三十年了，都没能完全恢复当时的最高点。但凡有人告诉你买房长期可以升值，你就得好好考虑一下他这么说的动机。}
\BookParagraph{所以说，买房不应该是为了投资升值，而是为了你女朋友提到的「安全感」。住在自己买的房子里会比住在租来的房子里感觉更安心、更幸福，这是人之常情。为了这样的「安全感」买房，我认为很值得。}
\BookParagraph{所以说，买房不是为了投资，而是在买一样儿给自己多一点儿幸福感的商品。这一点和买车、买手机其实差不多，唯一的区别是房子比车和手机贵得多，需要几十年收入的很大一部分慢慢才能还清。}
\BookParagraph{就是因为买房贵得多，所以就需要慎重地计算一下你的收入到底可以买多少钱的房子。这个道理其实很浅显：比如如果你只能买得起丰田汽车，你绝对不会去买什么「奔驰大 G」，因为你很明确地知道自己买不起。可是到了买房，很多人就跟你女朋友一样，觉得「咬咬牙也就买了」，不再去精确计算自己到底能买得起多少钱的房产。}
\BookParagraph{那买房如何才能量力而行，如何计算自己的收入可以负担得起多少钱的房子呢？我们可以参考加拿大联邦政府建议的计算方式，叫做 Gross Debt Service (GDS) ratio 和 Total Debt Service (TDS) ratio \footnote{\href{https://itools-ioutils.fcac-acfc.gc.ca/MQ-HQ/MQCalc-EAPHCalc-eng.aspx}{Mortgage Qualifier Tool}, last accessed on August 13, 2020.}。基于 GDS ratio，加拿大政府建议房屋按揭贷款加上每月生活必须支付的水电费（在加拿大还需要再加上房产税），不能超过税前总收入的 32\%。而基于 TDS ratio，则建议按揭贷款、水电费、买车等其他贷款的总和，不能超过税前总收入的 40\%。}
\BookParagraph{举例来说，如果你每月的税前收入是三万元，水电费每月两百元，那你每月的按揭贷款不能超过 9400 元。如果按照年利率百分之五、二十五年付清、首付百分之二十五来计算，那你要买的房子不能超过 215 万元，而买车的开销每月不能超过 2400 元。当然，由于加拿大和国内个人所得税的税率不一样，这些计算还可以再进一步调整。}
\BookParagraph{「咬牙」可以，千万要在你能负担的范围内去「咬牙」。其实呢，我倒是觉得「奔驰大 G」这样的豪华车很多人都可以「咬咬牙」去买，最不济就是买辆二手车再少买些其他生活必需品。而买房子一旦超过了 GDS ratio 和 TDS ratio 的上限「咬牙」去买，搞不好会因小失大，为了一点儿「安全感」，小半辈子的日子一直在缺钱的状态中度过，会过得很不开心。}
\BookParagraph{「房奴」有两种：幸福像花儿一样的「房奴」和捉襟见肘狼狈不堪的「房奴」。你希望做哪一种呢？}
\BookDate{2020 年 1 月，多伦多}
\BookArticle{理发}{理发}
\BookParagraph{理发店关门了。}
\BookParagraph{郭芳说，没事儿，我给你理发不就得了 —— 我自然而然地想起了二十多年前从北京带到美国的那把理发推子。}
\BookParagraph{那时候我们几个比较熟的穷学生，为了省钱，隔一段时间聚在一起互相理发的时候，用的就是我那把千里迢迢带过来的手动理发推子。有一个朋友长得很帅气，手艺也不错，非常乐意给别的同学理发，自己却从来不在这儿理。剩下我们几个，觉得形象虽然重要，但美金更重要，就凑合一下算了。}
\BookParagraph{你还别说，我的理发推子是家里传下来的。时不时上一点儿机油，在手上一收一放，只要用力恰到好处，弹簧及时弹开，咔嚓咔嚓地像是音符在头顶跳动，不但不会夹头发，还有一种丰收季节收割的快感。我的手艺比较随机，时好时坏，有一次一位同学的头发剪白了一大块儿，不知这位如今已是功成名就的大教授是否还记得。郭芳后来给我剪头发的时候，竟也经常陶醉于这样的快感之中，理着理着就理成了「板寸儿」，适合配上绿色的军装，美其名曰「短点儿凉快」。}
\BookParagraph{还记得我小时候在北大燕南园住的那会儿，夏日的周末蝉鸣声四起，要不是姥姥的冰激凌和冰镇西瓜，还真没法儿「心静自然凉」。外公虽然没有多少头发，却非常注意形象，理发这件事儿只有小姨父可以胜任。于是小姨父经常中午来到燕南园，白色的围裙一系，一小会儿的功夫，就给外公把头发理好了。然后一大家人坐上餐桌，热菜热汤端上来，天南地北地闲聊着，午后的时光过得也快些。}
\BookParagraph{当年从美国辗转搬家到加拿大，因为我的双门小跑车放不了太多东西，很多物件扔的扔，送人的送人，不过这把理发推子算是一样儿宝贝，记得还是随车带过来了的。到储物间找了一会儿没找到，只好到亚马逊上下了单。}
\BookParagraph{现在早已没有手动的理发推子了，都是电动的。其实越是电动，越有点儿担心什么地方随手就白了一块儿。不过现在都是视频会议，看不了太真切，大不了关了视频就成。}
\BookParagraph{「还是有媳妇好啊 —— 就算全世界都关门了，还有个人给你理发。」我跟郭芳感慨。}
\BookDate{2020 年 3 月，多伦多}
\BookArticle{好酒不怕巷子深}{好酒不怕巷子深}
\BookQuote{李葆春老师中午好。在电子诈骗中被骗的人，无论是出于贪还是侥幸的心理，您对这样的人、对被骗这种事怎么看？}
\BookQuote{撇开不讨论电子诈骗以及其他金钱骗局的低端、低技术含量和明显的诈骗性，也不谈普通人要有防诈骗的警惕心、不要明知故犯或贪心、要培养脚踏实地的性格或者政府应该积极打击诈骗犯，只谈那些不幸被骗的人，是不是被骗本身说明这类人确实人生的路会更难？}
\BookQuote{一辈子有可能不被骗吗？被骗过后会不再上当吗？人生又是否又会有重来的机会呢？为什么人性这么脆弱？（问题来自微博私信）}
\BookSubtitle{一}
\BookParagraph{七年前有一个夏天，我正好在西安出差，准备在当当网上邮购一些书带回加拿大。我用我当时的手机号码注册了一个账号，选好了几本书，就下单了。虽然那时我没怎么在国内网站上买过东西 —— 即使到现在我都没有淘宝的账号 —— 但我一直听说国内的电商极其发达，只需要一两天时间，书就可以轻松地快递上门。}
\BookParagraph{第二天，我接到当当网的客服打来的电话。客服首先跟我核对了一遍我的订单号、名字和快递收货的地址，然后客气地说：「李先生，不好意思，您的快递暂时缺货，需要先给您操作退款，等有货了再给您重新下单。」}
\BookParagraph{我回答：「行，那你就先退款吧。」}
\BookParagraph{「先生，您有没有 QQ？」}
\BookParagraph{「对不起，我没有 QQ。你要 QQ 干嘛？」}
\BookParagraph{「哦，是这样的先生，咱们这里查到您是用网银支付的，操作退款需要把钱原路退回给您的银行账号，需要您提供一些信息，所以必须用 QQ 联系您。」}
\BookParagraph{「我确实没有 QQ。你要什么信息？」}
\BookParagraph{「那实在不行要不然这样：我这里有个网站，给您短信发过去了，您登录填写一下您的信息，填好了告诉我，我在线等您。」}
\BookParagraph{网站我一下就打开了，里面提示我填写银行账号、收款人和密码。}
\BookParagraph{后来我回想起来，真的是有点儿悬。直到打开这个网站之前，我还没有意识到这其实是个骗局。主要原因是对方刚开始就报出了我的订单号、收件人姓名和详细地址，骗取了我的信任。但是这个骗局的设计稍微有些粗线条，有两个明显的破绽：第一是对方带有南方口音的普通话，第二是让我填写账号密码。}
\BookParagraph{对了，还有他后来愈发焦急的催促：「您填好了没有？填好了没有？」}
\BookParagraph{电信诈骗之所以这么多年来一直没有杜绝，非常符合经济学的原理。由于作案风险相对比较低，获利又相对比较高，就会有更多的人投入其中，组织和分工也会越来越精准严密。再加上其套路层出不穷，年年都可以推陈出新，令人防不胜防。}
\BookParagraph{所以我对在电信诈骗案中被骗的人非常同情 —— 我自己研究了半辈子计算机，也没有什么贪小便宜的心理，又历来小心谨慎，当年都照样差点儿上当。}
\BookSubtitle{二}
\BookParagraph{那将来会不会再次上当受骗呢？}
\BookParagraph{这得看你这个「受骗」如何定义了。}
\BookParagraph{如果只是电信诈骗中的骗局，我想大多数人多半儿可以「吃一堑长一智」，不接陌生人的电话，基本可以对这种骗局「免疫」。}
\BookParagraph{但如果你把什么「P2P 理财」啊、「币圈儿」啊、「直销」啊、「万能保险」啊、「投行权威分析」啊什么的全都算上，那很可能你今后还会再「受骗」，而且还不止一次。}
\BookParagraph{这是因为在如今的信息时代，每天主动读到和被动接受的信息和观点实在太多，哪些是正确的，哪些是错误的，哪些是事实真相，哪些则是以假乱真的宣传，不太可能有一种万灵的方法，可以保证自己的判断每次都准确无误。}
\BookParagraph{不过你大可放心，这样的「受骗」就如同一家企业在市场中的亏损一样，并不一定是一件坏事儿。盈利是一种市场信号，亏损也是一种信号。和盈利比起来，亏损甚至可以让一家企业更有效地重新优化整合自己的资源配置和产品策略，积累更宝贵的经验。}
\BookParagraph{一个人更是如此：赚钱的时候春风得意，而一旦因为判断失误听信了别人的「权威分析」上了当亏了钱，「痛定思痛」，总结出来的经验只会更加宝贵。}
\BookParagraph{所以说啊，人生完全不用「重新开始」，以后的好日子还长着呢。}
\BookSubtitle{三}
\BookParagraph{不过道理虽然这么说，要是换做是我，总还是会希望尽量多「盈利」少「亏损」，少总结宝贵的失败经验，心情会好很多。}
\BookParagraph{我有三条「锦囊妙计」，说不定可以派上用场。}
\BookParagraph{第一条是尽量不要人云亦云，盲目「跟风」。比如，如果大家都去买比特币或某种高回报的理财产品，你则要冷静下来持怀疑态度，问一个关键的问题：买房也好，比特币也好，P2P 理财也好，这种投资为什么可以盈利？风险多大？风险从何而来？如果经过冷静的独立思考，觉得对这些问题的回答逻辑严密、符合常理，再去跟风不迟。}
\BookParagraph{第二条是多考虑长期的影响，少考虑短期的效益。经济学家 Thomas Sowell 管只看短期效益的短视行为叫做「Stage One Thinking」。Sowell 指出，很多决策的制订都仅仅看到了短期的好处，没有考虑长远的负面影响。比如，如果有这么一种理财产品，头一年的回报百分之十，但长期来看没有保证，我们就一定要估算一下其长期的风险，连为其担保的公司倒闭的风险也要计算在内。}
\BookParagraph{最后一条是尽量避免「冲动消费」。商场或网站大减价的或者「直播带货」的产品，都属于别人给你推销的产品。这些绝大多数都不是你需要的产品，很可能完全不适合你。但是你因为一时兴起，说不定就下单买下了。买了以后没怎么用，就造成了有限资源的浪费。而「好酒不怕巷子深」，质量上乘、做工精细的产品，你排队抢购都不一定抢得到，不会天天大减价推销给你。}
\BookParagraph{比如说吧，如果我现在告诉你，我准备近期隆重推出一个名为「好酒不怕巷子深」的私人朋友圈，为期四周，每周两小时。届时将汇集我的三位老朋友作为嘉宾，有你早已耳熟能详的「大 V」，有学识渊博的学者，还有神龙见首不见尾的高人，没有什么主题，大家根据共同关心的话题，全程语音互动，想聊哪儿就聊哪儿。由于平台资源限制，名额有限，优惠价格在评论里找，请大家尽快报名。}
\BookParagraph{你可千万不要参加。}
\BookDate{2020 年 6 月，香港}
\BookArticle{比特币（三）}{比特币（三）}
\BookQuote{李老师您好。比特币仿佛又迎来了新的春天，我这样不太懂的人很难分辨购买比特币是投资还是投机，打算用小部分闲钱购买比特币等数字货币基金并长期持有。想知道您对比特币这类数字货币的未来怎么看？（问题来自微博私信）}
\BookSubtitle{一}
\BookParagraph{我一直认为，如果对一件事儿不是很了解，应该先花时间做一些调查研究，再去做决定。}
\BookParagraph{比如打算买一辆新车，是买丰田还是买特斯拉？虽然当今买特斯拉比较时髦，我经过一段时间仔细的研究，还是决定买一辆丰田。}
\BookParagraph{买车完全可以看心情，但是投资不同，由于关系到几十年以后的生活品质，花点儿时间做研究是相当合算的。我说的针对投资的研究，不是简单地看看微博上「大 V」关于特斯拉比特币的新鲜观点，而是从关于股票和债券的基本知识开始，读一读 Benjamin Graham 写的《The Intelligent Investor》或者 Warren Buffett 写的《Letters to Shareholders》，再追一追三十年前的股神 Peter Lynch 或者还继续活跃在投资第一线的 Michael Burry，你的研究工作才算是全面而深入。}
\BookParagraph{你发现没有，我所推崇的研究不是去被动地去刷每天推送给你的内容，而是自己主动地去筛选和寻找高质量的信息，甚至去读几十年前写的书。每天社交媒体的新闻很容易让你觉得机会再不抓住就一去不复返了 —— 英文叫 Fear of Missing Out (FOMO)。无论是买房、买特斯拉还是买比特币，一旦有了这种「FOMO」的焦虑，早晚会有「all in」的冲动。偏爱「赌一把大的」迅速达到「财富自由」，是我们大家都有的美好理想。}
\BookParagraph{三十多年前巴菲特有一次访谈中曾经提到过一个观点叫「tune out the noise」。说的就是在如今信息爆炸的时代，应该尽量少看短期的新闻和投资分析，而是去研究一家公司长期的价值。那时的 IBM 有点儿像现在的特斯拉，是高科技领域的新宠，而巴菲特没有买 IBM 的股票，一是因为他不懂高科技，无法准确地评估这家公司的价值；二是他完全不怕「Missing Out」：一个人再努力也不可能什么都懂，错过了自己不懂的机会有什么可遗憾的？}
\BookParagraph{所以，与其天天「FOMO」，我提倡「Fear of Mixing In」(FOMI)—— 一去不复返的机会没有抓住没关系，我更害怕随波逐流，人云亦云。}
\BookSubtitle{二}
\BookParagraph{几十年长期投资需要考虑的因素中，最重要的是风险。}
\BookParagraph{风险的种类有很多。比如如果只买一家公司的股票，会有投资过于集中的风险（concentration risk）。}
\BookParagraph{假设你花了一万块钱买了比特币或者特斯拉，将来如果有一天中央银行限制了比特币的流通，或者电池生产的原材料紧缺，你的一万块钱就有可能大幅度缩水。}
\BookParagraph{所以，现在分散风险的常规做法是买入 Exchange Traded Funds (ETF)，相当于你这一万块钱同时投资了几十、几百、甚至几千家公司的股票，涵盖了不同的市值、不同的产业、甚至不同的国家和地区，而不是只买一家公司的股票。}
\BookParagraph{花一万块钱去买比特币，自然也会有投资过于集中的风险。}
\BookParagraph{那如果我有十万块钱，只拿其中一万块去买比特币，不怕什么投资集中的风险，是不是就没问题了？}
\BookParagraph{咱找一家历史悠久的加拿大公司 —— Enbridge Gas —— 来横向比较一下。Enbridge Gas 的业务很多，主要的一项业务是石油运输管道（比如其著名的 Line 3 pipeline）。就好像我们每天都离不开手机一样，无论油价高低，石油总是要通过管道运输的。由于新建管道越来越困难，已有的管道也就可以旱涝保收，越来越重要。正因为这个原因，以 Enbridge Gas 现在每股 35 美元的价格，每年的分红是 7.36\%。}
\BookParagraph{如果你把每年投资 Enbridge Gas 的分红拿来继续买入这家公司的股票，三十年以后，你投资的一万块钱，可以涨到八万四千元。}
\BookParagraph{由于比特币不是公司，所以没有分红。三十年以后，比特币如今 46000 美元的价位，需要涨到三十九万美元，盈利才能超过投资 Enbridge Gas。}
\BookParagraph{比特币有没有可能涨到三十九万美元呢？当然有可能。不但三十九万美元有可能，三千九百万美元都有可能。但是它也有可能三十年后无人问津，一文不值。而 Enbridge Gas 的石油运输管道，三十年后一定还在继续运营。虽然投资比特币潜在的收益很高，但承担的风险更高。你可以选择投资比特币，但是你的选择不应该是人云亦云，而是完成了上述研究之后才得出的结论。}
\BookParagraph{我选择 Enbridge Gas —— 而且宁愿错过比特币这个说不定千载难逢的历史机遇。}
\BookSubtitle{三}
\BookParagraph{如果你打算买比特币的话，最好先确认一下你买的确实是货真价实的比特币。}
\BookParagraph{如果你去买基于比特币的基金或者 ETF，那么你这一万块买到的其实是一种衍生产品：你的比特币不属于你，而是属于基金公司 —— 最多算是基金公司替你代持。}
\BookParagraph{那你研究过没有，如果该基金公司的比特币哪天被人偷走了，或者基金公司自己倒闭了，怎么办？}
\BookParagraph{这就是所谓的 counterparty risk: 基金公司作为中间人替你去买比特币，那你就必须多承担一点儿风险。理论上来说，基金公司可以将其管理的比特币找保险公司投保，那你的风险就相应转化为保险公司是否会理赔的风险，风险就有可能大大降低。}
\BookParagraph{但是我的研究表明，绝大多数基金公司并没有为其管理的比特币买任何保险。既然你是打算长期持有比特币，这样的风险不可忽略。}
\BookParagraph{所以如果你想买的是货真价实的比特币的话，我强烈建议你仔细研究一下有关区块链技术和比特币交易的知识，把你买来的比特币用一个硬件钱包（hardware wallet）离线「冷藏」（cold storage）起来。}
\BookParagraph{这样一来，你就可以避开任何有关基金公司代持的风险，唯一的要求是你必须记住一条二十四个单词组成的短语，做为你拥有比特币的专属密码。}
\BookParagraph{这二十四个单词太重要了：谁有了这条短语，谁就有可能卖掉你的比特币。所以你不能存在电脑中，不能上传到云服务器，也不能写在纸条上放到保险箱里 —— 你必须死记硬背下来，保证未来的三十年中任何时候都记忆犹新。}
\BookParagraph{我在研究比特币的过程中，还真的为这小小的二十四个单词从何而来的细节犹豫了半天。}
\BookParagraph{最终我决定，如果我买比特币，应该把它存在一首我最为熟悉的宋词中：}
\BookParagraph{「遥想公瑾当年，小乔初嫁了，雄姿英发。羽扇纶巾，谈笑间、樯橹灰飞烟灭。故国神游，多情应笑我，早生华发。人间如梦，一樽还酹江月。」}
\BookDate{2021 年 2 月，香港}
\BookArticle{偶遇}{偶遇}
\BookParagraph{一直觉得，一个线上举办的学术会议不能全是「理工科」的感觉，至少应该在细节里融入一点儿文艺的氛围。比如会议网站的设计既要简单好用又要漂亮大气 —— 去年初春，为了这点儿理想，我自己设计实现了一个全新的会议网站 \footnote{\href{https://duetone.com}{Duetone Corp.}}。}
\BookParagraph{我组织的一个会里还有一个小小的细节：我希望在每一个 session 之前加入一个五分钟的倒计时。倒计时的动画是用 Keynote 设计的，将几百个动画步骤一秒一秒地逐次播放出来，再配上钢琴演奏的背景音乐。}
\BookParagraph{选择背景音乐时花了一点儿时间。既要让人放松心情，不能喧宾夺主，又要有一点儿小小的高潮，最终朴素地转入主题。我去年选择的曲目叫「Almost a Love Story」；今年选的二个曲子之一叫「I Vow to Thee, My Country」\footnote{\href{https://www.bilibili.com/video/BV1gz4y117Jk}{INFOCOM 2021 5-minute countdown}}。}
\BookParagraph{虽然我不懂音乐，更不懂钢琴，还是暗自希望，当那极少的几位早到了几分钟的幸运观众「偶遇」到这一段倒计时和背景音乐的时候，心情可以更加舒展愉快。}
\BookDate{2021 年 3 月，香港}
\BookArticle{煲电话粥}{煲电话粥}
\BookParagraph{昨儿早上跟一人在香港的「港漂」朋友打电话聊天儿。各种新鲜八卦和随之而来的感慨，一没留神儿聊 high 了，聊了一个多钟头。}
\BookParagraph{我在美国上学的时候有个好朋友，一口极其标准流利的普通话、广东话和英文，我一度羡慕不已。这位朋友管这种聊法儿叫「煲电话粥」—— 哎呀，太形象了。}
\BookParagraph{后来我发现，我是个名副其实的「煲电话粥」爱好者。打电话自然不如见面聊，但是也没有相应的限制 —— 孩子不会在一边追逐打闹，饭馆儿不会打烊，不用边聊边听旁人唱 K，也不用编写剧情剧本杀。}
\BookParagraph{我的这位朋友后来谈了个女朋友，几年后由于工作关系，不得不分居两地。他俩于是每天晚上都要「煲粥」煲到很晚，情投意合，恩恩爱爱。}
\BookParagraph{直到有一天，朋友发现能聊的话题越来越少，有时候需要搜肠刮肚地找「谈资」了。甚至好一点儿的「谈资」也最好不能一次用光，抻这点儿用，明天还能继续聊。}
\BookParagraph{我一度以为，谈恋爱没了谈资，也就不太容易继续谈下去了。可我这位朋友本着「坚持就是胜利」「笑到最后笑得最甜」的准则，终于不但完美地结了婚，还调动了工作搬到一起住，「lived happily ever after」。}
\BookParagraph{可惜的是，在现在这个社交网络和手机过于发达、碎片信息大爆炸的时代，已经很少有人还能乐意「煲电话粥」了。我想，如果真的有人忽地拨打电话过来跟我简单纯粹地闲聊，我一定会开心地四处寻找「谈资」和充电线，「因过竹院逢僧话，偷得浮生半日闲」。}
\BookDate{2022 年 5 月，多伦多}
\BookArticle{礼轻情意重}{礼轻情意重}
\BookParagraph{时间长了，很多小礼物已经不记得了。俗话说「礼轻情意重」，我觉得礼物越轻就越是仔细斟酌过，弥足珍惜，如果还能有一份惊喜，就更加难能可贵，会一直记得。}
\BookParagraph{十几年前一个学生送给我一个开瓶器，现在学生早已失去了联系，但每次喝葡萄酒，都会不经意地想起。一位朋友送给我的手机无线充电器，还有另外一位朋友送的太阳镜，收到时都是意料之外。感动之余，每次出远门都会随身携带，经常可以用到。}
\BookParagraph{刚刚发现香港两位好朋友送了我一年的 Telegram Premium 会员，非常开心。大概是因为前一阵子聚会时聊天，聊到了 Telegram 最近推出的 premium 会员制。人家一查 Telegram，发现我这么了解，竟然没有入手，一定是觊觎已久吧。}
\BookParagraph{想起多年前看刘震云的小说《一地鸡毛》。主人公小林为了找领导办事儿，斟酌再三，既要拿得出手又不能太贵，最后送了一箱可口可乐。}
\BookDate{2022 年 8 月，多伦多}
\BookArticle{Hit and Roll}{Hit and Roll}
\BookParagraph{冰壶常规赛季结束了，我们队排名第二。昨天季后淘汰赛（Playoff）开始，我打二垒，第一场对阵排名第三的队。}
\BookParagraph{冰壶运动一场分为八局，四个人一队，一垒（lead）和二垒（second）不如三垒（vice）和四垒（skip）重要，后者可以随时扭转乾坤，左右胜负。}
\BookParagraph{第六局打完，不但比分二比五落后，对方还拥有最后掷球权。英文把最后掷球权形象地称为「锤子」，颇有一锤定音之意。若是能从「锤子」手里得分，英文叫做「您就偷着乐吧」（steal）。}
\BookParagraph{我们队超常发挥，竟然「偷着乐」了两局，第八局打完，追成五比五平。}
\BookParagraph{这可是新鲜事儿。我打了五年冰壶，从来没有过淘汰赛平局需要加赛的经历。}
\BookParagraph{至关重要的第九局，开场对方三个冰壶在营垒（house）里面，形势极其不利。我们 skip 叫 Joel，属于想赢怕输的性格，离得老远都能看到，他的脸色已经很难看了。}
\BookParagraph{我打第四个冰壶之前，Joel 用冰刷比划了一下，用意是让我「hit and roll」—— 把对方冰壶打出去的同时，自己的冰壶侧滑一小段，躲在对方其他冰壶的后方。}
\BookParagraph{冰壶两大要素是「line」和「weight control」。「line」是要不偏不倚地对准，而「weight control」则是力道要合适，冰壶速度不快不慢。这个「hit and roll」要把对方冰壶打出去，所以不能太轻；自己又要侧滑一小段，所以又不能太重。}
\BookParagraph{近一年来我冰壶水平有很大进步，竟然在关键时刻打出了一个完美的「hit and roll」：把对方的冰壶打出去以后，自己的冰壶 —— 英文形象地叫做「shooter」—— 竟然真的缓缓侧滑了一小段，停在了「house」的中心。}
\BookParagraph{最后我们队因为我这一个冰壶险胜。}
\BookParagraph{看到 Joel 喜笑颜开的样子，我想起了五年前刚认识他的时候，他还不在我所在的队。当时比赛结束大家一起喝酒，Joel 听说我要去香港访问，也是像此时此刻这样的神情，眉飞色舞地跟我吹：「你知道吗？香港有一家馆子叫 Mott 32，就在汇丰银行总部的楼下，他家的烤鸭惊艳无比，真的是一绝！」}
\BookParagraph{我当时差点儿跟他说：「你知道吗？我是一个北京人。」}
\BookDate{2023 年 4 月，多伦多}
\BookArticle{微信群}{微信群}
\BookParagraph{我其实不怎么用微信，也基本上没有时间看微信群。这几天想，是不是也可以建一个我自己的微信群呢？}
\BookParagraph{我以前之所以不喜欢微信群，是因为这些群要不然众说纷纭，热闹非凡，看不过来；要不然就只有广告和过年过节的祝福，没有什么值得看的内容。两者的共同之处在于「信噪比」太低，不值得花时间浏览。}
\BookParagraph{而在我的想象之中，一个趋于完美的微信群应该是这样的：没有人重复其他人早先发过的内容，每一段话都有一定的信息量，大部分发言都是段落清晰的长文，大部分推荐的都是高质量值得收藏的内容，并附有一小段推荐的原因。}
\BookParagraph{有没有可能有这样的微信群存在呢？我觉得这个问题挺好玩儿的，准备自己做一个试验。}
\BookParagraph{我的微信群是个稍纵即逝（ephemeral）的免费临时群：它大多数时间一直是冷场，随时都有可能消失，进来的朋友也随时可以退出。您如果对这个试验感兴趣，可以微信、微博私信或微博评论联系我，我先邀请其中两百人，以后再酌情邀请更多的朋友。}
\BookParagraph{群的名字就叫「葆春和他的朋友们」吧。}
\BookDate{2024 年 3 月，多伦多}
\BookArticle{细节}{细节}
\BookParagraph{我发现这个微信群有一个好处：有些和产品相关的想法，我不太敢发在微博 —— 怕人说是「带货」—— 但是正好可以在群里讨论。}
\BookParagraph{这个关于「带货」的执念也挺神奇的。无论你说点儿什么， 但凡是跟某产品有关，就一定有人觉得你是「充了值」，是「带货」，不知道是不是因为不「充值」不「带货」的人现在已经绝迹了。}
\BookParagraph{有一个很小的细节，我一直在琢磨着，琢磨了很多年。}
\BookParagraph{这就是如何建立和使用一个「To Do List」—— 中文还真搞不清如何翻译 —— 「待办事务清单」？}
\BookParagraph{To Do List 的特点，是一共只有两个动作：加入一个新的 To Do，和打勾完成一个已有的 To Do。加入新的 To Do 时，可能会有一个需要完成的日期。无论是什么 app，都应该支持显示现在需要完成的事项。这两个动作需要同时可以在手机或电脑上完成，To Do List 两者需要完美同步。}
\BookParagraph{关键的问题在于，这两个动作，都必须非常容易完成，否则我时间长了坚持不下去。可全靠脑子记也不是一个办法，因为不可避免的会忘掉一些重要的事情。}
\BookParagraph{以前曾经试用过很多 app，都感觉不是很好用。有的 app 太复杂 —— 但凡需要学习怎么用的 app 都太复杂 —— 而有些 app 呢，比如一个名叫「Clear」的 app，则又太简单，To Do List 太长时，无法只显示马上要做的工作，也无法和电脑同步。最重要的问题是，绝大多数 apps —— 比如 iOS 自带的 Reminders —— 输入一个需要完成的日期太麻烦，需要在手机或电脑上进行好几步操作。}
\BookParagraph{最近终于找到了我觉得比较完美的 To Do List，是用一个万能的叫做「Obsidian」的 super app 实现的。Obsidian 本身免费，但是和电脑同步需要一小笔订阅费。用 Obsidian 输入 To Do List 的好处是非常简单，它支持 Markdown 格式，只需要输入 「- [] Proposal reviews due @due (2024-03-22)」一行字，几秒钟就可以同时记录待办事项和需要完成的日期。Obsidian 里有一个叫做 Cardboard 插件，可以完成打勾的动作，自动记录完成时间，并可以显示今天需要完成的事项。}
\BookParagraph{Obsidian 还有很多其他功能和优点，我有机会以后再慢慢「带货」。但是这个关于 To Do List 多年悬而未决的难题，竟然就这么轻而易举地解决了。有了这么一个普通但是很好用的 To Do List，我完成工作的效率有明显的提高。平常不怕忘记了什么重要的工作，也不如以前那么焦虑了。}
\BookParagraph{这篇文字，也是在 Obsidian 中写成的。它自动同步到我的手机，还可以在发出之前简单修改一下，堪称完美。}
\BookDate{2024 年 3 月，多伦多}
\BookArticle{直播带货}{直播带货}
\BookParagraph{我其实不太明白为什么大家都这么喜欢看带货的直播。}
\BookParagraph{我在美国读博士的时候，网络还不怎么发达，一种主要的娱乐方式还是看电视。当然那时的电视是有线电视（cable TV），普通的电视频道是包在基础月费中的，但是比较好看的电影频道（比如 HBO）需要另外付费才能看到。我一般是付费看几个月，看到好电影用录像机录下来，当时的录像带到现在还在家里保存着。}
\BookParagraph{包在基础月费中我觉得不怎么好看的一个频道，就是 The Shopping Network (TSN），是一个一天二十四小时连轴转的直播带货频道。别的频道还可以看看新闻，历史，甚至学学养花（Home and Garden TV) 和做饭（The Food Network），这个带货频道却有一种一直在看广告的感觉：那些需要另外付费的电影频道在播放电影时有是也会插播广告，但是至少忍耐一下就过去了。而这个 TSN 则是一直在放广告。}
\BookParagraph{我一向认为，广告属于市场推广的一部分，无论是以前的有线电视时代还是现在的互联网时代，是需要企业花钱在更吸引人的内容中插播，才能吸引观众来看的。一个二十四小时播放广告的频道，我一直不是很理解为什么会有人主动去看。}
\BookParagraph{有一种解释挺合理的：看直播带货是因为比正常市场价便宜一些，有时间的观众看了，又恰巧买到了自己喜欢的物品，可以省钱。可问题在于，到底是否真的比市场价便宜，不太容易核实；而物品是否真的跟主持人说的那样好用，也并不一定。}
\BookParagraph{相比直播，我非常喜欢看一些名不见经传的普通人用过一样东西以后写的详细评价。无论是手电筒、小刀、手表、衣服，我一直觉得辅以照片的文字是最好的媒介 —— 看文字比看视频速度快得多。文字写起来虽然慢一些，但是可以随时迅速阅读，写得好的独立评价，如果这样东西正好我需要，我甚至看完会直接去买。}
\BookParagraph{可惜的是，我看过的这些评价，清一色都是英文的。在中文的世界里，视频正在全面取代文字。好看的文字，即使是「带货」，也越来越少了。}
\BookDate{2024 年 3 月，多伦多}
\BookArticle{两地}{两地}
\BookQuote{想请教朋友们一个关于人生选择的问题，希望大家不吝赐教。问题是：我该为了家庭放弃编制吗？}
\BookQuote{背景：我是一所小学副科老师，算上公积金年入 30 万，自己专业水平一般，现在工作是我人生的天花板，放弃编制去社会上找工作估计是个噩梦。孩子爹学历还行，跟随团队在另外一个省的新型科研单位做人工智能，收入比我高，他过来我这里的话又很难找到对口的方向。}
\BookQuote{我们婚后一直异地，他回家频率是每月 1-2 天。现在孩子快读小学了，很多问题在长期异地中越来越凸显，所以我一直想着要不要辞职过去那边结束异地。我年龄快超龄了，再备考事业编也没啥希望，很难割舍下现在的工作。现在我到底该怎么办？（问题来自「葆春和他的朋友们」）}
\BookParagraph{无论如何选择，你不要选择辞职。}
\BookParagraph{假设有两份工作：一份工作非常稳定，没有什么风险，无论大环境如何变迁，都可以保证或多或少一份工资，但是收入比较低。}
\BookParagraph{还有一份工作，收入相对高一些。将来的收入甚至有可能更上一层楼，但是随时都有可能被解雇。再找工作的风险相对比较大，不一定可以找到和自己能力和专业背景匹配的工作。}
\BookParagraph{你觉得应该选哪一份工作呢？}
\BookParagraph{在你的具体事例中，结束长期异地的方案有两种：一种是你辞职，一种是他辞职。咱们假设没有其他因素影响你的选择 —— 比如地理位置哪个更优越 —— 应该如何选择呢？}
\BookParagraph{这个问题其实在另外一种人才市场中，有很多人做出过选择。}
\BookParagraph{这就是具备博士学位且成绩出色的应届博士研究生毕了业，是去学术界做研究还是工业界研究机构的选择。}
\BookParagraph{大多数毕业生，都会优先选择学术界，即使学术界的收入比工业界低得多。}
\BookParagraph{原因虽然各异，但是其中之一是因为学术界一旦几年后拿到终身教职，基本上可以算是「铁饭碗」。}
\BookParagraph{当然也有选择工业界里的研究岗位的 —— 甚至有学术界已经拿到终身教职的教授，跳槽去工业界发展的。但是这需要一定的激励机制，而最重要的激励当然就是收入。据我粗浅的估计，工业界的收入至少要到学术界收入的三倍 —— 也就是税后收入的两倍 —— 以上，工业界的研究岗位才开始有竞争力。甚至这还得是在跨国高科技大公司，收入相对稳定，被解雇的风险相对比较低的情形。}
\BookParagraph{细心的读者会发现，我刚才讲的是一个人自己面对两种不同工作的选择，你的情况则略有不同。你的选择中，还包括了辞职后是否能找到新工作的风险。你辞职后顺利找到新工作的概率，明显低于他辞职后找到新工作的概率，选择的天平在向他辞职倾斜。如果我们同时考虑工作本身的稳定性和顺利找到新工作的概率，他的工资可能需要比你的工资高好几倍，才有可能开始认真考虑你辞职这种选择。}
\BookParagraph{显而易见，以上分析都是假设你们两人可以以平等而客观的角度，逻辑清晰、按部就班、就事论事地讨论这两种方案的选择。讨论的过程中，更看重中长期的风险规划。因为一般来说，随着时间的推移，家庭开销会逐渐上升，未来的收入和失去该收入的风险，都需要仔细考虑。}
\BookParagraph{你可以看到，我的观点是纯粹地从收入角度来看的，而没有提及孩子。这是因为我假设，无论孩子在哪里上学，只要异地的问题解决了，结果完全一致。如果两地中有更好的选择，我们当然可以继续给选择的天平加上权重，做一个加权的优化问题。}
\BookParagraph{如果咱们可以客观地看这个优化问题，我觉得结论没有异议：与其选择你辞职，不如选择他辞职。}
\BookDate{2024 年 3 月，多伦多}
\BookArticle{纸质书}{纸质书}
\BookParagraph{昨天新买了一台相机。打开包装，发现没有纸质版的使用说明书，只给了一个电子版说明书的链接。}
\BookParagraph{上次买相机还是十四年前。那时，包装里附上了英法双语两大本说明书，我还仔细研读了好长一段时间。这回没有了说明书可以翻阅，我坐在办公室的大屏幕前，全屏打开了刚下载好的电子版。虽然字大得出奇，也许是脑子还没适应，总是觉得看不下去。}
\BookParagraph{正好，在「葆春和他的朋友们」微信群中，有人提到了纸质版的《幸福》。这本书是我五年前花了很多时间整理排版出来的。光是校对版就印了五六本，重头学了书籍排版技术，颇花了不少心思。可惜的是，那年只有机会寄出了两百本书，真的像当年微博里戏称的，成了「法拉利限量版」。这两年所有回国的航班都需要转机，由于怕转机丢失行李，每次回国都只带了随身行李，基本没有机会再带书了。}
\BookParagraph{四年前，我又花了大量时间，把所有书里书外的文字，都整理发布在个人网站（baochun.ca）。这样想看这些文字的朋友，至少能够看得到。后来网站的代码一直在维护升级，搜索功能改了两版，现在总算是稳定下来了。唯一美中不足之处是，由于需要打开思源宋体这个字体，第一次加载网站速度不够快，需要等很长一段时间。}
\BookParagraph{现在纸质书和 CD 唱片、录音机、胶卷相机、现金等大量智能手机出现之前的事物的命运一样，早已退居二线了。《幸福》虽然现在网站上有电子版，也没什么人真正有时间看纸质版了，我还是觉得纸质版有点儿电子版所没有的吸引力。不说别的，仅仅是中英文采用什么字体，什么字号，留白留多少，都能给读者带来不同的体验。我小时候父母在出版社工作，常年都负责编辑大学课本，从小就看着他们用图章来标注字体和字号的选择，也算是耳濡目染了。}
\BookParagraph{这本书当时印刷「法拉利限量版」的时候，还有一段小插曲。书虽然没什么人看，我却执意希望能够用传统的胶印印刷。试了好几家印刷厂，都因为起印的数目太大而作罢。最后功夫不负有心人，终于找到一家位于魁北克省的印刷厂，量比较少也可以胶印，印好之后用卡车从魁北克千里迢迢运到了多伦多。遗憾的是，这家印刷厂质量不太灵：好多书里都有油墨还没干透涂抹的痕迹，不过还不算特别严重，至少字还能看清。我遗憾了很久，不过后来转念一想，也许这种与生俱来的不完美，正是纸质版魅力所在的一部分。}
\BookParagraph{令人感动的是，直到五年后的现在，还有朋友记挂着这本书，给我发微博私信，询问什么时候能收到书。不能一一回复之余，还是希望将来能有机会，可以陆续再带一些书回国 —— 即使是不完美的「法拉利限量版」，也不能只有两百本是吧。}
\BookDate{2024 年 8 月，香港}
\BookFinish
